\cleardoublepage
\phantomsection
\addcontentsline{toc}{part}{Ordered and Associative Chiral Geometry}
\chapter*{Bar constructions and their comparisons}
\addcontentsline{toc}{chapter}{Bar constructions and their comparisons}
\begin{summarybox}
The three bars use transverse associative composition, associative chiral
composition, and the chiral Lie bracket. Mixed distributive laws compare
iterated bars when they extend to the full bisimplicial object and tensor
products preserve its realizations. Enveloping comparisons use a
Poincar\'e--Birkhoff--Witt isomorphism and strongly convergent bar filtrations.
\end{summarybox}
\part{The invariant object}

\chapter{Two directions of locality}
\section{The local experiment}
Let $X$ be a smooth complex algebraic curve and let $J\subset\R$ be an
oriented interval. Consider a sector of observables that is holomorphic
along $X$ and locally constant along $J$. Two insertions at $(z_1,t_1)$
and $(z_2,t_2)$ can approach one another in either direction. The limit
$z_1\to z_2$ tests the singular coefficients along the diagonal of $X^2$.
Composition as $t_1\to t_2$ uses the order selected by the orientation
of $J$. A theory carrying both operations must specify their compatibility.

\begin{principle}[Two-direction principle]
\label{cbc:loc:two-directions}
Holomorphic collision uses the coefficient system on powers of $X$ and
its chiral operation maps. Transverse composition uses an internal
$\Eone$ multiplication. Exchange compares orders through a specified
transport on the configuration space.
\end{principle}

The geometry determines which operation must be supplied:
\[
\begin{array}{c|c}
\text{carrier}&\text{operation}\\ \hline
\text{diagonal in }X^2&\text{singular chiral operation}\\
\text{oriented interval}&\text{ordered composition and associative bar}\\
\text{oriented plane}&\text{transport along braids}\\
\text{circle}&\text{Hochschild trace}\\
\text{stable curve}&\text{modular sewing}.
\end{array}
\]
The trace, transport and sewing maps belong to the chosen coefficient
theory. For ordered word coefficients, the convolution
\eqref{eq:words-cut-product} implements composition by consecutive cuts.
Its realization by observables requires the corresponding operation maps.

\section{Carrier and operation}
Write $D_X$ for the sheaf of differential operators on $X$. Let $\g$ be
a right $D_X$-module, let $\Delta:X\hookrightarrow X^2$ be the diagonal,
and let $j:X^2\setminus\Delta(X)\hookrightarrow X^2$ be its open complement.
A binary chiral operation on $\g$ is a morphism of right $D_{X^2}$-modules
\[
 \mu:j_*j^*(\g\boxtimes\g)\longrightarrow \Delta_*\g.
\]

Here the direct images are those of $D$-modules. A chiral Lie algebra
requires this operation to be skew-symmetric and to satisfy Jacobi,
with compositions taken in the chiral pseudo-tensor category
\cite[\S\S0.8, 3.1.1--3.1.2]{BD04}. Thus $\g$ specifies the carrier,
while $\mu$ and its identities specify the Lie operation. A unit is
additional data.

A factorization family records complexes of $D$-modules on the powers
$X^S$ for finite sets $S$. It also supplies coherent collision
identifications and structure maps on disjoint loci. The chosen algebra or coalgebra
structure determines the direction of these maps. The full operator
product retains both the family and its operation maps.

\section{Principal symbols and residues}
The simplest collision coefficient already shows why the carrier matters.
Regard $\Delta(X)$ as a Cartier divisor, and set
\[
 F_m=\Ocal_{X^2}(m\Delta)\subset\Ocal_{X^2}(*\Delta),\qquad m\geq0.
\]
These are sheaves of meromorphic functions with pole order at most $m$.
Let $\mathcal I_\Delta$ be the ideal of the diagonal. Its normal line is
$N_{\Delta/X^2}=(\mathcal I_\Delta/\mathcal I_\Delta^2)^\vee$.

\begin{proposition}[Principal symbols along the diagonal]
\label{cbc:loc:normal-symbol}
For $m\geq1$, the pole filtration has the exact sequence of
$\Ocal_{X^2}$-modules
\[
 0\longrightarrow F_{m-1}\longrightarrow F_m
 \longrightarrow\Delta_*N_{\Delta/X^2}^{\tensor m}
 \longrightarrow0.
\]
In this sequence, $\Delta_*$ denotes the direct image of a sheaf on the
closed divisor. The leading symbol is therefore a section of a normal
line power on $\Delta$.
\end{proposition}
\begin{proof}
Tensor the divisor sequence
$0\to\Ocal_{X^2}(-\Delta)\to\Ocal_{X^2}\to\Delta_*\Ocal_\Delta\to0$
with the line bundle $\Ocal_{X^2}(m\Delta)$.
The restriction of $\Ocal_{X^2}(\Delta)$ to $\Delta$ is the normal line.
This gives the asserted quotient and its ambient sheaf.
\end{proof}

Locally choose formal coordinates $(z,u)$ with $z=z_2$ and $u=z_1-z_2$.
The class of $a(z)u^{-m}$ is its leading normal symbol. If
$v=\alpha(z)u+O(u^2)$ with $\alpha(z)\ne0$, then
$u^{-m}=\alpha(z)^m v^{-m}+O(v^{-m+1})$. This is the transition rule
for $N_{\Delta/X^2}^{\tensor m}$. Coordinate changes can mix the lower
pole terms, so the associated graded sheaf alone does not split the
pole filtration.

A scalar residue uses a relative one-form. For the projection to the
second factor, write such a form as $f(u,z)\,du$. Its residue along
$\Delta$ is the coefficient of $u^{-1}du$. For a simple pole this pairs
the normal symbol with the relative cotangent line. The coefficient is
unchanged under a change of the transverse parameter, including its
Jacobian in the differential.

Indeed, hold $z$ fixed and expand in the new parameter $v$.
For $m\ne1$, the form $v^{-m}dv=d(v^{1-m})/(1-m)$ is a Laurent
derivative, whose coefficient of $u^{-1}du$ vanishes. For a simple pole,
\[
 \frac{dv}{v}=\frac{du}{u}+\frac{d(v/u)}{v/u}.
\]
The last term is regular because $v/u$ is a unit. Linearity on the
finite polar part proves the claimed invariance.

\begin{example}[A double pole]
\label{cbc:loc:residue}
Set $v=u+u^2$ near $u=0$. Then
\[
 v^{-2}=u^{-2}-2u^{-1}+O(1),\qquad
 v^{-2}dv=\bigl(u^{-2}+O(1)\bigr)du.
\]
The function acquires a simple-pole term, while the one-form has zero
residue in either coordinate. For $m\geq1$, normal derivatives detect
higher pole orders in the coordinate $u$:
\[
 \partial_u^{m-1}(u^{-1})=(-1)^{m-1}(m-1)!\,u^{-m}.
\]
Thus a collision coefficient retains its filtered principal parts and
their coordinate transformations as well as its leading symbols.
\end{example}

Collisions of arbitrary subsets require these carriers and operation
maps on every finite power of $X$. Their compatible organization is
the purpose of the Ran coefficient category.

\chapter{The Ran coefficient category}
\section{Decompositions and alignment}
Let $X$ be a nonempty smooth separated algebraic curve over $\C$.
Two chiral products can divide the same configuration into different blocks.
Their comparison must retain the common divisions and commute with collisions.
For two labeled points there are two ordered divisions into nonempty blocks.
Two independent divisions give four pairs, of which two agree.
This finite projection has a geometric realization on collision coefficients.

For a nonempty finite set $I$ and $r\geq1$, put $[r]=\{1,\ldots,r\}$ and
\[
 \Dec_r(I)=\{d:I\twoheadrightarrow[r]\},\qquad
 I_a=d^{-1}(a).
\]
Thus $d$ is an ordered decomposition into $r$ nonempty blocks.
The block labels are ordered, while $I$ has no chosen order.

For example, $\{1,3\}\sqcup\{2\}$ is a division of $\{1,2,3\}$.
It is not a consecutive cut in the order $1,2,3$.
The consecutive-cut product \eqref{eq:words-cut-product} uses its own carrier.
A comparison with the present product must account for this difference.

Write $\Dec(I)=\Dec_2(I)$. Let $\mathcal C$ be a stable $\infty$-category and let
$M_{d,e}\in\mathcal C$ be objects indexed by $d,e\in\Dec(I)$.
Finite sums in $\mathcal C$ are also products, so they have canonical
inclusions and projections.

\begin{proposition}[Aligned sector in the finite decomposition model]
\label{cbc:ran:finite}
The object $M=\bigoplus_{d,e}M_{d,e}$ has the split idempotent
\[
 e_{\mathrm{al}}\big|_{M_{d,e}}=
 \begin{cases}\id,&d=e,\\0,&d\ne e.\end{cases}
\]
Its image is $\bigoplus_d M_{d,d}$. The projection is natural under
bijections of $I$. Repeated alignment, in any parenthesization, selects
the summands whose decomposition labels all agree.
\end{proposition}
\begin{proof}
Let $q$ be projection onto the indicated summands and $s$ their inclusion.
Then $qs=\id$ and $e_{\mathrm{al}}=sq$, which proves idempotence and the
image statement. A bijection preserves equality of decomposition labels.

For three labels, the two composites impose $d=e$ and $e=f$.
Both therefore select exactly $d=e=f$.
For any binary tree of alignments, equality at each vertex forces equality
of all leaf labels. The canonical finite-sum maps give the same splitting.
\end{proof}

To apply this projection to collision coefficients, the summands must
form a family compatible with diagonal pullback.

\section{Coefficients and chiral convolution}
Coefficient categories below are $\C$-linear stable $\infty$-categories.
Tensor products are derived. Scalar coefficient complexes use the usual
cohomological Koszul signs.
For a separated scheme $Y$ of finite type, let $\Dcal(Y)$ be the presentable
stable $\infty$-category of complexes of right $D$-modules localized
at quasi-isomorphisms.

The correspondence tensor uses de Rham direct image, as in
\cite[\S1.4]{FG12}.
A correspondence $Y\xleftarrow{a}Z\xrightarrow{b}W$ acts by $b_*a^!$.
Cartesian composition supplies coherent base-change equivalences,
compatible with external products.
For an open embedding $j$, restriction $j^!=j^*$ is left adjoint to $j_*$.

Let $\mathsf{Fin}_{\mathrm{surj}}$ be the category of nonempty finite sets
and surjections. A surjection $p:I\twoheadrightarrow J$ gives the diagonal
\[
 \Delta_p:X^J\longrightarrow X^I,\qquad
 (x_j)_{j\in J}\longmapsto(x_{p(i)})_{i\in I}.
\]
The positive Ran coefficient category is
\[
 \Dcal_X=\lim_{I\in\mathsf{Fin}_{\mathrm{surj}}}\Dcal(X^I),
\]
where the transition functor for $p$ is $\Delta_p^!$.
An object $F$ consists of complexes $F_I$ and coherent equivalences
$\Delta_p^!F_I\simeq F_J$ \cite[Definition 2.1.2]{FG12}.
The evaluations $F\mapsto F_I$ jointly detect equivalences.

Diagonal pullback defines a pointwise symmetric monoidal product:
\[
 (F\tensor^!G)_I=F_I\tensor^!G_I.
\]
Its unit $\omega_{\Ran}$ has components $\omega_{X^I}$, the dualizing
$D$-modules. Extraordinary pullback preserves these units and products.

A decomposition $d\in\Dec_r(I)$ determines the disjointness locus
\[
 U_d=\{(x_i)\in X^I:x_i\ne x_j\text{ whenever }d(i)\ne d(j)\},
 \qquad j_d:U_d\hookrightarrow X^I.
\]
Points in different blocks must be distinct. Points within one block
may coincide.

For $F_1,\ldots,F_r\in\Dcal_X$, define the chart expression
\begin{equation}
 \left(\mathop{\tensor^{\ch}}_{a=1}^r F_a\right)_I
 =\bigoplus_{d\in\Dec_r(I)}
 j_{d,*}j_d^!\left(\mathop{\boxtimes}_{a=1}^r(F_a)_{I_a}\right).
 \label{eq:cbc:ran:chiral}
\end{equation}
The external product uses $X^I=\prod_a X^{I_a}$.
The sum retains every ordered decomposition, without dividing by permutations of the blocks.
These charts describe the disjoint-union correspondence
\[
 \Ran(X)\times\Ran(X)\xleftarrow{j}
 (\Ran(X)\times\Ran(X))_{\mathrm{disj}}
 \xrightarrow{u}\Ran(X).
\]
For $F,G\in\Dcal_X$, the resulting chiral convolution is
\begin{equation}\label{eq:cbc:ran:star-convention}
 F\tensor^{\ch}G=u_*j^!(F\boxtimes G),
\end{equation}
with finite-chart expression~\eqref{eq:cbc:ran:chiral}
\cite[\S2.3.3 and Lemma~2.3.4]{FG12}.

The symbol $u_*$ denotes de Rham direct image.
An operation defined using extraordinary direct image $u_!$ would require its own domain and a comparison with this tensor.
No such comparison is included in the definition.
All aligned-projector and factorization statements below use~\eqref{eq:cbc:ran:star-convention}.

These finite-chart constructions require finite sums and the indicated base changes.
For a completed construction, compatibility with its inverse limits is an additional hypothesis.
Preservation of arbitrary totalizations is not assumed.

\begin{lemma}[Collision pullback]
\label{cbc:ran:collision}
Let $p:I\twoheadrightarrow J$ and $d\in\Dec_r(I)$.
If $d$ is nonconstant on a fiber of $p$, then $\Delta_p^{-1}(U_d)$ is empty.
Otherwise there is a unique $\bar d\in\Dec_r(J)$ such that
$d=\bar d\circ p$, and
\[
 \Delta_p^{-1}(U_d)=U_{\bar d}.
\]
Under base change, the summand of~\eqref{eq:cbc:ran:chiral} indexed by $d$
therefore vanishes or becomes the summand indexed by $\bar d$ on $X^J$.
These identifications are coherent under composition of surjections.
\end{lemma}
\begin{proof}
If $p(i)=p(j)$ but $d(i)\ne d(j)$, the diagonal requires $x_i=x_j$,
while $U_d$ requires $x_i\ne x_j$. Their intersection is empty.
If each fiber has one color, assign that color to its point in $J$.
This defines $\bar d$, uniquely because $p$ is surjective.
Its fibers are nonempty because those of $d$ are nonempty.
The defining inequalities then give the asserted equality of opens.

For a surviving block, the restriction
$p:I_a\twoheadrightarrow\bar d^{-1}(a)$ is surjective.
The Ran compatibility of $F_a$ identifies its diagonal pullback with
$(F_a)_{\bar d^{-1}(a)}$. Base change for $j_{d,*}$ gives the claimed
summand identification. Descent of a color through two successive
quotients equals descent through their composite.
Coherence of the correspondence functor supplies the same compatibility
for these pullback equivalences.
\end{proof}

The lemma makes~\eqref{eq:cbc:ran:chiral} an object of $\Dcal_X$.
Its associativity can also be seen on the charts.
Expanding either $(F\tensor^{\ch}G)\tensor^{\ch}H$ or
$F\tensor^{\ch}(G\tensor^{\ch}H)$ gives ordered divisions
$I=I_1\sqcup I_2\sqcup I_3$ into nonempty blocks.
In both expansions the open locus separates all three blocks, and the
coefficient is the restriction of $F_{I_1}\boxtimes G_{I_2}\boxtimes H_{I_3}$.

Base change identifies the iterated direct images with the direct image
from this common open locus. Every longer parenthesization uses the same
ordered division into leaf blocks. The coherent correspondence functor
therefore gives the associativity coherences in every arity.
Permuting the block labels supplies the symmetry.
Thus $\tensor^{\ch}$ is a symmetric monoidal product without a unit
on the positive Ran category.

\section{The geometric aligned projector}
Fix $F,G,H,K\in\Dcal_X$ and put
\[
 M=(F\tensor^{\ch}G)\tensor^!(H\tensor^{\ch}K),\qquad
 N=(F\tensor^!H)\tensor^{\ch}(G\tensor^!K).
\]
For $d=(A,B)\in\Dec(I)$, write
\[
 E_d(F,G)=j_d^!(F_A\boxtimes G_B)\in\Dcal(U_d).
\]
For $d,e\in\Dec(I)$, set $W_{d,e}=U_d\cap U_e$.
Let $\ell_d:W_{d,e}\hookrightarrow U_d$ and
$\ell_e:W_{d,e}\hookrightarrow U_e$ be the inclusions, and let
$k_{d,e}:W_{d,e}\hookrightarrow X^I$ be their common composite.
Define
\[
 \mathcal K_{d,e}
 =\ell_d^!E_d(F,G)\tensor^!\ell_e^!E_e(H,K)
 \quad\text{in }\Dcal(W_{d,e}).
\]

\begin{lemma}[The mixed correspondence]
\label{cbc:ran:mixed}
There is a canonical equivalence
\begin{equation}
 M_I\simeq\bigoplus_{d,e\in\Dec(I)}k_{d,e,*}\mathcal K_{d,e}.
 \label{eq:cbc:ran:mixed}
\end{equation}
For $d=e=(A,B)$, its summand identifies with
\[
 j_{d,*}j_d^!\bigl((F_A\tensor^!H_A)\boxtimes
                         (G_B\tensor^!K_B)\bigr).
\]
Consequently the sum of the aligned summands is $N_I$.
\end{lemma}
\begin{proof}
Expand both chiral factors and distribute the pointwise tensor over
their finite sums. Let $\delta_I:X^I\to X^I\times X^I$ be the diagonal.
The summand indexed by $(d,e)$ is
\[
 \delta_I^!(j_d\times j_e)_*
       \bigl(E_d(F,G)\boxtimes E_e(H,K)\bigr).
\]
The fiber product of $j_d\times j_e$ with $\delta_I$ is $W_{d,e}$.
Its base-change equivalence is exactly $k_{d,e,*}\mathcal K_{d,e}$.
This proves~\eqref{eq:cbc:ran:mixed} on the actual coefficient objects.

When $d=e$, the intersection is $U_d$.
Compatibility of external and pointwise products rearranges the factors
as $(F_A\tensor^!H_A)\boxtimes(G_B\tensor^!K_B)$ on that open locus.
The middle exchange uses the symmetry of the fixed correspondence theory.
Summing these identifications over $d$ gives
\eqref{eq:cbc:ran:chiral} for $N_I$.
\end{proof}

\begin{definition}[Aligned correspondence]
\label{cbc:ran:correspondence}
The labeled mixed correspondence and its aligned subspace are
\[
 Y_I=\coprod_{d,e\in\Dec(I)}W_{d,e},\qquad
 Y_I^{\mathrm{al}}=\coprod_{d\in\Dec(I)}U_d.
\]
Let $a_I:Y_I^{\mathrm{al}}\hookrightarrow Y_I$ select the components
with equal labels. Let $t_I:Y_I\to X^I$ restrict to $k_{d,e}$ on
each component.
\end{definition}

The inclusion $a_I$ is finite and clopen because it selects components
of a finite disjoint union. The functor $a_{I,*}a_I^!$ selects the aligned
part of the object with components $\mathcal K_{d,e}$ on $Y_I$.
Pushing its projection and inclusion along $t_I$ gives maps on $M_I$.
The clopen decomposition belongs to the labeled correspondence $Y_I$.
Different labeled components can have the same open image in $X^I$.

\begin{theorem}[Global aligned projector]
\label{cbc:ran:global}
There are natural maps in $\Dcal_X$
\[
 M\xrightarrow{q_{\mathrm{al}}}N
   \xrightarrow{s_{\mathrm{al}}}M,
 \qquad q_{\mathrm{al}}s_{\mathrm{al}}=\id_N.
\]
The composite $e_{\mathrm{al}}=s_{\mathrm{al}}q_{\mathrm{al}}$ is the
idempotent selecting the aligned summands of
\eqref{eq:cbc:ran:mixed}. These maps commute with collision pullback,
permutations and ordered refinement. Repeated alignments give the same
common-decomposition map under the canonical associativity identifications.
The map $q_{\mathrm{al}}$ is an equivalence exactly when every off-aligned
summand $k_{d,e,*}\mathcal K_{d,e}$ vanishes, for every $I$ and $d\ne e$.
\end{theorem}
\begin{proof}
On $X^I$, project to the aligned summands in
\eqref{eq:cbc:ran:mixed}, then use Lemma~\ref{cbc:ran:mixed} to identify
their sum with $N_I$. The inverse identification followed by inclusion
defines $s_{\mathrm{al},I}$. The finite-sum splitting gives
$q_{\mathrm{al},I}s_{\mathrm{al},I}=\id$.

Let $p:I\twoheadrightarrow J$ be a collision map.
Lemma~\ref{cbc:ran:collision} describes every component of
$Y_I\times_{X^I}X^J$. A component vanishes if either decomposition
separates a fiber of $p$. Every surviving pair is uniquely
$(\bar d\circ p,\bar e\circ p)$ with
$\bar d,\bar e\in\Dec(J)$.
Its pullback is $U_{\bar d}\cap U_{\bar e}$.

This gives a canonical isomorphism
\[
 Y_I\times_{X^I}X^J\simeq Y_J.
\]
Surjectivity of $p$ gives $\bar d\circ p=\bar e\circ p$ exactly when
$\bar d=\bar e$. The same isomorphism therefore identifies the pullback
of $Y_I^{\mathrm{al}}$ with $Y_J^{\mathrm{al}}$.

These are Cartesian comparisons of the labeled correspondences,
including their clopen inclusions. Base change for the coefficient
objects makes their projections and inclusions commute with $\Delta_p^!$.
For composed collisions, the unique descended labels and the Cartesian
comparisons compose. Coherence of the correspondence functor makes
these maps a transformation of the full Ran diagram.
They consequently define the asserted maps in its limit $\Dcal_X$.

For repeated pointwise products, each summand has one decomposition
label for each factor. Every parenthesized alignment keeps precisely
the tuples with all labels equal, as in Proposition~\ref{cbc:ran:finite}.

For chiral refinement, expand each binary tree into its ordered leaf
blocks. Equality at every vertex is equivalent to equality of the
resulting leaf decompositions. Both routes thus use the same open
locus, coefficient rearrangement and finite-sum projection.
Their graded permutations agree, and their pull--push identifications
are coherent by Cartesian composition. Permutations simply relabel
these constructions.

Finally, the fiber of $q_{\mathrm{al},I}$ is the sum of its off-aligned
summands. It vanishes exactly when each summand vanishes, since each
is a retract of that finite sum. The chart evaluations jointly detect
equivalences, which proves the last assertion.
\end{proof}

\begin{example}[Two disjoint points]
\label{cbc:ran:two-points}
Take $I=\{1,2\}$ and let $U=X^2\setminus\Delta(X)$.
The decompositions $d_+=(\{1\},\{2\})$ and
$d_-=(\{2\},\{1\})$ have the same disjointness locus $U$.
For $F=G=H=K=\omega_{\Ran}$, restriction to $U$ gives
\[
 M_I|_U\simeq\omega_U^{\oplus4},\qquad
 N_I|_U\simeq\omega_U^{\oplus2}.
\]
In the order $(+,+),(+,-),(-,+),(-,-)$, the projection is
\[
 q_{\mathrm{al}}|_U=
 \begin{pmatrix}1&0&0&0\\0&0&0&1\end{pmatrix},
 \qquad
 e_{\mathrm{al}}|_U=\operatorname{diag}(1,0,0,1).
\]
The two off-aligned copies survive on this open locus.
Thus the aligned comparison is generally a split projection, even
when its four inputs are the pointwise unit.
\end{example}

The same construction applies to $r$ ordered blocks.
Write $q_{\mathrm{al}}^{(r)}$ for the projection that pairs equal
elements of $\Dec_r(I)$. Its target pairs the coefficient factors
within each block.

To display its scalar signs, let $V_a,W_a$ be coefficient complexes.
Take block factors $V_a\tensor\omega_{X^{I_a}}$ and
$W_a\tensor\omega_{X^{I_a}}$. For homogeneous coefficients of degrees
$p_a$ and $q_a$, the rearrangement contributes the sign
\[
 (-1)^{\sum_{a<b}q_a p_b}.
\]
This is the sign for interleaving the two ordered lists of factors.
The refinement argument in Theorem~\ref{cbc:ran:global} applies in every arity.

\section{Factorization coalgebras and pointwise products}
A chiral commutative coalgebra is a noncounital commutative coalgebra
object in $(\Dcal_X,\tensor^{\ch})$.
It has coherent coproducts
$\delta_r:F\to\mathop{\tensor^{\ch}}_{a=1}^r F$ for $r\geq1$.
For $d\in\Dec_r(I)$, projection to the $d$-summand and open adjunction
give the map
\begin{equation}
 j_d^!F_I\longrightarrow
 j_d^!\left(\mathop{\boxtimes}_{a=1}^rF_{I_a}\right).
 \label{eq:cbc:ran:factorization}
\end{equation}
The coalgebra is \emph{factorizing} when these maps are equivalences
for every $I$, $r$ and $d$.
Write $\Fact_D(X)$ for the full subcategory of such coalgebras.
This is the positive convention of \cite[\S2.4]{FG12}.

\begin{proposition}[Pointwise closure]
\label{cbc:ran:pointwise}
The pointwise tensor product defines a symmetric monoidal structure
on $\Fact_D(X)$. Its unit is $\omega_{\Ran}$ with its canonical
factorization coproducts.
\end{proposition}
\begin{proof}
For chiral commutative coalgebras $F$ and $G$, define the $r$-fold
coproduct on $F\tensor^!G$ by
\[
 F\tensor^!G
 \xrightarrow{\delta_r^F\tensor^!\delta_r^G}
 \left(\mathop{\tensor^{\ch}}_{a=1}^rF\right)
 \tensor^!
 \left(\mathop{\tensor^{\ch}}_{a=1}^rG\right)
 \xrightarrow{q_{\mathrm{al}}^{(r)}}
 \mathop{\tensor^{\ch}}_{a=1}^r(F\tensor^!G).
\]
The common-decomposition maps commute with refinements and permutations.
Hence these maps satisfy the coherent commutative-coalgebra identities.
Repeated pointwise products use the same common-decomposition projections,
so their associativity and symmetry constraints are coalgebra maps.

On $U_d$, the $d$-component of this coproduct is the pointwise tensor
of the two maps~\eqref{eq:cbc:ran:factorization}, followed by the graded
rearrangement of external factors. If $F$ and $G$ factorize, every
one of these maps is an equivalence. Their product therefore factorizes.

For the unit, the identifications
$j_d^!\omega_{X^I}=\omega_{U_d}
 \simeq j_d^!(\boxtimes_a\omega_{X^{I_a}})$
give the component maps by adjunction.
Their finite sum defines each coproduct on $\omega_{\Ran}$.
The unit identifications are compatible with all correspondences,
so these coproducts satisfy the same refinement identities.
Under $F\tensor^!\omega_{\Ran}\simeq F$, the common-decomposition
formula recovers each coproduct of $F$.
Thus $\omega_{\Ran}$ is the unit in the coalgebra category as well.
\end{proof}

The chiral product also has a finite-arity filtration.
For $n\geq1$, let $\Dcal_{X,\leq n}$ be the same compatible-family
category restricted to sets $I$ with $1\leq|I|\leq n$.
Restriction gives
\[
 \Dcal_X\simeq\lim_n\Dcal_{X,\leq n}.
\]
The chart formula~\eqref{eq:cbc:ran:chiral} uses only smaller blocks,
so the chiral products descend to these categories and restriction
preserves them. Colimits of compatible families are computed chartwise:
each diagonal pullback preserves colimits.
Thus restriction preserves colimits as well.

Suppose $F\in\Dcal_{X,\leq n}$ vanishes on sets of size less than $n$.
For any $G\in\Dcal_{X,\leq n}$, every binary decomposition of a set
of size at most $n$ has both blocks smaller than $n$.
The chart formula therefore gives $F\tensor^{\ch}G=0$.
For $n=1$ the product is zero because no binary decomposition exists.

A limit of stable presentable tensor categories without units, with
these restriction and vanishing properties, is called a
\emph{pronilpotent tensor category}
\cite[Definition 14.10]{Heuts}.
This arity filtration supplies the setting for the operadic Koszul
equivalence of \cite[Theorem 14.13]{Heuts}.
It concerns the positive chiral product.
The pointwise product has the nonzero unit $\omega_{\Ran}$.

Let $\Delta_{\mathrm{main},*}:\Dcal(X)\to\Dcal_X$ be the left adjoint
to singleton evaluation, the direct image from the main diagonal
$X\to\Ran(X)$ \cite[\S2.1.5]{FG12}. An object is supported on this
diagonal when it belongs to the essential image of that functor.

The chiral Chevalley functor computes Lie homology in
$(\Dcal_X,\tensor^{\ch})$. Francis--Gaitsgory's theorem identifies
chiral Lie objects in this category with chiral commutative coalgebras
\cite[Theorems 1.2.4 and 5.1.1]{FG12}.
It restricts to factorization coalgebras exactly on Lie objects
supported on the main diagonal $X\to\Ran(X)$
\cite[Theorem 5.2.1]{FG12}.
The pointwise product constructed above permits an additional
transverse associative multiplication on a factorization coalgebra.

\section{The empty configuration and the two units}
On the singleton chart, $F\tensor^{\ch}G$ is zero because
$\Dec_2(\{1\})$ is empty. The positive chiral category therefore
has no tensor unit: the singleton component of $\omega_{\Ran}$ is nonzero.
An empty-configuration component supplies a unit.

Adjoin this component by setting
\[
 \Dcal_X^+=\operatorname{Vect}_{\C}\times\Dcal_X,
 \qquad \widehat F=(V,F),
\]
where $\operatorname{Vect}_{\C}$ denotes complexes of complex vector spaces.
The complex $V$ is the coefficient of the empty configuration.

The unital chiral product is
\[
 (V,F)\tensor^{\ch}(W,G)
 =\bigl(V\tensor W,
       (V\tensor G)\oplus(W\tensor F)\oplus(F\tensor^{\ch}G)\bigr).
\]
Here tensoring a complex with a Ran object uses its $\C$-linear structure.
The first two nonempty components correspond to the two possible empty blocks.
On each finite set, the expanded chart formula sums over all maps
$I\to[r]$, allowing empty fibers with their specified coefficients.
Ordered refinements prove its associativity exactly as before.

The unit is $\mathbbm1_{\ch}=(\C,0)$, as direct substitution shows.
The pointwise product on $\Dcal_X^+$ is instead
\[
 (V,F)\tensor^!(W,G)=(V\tensor W,F\tensor^!G),\qquad
 \mathbbm1_!=(\C,\omega_{\Ran}).
\]
The aligned construction extends to empty blocks: a surjection of
finite sets still reflects equality of their block-label maps.
The two units specify distinct targets for augmentations of the
two associative structures.

On $\Dcal_X^+$, both products preserve colimits separately in each variable.
For an inverse system $(F_n)$, passage to a completed tensor product
requires the comparison
\[
 (\lim_n F_n)\tensor^?G\longrightarrow
       \lim_n(F_n\tensor^?G),\qquad ?\in\{!,\ch\},
\]
to be an equivalence for the system in question.
This limit compatibility is an additional hypothesis of the
corresponding completed bar construction.

\chapter{The two noncommutative chiral objects}
\label{cbc:two:chapter}
Two operators can be composed at the same chiral support or on disjoint
chiral configurations. The first composition can record their order
along a transverse line. The second uses chiral convolution.
Their domains differ even before associativity is imposed.
We shall construct both products on one factorizing coefficient family
and vary them independently.

\section{The two products and their units}
Retain the smooth complex curve $X$ and the positive Ran category $\Dcal_X$.
Its empty-configuration extension is $\Dcal_X^+$.
Write $\starCh$ for the chiral product $\tensor^{\ch}$.
The notation $\Alg_{\Eone}^{\mathrm{aug}}(\mathcal C)$ means associative
unital algebra objects of a monoidal $\infty$-category $\mathcal C$,
with an algebra map to its tensor unit.
Associativity includes its coherent higher homotopies.

\begin{definition}[Transversely ordered chiral algebra]
\label{cbc:two:transverse}
A transversely ordered chiral algebra is an object
\[
 A^\perp\in\Alg_{\Eone}^{\mathrm{aug}}(\Fact_D(X),\tensor^!).
\]
Its multiplication $m^\perp:A^\perp\tensor^!A^\perp\to A^\perp$
and its unit $\omega_{\Ran}\to A^\perp$ are maps of factorization
coalgebras. We also call it an \emph{ordered chiral algebra}.
\end{definition}

The superscript $\perp$ records an independent oriented direction.
When a field theory supplies locally constant fusion along that line,
the order of transverse intervals determines the order of multiplication.
The definition itself specifies the coefficient object and its maps.

\begin{definition}[Associative chiral algebra]
\label{cbc:two:chiral}
An associative chiral algebra is an object
\[
 \widehat A^{\ch}_{\mathsf{Ass}}
 \in\Alg_{\Eone}^{\mathrm{aug}}(\Dcal_X^+,\starCh).
\]
For an augmentation whose empty component is the identity of $\C$,
write $\widehat A=(\C,A)$. Its augmentation ideal is the positive object
$A$, with an associative product $m^{\ch}:A\starCh A\to A$.
The chiral unit is $\mathbbm1_{\ch}=(\C,0)$.
\end{definition}

The commutator $m^{\ch}-m^{\ch}\circ\tau$, where $\tau$ is the
graded symmetry, satisfies the chiral Jacobi identity.
It gives a Lie object in $(\Dcal_X,\starCh)$ on the positive ideal.
For strict differential graded products this follows by expanding Jacobi
and cancelling the associative triple products with their Koszul signs.
For coherent products it is the induced map from the Lie operad
to the associative operad.

The pointwise unit in $\Dcal_X^+$ is instead
$\mathbbm1_!=(\C,\omega_{\Ran})$.
Thus a common object can have two augmentations with different targets
and different ideals. No identification of those ideals is implicit.

\section{A factorizing family supported at a marked point}
Fix $x\in X(\C)$ and a finite-dimensional complex vector space $V$
in degree zero. Let $i_I:\{x^I\}\hookrightarrow X^I$ be the closed
embedding of the point whose coordinates all equal $x$. Define
\[
 M(V)_I=i_{I,*}V,\qquad
 L(V)=M(V)\oplus\bigl(M(V)\starCh\omega_{\Ran}\bigr),\qquad
 N(V)=\omega_{\Ran}\oplus L(V).
\]
Closed direct image is taken in the category of $D$-modules.
In particular, $i_{I,*}V$ retains the normal derivatives at $x^I$.
Proper base change equips $M(V)$ with its Ran collision equivalences.

The other two objects use finite sums on each chart.

On two distinct points, at most one point can equal $x$.
Consequently two marked contributions cannot occur in two separated
groups of points. This observation makes $N(V)$ factorize despite its
description as a sum. General componentwise sums have no such property.

For a nonempty subset $S\subseteq I$, let
$j_S:U_S\hookrightarrow X^I$ be the locus separating $S$ from
$I\setminus S$. If $S=I$, take $U_S=X^I$ and $j_S=\id$.
Put
\[
 E_{I,S}(V)=j_{S,*}j_S^!
 \bigl(i_{S,*}V\boxtimes\omega_{X^{I\setminus S}}\bigr).
\]
For $S=I$ this expression means $i_{I,*}V$.
The chiral chart formula gives the decomposition
\begin{equation}
 L(V)_I=\bigoplus_{\varnothing\ne S\subseteq I}E_{I,S}(V).
 \label{eq:cbc:two:marked-sum}
\end{equation}
Before open direct image, its $S$-coefficient sets every $S$ coordinate
equal to $x$ and requires every complementary coordinate to differ from $x$.

\begin{lemma}[Marked-subset comparisons]
\label{cbc:two:subsets}
Let $p:I\twoheadrightarrow J$ be a collision map.
The pullback $\Delta_p^!E_{I,S}(V)$ vanishes unless
$S=p^{-1}(S')$ for a nonempty subset $S'\subseteq J$.
In that case it is canonically $E_{J,S'}(V)$.
For finite-dimensional complex vector spaces $V,W$, there are natural equivalences
\begin{equation}
 L(V)\tensor^!L(W)\simeq L(V\tensor W).
 \label{eq:cbc:two:marked-tensor}
\end{equation}
Both comparisons commute with collisions, relabeling and repeated products.
\end{lemma}
\begin{proof}
If a fiber of $p$ meets both $S$ and its complement, the collision
equates two coordinates that $U_S$ separates.
Its inverse image is empty, so open base change gives zero.
Otherwise membership in $S$ descends uniquely to $J$.
Closed and open base change identify the surviving coefficient with
$E_{J,S'}(V)$. Descent of a subset through successive surjections agrees
with descent through their composite.

Expand the left side of~\eqref{eq:cbc:two:marked-tensor} using
\eqref{eq:cbc:two:marked-sum} twice.
The mixed-correspondence calculation of Lemma~\ref{cbc:ran:mixed}
places the $(S,T)$ coefficient on $U_S\cap U_T$.
If $S\ne T$, choose an index in one subset outside the other.
The two closed coefficients force it and an index of the other subset
to equal $x$. The latter subset's open condition separates them.
Thus this coefficient is zero.

When $S=T$, the closed embedding $i_S$ is the same in both factors.
Apply diagonal base change to $i_S\times i_S$.
It identifies their extraordinary tensor with $i_{S,*}(V\tensor W)$.
The remaining factors are dualizing units.
Open base change therefore gives $E_{I,S}(V\tensor W)$.
Summing over $S$ proves~\eqref{eq:cbc:two:marked-tensor}.

For any number of factors, the identical argument keeps precisely
the tuples with all marked subsets equal.
The coefficient there is the tensor product of all the vector spaces.
The coherent correspondence functor identifies these maps for every
parenthesization. Surjections reflect equality of their inverse-image
subsets, which proves compatibility with collisions as well.
\end{proof}

\begin{proposition}[The common factorizing carrier]
\label{cbc:two:factorization}
The object $N(V)$ has a canonical structure in $\Fact_D(X)$.
Its inclusion $\omega_{\Ran}\to N(V)$ and projection
$N(V)\to\omega_{\Ran}$ are maps of factorization coalgebras.
Suppose $V$ has an associative multiplication $\mu:V\tensor V\to V$.
Every $t\in\C$ then gives an augmented pointwise algebra on $N(V)$.
Its ideal product on $L(V)$ is $tL(\mu)$ under
\eqref{eq:cbc:two:marked-tensor}.
\end{proposition}
\begin{proof}
Fix a separated decomposition $I=\bigsqcup_{a=1}^r I_a$.
An $S$-coefficient of $L(V)_I$ survives restriction only if $S$
lies in a single part $I_a$.
Two different parts cannot both contain a coordinate equal to $x$.
For $S\subseteq I_a$, base change identifies this restriction with
the corresponding coefficient of
\[
 L(V)_{I_a}\boxtimes\mathop{\boxtimes}_{b\ne a}\omega_{X^{I_b}}.
\]
All complementary coordinates in other parts differ from $x$ already
by separation from $S$. Hence these identifications exhaust the
restricted sum~\eqref{eq:cbc:two:marked-sum}.

Now expand $\boxtimes_a N(V)_{I_a}$.
Every term containing two $L(V)$ factors vanishes on the separated
open by the same argument. Its surviving terms are the dualizing
unit and the single marked-part terms just identified.
This gives the required factorization equivalence.
Open adjunction supplies the coproduct components on $X^I$.

Under refinement, a surviving marked subset has a unique containing
final part. Each iterated comparison uses this part and the same
closed and open base-change maps.
They therefore give coherent coproducts in all arities.
Relabeling merely permutes the parts. The scalar summand is preserved
by these maps, which proves the unit and augmentation assertions.

For the multiplication, use the ordinary unital extension of the
possibly nonunital product $t\mu$ on the ideal $L(V)$.
The comparisons~\eqref{eq:cbc:two:marked-tensor} respect all iterated
tensors, so associativity reduces to associativity of $t\mu$.
On a separated open at most one part carries an ideal contribution.
The same multiplication acts in that part, while the others have
their scalar products. Thus the multiplication respects every
factorization map and its refinements.
\end{proof}

For $V=M_2(\C)$ with matrix multiplication, extraordinary pullback
to the singleton $x$ identifies this algebra with $\C\oplus M_2(\C)$
and product
\begin{equation}
 (a,v)(b,w)=(ab,aw+bv+t\,vw).
 \label{eq:cbc:two:matrix}
\end{equation}
The unit is $(1,0)$ and the augmentation is $(a,v)\mapsto a$.
At $t=0$ the ideal has square zero.
For matrix units, the commutator of $(0,E_{12})$ and $(0,E_{21})$
is $(0,t(E_{11}-E_{22}))$, so the product is noncommutative for $t\ne0$.

\section{An independently variable chiral product}
Write
\[
 P(V)=M(V)\starCh\omega_{\Ran},\qquad
 N(V)=\omega_{\Ran}\oplus M(V)\oplus P(V).
\]
The summands and their projections are Ran morphisms by
Lemma~\ref{cbc:two:subsets}.
Let $q:N(V)\to M(V)$ and $\epsilon:N(V)\to\omega_{\Ran}$
be the indicated projections, and let $i:P(V)\to N(V)$ be inclusion.
For $\lambda\in\C$, define
\begin{equation}
 \mu_\lambda=\lambda\,i\circ(q\starCh\epsilon):
 N(V)\starCh N(V)\longrightarrow N(V).
 \label{eq:cbc:two:chiral-product}
\end{equation}

\begin{theorem}[Independent products on one factorizing family]
\label{cbc:two:independence}
For every $t,\lambda\in\C$, the object
$\widehat N=(\C,N(M_2(\C)))$ carries both the pointwise algebra
of Proposition~\ref{cbc:two:factorization} and an augmented chiral algebra
whose positive product is~\eqref{eq:cbc:two:chiral-product}.
Its factorization coalgebra is fixed as $t$ and $\lambda$ vary.
Fixing either multiplication leaves the other parameter free.
\end{theorem}
\begin{proof}
Since $qi=0$ and $\epsilon i=0$, both nested positive triple products
of $\mu_\lambda$ are zero.
In the left-associated product the outer $q$ kills the inner image.
In the right-associated product the outer $\epsilon$ kills it.
This proves associativity. Adjoining the empty component gives its
unital extension with unit $(\C,0)$ and augmentation to that unit.
The split projections make every nonunit product of length at least
three vanish, giving the higher associative coherences as well.

The positive product is nonzero for $\lambda\ne0$.
On a two-point chart $(x,y)$ with $y\ne x$, take the first input
in $M(V)$ and the second in the positive dualizing summand.
The corresponding ordered decomposition maps identically, multiplied
by $\lambda$, to its summand of $P(V)$.
Reversing these two inputs gives zero because $q$ kills the dualizing
summand. Thus this chiral product is also noncommutative when
$\lambda\ne0$ and $V\ne0$.

The pointwise multiplication extends to the empty component by scalar
multiplication. Its augmentation has target $(\C,\omega_{\Ran})$
and ideal $(0,L(V))$, whereas the chiral augmentation has ideal
$(0,N(V))$.
Neither formula for a multiplication uses the other parameter.
The factorization maps in Proposition~\ref{cbc:two:factorization}
use only the marked subsets and $V$, so they remain fixed.

For each product, the zero parameter makes the augmentation ideal's
multiplication zero, while a nonzero parameter gives a nonzero map.
These properties are invariant under equivalence of augmented algebras.
Thus the zero and nonzero choices are inequivalent in the respective
algebra categories.
\end{proof}

\section{Unitization and factorization}
The marked-point construction has a general algebraic part: adjoining a unit.
Its geometric factorization uses the additional support calculation.

\begin{proposition}[Split unitization]\label{cbc:two:unitization}
Let $(\mathcal C,\otimes,\mathbbm1)$ be a stable unital monoidal category whose tensor preserves finite sums separately.
Every nonunital associative algebra object $M$ has an augmented unitization $\mathbbm1\oplus M$.
On the four summands of its tensor square, multiplication is the unit multiplication, the two unit actions, and the product of $M$.
The unit is the inclusion of $\mathbbm1$, and the augmentation is its projection.
For any object $M$, choosing the zero nonunital algebra structure gives the square-zero unitization.
\end{proposition}
\begin{proof}
For each arity $r$, expand $(\mathbbm1\oplus M)^{\otimes r}$ over subsets of the $r$ input positions.
Remove unit factors by the monoidal unit constraints.
On the remaining ordered $M$ factors apply the given associative operation, and include its value in $M\subset\mathbbm1\oplus M$.
When the subset is empty, use the unit summand.

For a composite operation, removing the unit positions in stages gives the same ordered $M$ positions as removing them at once.
The associativity coherences of $M$ and the monoidal unit coherences therefore supply all substitution and unit identities.
Projection to $\mathbbm1$ respects these operations, since every term with an $M$ input has zero scalar component.
\end{proof}

This construction applies in the ambient coefficient categories.
To obtain an object of $\Alg_{\Eone}^{\mathrm{aug}}(\Fact_D(X),\tensor^!)$, one must also supply a factorization coalgebra on $\omega_{\Ran}\oplus M$.
Its multiplication, unit and augmentation must be coalgebra maps.
Componentwise direct sums alone need not admit this structure.

For example, $\omega_{\Ran}\oplus\omega_{\Ran}$ has two dualizing summands on both the singleton chart and every separated configuration chart.
Factorization over $r$ separated singleton blocks would require $2^r$ summands on the latter chart.
For $r=2$, the required rank is four, whereas the given rank is two.
Thus the factorization equivalence already fails on two distinct points.

The marked-point family supplies the extra factorization maps on an actual common carrier and proves independence of the two products.
A mixed distributive law additionally requires its full coherent data.
The family uses a fixed point of $X$, so translation-equivariant field realizations require further structure.

\section{The mixed compatibility equation}
Independence leaves a separate question: when do the two products
commute with one another?
For a common object $\widehat A\in\Dcal_X^+$, set
\[
 Q=(\widehat A\tensor^!\widehat A)\starCh
                  (\widehat A\tensor^!\widehat A).
\]
The aligned section constructed in Theorem~\ref{cbc:ran:global},
extended to empty blocks, is a map
\[
 s_{\mathrm{al}}:Q\longrightarrow
 (\widehat A\starCh\widehat A)\tensor^!
 (\widehat A\starCh\widehat A).
\]
It allows two composites with the same source and target:
\begin{equation}
 m^\perp\circ(m^{\ch}\tensor^!m^{\ch})\circ s_{\mathrm{al}}
 \quad\text{and}\quad
 m^{\ch}\circ(m^\perp\starCh m^\perp):Q\longrightarrow\widehat A.
 \label{eq:cbc:two:mixed}
\end{equation}

\begin{definition}[Mixed associative chiral datum]
\label{cbc:two:mixed}
A mixed associative chiral datum consists of a common object with
the two unital associative products and a factorization coalgebra
structure on its positive part.
The pointwise multiplication and unit are coalgebra maps.
It also includes a coherent aligned distributive law.
The latter includes a homotopy between the maps
\eqref{eq:cbc:two:mixed}, compatible with both units, every iterated
product, relabeling and collision. Its higher homotopies are part
of the datum.
\end{definition}

A binary equation and the two triple-product faces check only the
first coherence conditions. A construction by a mixed operad or by
iterated bars must supply its remaining diagrams as well.

The family above detects a failure already on a unit slice of
\eqref{eq:cbc:two:mixed}.
Let $u$ be the pointwise unit and choose $a$ in a marked $M(V)$ coefficient.
Let $b$ lie in the positive dualizing coefficient at a separated point.
Evaluate the two composites on
$(a\tensor^!u)\starCh(u\tensor^!b)$.
The first is zero because $m^{\ch}(u,b)=0$.
The second is $m^{\ch}(a,b)$, which is nonzero for $\lambda\ne0$.

Extraordinary pullback to this configuration gives zero for the first
map and $\lambda\id_V$ for the second, on $V$ in degree zero.
A homotopy therefore cannot equate them.
This particular family with $\lambda\ne0$ therefore fails the aligned
mixed law. The calculation imposes no obstruction on other choices
of carrier and operations.

\section{Formal fields and their comparison problem}
Choose a formal coordinate $z$.
For a complex vector space $V$, a field on $V$ is an element of
$\operatorname{Hom}_{\C}(V,V((z)))$.
Its pole bound may depend on the vector to which it is applied.
It need not belong to $(\operatorname{End}_{\C}V)((z))$, where one
bound would work on every vector.

\begin{definition}[Nonlocal vertex algebra]
\label{cbc:two:nonlocal}
A nonlocal vertex algebra consists of a vector space $V$, a vector
$\mathbf1\in V$, and a linear map
\[
 Y:V\longrightarrow\operatorname{Hom}_{\C}(V,V((z))),
 \qquad a\longmapsto Y(a,z),
\]
with the vacuum and creation properties
\[
 Y(\mathbf1,z)a=a,\qquad
 Y(a,z)\mathbf1\in V[[z]],\qquad
 Y(a,z)\mathbf1\big|_{z=0}=a.
\]
It also satisfies weak associativity: for every $a,b,c\in V$ there
is $N\geq0$ such that
\begin{equation}
 (z+w)^N\iota_{z,w}Y(a,z+w)Y(b,w)c
   =(z+w)^NY(Y(a,z)b,w)c.
 \label{eq:cbc:two:weak-assoc}
\end{equation}
Here $\iota_{z,w}$ expands in nonnegative powers of $w$ relative to $z$.
The factor is chosen so that the indicated substitution and equality
of formal distributions are defined.
\end{definition}

These are the axioms of \cite[Definition 2.3]{LiNonlocal}.
Creation already makes $a\mapsto Y(a,z)$ injective.
The translation operator
$Ta=\partial_zY(a,z)\mathbf1|_{z=0}$ satisfies
\[
 Y(a,z)\mathbf1=e^{zT}a,\qquad
 [T,Y(a,z)]=Y(Ta,z)=\partial_zY(a,z)
\]
by \cite[Lemma 2.5]{LiNonlocal}.
Thus weak associativity is an axiom here. A reconstruction theorem
starts with fields on a different space and proves these axioms for
the space they generate.

For example, let $B$ be a unital associative algebra with derivation $T$.
The formula
\[
 Y(a,z)b=(e^{zT}a)b
\]
defines a nonlocal vertex algebra with no poles.
Indeed, $e^{wT}$ preserves multiplication, and both sides of
\eqref{eq:cbc:two:weak-assoc} with $N=0$ equal
$(e^{(z+w)T}a)(e^{wT}b)c$.
The vacuum and creation properties follow from $T1=0$.
Taking $B=M_2(\C)$ and $T=0$ gives noncommuting constant fields.

\emph{$S$-locality} over $\C$ adds a finite exchange relation.
For each $a,b$ it requires vectors $a_i,b_i$, series
$f_i(z)\in\C((z))$ and an integer $k\geq0$ such that
\begin{equation}
 (z-w)^kY(a,z)Y(b,w)
 =(z-w)^k\sum_i\iota_{w,z}f_i(w-z)Y(b_i,w)Y(a_i,z).
 \label{eq:cbc:two:s-local}
\end{equation}
A nonlocal vertex algebra with this property is a weak quantum
vertex algebra \cite[Remark 5.6]{LiNonlocal}, \cite[\S3]{LiQVA}.
The expansion convention in~\eqref{eq:cbc:two:s-local} is part of
the identity. A chosen exchange operator must reproduce it.

There is also a formal deformation convention.
Let $\hbar$ be a formal parameter and let $V=V_0[[\hbar]]$ be
a topologically free $\C[[\hbar]]$-module.
Here $V_0$ is a complex vector space.
Write $V\widehat\otimes V=(V_0\otimes V_0)[[\hbar]]$ and define
\[
 V_\hbar((z))=\varprojlim_m(V/\hbar^mV)((z)).
\]
Thus a Laurent series has a finite pole bound modulo each power
of $\hbar$, with no common bound required as $m$ grows.

A topological state--field correspondence consists of a continuous
$\C[[\hbar]]$-linear map $Y:V\widehat\otimes V\to V_\hbar((z))$
and a $\C[[\hbar]]$-linear operator $T$.
The map $Y$ has the vacuum and creation properties above, and $T$ satisfies
\[
 T\mathbf1=0,\qquad
 \partial_zY(z)=TY(z)-Y(z)(1\otimes T)=Y(z)(T\otimes1).
\]
A braiding is a continuous $\C[[\hbar]]$-linear map
\[
 S(z):V\widehat\otimes V\longrightarrow
 (V_0^{\otimes2}\otimes_{\C}\C((z)))[[\hbar]],
 \qquad S(z)=1+O(\hbar).
\]
Its topological locality condition is that, for every $a,b$ and $m$,
there is $k\geq0$ such that for all $c$
\begin{equation}
 \begin{aligned}
 &(z-w)^kY(z)(1\otimes Y(w))
       \bigl(S(z-w)(a\otimes b)\otimes c\bigr)\\
 &\hspace{2em}=(z-w)^kY(w)(1\otimes Y(z))(b\otimes a\otimes c)
       \pmod{\hbar^m}.
 \end{aligned}
 \label{eq:cbc:two:top-local}
\end{equation}
The Laurent coefficients are expanded in the order used by each
field product, with the factor clearing the required poles.

A \emph{topological quantum vertex algebra} in the convention of
\cite[Definitions 3.5 and 3.12]{DSGKQVA} is this data satisfying
\eqref{eq:cbc:two:top-local} and weak associativity modulo every
$\hbar^m$. The exponent in weak associativity may depend on $a,b,c,m$.
This convention states associativity explicitly.

There is a useful sufficient exchange condition.
Let $S^{ij}$ act on tensor factors $i,j$, leaving the third fixed.
If the topological state--field correspondence satisfies
\eqref{eq:cbc:two:top-local} and the exact hexagon identity
\begin{equation}
 S(w)(Y(z)\otimes1)
  =(Y(z)\otimes1)S^{23}(w)\iota_{w,z}S^{13}(z+w),
 \label{eq:cbc:two:hexagon}
\end{equation}
then it is a topological quantum vertex algebra
\cite[Propositions 3.8 and 3.13]{DSGKQVA}.

Indeed, locality gives the quasi-associative identity with the
extra factor $S^{23}(w)S^{13}(z+w)$.
Substituting~\eqref{eq:cbc:two:hexagon} puts that factor on both
sides. Modulo each $\hbar^m$, one clears its poles and cancels its
inverse, which exists because $S=1+O(\hbar)$.
The resulting identity is weak associativity modulo $\hbar^m$.

The converse hexagon follows only after composition with the field
map unless an additional injectivity statement is proved
\cite[Proposition 3.14]{DSGKQVA}.
Likewise, operator unitarity and the Yang--Baxter equation must be
distinguished from the identities after evaluation by fields.
Injectivity of $a\mapsto Y(a,z)$ alone does not give the required
injectivity on tensor products of states.

These field axioms specify possible recipients of a formal-disc realization.
The geometric-to-field map must supply the state space, vacuum and state--field operation, with the following compatibility.

\begin{proposition}[Conditional formal-disc field realization]
\label{cbc:two:field-realization}
Suppose an associative chiral algebra has a specified formal-disc realization with state space $V$, vacuum $\mathbf1$ and field map $Y$.
Assume this map satisfies the vacuum, creation and lower-truncation conditions of Definition~\ref{cbc:two:nonlocal}.
Suppose the comparison with geometric iterated products proves~\eqref{eq:cbc:two:weak-assoc} for every triple of states.
Then $(V,Y,\mathbf1)$ is a nonlocal vertex algebra.

In the topological convention above, a topologically free state--field system satisfying~\eqref{eq:cbc:two:top-local} and~\eqref{eq:cbc:two:hexagon}, with $S=1+O(\hbar)$, is a topological quantum vertex algebra.
The stated translation, vacuum and completion conventions are part of this hypothesis.
\end{proposition}
\begin{proof}
The first assertion follows from the defining nonlocal vertex axioms.
Creation also makes $a\mapsto Y(a,z)$ injective by evaluation on the vacuum at $z=0$.
The second assertion is the locality--hexagon implication proved above, using the inverse of $S$ modulo each power of $\hbar$.
These implications use the specified realization or state--field system as input.
\end{proof}

The condition $S=1+O(\hbar)$ and the $\hbar$-adic completions describe this formal deformation class.
An arbitrary unitary rational Yang--Baxter operator does not include those data.
Its comparison with this class must specify the deformation parameter, completion and exchange normalization.

\begin{problem}[The reverse geometric construction]\label{cbc:two:reverse}
For a given nonlocal or quantum vertex algebra, construct a corresponding formal-disc coefficient system.
Fix the category of field systems and the geometric coefficient category, with their morphisms, supports and completions.
Construct the coefficients on every finite power of the disc, their collision transitions, relabeling maps, factorization comparisons and associative chiral operations.
Supply every required coherence, and identify the recovered state space, vacuum and field operations with the original field system.
The empty chiral unit and the chosen vacuum require their own comparison.

An equivalence of categories further requires both functors and coherent natural equivalences from both composites to the respective identity functors.
Determine the admissible field class on which these constructions exist, retaining the question for the full stated class.
Weak associativity and injectivity of the single-state field map do not alone construct these geometric objects or inverse comparisons.
\end{problem}

\section{The geometry supplying each operation}
A pole has more data than its residue.
For a normal coordinate $z$, the successive coefficients of
$\sum_{j\geq1}a_jz^{-j}$ are its principal parts of successive orders.
Keeping all these coefficients preserves every normal pole symbol.
It does not specify a transverse ordering.

Ordered disjoint intervals in a line have contractible configuration
chambers, one for each order. Their composition supplies the coherent
associative operation when the theory is locally constant in that
direction. Moving points in a plane also supplies braid paths.
An exchange operation requires an assignment to those paths and
its compatibility with fusion.
Principal parts, interval composition and braided exchange therefore
come from distinct supplied geometric data.

\chapter{Transversely ordered chiral algebras}
\label{cbc:ord:chapter}
Fuse two operators, then separate their chiral supports.
Alternatively, separate first and fuse in each cluster.
Transverse locality requires these operations to agree, including
their associativity homotopies.
We first express this condition on the coefficient family.
We then construct its interval origin and specify the comparison
needed to obtain algebraic coefficients from a field theory.

\section{Locality on the coefficient family}
By Definition~\ref{cbc:two:transverse}, an ordered chiral algebra is
an augmented $\Eone$-algebra in $(\Fact_D(X),\tensor^!)$.
Write $A_I$ for its coefficient on $X^I$.

An operation of arity $r$ is parametrized by the little-interval
space $\Eone(r)$. Its component will be denoted by
\[
 \mu_{r,I}(\alpha):A_I^{\tensor^!r}\longrightarrow A_I,
 \qquad \alpha\in\Eone(r).
\]
The maps vary coherently with $\alpha$ and with operadic substitution.
At arity zero, the source is the pointwise unit $\omega_{X^I}$.
Write $u_I:\omega_{X^I}\to A_I$ for this unit and
$\epsilon_I:A_I\to\omega_{X^I}$ for the augmentation.

Let $d=(I_1,\ldots,I_k)$ be an ordered decomposition of $I$.
The open inclusion $j_d:U_d\hookrightarrow X^I$ separates its parts.
The factorization structure supplies
\[
 \delta_d:j_d^!A_I\xrightarrow{\sim}
 j_d^!\left(\mathop{\boxtimes}_{a=1}^k A_{I_a}\right).
\]
For $r=2$ and $k=2$, compatibility of the binary multiplication
$m_I=\mu_{2,I}(\alpha)$ with separation is
\begin{equation}
 \delta_d\circ j_d^!m_I
 =j_d^!(m_{I_1}\boxtimes m_{I_2})\circ\sigma_d
       \circ(\delta_d\tensor^!\delta_d).
 \label{eq:cbc:ord:locality}
\end{equation}
Here the tensor comparisons for $j_d^!$ are understood.
The map $\sigma_d$ groups the first cluster's two factors together,
then the second cluster's factors.

For homogeneous coefficients, this regrouping reads
\[
 (a\tensor b)\tensor(c\tensor d)
 \longmapsto
 (-1)^{|b||c|}(a\tensor c)\tensor(b\tensor d).
\]
Thus the product on two separated clusters is
$(-1)^{|b||c|}ac\tensor bd$.
If $b$ and $c$ are odd, omitting the middle crossing reverses its sign.
The sign records the actual permutation of coefficients.

For general $r,k$, let $\sigma_{d,r}$ transpose the array of
$r$ inputs and $k$ clusters, with its Koszul signs.
The full locality equation is
\begin{equation}
 \delta_d\circ j_d^!\mu_{r,I}(\alpha)
 =j_d^!\left(\mathop{\boxtimes}_{a=1}^k
          \mu_{r,I_a}(\alpha)\right)
       \circ\sigma_{d,r}\circ\delta_d^{\tensor^!r}.
 \label{eq:cbc:ord:all-locality}
\end{equation}
The same parameter $\alpha$ acts in every cluster.
For $r=0$, this is the unit equation.
There is also the corresponding equation for $\epsilon_I$.
These equations include homotopies compatible with every substitution,
collision, permutation and refinement of $d$.

\begin{theorem}[Coherent locality recognition]
\label{cbc:ord:locality}
An augmented ordered chiral algebra is equivalent to a factorization
family $A$ with coherent augmented $\Eone$-algebra structures on its
coefficients.
The collision equivalences and every $\delta_d$ must be equivalences
of augmented $\Eone$-algebras, coherently under composition and refinement.
On the separated locus, this condition is precisely
\eqref{eq:cbc:ord:all-locality} with its unit, augmentation and
higher homotopies.
\end{theorem}
\begin{proof}
Evaluation on each finite power preserves the pointwise tensor and
its unit. It therefore carries an internal $\Eone$-algebra to the
stated component algebras.
Its operations are morphisms of factorization coalgebras, so they
intertwine every collision and separation map.
The tensor product of the component algebras uses $\sigma_{d,r}$,
which gives~\eqref{eq:cbc:ord:all-locality}.

Conversely, fix an arity and a parameter in $\Eone(r)$.
The component operations assemble to a Ran morphism because they
respect all collision equivalences.
Equation~\eqref{eq:cbc:ord:all-locality} makes this morphism a map
of factorization coalgebras.
The prescribed coherent compatibilities perform the same assembly
for the entire parameter space and every operadic substitution.
They therefore give an $\Eone$-action on $A$ in $\Fact_D(X)$.
The arity-zero operations and the maps $\epsilon_I$ assemble to its
unit and augmentation.

A morphism between either kind of data is the same compatible family
of component algebra maps, with its homotopies.
Both constructions preserve these mapping spaces and undo each other.
Hence they give an equivalence of the stated categories.
\end{proof}

In a strict differential graded model, strictly associative binary
products and strict unit equations determine the iterated products.
If all transition and factorization maps preserve them, induction
proves every higher locality equation.
For a general $\Eone$-algebra, the higher homotopies are additional
data. The binary square alone does not construct them.
The marked family of Proposition~\ref{cbc:two:factorization} illustrates
\eqref{eq:cbc:ord:all-locality}.
Only one separated cluster can contain its marked ideal, and the
same matrix product acts in that cluster before or after separation.

\section{Fusion along an interval}
Fix a symmetric monoidal coefficient $\infty$-category $\mathcal V$.
Let $\mathcal H$ be a specified symmetric monoidal $\infty$-category
of holomorphic factorization objects on $X^{\mathrm{an}}$.
Write $\odot$ for its tensor product and $\mathbf1_{\mathcal H}$
for its unit.
Assume it is a full subcategory of $\mathcal V$-valued prefactorization objects with
the pointwise tensor product.
Thus $(B\odot C)(U)=B(U)\tensor C(U)$ on each chiral open $U$.
The chosen holomorphic structure and this closure are hypotheses
on $\mathcal H$.

Let $\F$ be a unital $\mathcal V$-valued factorization algebra on
$X^{\mathrm{an}}\times\R$.
For an open interval $J$, define its interval coefficient by
\[
 B_J(U)=\F(U\times J).
\]
Require $B_J\in\mathcal H$ and require every inclusion of intervals
to induce an equivalence in $\mathcal H$.
For disjoint intervals $J_1<\cdots<J_r$ inside $J$, factorization gives
\begin{equation}
 B_{J_1}\odot\cdots\odot B_{J_r}\longrightarrow B_J.
 \label{eq:cbc:ord:interval}
\end{equation}
The empty input supplies the unit map.
A compatible augmentation of this interval family is a further
map to $\mathbf1_{\mathcal H}$ that preserves every operation.

The rectangular maps in~\eqref{eq:cbc:ord:interval} respect chiral
factorization. To see this, choose disjoint chiral opens
$U_1,\ldots,U_k\subset U$.
The rectangles $U_a\times J_i$ form one disjoint family.
Fusing first across the chiral opens and then across intervals,
or in the other order, gives its same factorization map.
The coefficient permutation is exactly $\sigma_{d,r}$.
The field theory's composition homotopies give all refinements of
this identity, so these are morphisms in $\mathcal H$.

The ordered configurations of $r$ affine subintervals of $(0,1)$
form the convex region
\[
 0<a_1<b_1<a_2<b_2<\cdots<a_r<b_r<1.
\]
Each other component is obtained by permuting the input labels.
The contractibility of these components explains associativity:
different nested fusions lie in the same ordered component.
Their coherent paths, together with substitution, give the
little-interval operations. They do not exchange the order of labels.

\begin{proposition}[The transverse interval action]
\label{cbc:ord:interval}
Under the preceding hypotheses, restriction to any fixed interval
$J_0$ gives $B_{J_0}$ a coherent unital $\Eone$-algebra structure
in $\mathcal H$.
A compatible augmentation of the interval family makes this an
augmented algebra.
The construction is natural in the field theory and in its
compatible maps.
\end{proposition}
\begin{proof}
Consider the colored operad whose colors are open intervals in $\R$.
Its operations are inclusions of disjoint intervals into one interval.
The rectangular maps define an algebra over this operad in
$\mathcal H$, by their compatibility with chiral factorization.
Its unary interval inclusions act by equivalences.

The operadic localization theorem
\cite[Theorem 5.4.5.9]{LurieHA} therefore identifies it with an
$\Eone$-algebra, using
\cite[Example 5.4.5.3]{LurieHA} for $M=\R$.
The underlying object is its value on $J_0$.
This theorem retains all substitution homotopies. Choosing separate
inverses to the unary maps would not by itself specify them.
It also identifies the spaces of compatible maps.
Consequently it transports an augmentation to the tensor unit
and proves the asserted naturality.
\end{proof}

In degree-zero vector spaces, the construction reduces to ordered
interval multiplication in \cite[2014 preprint, \S3.2]{CG1}.
The unit comes from the empty input.
An augmentation instead maps the resulting algebra to its scalar
coefficient object and must be supplied separately.

\section{Algebraization and product reconstruction}
The interval action lives in $\mathcal H$.
Fix a full subcategory $\mathcal C_D\subseteq\Fact_D(X)$ that
contains $\omega_{\Ran}$ and is closed under $\tensor^!$.
It specifies the admitted algebraic coefficients.

To obtain such coefficients from the field theory, suppose there
is a strong symmetric monoidal functor
\begin{equation}
 \Phi:(\mathcal H,\odot,\mathbf1_{\mathcal H})
       \longrightarrow(\mathcal C_D,\tensor^!,\omega_{\Ran}).
 \label{eq:cbc:ord:comparison}
\end{equation}
Strong monoidality includes coherent equivalences
\[
 \phi_r:\Phi(B)^{\tensor^!r}\xrightarrow{\sim}\Phi(B^{\odot r}),
 \qquad
 \phi_0:\omega_{\Ran}\xrightarrow{\sim}\Phi(\mathbf1_{\mathcal H}).
\]
Landing in $\Fact_D(X)$ requires actual collision and separated
factorization maps on every finite power of $X$.
The functor must preserve their coherent compositions.

\begin{theorem}[Holomorphic--topological recognition under coefficient comparison]
\label{cbc:ord:recognition}
Let $\F$ satisfy the hypotheses of
Proposition~\ref{cbc:ord:interval}, with a compatible augmentation.
Given the comparison~\eqref{eq:cbc:ord:comparison}, the coefficient
$A=\Phi(B_{J_0})$ is an augmented ordered chiral algebra.
Its multiplication is transverse interval fusion transported by $\Phi$.
\end{theorem}
\begin{proof}
Write $\mu_r^B$ for the interval operations on $B=B_{J_0}$.
The transported operations, unit and augmentation are
\[
 \mu_r^A=\Phi(\mu_r^B)\circ\phi_r,\qquad
 u^A=\Phi(u^B)\circ\phi_0,\qquad
 \epsilon^A=\phi_0^{-1}\circ\Phi(\epsilon^B).
\]
The monoidal coherence of $\Phi$ carries every substitution and
unit diagram for $B$ to the corresponding diagram for $A$.
The same holds for the augmentation.
All maps lie in $\Fact_D(X)$, so the locality equations and their
higher homotopies follow from Theorem~\ref{cbc:ord:locality}.
\end{proof}

A regular Riemann--Hilbert correspondence can enter the construction
of $\Phi$ after the chosen analytic coefficients have been identified
with the appropriate constructible complexes.
One must then verify the tensor normalization, units, diagonal
restrictions, external products and separated direct images.
An algebraic target also requires an algebraic realization and the
relevant regularity conditions.
Holomorphic dependence of a field theory alone supplies none of
these comparison maps.
The distinction between analytic factorization data and algebraic
Ran coefficients is already explicit in \cite[2014 preprint, \S3.3.1]{CG1}.

The converse recognition question concerns reconstruction on the
whole product space.
Let $\mathcal R(W)$ be the category of finite disjoint unions of
rectangles $U_a\times J_a$ contained in an open
$W\subset X^{\mathrm{an}}\times\R$.
Here each $U_a$ is an analytic disc, each $J_a$ an open interval,
and morphisms are inclusions of unions. The empty union is allowed.

Assume the following colimits exist in $\mathcal V$.
Call a theory \emph{product-generated} when its canonical maps
\begin{equation}
 \mathop{\operatorname{colim}}_{\bigsqcup_a(U_a\times J_a)\in\mathcal R(W)}
       \bigotimes_a\F(U_a\times J_a)
       \longrightarrow\F(W)
 \label{eq:cbc:ord:product-generation}
\end{equation}
are equivalences for every $W$.
The colimit is derived.

Starting from an internal $\Eone$-algebra, an inverse must first
construct the diagram on the left of
\eqref{eq:cbc:ord:product-generation}.
Its maps must handle every rectangular inclusion, including families
whose chiral projections overlap.
One must prove their coherent compatibility and factorization descent,
then identify both composites with the original objects.

Write $\Fact_{\mathrm{hol,lc}}^{\mathrm{prod,aug}}$ for the category
of admitted augmented product-generated theories with interval coefficients
in $\mathcal H$.
If this product reconstruction is inverse to interval restriction,
and $\Phi$ is a monoidal equivalence,
the corresponding categories are equivalent:
\[
 \Fact_{\mathrm{hol,lc}}^{\mathrm{prod,aug}}
       (X^{\mathrm{an}}\times\R)
 \simeq
 \Alg_{\Eone}^{\mathrm{aug}}(\mathcal C_D,\tensor^!).
\]
Without these comparisons, the proved construction is the forward
interval action and its transport through~\eqref{eq:cbc:ord:comparison}.
The general inverse remains a separate reconstruction problem.

\chapter{Block-ordered operations}
Rectangular locality compares two compositions: row multiplication followed by column factorization, and factorization followed by multiplication. An operadic presentation must encode these comparisons together with their higher homotopies. A forest calculation has a further requirement: every edge contraction must act on the actual collision coefficients. We describe the operadic comparison first and then calculate its possible contraction differential.

\section{Chiral operations and ordered profiles}
An \emph{ordered profile} is a list $(I_1,\ldots,I_r)$ of disjoint nonempty finite label sets. Its coefficient space is
\[
 X^{I_1}\times\cdots\times X^{I_r}.
\]
The list records the transverse order of the blocks. Diagonals inside a factor record collisions of its labels. These are different data: a profile specifies neither the coefficient maps for composing adjacent blocks nor the maps on the collision strata.

The chiral operations on a smooth curve already have a precise pseudotensor description. For a nonempty finite set $I$, let $j_I:U_I\hookrightarrow X^I$ be the inclusion of the locus of pairwise distinct points, and let $\Delta_I:X\hookrightarrow X^I$ be the small diagonal. For right $D_X$-modules $M_i$ and $N$, the space of chiral operations is
\begin{equation}\label{cbc:ord:pseudotensor}
 P_I^{\mathrm{ch}}(\{M_i\},N)
 =\operatorname{Hom}_{D_{X^I}}
   \left(j_{I*}j_I^*\boxtimes_{i\in I}M_i,\Delta_{I*}N\right).
\end{equation}

Here the operations are those of the abelian category of right $D$-modules. Substitution along $J\twoheadrightarrow I$ first applies the operations on its fibers, on the open locus where different fibers are separated. The $I$-operation then composes their outputs. Diagonal base change identifies the intermediate targets, and composition of these maps is associative.

This is the chiral pseudotensor structure of Beilinson--Drinfeld~\cite[\S\S3.1.1--3.1.2]{BD04}. Its nonempty operation spaces do not by themselves specify a nullary operation.

An operad acting in this pseudotensor category specifies additional algebraic identities. A presentation of the entire derived factorization category must also account for its collision transitions, separated-factorization equivalences and mapping spaces. In particular, the existence of~\eqref{cbc:ord:pseudotensor} is distinct from a symmetric monoidal equivalence between that factorization category and algebras over a chosen infinity-operad.

\section{The operadic exponential law}
Let $P$ and $Q$ be small infinity-operads, and let $\Vcal$ be a symmetric monoidal infinity-category. The category $\Alg_Q(\Vcal)$ has the \emph{pointwise tensor}. It tensors underlying objects and applies each operation to both factors, using the symmetry to group the inputs. The tensor unit has its canonical $Q$-algebra structure. This is the symmetric monoidal structure of~\cite[Construction~3.2.4.1, Proposition~3.2.4.3 and Example~3.2.4.4]{LurieHA}.

An operadic bifunctor is a functor
\[
 P^\otimes\times Q^\otimes\longrightarrow\Vcal^\otimes
\]
over the smash product of pointed finite sets, preserving inert morphisms in both variables. Its two families of operations satisfy interchange, including all coherent substitution relations. The Boardman--Vogt tensor product $P\otimes_{\mathrm{BV}}Q$ represents these bifunctors.

\begin{theorem}[Operadic exponential law]\label{cbc:ord:operadic-exp}
Let $P$ and $Q$ be small infinity-operads and $\Vcal$ a symmetric monoidal infinity-category. Equip $\Alg_Q(\Vcal)$ with its pointwise tensor. There is a natural equivalence
\[
 \Alg_{P\otimes_{\mathrm{BV}}Q}(\Vcal)
 \simeq
 \Alg_P\bigl(\Alg_Q(\Vcal)\bigr).
\]
The equivalence identifies operation maps, their coherent substitutions and the full mapping spaces of algebras.
\end{theorem}
\begin{proof}
By the defining universal property of the operadic tensor product, the left side is the infinity-category of the displayed bifunctors~\cite[Definition~2.2.5.3]{LurieHA}. Fixing the $P$-variable gives a $Q$-algebra. An operation in $P$ becomes a map from the pointwise tensor of these $Q$-algebras to another such algebra. Compatibility with the $Q$-operations is exactly the bifunctor's interchange. Currying therefore gives a $P$-algebra in $\Alg_Q(\Vcal)$.

Conversely, evaluation of a $P$-algebra in this pointwise algebra category gives the bifunctor. The algebra-category construction in~\cite[Construction~3.2.4.1]{LurieHA} characterizes these two operations on every simplicial parameter space. Currying and evaluation are inverse there and preserve inert maps by Proposition~3.2.4.3. They therefore identify transformations and all their higher homotopies as well as objects, proving the asserted equivalence of infinity-categories.
\end{proof}

To apply this law to the preceding chapter, one needs a specified symmetric monoidal equivalence
\begin{equation}\label{cbc:ord:operadic-comparison}
 \Phi:(\Fact_D(X),\tensor^!)
 \xrightarrow{\ \sim\ }
 \bigl(\Alg_{Q_X}(\Vcal),\otimes_{\mathrm{pt}}\bigr).
\end{equation}
Here $Q_X$ is the chosen infinity-operad, and $\otimes_{\mathrm{pt}}$ denotes the pointwise tensor just defined. The equivalence must identify the tensor units and the specified augmentation targets. Applying $\Phi$ to every ordered operation and then using Theorem~\ref{cbc:ord:operadic-exp} gives
\begin{equation}\label{cbc:ord:operadic-ordered}
 \Alg_{\Eone}^{\mathrm{aug}}(\Fact_D(X),\tensor^!)
 \simeq
 \Alg_{\Eone\otimes_{\mathrm{BV}}Q_X}^{\mathrm{aug}}(\Vcal).
\end{equation}

On the right, augmentation means a map to the algebra corresponding to the tensor unit on the left. Applying a monoidal inverse to $\Phi$ gives the inverse equivalence. Thus~\eqref{cbc:ord:operadic-ordered} is conditional on the full comparison~\eqref{cbc:ord:operadic-comparison}, not merely on the agreement of its objects or its binary operation spaces. Under that hypothesis one may write $\mathsf{OrdCh}_X=\Eone\otimes_{\mathrm{BV}}Q_X$. The intrinsic ordered algebra does not require this presentation.

\section{The sign of a mixed contraction}
Consider a forest $F$ with one transverse edge $e$ and one chiral edge $f$, oriented by $e\wedge f$. Write $F/e$ for the forest obtained by contracting $e$. The degree-zero map $c_e^F$ takes coefficients of $F$ to those of $F/e$. Use the same notation for the other contractions. For a coefficient $x$ of degree zero, the contraction boundary is
\[
 \delta(x\otimes e\wedge f)
 =c_e^Fx\otimes f-c_f^Fx\otimes e.
\]

The second boundary is
\begin{equation}\label{cbc:ord:two-edge}
 \delta^2(x\otimes e\wedge f)
 =\bigl(c_f^{F/e}c_e^F-c_e^{F/f}c_f^F\bigr)x.
\end{equation}
The square vanishes when its two coefficient composites agree. Its sign comes from the two edge orientations. Its commutativity is a separate property of the coefficients.

We now give the general construction with the grading included. Work over a field $k$ of characteristic zero and use cohomological differentials of degree one. Let $\mathcal F$ be a collection of forests closed under the chosen edge contractions. For each $F\in\mathcal F$, its finite internal-edge set is a disjoint union
\[
 E(F)=E_\perp(F)\sqcup E_{\mathrm{ch}}(F).
\]

The two names distinguish directions of contraction, rather than the object colors of a colored operad. Contracting $e$ gives $F/e$ with the remaining named edges and their directions unchanged. Require the two targets $F/e/f$ and $F/f/e$ to be identified for every pair of distinct edges.

\begin{definition}[Coefficient contraction system]\label{cbc:ord:contraction-system}
A coefficient contraction system assigns a cochain complex $C(F)$ to each forest. Write $d_C$ for its differential. Each edge contraction has a degree-zero cochain map
\[
 c_e^F:C(F)\longrightarrow C(F/e).
\]
For distinct edges these maps satisfy, under the specified target identifications,
\begin{equation}\label{cbc:ord:contraction-square}
 c_f^{F/e}c_e^F=c_e^{F/f}c_f^F
 \qquad(e\ne f).
\end{equation}
The identifications are required to agree for longer contraction sequences.
\end{definition}

Put the orientation line
\[
 \operatorname{or}(F)
 =\det(kE_\perp(F))\otimes\det(kE_{\mathrm{ch}}(F))
\]
in degree $-|E(F)|$. Its product orientation places all transverse edges before all chiral edges. Form the total graded space
\begin{equation}\label{cbc:ord:forest-total}
 K^n=\bigoplus_{F\in\mathcal F}
 C^{n+|E(F)|}(F)\otimes\operatorname{or}(F).
\end{equation}
The direct sum uses finite support, with no implicit completion. For an ordered wedge $e_1\wedge\cdots\wedge e_N$, define exterior contraction by
\[
 \iota_{e_i}(e_1\wedge\cdots\wedge e_N)
 =(-1)^{i-1}e_1\wedge\cdots\widehat e_i\cdots\wedge e_N.
\]
This determines a map of determinant lines independently of the displayed basis order.

For homogeneous $x\in C^q(F)$ and $\omega\in\operatorname{or}(F)$, define
\begin{align}
 d_{\mathrm{int}}(x\otimes\omega)&=d_Cx\otimes\omega,
 \label{cbc:ord:forest-internal}\\
 d_\perp(x\otimes\omega)
 &=(-1)^q\sum_{e\in E_\perp(F)}c_e^Fx\otimes\iota_e\omega,
 \label{cbc:ord:forest-transverse}\\
 d_{\mathrm{ch}}(x\otimes\omega)
 &=(-1)^q\sum_{f\in E_{\mathrm{ch}}(F)}c_f^Fx\otimes\iota_f\omega.
 \label{cbc:ord:forest-chiral}
\end{align}
All three maps have total degree one. In the product orientation the chiral-edge sign includes $(-1)^{|E_\perp(F)|}$, because its contraction passes the transverse determinant factor.

\begin{theorem}[Two-color forest differential]\label{cbc:ord:forest-d}
A coefficient contraction system gives a differential
\[
 d=d_{\mathrm{int}}+d_\perp+d_{\mathrm{ch}}
\]
on~\eqref{cbc:ord:forest-total}. Its summands satisfy
\[
 d_\perp^2=d_{\mathrm{ch}}^2=0,
 \qquad d_\perp d_{\mathrm{ch}}+d_{\mathrm{ch}}d_\perp=0,
 \qquad d_{\mathrm{int}}d_\alpha+d_\alpha d_{\mathrm{int}}=0
 \quad(\alpha=\perp,\mathrm{ch}).
\]
\end{theorem}
\begin{proof}
For distinct edges, exterior contractions anticommute:
\[
 \iota_f\iota_e=-\iota_e\iota_f.
\]
The two coefficient composites agree by~\eqref{cbc:ord:contraction-square}. Contracting a fixed pair in the two orders therefore gives opposite terms. Pairing all such terms proves the same-direction squares and the mixed anticommutator. Contracting an edge twice is impossible.

For the internal differential and contraction of $e$, the two coefficient terms are $d_Cc_e x$ and $c_ed_Cx$. They agree because $c_e$ is a cochain map. Their totalization signs are $(-1)^q$ and $(-1)^{q+1}$, so they cancel. Finally $d_{\mathrm{int}}^2=0$ follows from the coefficient differential. Each forest has finitely many edges, so these calculations define maps on the direct sum and prove $d^2=0$ there.
\end{proof}

For the constant coefficient system $C(F)=k$, take $c_e=c_f=\id$. Formula~\eqref{cbc:ord:two-edge} is $1-1=0$. Giving the chiral contraction its sign within $E_{\mathrm{ch}}$ alone would instead give $1+1=2$. Conversely, with the correct orientation but noncommuting coefficient maps,~\eqref{cbc:ord:two-edge} is their nonzero commutator. Thus the orientation calculation neither supplies the coefficient square nor replaces its required homotopy in a coherent model.

\section{Resolution of an operadic model}
A Boardman--Vogt resolution carries more data than the contraction complex above. Let $\mathcal M$ be a cofibrantly generated symmetric monoidal model category with cofibrant tensor unit $\mathbbm1_{\mathcal M}$. Fix a set of object colors $\mathcal C$.

An \emph{interval object} $H$ has endpoints $0,1:\mathbbm1_{\mathcal M}\to H$, a weak equivalence $\epsilon:H\to\mathbbm1_{\mathcal M}$, and an associative operation $\vee:H\otimes H\to H$. The element $0$ is a unit for $\vee$, and $1$ is absorbing. The map $\epsilon$ preserves this operation, and both endpoint composites with $\epsilon$ are the identity. The endpoint map $(0,1):\mathbbm1_{\mathcal M}\sqcup\mathbbm1_{\mathcal M}\to H$ is required to be a cofibration. The usual interval $[0,1]$ with maximum is the topological example.

Let $R$ be a $\mathcal C$-colored operad in $\mathcal M$. It is \emph{well pointed} when each operadic unit is a cofibration. It is \emph{$\Sigma$-cofibrant} when its underlying collection is cofibrant with its input-permutation actions.

The construction $W(H,R)$ uses trees with $R$-operations at vertices and lengths in $H$ on internal edges. It identifies a zero-length edge with its operadic contraction. Removing a vertex labeled by an operadic unit combines two adjacent internal lengths by $\vee$, in their order along the tree. If one adjacent edge is external, it removes the other length by $\epsilon$. Tree symmetries give further identifications, and grafting introduces edges of length~$1$.

The full interval object and these identifications are part of the construction. In a cellular interval model, the boundary of an interval edge has both endpoint faces.

For well-pointed $\Sigma$-cofibrant $R$, the operadic augmentation
\begin{equation}\label{cbc:ord:w-resolution}
 \epsilon_R:W(H,R)\longrightarrow R
\end{equation}
is a weak equivalence from a cofibrant operad. This is Berger--Moerdijk's resolution theorem~\cite[Theorem~3.5]{BergerMoerdijk2007}. Its uncolored construction and proof are given in~\cite[\S\S4--5]{BergerMoerdijk2006}. Here weak equivalence is tested on every operation object. Cofibrancy means the lifting property against maps that are trivial fibrations on those objects. The resolution theorem does not require a transferred model structure on all colored operads.

An operad is \emph{admissible} if its algebras have the transferred model structure. Its weak equivalences and fibrations are detected on the underlying objects. Suppose in addition that $\mathcal M$ is left proper and that both operads in~\eqref{cbc:ord:w-resolution} are admissible, well pointed and $\Sigma$-cofibrant. The comparison theorem of~\cite[Theorem~4.1]{BergerMoerdijk2007} then makes extension and restriction along $\epsilon_R$ a Quillen equivalence. After inverting weak equivalences, it identifies algebras over the resolution with algebras over $R$.

To model ordered chiral algebras by this route, one must supply~\eqref{cbc:ord:operadic-comparison}. A second comparison must identify $R$-algebras with $(\Eone\otimes_{\mathrm{BV}}Q_X)$-algebras and preserve the specified augmentation. The derived change-of-operad functor, that comparison, and~\eqref{cbc:ord:operadic-ordered} then compose to the required equivalence. A concrete forest description must further identify its collision coefficients, substitution maps, interval faces and unit relations with this resolution. Theorem~\ref{cbc:ord:forest-d} proves the contraction identity once its coefficient maps are given. Constructing those maps and the comparison is the geometric problem that remains.


\part{Bar, cobar, and the boundary state}


\chapter{Three bar constructions}\label{cbc:bars:chapter}
A bar construction records a multiplication through its strings of composable inputs. Changing the multiplication changes the bar, even when the underlying coefficient object stays the same. Transverse composition uses the pointwise tensor, associative chirality uses chiral convolution, and a chiral Lie bracket gives Chevalley chains. We specify these three constructions and then construct the comparison that exists on associative envelopes.

\section{The products and their units}
Recall the positive Ran coefficient category $\Dcal_X$ and its extension $\Dcal_X^+=\operatorname{Vect}_{\C}\times\Dcal_X$ by an empty-configuration coefficient. The pointwise tensor on $\Dcal_X$ has unit
\[
 \mathbf1_\perp=\omega_{\Ran}.
\]
The unital chiral convolution $\starCh=\tensor^{\ch}$ on $\Dcal_X^+$ has the different unit
\[
 K=\mathbf1_{\ch}=(\C,0).
\]
The positive chiral product on $\Dcal_X$ has no unit. Its $r$-fold product vanishes on the chart $X^I$ when $r>|I|$, because $I$ has no decomposition into $r$ nonempty blocks.

Let $A^\perp$ be an associative algebra for $\tensor^!$, with augmentation to $\mathbf1_\perp$. An ordered chiral algebra supplies such an object after forgetting its factorization structure. Let $A^{\ch}$ be an associative algebra for $\starCh$, with augmentation to $K$. Finally, let $L$ be a Lie object for the positive chiral product. These are three specified inputs, with no identification between their operations assumed.

For an augmented associative algebra $A\to\mathbf1$ in a monoidal category $(\mathcal C,\odot,\mathbf1)$, its two-sided simplicial bar has terms
\[
 B_n(\mathbf1,A,\mathbf1)=\mathbf1\odot A^{\odot n}\odot\mathbf1.
\]
Its interior faces multiply adjacent entries. Its two end faces apply the augmentation, and its degeneracies insert the algebra unit. The coherent associative and augmentation identities make these maps a simplicial object. When its realization exists, write
\[
 B_\odot A=|B_\bullet(\mathbf1,A,\mathbf1)|
          =\mathbf1\odot_A^{\mathbb L}\mathbf1.
\]

Use cohomological grading and let $sV=V[1]$, so $|sa|=|a|-1$. Choose a differential graded presentation in which the tensors compute the indicated derived tensors and direct-sum normalization computes realization. Write $d_L$ for the differential of the chosen presentation of $L$.

For a positive object $V$, $\Sym_{\ch}^r(V)$ means the coinvariants of $V^{\starCh r}$ under the signed permutation action. Characteristic zero makes these finite-group coinvariants exact. Write $\Sym_{\ch}(V)=K\oplus\bigoplus_{r\geq1}\Sym_{\ch}^r(V)$.

On $\Sym_{\ch}(sL)$, the coproduct is the unshuffle. Choose a subset of inputs for the first factor and place its complement in the second, with the suspended Koszul sign. Let $Q$ be the coderivation with $Q|_K=0$ and components to the cogenerators
\begin{equation}\label{cbc:bars:chevalley-components}
 Q_1(sx)=-s(d_Lx),\qquad
 Q_2(sx\,sy)=(-1)^{|x|}s[x,y].
\end{equation}
All components with more than two inputs are zero. The second formula is symmetric on the suspended inputs by graded antisymmetry of the bracket.

\begin{definition}[Transverse, associative-chiral, and Lie-chiral bars]
\label{cbc:bars:three}
The two associative bars are
\begin{align*}
 \BarPerp(A^\perp)
 &=\mathbf1_\perp\tensor_{A^\perp}^{\mathbb L,!}\mathbf1_\perp,\\
 \BarAssCh(A^{\ch})
 &=K\starCh_{A^{\ch}}^{\mathbb L}K.
\end{align*}
Their ambient coefficient categories are respectively $\Dcal_X$ with $\tensor^!$ and $\Dcal_X^+$ with $\starCh$. The Lie-chiral bar is the noncounital chiral Chevalley coalgebra
\[
 \BarLieCh(L)=\operatorname{CE}_{\ch}(L)
 =\left(\bigoplus_{r\geq1}\Sym_{\ch}^r(L[1]),Q\right).
\]
For comparison with a counital associative bar, put
\[
 \operatorname{CE}_{\ch}^+(L)=K\oplus\operatorname{CE}_{\ch}(L).
\]
\end{definition}

All these sums are direct sums. On each finite configuration the positive chiral sums have only finitely many nonzero terms. Pointwise tensor powers have no corresponding configuration bound. Any product completion of a transverse bar therefore requires its own convergence and continuity hypotheses.

\section{Normalization and the differential}
The augmentation is split by the unit. Write $d_A$ for the coefficient differential and $\bar A$ for the augmentation kernel in this presentation. Write
\[
 [a_1|\cdots|a_n]=sa_1\odot\cdots\odot sa_n.
\]
The empty word denotes $\mathbf1$.

On the tensor coalgebra
\[
 T^c_\odot(s\bar A)=\mathbf1\oplus\bigoplus_{n\geq1}(s\bar A)^{\odot n},
\]
the coproduct cuts a word at every position. The differential is the coderivation whose components to words of length one are
\begin{equation}\label{cbc:bars:components}
 b_1(sa)=-s(d_Aa),\qquad
 b_2(sa\odot sb)=(-1)^{|a|}s(ab).
\end{equation}
Inserting either component after a prefix gives the Koszul sign of that prefix of suspended factors. Thus multiplication in slots $i,i+1$ has sign
\[
 (-1)^{\sum_{j<i}(|a_j|-1)+|a_i|}.
\]

\begin{theorem}[Typed bar theorem]\label{cbc:bars:typed}
Under the stated differential graded presentation hypotheses, the two associative bars are represented in their respective ambient coefficient categories by
\[
 \BarPerp(A^\perp)\simeq
 \bigl(T^c_{\tensor^!}(s\bar A^\perp),b_1+b_2\bigr),
 \qquad
 \BarAssCh(A^{\ch})\simeq
 \bigl(T^c_{\starCh}(s\bar A^{\ch}),b_1+b_2\bigr).
\]
They are coaugmented counital coassociative coalgebras with deconcatenation. The Lie-chiral bar is the noncounital cocommutative coalgebra with differential~\eqref{cbc:bars:chevalley-components}. Each differential squares to zero.
\end{theorem}
\begin{proof}
In the normalized associative bar, a word containing a unit is degenerate. Splitting the augmentation therefore leaves exactly $\bar A^{\odot n}$ in simplicial degree $n$. Totalization suspends every factor. With the convention~\eqref{cbc:bars:components}, the two end faces vanish on these reduced words and the interior faces give $b_2$.

The square of $b_1$ is induced by $d_A^2$. The relation $b_1b_2+b_2b_1=0$ is the graded Leibniz identity. In $b_2^2$, contractions of disjoint adjacent pairs cancel with their Koszul signs. The remaining terms contract three consecutive factors in the two orders, and cancel by associativity. For degree-zero inputs this last identity is already visible in
\[
 b_2[a|b|c]=[ab|c]-[a|bc],\qquad
 b_2^2[a|b|c]=[(ab)c]-[a(bc)]=0.
\]

Deconcatenation is coassociative because successive cuts enumerate the same three subwords. A contraction and a surviving cut act on one side of the cut, so the differential is a coderivation. The empty word supplies coaugmentation and counit.

For the Chevalley coalgebra, the unary, binary and ternary components of $Q^2$ are respectively $d_L^2$, the derivation identity and graded Jacobi. They vanish. All longer components are determined by the coderivation identity. This proves the Lie assertion with its unshuffle coproduct.
\end{proof}

Normalization here takes place in the stable ambient coefficients. Its individual suspended summands need not be factorization objects. The realization criterion below addresses that additional structure on the whole bar.

The positive Chevalley construction is the functor in Francis--Gaitsgory's Lie--commutative chiral duality~\cite[\S4.3 and Theorem~5.1.1]{FG12}. It factorizes exactly when $L$ is supported on the main diagonal $X\to\Ran(X)$~\cite[Theorem~5.2.1]{FG12}. The equivalence thus identifies a diagonally supported chiral Lie object with its factorization coalgebra. An additional associative multiplication or a second associative algebra is separate data.

\section{The bar of an associative envelope}
Let $U_{\ch}^{\mathsf{Ass}}(L)$ denote the derived associative envelope of $L$ in $(\Dcal_X^+,\starCh,K)$. Among augmented algebras with positive augmentation ideal, it is left adjoint to taking that ideal's commutator Lie object. Its augmentation kills $L$, and its augmentation ideal has zero empty-configuration component. In a cofibrant strict presentation, its defining relation is
\[
 xy-(-1)^{|x||y|}yx=[x,y].
\]
This is an associative envelope for the chiral tensor. The chiral envelope from star-Lie objects to chiral Lie objects in~\cite[\S6]{FG12} is a different functor.

Define a map $\kappa:\operatorname{CE}_{\ch}^+(L)\to\BarAssCh(U_{\ch}^{\mathsf{Ass}}(L))$ by the suspended permutation sum
\begin{equation}\label{cbc:bars:pbw-map}
 \kappa(sx_1\cdots sx_r)
 =\sum_{\sigma\in S_r}
   \varepsilon(\sigma;sx_1,\ldots,sx_r)
   [x_{\sigma(1)}|\cdots|x_{\sigma(r)}],
\end{equation}
where $\varepsilon$ is the Koszul sign of the displayed permutation. Set $\kappa|_K=\id_K$. The sum has no factorial denominator. For degree-zero $x,y$, its first nontrivial component and its reduced coproduct are
\begin{align*}
 \kappa(sx\,sy)&=[x|y]-[y|x],\\
 \overline\Delta\kappa(sx\,sy)&=[x]\starCh[y]-[y]\starCh[x].
\end{align*}
The same expression is obtained by applying $\kappa\starCh\kappa$ to the unshuffle coproduct. Dividing only the length-two component by $2$ halves $\overline\Delta\kappa$ but leaves $(\kappa\starCh\kappa)\overline\Delta$ unchanged.

\begin{theorem}[Filtered chiral PBW comparison]\label{cbc:bars:pbw}
For a positive chiral Lie object $L$, the map~\eqref{cbc:bars:pbw-map} gives a natural equivalence of coaugmented counital coassociative coalgebras
\[
 \kappa:\operatorname{CE}_{\ch}^+(L)
 \xrightarrow{\ \sim\ }
 \BarAssCh\bigl(U_{\ch}^{\mathsf{Ass}}(L)\bigr).
\]
The source has its unshuffle coproduct, and the target has deconcatenation. The comparison preserves the filtration by the total number of original $L$-inputs. This filtration is finite on every configuration chart.
\end{theorem}
\begin{proof}
Cutting a permutation in~\eqref{cbc:bars:pbw-map} amounts to choosing a subset of its inputs, ordering that subset and ordering its complement. These are exactly the terms of the unshuffle coproduct followed by $\kappa\starCh\kappa$, with the same signs. Hence $\kappa$ preserves coproducts, coaugmentation and counit.

For two homogeneous inputs, the differential identity is
\[
 b_2\kappa(sx\,sy)
 =(-1)^{|x|}s\bigl(xy-(-1)^{|x||y|}yx\bigr)
 =(-1)^{|x|}s[x,y].
\]
In any length, pair the permutations in which the two factors being multiplied occur in opposite adjacent orders. Their difference is the displayed bracket relation. Summing over the pairs gives the bracket part of the chain-map equation. The unary parts commute by naturality of the internal differential. Thus $\kappa$ is a chain coalgebra map.

Filter the envelope by the number of original $L$-factors. The operadic PBW identity identifies its associated graded with $\Sym_{\ch}(L)$. The required identity is the freeness of the associative operad as a right Lie-module, with symmetric operations as generators~\cite[Theorem~4.2 and Proposition~4.4]{DotsenkoTamaroffPBW}. It is an identity of multilinear operations over a characteristic-zero field, so it can be evaluated using the exact, colimit-preserving chiral tensor. A free simplicial Lie resolution computes the derived envelope. Realization preserves this identity because finite tensor powers and finite symmetric-group coinvariants preserve geometric realizations.

Give the Chevalley source its symmetric-degree filtration and the bar target the sum of the PBW degrees of its factors. The associated graded comparison is
\begin{equation}\label{cbc:bars:abelian-koszul}
 \Sym_{\ch}(sL)\longrightarrow B_{\starCh}(\Sym_{\ch}L),
\end{equation}
with only the internal differential on the left.

To check this comparison, put $S=\Sym_{\ch}L$ and consider the free Koszul resolution
\[
 S\starCh\Sym_{\ch}(sL).
\]
In addition to its internal differential, give it the derivation $d_K(sx)=-x$, $d_K(x)=0$. The degree-minus-one derivation $h(x)=-sx$, $h(sx)=0$ satisfies
\[
 d_Kh+hd_K=N,
\]
where $N$ counts total polynomial weight. The internal differential anticommutes with $h$. Thus $h/N$ contracts every positive-weight part, and the augmentation has homology $K$.

The filtration by suspended degree has subquotients induced from coefficient objects by the free $S$-module functor. On each configuration this filtration is finite. The chosen presentation computes the tensors in these subquotients, so this resolution computes the derived tensor with $K$. The normalized free bar resolution does so as well, by its extra-degeneracy contraction.

Use the sign convention in which the bar resolution's action term sends $1\otimes[x]$ to $-x$. Extending~\eqref{cbc:bars:pbw-map} $S$-linearly then gives the augmentation-preserving comparison between the resolutions. Its action terms match $d_K(sx)=-x$, and its interior terms cancel in adjacent pairs by graded commutativity. Tensoring this comparison with $K$ proves~\eqref{cbc:bars:abelian-koszul}.

On the chart $X^I$, a term of PBW weight $r$ contains $r$ positive original inputs and vanishes for $r>|I|$. Every filtration just used is therefore finite on that chart. The cone of $\kappa$ has acyclic associated graded. Starting at weight zero and successively extending by its graded pieces proves that the cone is acyclic. The chart evaluations are jointly conservative, which proves the claimed equivalence.
\end{proof}

The factorial in ordinary PBW symmetrization and its absence in~\eqref{cbc:bars:pbw-map} have different purposes. Averaging identifies symmetric powers with the associated graded envelope. The unnormalized sum intertwines unshuffle with deconcatenation.

For comparison, let $k$ be a field of characteristic zero and let $A=k[x]$ with augmentation $x\mapsto0$. Its augmentation ideal, regarded only as a Lie algebra, is abelian. The Chevalley class $s(x^2)$ is nonzero. Its image in the associative bar is $[x^2]=b_2[x|x]$, a boundary. Thus replacing the envelope hypothesis by the commutator Lie algebra of an arbitrary associative algebra changes the comparison.

\begin{definition}[PBW-admissible completion]\label{cbc:bars:pbw-admissible}
Consider a chiral Lie object and its envelope in a specified completed tensor category. Call this datum PBW-admissible when the following conditions hold. The envelope and the continuous comparison~\eqref{cbc:bars:pbw-map} exist. The PBW morphism identifies the associated graded envelope with the symmetric algebra. The induced associated graded comparison is~\eqref{cbc:bars:abelian-koszul}.

The filtrations on both complexes and their comparison cone have strongly convergent spectral sequences. Their limiting pages identify with the associated graded objects of complete, separated, exhaustive filtrations on cohomology.
\end{definition}

For a PBW-admissible completion the same comparison is a quasi-isomorphism. Indeed, its cone has zero associated graded cohomology by~\eqref{cbc:bars:abelian-koszul}. Strong convergence makes the associated graded of its cohomology zero, and separatedness and exhaustiveness make that cohomology zero. The finite-configuration proof above supplies these convergence requirements on positive Ran coefficients. It does not supply continuity for a different infinite product completion.

\section{Realization and factorization}
Factorization compares the coefficient at a separated configuration with the external tensor product of its coefficients at the separate configurations. This condition is nonlinear. In a rank-one example, the singleton and separated two-point coefficients are both $k$. The componentwise sum of two copies has coefficient $k^2$ at both configurations, whereas factorization would require $k^2\otimes_k k^2\cong k^4$ at two points.

A componentwise shift has the same difficulty. Its two-point coefficient is $k[1]$, whereas the tensor product of its singleton coefficients is $k[2]$. Thus stability of the ambient coefficients does not make factorization objects closed under componentwise sums or shifts.

Geometric realization has a different colimit property. An indexing infinity-category $J$ has a cofinal diagonal $J\to J\times J$ when restriction along it preserves the colimit of every diagram on $J\times J$. Call $J$ \emph{sifted} when it is nonempty and its diagonal is cofinal. Iteration gives the same statement for $J\to J^m$, for every finite $m\geq1$. The category $\Delta^{\mathrm{op}}$ indexing simplicial objects is sifted~\cite[Definition~5.5.8.1 and Lemma~5.5.8.4]{LurieHTT}.

\begin{proposition}[Sifted realization of factorization objects]
\label{cbc:bars:factorization-realization}
Fix coefficient categories on configuration charts, with specified restriction and collision functors and their coherent identifications. Suppose these categories admit $J$-indexed colimits, and these functors preserve them. Suppose external tensor products preserve those colimits separately in each variable. If $J$ is sifted, the componentwise colimit of a $J$-diagram of factorization objects and factorization morphisms is a factorization object.
\end{proposition}
\begin{proof}
Write $F_I=\operatorname*{colim}_{j\in J}(F_j)_I$. Collision maps and their coherences pass to the colimit by the stated preservation hypothesis. For a decomposition $I=I_1\sqcup\cdots\sqcup I_m$, let $u$ denote restriction to its disjointness locus. Restriction, factorization, cofinality and separate preservation give
\[
\begin{aligned}
 uF_I
 &\simeq\operatorname*{colim}_{j\in J}u(F_j)_I\\
 &\simeq\operatorname*{colim}_{j\in J}
        \mathop{\boxtimes}_{a=1}^m(F_j)_{I_a}\\
 &\simeq\operatorname*{colim}_{(j_1,\ldots,j_m)\in J^m}
        \mathop{\boxtimes}_{a=1}^m(F_{j_a})_{I_a}
 \simeq\mathop{\boxtimes}_{a=1}^m F_{I_a}.
\end{aligned}
\]
All comparisons are natural in the diagram and the decomposition, so they preserve refinement and permutation coherences. If an empty-configuration unit is specified, its constant diagram has the same unit as colimit: a sifted category has contractible groupoid completion~\cite[Proposition~5.5.8.7]{LurieHTT}.
\end{proof}

For a simplicial bar, every term, face and degeneracy must first be a morphism in the specified factorization category. Restriction and collision functors must also preserve realization. Under these hypotheses the proposition constructs the whole realized factorization object. Its normalized summands remain ambient coefficients, and any coproduct must separately be compatible with the factorization equivalences.

\section{Two orders of realization}
The two associative bars can be iterated only when their remaining structures and augmentations are preserved. A mixed distributive law must therefore extend to the complete simplicial diagrams, including end faces, units and all substitution homotopies.

\begin{definition}[Bar-admissible mixed datum]
\label{cbc:bars:mixed-admissible}
Fix a coefficient infinity-category $\mathcal C$. A mixed associative datum is bar-admissible when it carries a bisimplicial object
\[
 B_{\bullet,\bullet}(A):
 \Delta^{\mathrm{op}}\times\Delta^{\mathrm{op}}\longrightarrow\mathcal C
\]
whose two directions are identified with the specified two-sided bar constructions. The identifications include their augmentation and unit maps and every coherence. Both monoidal products are required to preserve the realizations used to form the structured iterated bars.
\end{definition}

This condition includes the compatibility of the two distinct units. A comparison of binary multiplications alone does not specify the displayed bisimplicial object.

\begin{theorem}[Mixed Fubini theorem]\label{cbc:bars:fubini}
Let $A$ be bar-admissible, and suppose its coefficient category $\mathcal C$ admits the displayed realizations. Then the two iterated realizations and the diagonal realization are naturally equivalent:
\begin{align*}
 \left|p\longmapsto\left|q\longmapsto B_{p,q}(A)\right|\right|
 &\simeq\left|n\longmapsto B_{n,n}(A)\right|,\\
 \left|q\longmapsto\left|p\longmapsto B_{p,q}(A)\right|\right|
 &\simeq\left|n\longmapsto B_{n,n}(A)\right|.
\end{align*}
\end{theorem}
\begin{proof}
Mapping either iterated colimit to a test object gives the limit, in the opposite order, of the same diagram of mapping spaces. Limits over the two indexing categories commute. The universal property of colimits therefore identifies both iterated realizations with
\[
 \operatorname*{colim}_{(p,q)\in\Delta^{\mathrm{op}}\times\Delta^{\mathrm{op}}}
 B_{p,q}(A).
\]

The diagonal $\Delta^{\mathrm{op}}\to\Delta^{\mathrm{op}}\times\Delta^{\mathrm{op}}$ is cofinal. Apply the cofinality criterion of~\cite[Theorem~4.1.3.1]{LurieHTT}. For $([p],[q])$, the relevant comma category is opposite to the category of monotone strings in $[p]\times[q]$.

Append the largest vertex $(p,q)$ to each string. Including the old string gives a natural transformation from the identity to this extension functor. Selecting the appended vertex gives another from the constant largest-vertex string. The nerve is therefore contractible. Cofinality identifies the displayed colimit with the diagonal realization, proving both equivalences.

The comparisons are natural in the entire bisimplicial diagram. Any additional operations retained through realization are consequently preserved. The monoidal preservation assumptions ensure that the iterated colimits are the intended structured bars. No tensor assumption is needed for the colimit identity once the bisimplicial object itself is given.
\end{proof}

In the remainder of this part,
\[
 \Barc_X^{\ord}=\BarPerp
\]
always denotes the transverse bar. The envelope comparison relates an associative chiral bar to a Lie-chiral coalgebra in the same chiral tensor. Fubini compares two orders of an already specified mixed construction. Neither comparison identifies the transverse and chiral products or constructs an additional associative algebra by itself.

\chapter{The internal ordered bar}
The transverse bar records the associative multiplication while retaining chiral factorization. Its normalized complex uses ordered words. The factorization maps compare entire bar complexes on separated configurations. The suspended shuffle connects these descriptions.

Work in the positive Ran coefficient category $\Dcal_X$ with its pointwise tensor $\tensor^!$ and unit $\mathbbm1=\omega_{\Ran}$. By Proposition~\ref{cbc:ran:pointwise}, the same product and unit belong to $\Fact_D(X)$. If an empty-configuration coefficient is adjoined, the pointwise unit throughout is instead $(\C,\omega_{\Ran})$.

Let $A$ be an associative algebra in $(\Fact_D(X),\tensor^!)$, with an algebra augmentation $\epsilon:A\to\mathbbm1$. Assume the restriction and collision functors preserve geometric realizations, as do the external and pointwise tensor products separately. The two-sided simplicial bar has terms
\[
 B_r(\mathbbm1,A,\mathbbm1)
 =\mathbbm1\tensor^!A^{\tensor^!r}\tensor^!\mathbbm1.
\]
Its inner faces multiply adjacent factors, its end faces use $\epsilon$, and its degeneracies insert the algebra unit. These are factorization morphisms, including their simplicial coherences. Proposition~\ref{cbc:bars:factorization-realization} therefore constructs the realization
\[
 BA=\Barc_X^{\ord}(A)
   :=\left|B_\bullet(\mathbbm1,A,\mathbbm1)\right|
\]
as a factorization object.

\section{Normalized words and their cuts}
Choose a symmetric monoidal dg presentation of the ambient coefficients. Require the displayed tensors to compute derived tensors and direct-sum normalization to compute realization. The following chain formulas refer to this presentation. Write $\bar A=\ker(\epsilon)$. The unit splits $A=\mathbbm1\oplus\bar A$. Let $s=[1]$ lower cohomological degree by one, so $|sa|=|a|-1$.

The tensor coalgebra on $s\bar A$ is the ambient graded object
\[
 T^c(s\bar A)=\mathbbm1\oplus
       \bigoplus_{n\geq1}(s\bar A)^{\tensor^!n}.
\]
Write $[a_1|\cdots|a_n]$ for a tensor of suspended homogeneous inputs, and $[]$ for the length-zero component $\mathbbm1$. These expressions specify coefficient maps in the dg presentation. No choice of global pure-tensor generators is needed. Its deconcatenation is
\begin{equation}\label{cbc:internal:deconcatenation}
 \Delta_{\mathrm{dec}}[a_1|\cdots|a_n]
 =\sum_{i=0}^n[a_1|\cdots|a_i]\tensor^!
                    [a_{i+1}|\cdots|a_n].
\end{equation}
The counit projects to length zero, and the coaugmentation includes that component.

Define the degree-one maps
\[
 b_1(sa)=-s(da),\qquad
 b_2(sa\tensor^!sb)=(-1)^{|a|}s(ab).
\]
Extending them as a coderivation means applying each map at every consecutive position, with the sign obtained by passing the preceding suspended factors. Explicitly, if $e_i=\sum_{j<i}(|a_j|-1)$, then
\begin{align*}
 b[a_1|\cdots|a_n]
 &=\sum_{i=1}^n(-1)^{e_i+1}
      [a_1|\cdots|da_i|\cdots|a_n]\\
 &\quad+\sum_{i=1}^{n-1}(-1)^{e_i+|a_i|}
      [a_1|\cdots|a_i a_{i+1}|\cdots|a_n].
\end{align*}
For closed degree-zero inputs, this gives
\[
 b[a|c]=[ac],\qquad
 b[a|c|e]=[ac|e]-[a|ce].
\]
The second formula displays the associativity relation in $b^2$.

\begin{theorem}[Normalized bar]\label{cbc:internal:normalized}
Under the stated dg and realization hypotheses, the underlying transverse bar is naturally represented by
\[
 BA\simeq\bigl(T^c(s\bar A),b\bigr).
\]
Deconcatenation makes it a coaugmented counital coassociative dg coalgebra in the pointwise coefficient tensor. Its word-length filtration
\[
 F_nBA=\bigoplus_{0\leq r\leq n}(s\bar A)^{\tensor^!r}
\]
is exhaustive, and the reduced coproduct into $n+1$ positive factors vanishes on $F_nBA$. This is the conilpotence asserted here.
\end{theorem}
\begin{proof}
In each simplicial degree, split every factor as $\mathbbm1\oplus\bar A$. The degeneracy subobject consists of summands with a unit factor. Its quotient is the tensor power of $\bar A$, shifted by the simplicial degree. The end faces vanish on this quotient, and the remaining faces and internal differential give the displayed $b$ with the stated suspension convention.

The unary part squares to zero because $d^2=0$. Its commutator with the binary part is zero by the graded Leibniz rule for $d$. Two disjoint multiplications cancel in opposite orders, whereas the two overlapping multiplications cancel by associativity. Thus $b^2=0$.

Two successive cuts enumerate the same three consecutive subwords in either order, proving coassociativity. In $\Delta_{\mathrm{dec}}b$, each cut surviving a multiplication lies wholly before or after the merged pair. These are exactly the terms of $(b\tensor^!1+1\tensor^!b)\Delta_{\mathrm{dec}}$, with the tensor sign supplied by the degree of the prefix. The unary terms satisfy the same identity. Hence deconcatenation is a chain map.

The differential preserves $F_n$, since its components preserve or lower length. Direct-sum normalization gives $BA=\operatorname*{colim}_nF_nBA$. A reduced coproduct into $n+1$ factors would divide a word of length at most $n$ into $n+1$ nonempty subwords. No such division exists. This proves the filtration form of conilpotence in the coefficient category.
\end{proof}

The pointwise tensor permits arbitrarily many factors on a fixed configuration. Conilpotence follows from the exhaustive word filtration, with no uniform configuration bound and no word-length completion. The filtration terms and their shifted summands are ambient coefficients. The whole realization carries the factorization structure.

\section{The shuffle comparison}
Write $\otimes$ for the tensor of a chosen symmetric monoidal dg coefficient model. Let $P,Q$ be augmented dg algebras in it, satisfying the derived tensor and realization hypotheses above. The multiplication on their tensor product is
\[
 (p\otimes q)(p'\otimes q')
   =(-1)^{|q||p'|}pp'\otimes qq'.
\]

Embed $P$ and $Q$ by $p\mapsto p\otimes1$ and $q\mapsto1\otimes q$. A shuffle of $r$ inputs from $P$ and $t$ from $Q$ is an interleaving preserving the order within each group. Define
\begin{equation}\label{cbc:internal:shuffle}
 \operatorname{sh}:BP\otimes BQ\longrightarrow B(P\otimes Q)
\end{equation}
by summing over these shuffles. Moving $1\otimes q_j$ before $p_i\otimes1$ contributes $(-1)^{(|p_i|-1)(|q_j|-1)}$, the sign for the suspended inputs. Empty words act as units for this formula. For degree-zero $p,q$,
\[
 \operatorname{sh}([p]\otimes[q])
 =[p\otimes1|1\otimes q]-[1\otimes q|p\otimes1].
\]

The tensor-product coalgebra also uses a sign. If
$\Delta u=\sum u'\otimes u''$ and $\Delta v=\sum v'\otimes v''$,
with homogeneous word factors, then
\begin{equation}\label{cbc:internal:tensor-coproduct}
 \Delta(u\otimes v)
 =\sum(-1)^{|u''||v'|}
       (u'\otimes v')\otimes(u''\otimes v'').
\end{equation}
The degrees here are the total suspended word degrees. In the binary example, the term $[1\otimes q]\otimes[p\otimes1]$ in the reduced coproduct has coefficient $-1$. This is precisely the crossing sign in~\eqref{cbc:internal:tensor-coproduct}.

\begin{proposition}[Shuffle and factorization]\label{cbc:internal:shuffle-factorization}
The shuffle~\eqref{cbc:internal:shuffle} is a natural quasi-isomorphism of coaugmented counital dg coalgebras. The iterated shuffles satisfy the unit, associativity and symmetry identities.

For geometric coefficients, assume the chosen restriction, collision and external-product comparisons commute coherently with the normalized bars and shuffles. Then the shuffles give the factorization equivalences of $BA$, and deconcatenation is a morphism of factorization objects after passage to the derived category.
\end{proposition}
\begin{proof}
Consider two neighboring suspended inputs from different factors. The two shuffles that interchange them contribute mixed multiplication terms with opposite signs, because
\[
 (p\otimes1)(1\otimes q)
   =(-1)^{|p||q|}(1\otimes q)(p\otimes1).
\]
Their cancellation leaves only adjacent multiplications within the same factor. These give the multiplication part of the differential of $BP\otimes BQ$. The internal differentials commute with the graded permutations, with the tensor differential sign. Thus $b\operatorname{sh}=\operatorname{sh}b$.

A cut in a shuffled word determines cuts in both original words and shuffles of the two prefixes and the two suffixes. Conversely, those four choices determine the cut shuffled word. The resulting permutation first moves the $Q$ prefix across the $P$ suffix, giving $(-1)^{|u''||v'|}$, and then shuffles each pair. This proves
\[
 \Delta_{\mathrm{dec}}\operatorname{sh}
   =(\operatorname{sh}\otimes\operatorname{sh})\Delta_{BP\otimes BQ}.
\]
The empty-word terms also give compatibility with counit and coaugmentation.

For the quasi-isomorphism, form the bisimplicial coefficient object
\[
 (r,t)\longmapsto B_r(\mathbbm1,P,\mathbbm1)
                  \otimes B_t(\mathbbm1,Q,\mathbbm1).
\]
Its iterated realization is $BP\otimes BQ$, because the tensor preserves realization separately. Its diagonal is the simplicial bar of $P\otimes Q$: moving the $P$ and $Q$ factors together gives exactly the displayed graded multiplication.

The diagonal and iterated realizations agree by Theorem~\ref{cbc:bars:fubini}. On normalized chains the comparison is the shuffle.

Indeed, the monotone paths across an $r$ by $t$ rectangle triangulate the product of the two simplices. Each path inserts the degeneracies of its interleaving. Its orientation, together with the internal graded rearrangement, is the suspended shuffle sign. Shared faces cancel, leaving the product boundary. This identifies the written chain map with the realization equivalence.

For several factors, either parenthesization enumerates each interleaving preserving all internal orders exactly once. Its sign is the Koszul sign of the same permutation. The same observation after permuting whole groups proves symmetry, and the empty group proves the unit identity.

Now let $I=I_1\sqcup I_2$, and let $j$ include the locus where the two groups of points are disjoint. The algebra factorization identifies $j^!A_I$ with $j^!(A_{I_1}\boxtimes A_{I_2})$. Applying the shuffle gives
\[
 j^!(BA_{I_1}\boxtimes BA_{I_2})
       \xrightarrow{\ \operatorname{sh}\ }j^!BA_I.
\]
Its inverse equivalence is the factorization map. The coalgebra identity gives compatibility with the pointwise deconcatenation, the several-factor identities give refinement and permutation coherence, and the assumed functor comparisons give collision coherence. These are the required compatibilities for a factorization morphism.
\end{proof}

The geometric hypothesis concerns coherent comparisons of functors and their operations. A stable ambient category alone does not supply them. Under that hypothesis, the equivalence concerns the whole bar, even though individual normalized summands need not factorize.

\section{Three coproducts}
There are two coalgebra structures on the transverse bar, with different tensor targets. Factorization separates the support:
\[
 j^!BA_I\xrightarrow{\ \sim\ }
       j^!(BA_{I_1}\boxtimes BA_{I_2}).
\]
These are the disjointness components of its chiral commutative coproduct. Deconcatenation instead repeats the same configuration in a pointwise tensor:
\[
 BA_I\xrightarrow{\ \Delta_{\mathrm{dec}}\ }
       BA_I\tensor^!BA_I.
\]
The shuffle proposition proves compatibility between these structures.

A third coproduct belongs to an acting bialgebra $H$. It is an algebra map $\Delta_H:H\to H\otimes H$, with an algebra counit, satisfying coassociativity and the counit identities. An antipode makes $H$ a Hopf algebra.

Let $M,N$ be $H$-representations. To act on $M\otimes N$, first apply $\Delta_H$ to the acting factor. Move the second copy of $H$ past $M$ by symmetry, then apply the two representation actions. A completed tensor product requires a specified topology and continuous maps throughout.

\begin{proposition}[Coproduct separation]\label{cbc:internal:three-coproducts}
Factorization, deconcatenation and an acting Hopf coproduct are related by comparisons in their specified tensor recipients. In strict dg presentations, let a strong monoidal dg functor $\Phi$ carry $H$ to the pointwise coefficient category. Let $f:\Phi H\to BA$ be a chain map preserving counit and coaugmentation. Then $f$ is a coalgebra comparison exactly when
\[
 \Delta_{\mathrm{dec}}f
   =(f\tensor^!f)\Phi(\Delta_H),
\]
using the monoidal identification of the target of $\Phi(\Delta_H)$.

For coherent coalgebras, this condition includes the compatible homotopies for every iterated coproduct. A comparison with the chiral coproduct requires transport to the chiral tensor recipient and compatibility on every disjointness locus. Retaining a counit requires the unitalized recipient.
\end{proposition}
\begin{proof}
The displayed equation is the coproduct condition for a strict dg coalgebra morphism in the pointwise tensor. Iterating it and using coassociativity gives every higher coproduct identity. With the stipulated unit maps, it therefore preserves the coaugmented counital coalgebra structure. For coherent coalgebras these iterations are part of the morphism datum. In the chiral recipient, restricting the corresponding coproduct components gives the disjointness equations.

For three representations $M,N,P$, applying $\Delta_H$ twice gives the two candidate actions on their tensor product. Coassociativity identifies them, and the counit gives the unit representation. This explains the representation tensor product supplied by $H$ without identifying its carrier with a Ran coefficient or a bar word.
\end{proof}

\section{Interval factorization homology}
Let $(\mathcal C,\otimes,\mathbbm1)$ be a symmetric monoidal infinity-category admitting sifted colimits, with tensor preserving them separately. An oriented interval coefficient system assigns an associative algebra $A$ to each interior interval. A left endpoint neighborhood receives a right $A$-module $R$, and a right endpoint neighborhood receives a left $A$-module $L$. Ordered embeddings act by multiplication and the module actions, with their coherences.

In the unital disk convention, nullary endpoint operations also specify maps $\mathbbm1\to R$ and $\mathbbm1\to L$. These are part of the coefficient system~\cite[Definition~2.30]{AFT}. Factorization homology is the colimit of its tensor values over disjoint unions of such basic intervals embedded in $[0,1]$.

\begin{proposition}[The endpoint bar]\label{cbc:internal:interval}
For this coefficient system there is a natural equivalence
\[
 \int_{([0,1],\{0,1\})}(A;R,L)
 \simeq\left|r\longmapsto R\otimes A^{\otimes r}\otimes L\right|
 \simeq R\otimes_A^{\mathbb L}L.
\]
The superscript records the derived relative tensor product. In an infinity-category this is the relative tensor product itself.
\end{proposition}
\begin{proof}
Replace the interval disk category by its equivalent category of finite ordered configurations with optional endpoint components. The full subcategory containing both endpoint components is the simplex category in the opposite variance. Adjoining a missing first or last endpoint component is left adjoint to its inclusion, so the inclusion is final. This is the interval finality argument of~\cite[Observations~2.31--2.32 and Proposition~2.34]{AFT}.

The restricted diagram has value $R\otimes A^{\otimes r}\otimes L$ in degree $r$. Its first and last faces apply the module actions, its inner faces multiply in $A$, and its degeneracies insert the unit. Finality proves the first equivalence. The realization of this two-sided bar computes the relative tensor product by~\cite[Theorem~4.4.2.8 and Example~4.4.2.11]{LurieHA}. Its tensor-preservation hypotheses are exactly those imposed above.
\end{proof}

For modules without endpoint point maps, the diagram with both endpoints present still defines the relative interval bar and gives $R\otimes_A^{\mathbb L}L$. To use the full unital disk category, the endpoint point maps must be included in the datum.

Apply this construction internally to the pointwise coefficient tensor. The augmentation $A\to\mathbbm1$ makes $R=L=\mathbbm1$ into the two endpoint modules, with identity point maps. The endpoint bar is then $BA=\Barc_X^{\ord}(A)$. The earlier functor hypotheses retain its chiral factorization. This identifies the interval coefficient object. A numerical amplitude or prepared boundary state requires its own evaluation or state map in addition to the augmentation.

\chapter{Cobar and twisting cochains}
A coalgebra map to a tensor bar is determined by its projection to bar length one.
The chain-map condition on that projection is a quadratic equation.
Tensor cobar represents the same equation by an algebra map in the other direction.

\section{The tensor cobar differential}
Fix a commutative ring $k$ and use cohomologically graded $k$-modules, with tensor product over $k$.
Let $A$ be a unital augmented dg associative algebra, with augmentation ideal $\bar A$.
Let $C$ be a counital coaugmented dg coassociative coalgebra.
Write $1$ for its coaugmentation and put $\bar C=\ker(\epsilon_C:C\to k)$.
On $\bar C$, the reduced coproduct is
\[
 \bar\Delta(c)=\Delta(c)-1\otimes c-c\otimes1.
\]

Let $\bar\Delta_r:\bar C\to\bar C^{\otimes r}$ be its iterates, with $\bar\Delta_1=\id$.
Coassociativity makes the choice of parenthesization immaterial.
Assume $C$ is conilpotent: every $c\in\bar C$ is killed by $\bar\Delta_r$ for all sufficiently large $r$.

Retain the bar suspension $|sa|=|a|-1$ of Theorem~\ref{cbc:internal:normalized}.
Write $t=s^{-1}$, so $|tc|=|c|+1$.
Both tensor constructions below use direct sums, with the empty tensor equal to $k$.
In a dg coefficient presentation of the transverse theory, the tensor being modeled is the pointwise product $\tensor^!$.

\begin{definition}[Tensor cobar]\label{cbc:cobar:definition}
The ordinary tensor cobar is the augmented algebra
\[
 \Omega C=T(t\bar C)=\bigoplus_{r\geq0}(t\bar C)^{\otimes r}.
\]
Write $\bar\Delta c=\sum c'\otimes c''$ for a homogeneous expansion of the reduced coproduct.
Its differential is the degree-one derivation determined by
\begin{equation}\label{eq:cbc:cobar:differential}
 d_\Omega(tc)=-t(d_Cc)
       -\sum(-1)^{|c'|}(tc')(tc'').
\end{equation}
The augmentation projects to tensor length zero.
\end{definition}

A primitive element $p$ has $\bar\Delta p=0$, so $d_\Omega(tp)=-t(d_Cp)$.
If $p,q,c$ are closed and $\bar\Delta c=p\otimes q$, then
\[
 d_\Omega(tc)=-(-1)^{|p|}(tp)(tq).
\]
The two factors retain their order.

\begin{proposition}\label{cbc:cobar:square}
The derivation~\eqref{eq:cbc:cobar:differential} squares to zero and makes $\Omega C$ an augmented dg algebra.
\end{proposition}
\begin{proof}
Write $d_\Omega=d_1+d_2$ for its internal and quadratic terms.
The square of an odd derivation is a derivation, so it suffices to check the generators.
The identity $d_1^2=0$ follows from $d_C^2=0$.

For a term $c_1\otimes c_2\otimes c_3$ of the iterated reduced coproduct, put $p=|c_1|$ and $q=|c_2|$.
Splitting the first factor after the first cut contributes the sign
$(-1)^{p+q+1}(-1)^{p+1}=(-1)^q$.
Splitting the second contributes
$(-1)^{p+1}(-1)^{p+1}(-1)^{q+1}=(-1)^{q+1}$.
Coassociativity identifies their coefficients, so they cancel.
Thus $d_2^2=0$.

The chain-map identity for the reduced coproduct is
\[
 \bar\Delta(d_Cc)
 =\sum\bigl(d_Cc'\otimes c''+(-1)^{|c'|}c'\otimes d_Cc''\bigr).
\]
Substitution in~\eqref{eq:cbc:cobar:differential} gives
\[
 d_1d_2(tc)=\sum\bigl((-1)^{|c'|}t(d_Cc')(tc'')-(tc')t(d_Cc'')\bigr)
             =-d_2d_1(tc).
\]
This proves the square-zero identity.
The differential has no scalar term, so the augmentation is a chain map.
\end{proof}

\section{The twisting equation}
For homogeneous graded maps $f,g:\bar C\to\bar A$, define the convolution product and differential by
\[
 (f\star g)(c)=\sum(-1)^{|g||c'|}f(c')g(c''),
 \qquad \partial f=d_Af-(-1)^{|f|}fd_C.
\]
Associativity and coassociativity make $\star$ associative.
The chain-map identities for multiplication and coproduct give
$\partial(f\star g)=(\partial f)\star g+(-1)^{|f|}f\star(\partial g)$.

\begin{definition}[Twisting cochain]\label{cbc:cobar:twisting}
A twisting cochain is a degree-one map $\tau:\bar C\to\bar A$ satisfying
\begin{equation}\label{eq:cbc:cobar:MC}
 \partial\tau+\tau\star\tau=0,
 \qquad
 d_A\tau(c)+\tau(d_Cc)+\sum(-1)^{|c'|}\tau(c')\tau(c'')=0.
\end{equation}
Write $\operatorname{Tw}(C,A)$ for this set.
Equivalently, extend $\tau$ to $C\to A$ by zero on the coaugmentation and require its augmentation component to vanish.
\end{definition}

This is the associative Maurer--Cartan equation.
It involves no division by two and is defined over the chosen ring $k$.

\section{The two universal extensions}
The tensor bar $BA=T^c(s\bar A)$ has the deconcatenation coproduct and the differential of Theorem~\ref{cbc:internal:normalized}.
In particular, its components with output of length one are
\[
 d_B(sa)=-s(d_Aa),\qquad d_{B,2}(sa\otimes sb)=(-1)^{|a|}s(ab).
\]
Let $p_A:BA\to s\bar A$ be the length-one projection.

\begin{theorem}[Bar--cobar adjunction]\label{cbc:cobar:adjunction}
For the augmented dg algebras and coaugmented conilpotent dg coalgebras specified above, there are natural bijections
\begin{equation}\label{eq:cbc:cobar:adjunction}
 \operatorname{Hom}_{\mathrm{dgAlg}^{\mathrm{aug}}}(\Omega C,A)
 \cong\operatorname{Tw}(C,A)
 \cong\operatorname{Hom}_{\mathrm{dgCoalg}^{\mathrm{coaug}}}(C,BA).
\end{equation}
Thus ordinary tensor cobar is left adjoint to tensor bar.
\end{theorem}
\begin{proof}
Given $\tau$, send $tc$ to $\tau(c)$ and extend multiplicatively, sending the empty tensor to $1_A$.
This defines the unique augmented graded algebra map $g_\tau:\Omega C\to A$ with the stated generator values.
Its chain-map condition on a generator is precisely~\eqref{eq:cbc:cobar:MC}.
Since both differentials are derivations, that condition suffices on every tensor.
Restriction to the generators gives the inverse bijection.

For the coalgebra extension, set $\pi_C(c)=c-\epsilon_C(c)1$ and define
\begin{equation}\label{eq:cbc:cobar:extension}
 f_\tau(c)=\epsilon_C(c)1+
       \sum_{r\geq1}(s\tau)^{\otimes r}\bar\Delta_r(\pi_C(c)).
\end{equation}
The map $s\tau$ has degree zero.
Conilpotence makes this a finite sum on every input, hence an element of the direct-sum bar.
Cutting its length-$r$ term after position $i$ gives the $i,r-i$ component of $(f_\tau\otimes f_\tau)\Delta_C$.
This proves the coalgebra identity, including the two counital cuts.

Conversely, any coaugmented coalgebra map to $T^c(s\bar A)$ has these components, determined by its length-one projection.
Thus~\eqref{eq:cbc:cobar:extension} is the unique extension of $s\tau$.

To check its differential, put $D=d_Bf_\tau-f_\tau d_C$.
The coderivation identities imply
\[
 \Delta_BD=(D\otimes f_\tau+f_\tau\otimes D)\Delta_C.
\]
Iterating this identity expresses the length-$r$ component of $D$ as the sum of tensors with one factor $p_AD$ and all other factors $p_Af_\tau$.
Consequently $D=0$ if and only if $p_AD=0$.

The two displayed bar components give
\[
 p_AD=-s\bigl(d_A\tau+\tau d_C+\tau\star\tau\bigr).
\]
This is zero exactly when $\tau$ is a twisting cochain.
The inverse map sends a dg coalgebra map $f$ to $t p_Af$.
Both extension formulas commute with maps of $C$ and $A$, proving naturality.
\end{proof}

The usual field version appears in~\cite[Theorem~2.2.9]{LV}.
The proof above uses the fixed cohomological signs and works over the chosen commutative ring, without flatness or a completed tensor algebra.

The universal twisting maps are
\[
 \tau_C:C\longrightarrow\Omega C,\qquad \tau_C(c)=t\pi_C(c),
 \qquad
 \tau_A:BA\longrightarrow A,\qquad \tau_A=t p_A.
\]
They correspond to the identity maps of $\Omega C$ and $BA$.
The unit is $\eta_C=f_{\tau_C}:C\to B\Omega C$.
The counit is $\epsilon_A=g_{\tau_A}:\Omega BA\to A$.
It sends $t(sa)$ to $a$ and generators of larger bar length to zero.

In $\epsilon_{\Omega C}\circ\Omega\eta_C$, the length-one term of~\eqref{eq:cbc:cobar:extension} returns each generator $tc$.
In $B\epsilon_A\circ\eta_{BA}$, only the decomposition into singleton bar blocks survives, returning the original bar tensor.
These are the two triangular identities.

For an internal adjunction on factorization objects, both extensions must preserve the chosen geometric data.
One must construct the factorization structure on cobar and verify compatibility with collisions, separated products, permutations and all their coherences.
The twisting maps and both inverse correspondences must preserve those structures.
The ordinary mapping-set adjunction above then restricts wherever these conditions hold.
In that recipient, write $\Cobarc_X^{\ord}(C)=\Omega C$ and $\Barc_X^{\ord}(A)=BA$.

An adjunction on derived mapping spaces requires a compatible derived enhancement.
A completed cobar also requires its own topology, convergence and comparison maps.
These are additional constructions, not consequences of the ordinary mapping-set bijections.

\chapter{Complete bar--cobar reconstruction}
The tensor counit reconstructs an algebra from its bar coalgebra.
Positive weights permit a completion of this reconstruction.
The quadratic criterion asks when diagonal cycles contain the homology of the full bar.

\section{The counit and the gap contraction}
Let $k$ be a commutative ring and let $A=k1\oplus I$ be an augmented associative algebra in cohomological degree zero.
All tensors are over $k$.
Retain the suspensions $|sa|=|a|-1$ and $t=s^{-1}$ from the preceding chapter.
Write
\[
 D=\Omega BA,\qquad c(a_1,\ldots,a_n)=t(sa_1\otimes\cdots\otimes sa_n),
 \qquad a_i\in I.
\]
The generator $c(a_1,\ldots,a_n)$ has degree $1-n$.

The cobar differential becomes
\begin{equation}\label{eq:cbc:reconstruction:blocks}
 dc(a_1,\ldots,a_n)=\sum_{i=1}^{n-1}(-1)^{i-1}
 \bigl(c(a_1,\ldots,a_i)c(a_{i+1},\ldots,a_n)
       -c(a_1,\ldots,a_i a_{i+1},\ldots,a_n)\bigr).
\end{equation}
Extend it by the graded Leibniz rule.
The counit $\epsilon:D\to A$ sends $c(a)$ to $a$ and longer generators to zero.
In particular, $dc(a,b)=c(a)c(b)-c(ab)$, whose image under $\epsilon$ is zero.

A product of generators is a \emph{block word}.
Flatten it to its $N$ original entries.
Each gap between consecutive entries is either internal to a block or an outer cut between two blocks.
If $S\subseteq\{1,\ldots,N-1\}$ is the set of outer cuts, its degree is $|S|-(N-1)$.

For an internal gap $i$, put
\[
 u_i(S)=\#\{j<i:j\notin S\}.
\]
The splitting term inserts the cut $i$ with sign $(-1)^{u_i(S)}$.
The multiplication term has the opposite sign and lowers the number of entries.
These signs follow from~\eqref{eq:cbc:reconstruction:blocks}, including the degrees of all preceding blocks.

\begin{theorem}[Ordinary reconstruction]\label{cbc:reconstruction:ordinary}
The counit $\epsilon:\Omega BA\to A$ is a quasi-isomorphism of augmented dg algebras.
No flatness assumption on $I$ is required.
\end{theorem}
\begin{proof}
Filter $D$ increasingly by the number of original entries.
Every element has finite support in this filtration.
On its $N$th associated graded piece only splitting remains.
For $N\geq1$, this complex is
\[
 I^{\otimes N}\otimes E^{\otimes(N-1)},\qquad
 E^{-1}=ke_-,\quad E^0=ke_+,\quad de_-=e_+.
\]
The states $e_-$ and $e_+$ mean an internal gap and an outer cut, respectively.
The tensor differential counts the preceding $e_-$ factors, giving exactly $u_i(S)$.

For $N>1$, define $h$ by changing the first gap from $e_+$ to $e_-$ and by sending a word with first gap $e_-$ to zero.
On that first factor, $dh+hd=\id$.
The differentials on the remaining factors cancel between the two composites because $h$ has degree $-1$.
Thus this identity holds after tensoring with $I^{\otimes N}$, without any tensor-exactness hypothesis.
The pieces $N=0,1$ are $k,I$, both in degree zero.

For a negative-degree cycle, apply $h$ to its highest-entry component and lift the resulting primitive to $D$.
Subtract its differential.
The largest entry number strictly decreases, so repetition terminates and proves negative acyclicity.
In degree zero, the boundaries impose exactly $c(a)c(b)=c(ab)$.
Every element reduces to a scalar plus $c(a)$, with $a\in I$, and its image in $k1\oplus I$ makes this expression unique.
Hence $H^0(D)=A$, through the counit.
\end{proof}

For example, take $A=k[x]/(x^3)$ and give $x$ weight one.
The bar complex of weight three is
\[
 k[x|x|x]\xrightarrow{\ d\ }
 k[x^2|x]\oplus k[x|x^2],\qquad
 d[x|x|x]=[x^2|x]-[x|x^2],
\]
in degrees $-3,-2$.
It has one copy of $k$ in degree $-2$.
Nevertheless the full tensor counit reconstructs $A$ by the theorem.

\section{Positive-weight completion}
Suppose now that $A_{\mathrm{alg}}=\bigoplus_{w\geq0}A_w$, with $A_0=k$, multiplication adding weights and augmentation projecting onto $A_0$.
Let $D_w$ be the weight-$w$ subcomplex of $\Omega B(A_{\mathrm{alg}})$.
Define
\[
 \widehat A=\prod_{w\geq0}A_w,\qquad
 \widehat D=\prod_{w\geq0}D_w.
\]
Products of complexes are degreewise.
The weight topology has neighborhoods $F^p=\prod_{w\geq p}(-)_w$ of zero.
Multiplication is continuous: each weight receives only finitely many products of nonnegative weights.

\begin{theorem}[Complete reconstruction]\label{cbc:reconstruction:complete}
The componentwise counit $\widehat D\to\widehat A$ is a quasi-isomorphism of complete dg algebras.
\end{theorem}
\begin{proof}
Every reduced entry has positive weight, so an entry list of total weight $w$ has length at most $w$.
The preceding filtration and contraction therefore apply within each weight and give
$H^j(D_w)=0$ for $j<0$ and $H^0(D_w)=A_w$.
The product differential has the product of the component kernels as its kernel and the product of the component images as its image.
Consequently
\[
 H^j\left(\prod_wD_w\right)=\prod_wH^j(D_w).
\]
This proves the assertion.
Continuity follows from preservation of weight.
\end{proof}

For $k[x]$ with weight one on $x$, the ordinary counit recovers $k[x]$,
whereas this weight completion recovers $k[[x]]$.
A completed geometric bar--cobar equivalence needs a specified completed tensor product,
coalgebra category, and comparison maps compatible with collision and inverse limits.
An operadic nilcompletion or a boundary double centralizer is not identified
with the ordinary tensor counit by this weight calculation.
The coalgebra-unit and geometric duality statements require their own convergence and comparison arguments.

\section{The quadratic coalgebra and its resolutions}
Let $V$ be a $k$-module in cohomological degree zero and weight one, let $R\subseteq V\otimes V$, and set
\[
 A=T_k(V)/(R),\qquad A_0=k.
\]
Assume each $A_w$ is flat over $k$.
Put $K_0=k$, $K_1=V$, and, for $n\geq2$, define
\begin{equation}\label{eq:cbc:quadratic:intersection}
 K_n=\bigcap_{i+j=n-2}V^{\otimes i}\otimes R\otimes V^{\otimes j}
       \ \subseteq V^{\otimes n}.
\end{equation}
These are embedded submodules: the surrounding tensor factors are flat.
Assume also that every $K_n$ is flat.
Both flatness conditions hold automatically over a field.

The quadratic dual coalgebra is
\[
 C=A^{\mathrm{q!}}=\bigoplus_{n\geq0}s^nK_n,
\]
with zero differential and deconcatenation coproduct.
Its $n$th summand has weight $n$ and degree $-n$.

To verify the coproduct target, express each $K_n$ as the common kernel of the adjacent quotient maps with target containing $A_2=(V\otimes V)/R$.
Tensoring these kernel sequences with the flat surrounding factors preserves their kernels.
For a cut after $i$ entries, the constraints on the first $i$ and last $n-i$ entries then place the cut tensor in $K_i\otimes K_{n-i}$.
Coassociativity and the counit restrict from the tensor coalgebra.
Positive weight implies conilpotence.

The weight-$n$, degree-$-n$ term of $BA$ is $(sV)^{\otimes n}$.
Its different adjacent multiplications land in distinct weight-composition summands.
Their common kernel is exactly $s^nK_n$, and there is no incoming differential in that degree.
Thus inclusion defines a dg coalgebra map
\[
 i:C\longrightarrow BA,
\]
identifying $C_n$ with the diagonal bar cycles.

Define $\kappa:C\to A$ by $\kappa(sv)=v$ and by zero on every other weight.
It has degree one.
The only possible quadratic term in its twisting equation is multiplication of $R$, which is zero in $A$.
Hence it is a twisting cochain and induces $q:\Omega C\to A$.
The corresponding coalgebra extension is $i$.

The left and right Koszul complexes are
\[
 K_L=A\otimes C,\qquad K_R=C\otimes A.
\]
For a homogeneous coproduct $\Delta c=\sum c'\otimes c''$, their differentials are
\begin{equation}\label{eq:cbc:quadratic:twisted}
 d_L(a\otimes c)=-\sum a\kappa(c')\otimes c'',\qquad
 d_R(c\otimes a)=\sum(-1)^{|c'|}c'\otimes\kappa(c'')a.
\end{equation}
The twisting equation makes both differentials square to zero.
Their augmentations to $k$ use the augmentations of both factors.
When their augmentations are quasi-isomorphisms, these are Priddy's Koszul resolutions~\cite[Introduction]{Priddy}.

\begin{theorem}[Quadratic Koszul criterion]\label{cbc:quadratic:criterion}
Under the stated flatness hypotheses, the following conditions are equivalent:
\begin{enumerate}
\item $i:C\to BA$ is a quasi-isomorphism.
\item $q:\Omega C\to A$ is a quasi-isomorphism.
\item both augmented complexes $K_L\to k$ and $K_R\to k$ are quasi-isomorphisms.
\item $\operatorname{Tor}^A_i(k,k)_w=0$ whenever $w\ne i$.
\end{enumerate}
Either one of the two assertions in (iii) suffices.
No finiteness assumption on duals is needed for these equivalences.
\end{theorem}
\begin{proof}
First compare the twisted complexes with the augmented bar resolutions.
For $P_L=A\otimes BA$, the differential is the bar differential plus
$-a_0a_1\otimes[a_2|\cdots|a_n]$.
Writing $\bar a_0=a_0-\epsilon(a_0)1$, its augmented contraction is
\[
 h_L\bigl(a_0\otimes[a_1|\cdots|a_n]\bigr)
       =-1\otimes[\bar a_0|a_1|\cdots|a_n].
\]
The adjacent multiplication involving $\bar a_0$ supplies the identity term.
The remaining multiplications cancel in $dh_L+h_Ld$.
Thus $dh_L+h_Ld=\id-\iota\epsilon$, where $\iota:k\to P_L$ is the unit.

For $P_R=BA\otimes A$, the last action has sign $(-1)^{n-1}$ and the corresponding contraction is
\[
 h_R\bigl([a_1|\cdots|a_n]\otimes a_0\bigr)
       =(-1)^n[a_1|\cdots|a_n|\bar a_0]\otimes1.
\]
The same cancellation proves its contraction identity.

The maps $\id_A\otimes i:K_L\to P_L$ and $i\otimes\id_A:K_R\to P_R$ respect differentials and augmentations.
At total weight $w$, filter either complex decreasingly by the weight $p$ in its $A$-factor.
The twisting term increases $p$, and the associated graded comparison is $A_p\otimes i_{w-p}$, with the factors reversed on the right.
Flatness of $A_p$ proves that (i) implies either assertion in (iii).

Conversely, suppose one twisted complex resolves $k$.
Its comparison with the corresponding $P$ is then a quasi-isomorphism.
Induct on $w$.
The assertion for $w=0$ is the identity of $k$.

In the cone at weight $w$, the subcomplex with positive $A$-weight has a finite filtration with acyclic quotients $A_p\otimes\operatorname{cone}(i_{w-p})$, by induction.
Its quotient is $\operatorname{cone}(i_w)$.
The whole cone and that subcomplex are acyclic, so the quotient is acyclic.
This proves (i), from either side of (iii).

For comparison with (ii), every weight component of $C$, $BA$, $\Omega C$ and $A$ is a bounded complex of flat $k$-modules.
At a fixed weight, filter bar increasingly and cobar decreasingly by tensor length.
These finite filtrations have tensor products of the input complexes as their associated graded pieces.
Tensor products of bounded flat complexes preserve quasi-isomorphisms: apply the finite degree filtration to the cone of each tensor-factor map.
Thus $B$ and $\Omega$ preserve the quasi-isomorphisms used here.
Since $q=\epsilon\circ\Omega i$, the ordinary counit theorem gives (i)$\Rightarrow$(ii).

For the converse, consider the unit $\eta_C:C\to B\Omega C$.
Flatten a term to $N$ original reduced $C$-entries and filter decreasingly by $N$.
Cobar splitting increases $N$.
On the associated graded, only multiplication of neighboring cobar blocks remains.
It changes an outer bar cut into an inner product gap.

For homogeneous entries $c_1,\ldots,c_N$ and a cut set $S$, its coefficient at $i\in S$ is
\[
 (-1)^{\sum_{j\leq i}|c_j|+u_i(S)+1}.
\]
Contract this reversed cube by inserting the first cut, with coefficient $(-1)^{|c_1|+1}$, when it is absent, and by zero when it is present.
The first differential gives the identity.
The other terms cancel with their graded signs.

All $N>1$ pieces are therefore contractible.
The unit is the identity on the surviving empty and singleton pieces.
At weight $w$ one has $N\leq w$, so the finite filtration proves that $\eta_C$ is a quasi-isomorphism.
If (ii) holds, $Bq$ is a quasi-isomorphism, and $Bq\circ\eta_C=i$ proves (i).

Finally, $P_L$ is a flat $A$-module resolution of $k$, since $I$ and all its tensor powers are flat over $k$.
Tensoring with $k$ over $A$ gives $BA$, so $H^{-i}(BA)_w=\operatorname{Tor}^A_i(k,k)_w$.
The diagonal homology is already $C_w$ by the kernel calculation, and terms with bar length greater than weight vanish.
Hence (i) is equivalent to (iv).
\end{proof}

Over a field this is the usual quadratic criterion~\cite[Theorem~3.4.6]{LV}.
The proof above isolates the flatness needed over a general ring.

\section{A torsion boundary for the Tor comparison}
Flatness cannot simply be omitted when ordinary bar homology is used for absolute Tor.
Take $k=\mathbb Z$ and
\[
 A=\mathbb Z[x]/(2x,x^2),\qquad |x|=0,\quad\operatorname{wt}(x)=1.
\]
Its augmentation ideal is $(\mathbb Z/2)x$ with zero product.
Thus its ordinary bar has zero differential and one copy of $\mathbb Z/2$ in each positive bidegree $(-n,n)$.

The following free $A$-resolution begins differently.
Let $F_0=A$, $F_1=Ae$, $F_2=Af\oplus Ag$, and $F_3=Ah\oplus Ai\oplus Aj$, with
\[
 de=x,\qquad df=2e,\quad dg=xe,
 \qquad dh=xf,\quad di=2g,\quad dj=xg.
\]
The weights of $e,f,g,h,i,j$ are $1,1,2,2,2,3$.

The kernel in $F_3$ is generated by
\[
 2h,\quad xh,\quad xi,\quad2j,\quad xj.
\]
Take $F_4$ free on generators mapping to these five elements.

To check exactness through $F_3$, write each coefficient uniquely as $a+bx$, with $a\in\mathbb Z$ and $b\in\mathbb Z/2$.
The successive kernels are $(2e,xe)$, $(xf,2g,xg)$ and the displayed five-generator module.
This also verifies all differential composites.
Extend by free covers of the later kernels.

After tensoring with $\mathbb Z$ over $A$, the weight-two part around homological degree three is
\[
 \mathbb Z\xrightarrow{\ (2,0)\ }
 \mathbb Zh\oplus\mathbb Zi
 \xrightarrow{\ (0,2)\ }\mathbb Zg.
\]
Therefore
\[
 \operatorname{Tor}^A_3(\mathbb Z,\mathbb Z)_2=\mathbb Z/2.
\]
The ordinary bar has no such off-diagonal class.
Its terms provide a relative resolution, but are not flat $A$-modules and do not compute this absolute Tor.

Here the quadratic coalgebra is the ordinary bar itself, so its inclusion is the identity and both twisted complexes contract.
The tensor counit is also a quasi-isomorphism, while the Tor-purity condition fails.
A flat or derived replacement supplies the missing comparison.
Finiteness of ordinary coefficient duals is insufficient.

\section{Coefficient duality}
The coalgebra comparison and its dual are separate constructions.
Return to $A=T_k(V)/(R)$ and its quadratic coalgebra $C$.
Assume the equivalent quadratic conditions and suppose, in addition, that every $A_w$ and $K_w$ is finitely generated projective over $k$.
Define the degreewise products
\[
 (BA)^\vee=\prod_{w\geq0}\operatorname{Hom}_k((BA)_w,k),
 \qquad
 C^\vee=\prod_{w\geq0}\operatorname{Hom}_k(C_w,k).
\]
For homogeneous cochains, the convolution product is
\[
 (f\star g)(c)=\sum(-1)^{|g||c'|}f(c')g(c''),
 \qquad \Delta c=\sum c'\otimes c''.
\]
Each output weight involves finitely many weight decompositions, so the product is continuous.

\begin{proposition}\label{cbc:quadratic:coefficient-dual}
The dual of $i$ is a quasi-isomorphism of complete augmented dg algebras
$i^\vee:(BA)^\vee\to C^\vee$.
Moreover $(BA)^\vee$ is a dg algebra model for $\RHom_A(k,k)$.
\end{proposition}
\begin{proof}
At each weight, the cone of $i$ is an acyclic bounded complex of finitely generated projectives.
Such a complex is contractible: split off its highest nonzero degree and proceed downwards.
Dualizing preserves that contraction.
Taking products therefore proves the quasi-isomorphism, and the coalgebra identity for $i$ gives compatibility with convolution.

For the derived endomorphisms, use $P=A\otimes BA$ with the left twisted differential above.
It is a bounded-above complex of projective $A$-modules resolving $k$, hence computes derived Hom.
Its right $BA$-coaction is $\rho=\id_A\otimes\Delta$.
Lift a homogeneous $f\in(BA)^\vee$ to the $A$-linear endomorphism
\[
 L_f=(\id_P\otimes f)\rho,
 \qquad
 L_f(a\otimes c)=\sum(-1)^{|f||c'|}a\otimes c' f(c'').
\]
Coassociativity gives $L_fL_g=L_{f\star g}$, and the chain-coaction identity gives $[d_P,L_f]=L_{\partial f}$.

Postcomposition with $P\to k$ is a quasi-isomorphism
$\operatorname{End}_A(P)\to\operatorname{Hom}_A(P,k)$ because $P$ is bounded above with projective terms.
Its composite with $L$ is the identity under $\operatorname{Hom}_A(P,k)=\operatorname{Hom}_k(BA,k)$.
Thus $L$ is a quasi-isomorphism of dg algebras, proving the asserted model.
\end{proof}

If $k$ is a field and $V$ is finite-dimensional, the graded pairing identifies
\[
 C^\vee\cong\prod_{n\geq0}
 \bigl(T(V^*[-1])/(R^\perp)\bigr)_n.
\]
In each finite-dimensional weight, the annihilator of the intersection in~\eqref{eq:cbc:quadratic:intersection} is the sum of the corresponding orthogonals.
That sum is the degree-$n$ part of the ideal generated by $R^\perp$.
This proves the quadratic presentation of the dual, including its multiplication.

Over a general ring the preceding proposition retains the actual coalgebra dual.
An orthogonal-relations presentation requires its own annihilator comparison.
Neither coefficient duality nor the ordinary tensor reconstruction constructs the geometric chiral comparison or a boundary-junction reconstruction.

\chapter{Modules, comodules, and Morita theory}
The tensor counit reconstructs an algebra.
Reconstruction of its representations also requires comparison maps for modules and comodules.
Positive weights make the two bar contractions finite in each component.
Morita reconstruction asks a different question: when does a chosen object's endomorphism algebra recover its entire category?

\section{The weighted representation categories}
Let $k$ be a field and let $A_{\mathrm{alg}}=\bigoplus_{w\geq0}A_w$ be an augmented associative algebra in degree zero.
Assume $A_0=k$, multiplication adds weights, and augmentation projects onto $A_0$.
Put $C_{\mathrm{alg}}=B(A_{\mathrm{alg}})$ and
\[
 A=\prod_{w\geq0}A_w,\qquad
 C=\prod_{w\geq0}(C_{\mathrm{alg}})_w.
\]
The bar has degree $-r$ on a word $[a_1|\cdots|a_r]$.
Its differential has the convention $d_C[a|b]=[ab]$.
All products of complexes below are degreewise.

A complete weight-graded left $A$-module is a product
\[
 M=\prod_{w\geq w_0}M_w
\]
for an integer $w_0$, with differential preserving weight and unital associative actions
$A_u\otimes M_v\to M_{u+v}$.
The weight topology has neighborhoods $\prod_{w\geq p}M_w$ of zero.
Morphisms preserve weight and commute with the action and differential.
Such maps are continuous.

For $D=A$ or $C$, define the completed tensor by
\[
 (D\widehat\otimes M)_w=\bigoplus_{u+v=w}D_u\otimes_k M_v.
\]
The sum is finite, since $u\geq0$ and $v\geq w_0$.
A complete weight-graded left $C$-comodule is a similarly bounded product $N$ with a continuous counital coassociative coaction
\[
 \rho:N\longrightarrow C\widehat\otimes N
\]
commuting with its differential.
Its morphisms preserve weight and coaction.
No coaugmentation of $N$ is chosen.

Each reduced coaction lowers the weight of the final $N$-factor by at least one.
It therefore vanishes after more than $w-w_0$ iterations on $N_w$.
Write $D_{\mathrm{wt}}(A)$ and $D_{\mathrm{wt}}(C)$ for the localizations at maps that are quasi-isomorphisms in every weight.
These weak equivalences are specified on the underlying complexes.

\section{The twisting adjunction}
The universal twisting map $\tau:C\to A$ has degree one.
It sends $[a]$ to $a$ and vanishes on the empty word and on words of length greater than one.
Define
\[
 F(M)=C\widehat\otimes_\tau M,\qquad
 G(N)=A\widehat\otimes_\tau N.
\]
The underlying tensors are the ones just defined.
Their coaction and action come from $\Delta_C$ and left multiplication in $A$, respectively.

For $\Delta(c)=\sum c_{(1)}\otimes c_{(2)}$ and $\rho(n)=\sum n_{(-1)}\otimes n_{(0)}$, set
\begin{equation}\label{eq:cbc:modules:twists}
\begin{aligned}
 d_F(c\otimes m)={}&d_Cc\otimes m+(-1)^{|c|}c\otimes d_Mm
   +\sum(-1)^{|c_{(1)}|}c_{(1)}\otimes\tau(c_{(2)})m,\\
 d_G(a\otimes n)={}&a\otimes d_Nn
   -\sum a\tau(n_{(-1)})\otimes n_{(0)}.
\end{aligned}
\end{equation}
Every sum is finite in each weight.
Expanding either square leaves the twisting equation for $\tau$, followed by the action or coaction.
Thus both differentials square to zero.
The same identities make the coaction of $F(M)$ and the action of $G(N)$ chain maps.

For instance,
\[
 d_F([a]\otimes m)=-[a]\otimes d_Mm+[\,]\otimes am.
\]
The minus sign makes the two resulting copies of $[\,]\otimes a\,d_Mm$ cancel in $d_F^2$.
The coefficient $a$ has degree zero, while its bar suspension has degree $-1$.

\begin{proposition}[Twisting adjunction]\label{cbc:modules:adjunction}
There is an adjunction $G\dashv F$.
Its counit and unit are
\begin{equation}\label{eq:cbc:modules:comparisons}
 \epsilon_M(a\otimes c\otimes m)=a\epsilon_C(c)m,
 \qquad
 \eta_N(n)=\sum n_{(-1)}\otimes1\otimes n_{(0)}.
\end{equation}
\end{proposition}
\begin{proof}
A degree-zero graded $A$-linear map $G(N)\to M$ is determined by a weight-preserving map $f:N\to M$.
Its value on $a\otimes n$ is $af(n)$.
It is a chain map precisely when
\begin{equation}\label{eq:cbc:modules:hom-condition}
 d_Mf(n)-f(d_Nn)=-\sum\tau(n_{(-1)})f(n_{(0)}).
\end{equation}

A graded $C$-colinear map $N\to F(M)$ is determined by its composite with $\epsilon_C\otimes\id_M$.
For the resulting map $f$, its value is
\[
 n\longmapsto\sum n_{(-1)}\otimes f(n_{(0)}).
\]
Coassociativity and~\eqref{eq:cbc:modules:twists} show that this is a chain map exactly when~\eqref{eq:cbc:modules:hom-condition} holds.
The free and cofree extension formulas are inverse and natural, proving the adjunction.

Applying them to identity maps gives~\eqref{eq:cbc:modules:comparisons}.
The counit is $A$-linear, and the unit is $C$-colinear.
In the two composites
\[
 \epsilon_{G(N)}G(\eta_N),\qquad F(\epsilon_M)\eta_{F(M)},
\]
the intervening coalgebra factor is removed by $\epsilon_C$.
The coalgebra and comodule counit identities reduce these composites to the respective identity maps.
\end{proof}

\section{The two bar contractions}
The composite $G F(M)$ has underlying tensor $A\widehat\otimes C\widehat\otimes M$.
Write its elements as $a_0[a_1|\cdots|a_r]m$ and put $\bar a=a-\epsilon_A(a)1$.
Its differential consists of the bar differential, the signed internal differential of $M$, and the two endpoint terms
\[
 -a_0a_1[a_2|\cdots|a_r]m,
 \qquad
 (-1)^{r-1}a_0[a_1|\cdots|a_{r-1}](a_rm).
\]
Endpoint terms are absent when $r=0$.

Define a degree-$-1$ map
\begin{equation}\label{eq:cbc:modules:left-contraction}
 h_L\bigl(a_0[a_1|\cdots|a_r]m\bigr)
   =-1[\bar a_0|a_1|\cdots|a_r]m.
\end{equation}
Let $i_M(m)=1[\,]m$.
Then
\[
 d h_L+h_Ld=\id-i_M\epsilon_M.
\]

Indeed, the new first endpoint returns the reduced part of the original leading coefficient.
The first endpoint followed by $h_L$ supplies its scalar part when $r>0$.
All other adjacent multiplications cancel in pairs, including the last module action.
At $r=0$ the remaining term is $a_0[\,]m-1[\,](a_0m)$.

The internal differential of $M$ cancels because insertion changes the bar-length sign.
This proves the identity in every degree.
The contraction is $k$-linear and preserves weight, though $i_M$ need not be $A$-linear.

The second contraction is on
\[
 Q=C\widehat\otimes_\tau A.
\]
Its last endpoint is $(-1)^{r-1}[a_1|\cdots|a_{r-1}]a_ra$.
Define
\begin{equation}\label{eq:cbc:modules:right-contraction}
 h_R\bigl([a_1|\cdots|a_r]a\bigr)
   =(-1)^r[a_1|\cdots|a_r|\bar a]1.
\end{equation}
With $j(1)=[\,]1$ and $p(c\otimes a)=\epsilon_C(c)\epsilon_A(a)$, one has
\[
 d h_R+h_Rd=\id-jp.
\]

The new last endpoint gives the identity term, its scalar remainder gives $jp$, and the preceding multiplications cancel.
This contraction also preserves weight.
Consequently $Q_0=k$ and every positive-weight component of $Q$ is contractible.

\begin{theorem}[Koszul module equivalence in positive weights]\label{cbc:modules:equivalence}
The functors $F,G$ preserve weightwise quasi-isomorphisms and induce inverse equivalences
\[
 D_{\mathrm{wt}}(A)\simeq D_{\mathrm{wt}}(C).
\]
Their unit and counit are~\eqref{eq:cbc:modules:comparisons}.
\end{theorem}
\begin{proof}
The left contraction proves that $\epsilon_M$ is a quasi-isomorphism in each weight.
For the unit, the underlying tensor of $FG(N)$ is $C\widehat\otimes A\widehat\otimes N$.
At total weight $w$, filter it increasingly by the weight $v$ of its final $N$-factor.
There are only the values $w_0\leq v\leq w$.

The reduced-coaction term lowers $v$, while every other differential term preserves it.
The associated graded component at $v$ is $Q_{w-v}\otimes_k N_v$ with its tensor differential.
Give $N_w$ the filtration concentrated at its largest value $v=w$.
The unit is filtered: its leading term is $[\,]\otimes1\otimes n$, and its other terms have lower final weight.

For $v<w$, the right contraction makes the target component contractible even after tensoring with $N_v$.
For $v=w$, the associated graded unit is the identity of $N_w$.
The cone therefore has a finite filtration with acyclic successive quotients.
It is acyclic, so $\eta_N$ is a quasi-isomorphism in weight $w$.

To check preservation of weak equivalences by $F$, filter by the final module weight decreasingly.
The endpoint action increases that weight and disappears in the associated graded.
For $G$, use the increasing filtration by final comodule weight, since its reduced coaction lowers that weight.
The associated graded maps tensor the given quasi-isomorphisms with $C_u$ or $A_u$ over the field $k$.
Such tensors preserve quasi-isomorphisms, even for unbounded complexes of vector spaces.
Both filtrations are finite at each total weight, so the functors preserve the specified weak equivalences.

They consequently descend to the localizations.
There their unit and counit are isomorphisms, and the triangle identities prove that the descended functors are inverse.
\end{proof}

The regular module and the augmentation module have different images.
The right contraction identifies $F(A)$ with the trivial comodule $k$ through the coaugmentation $k\to F(A)$.
In contrast, $F(k)=C$ with its regular coaction, since the reduced algebra acts trivially on $k$.
Thus the equivalence does not identify the regular representation with the augmentation representation.

\section{Other completions and two-sided comparisons}
The lower weight bound is part of the theorem.
It bounds both the bar length and the successive reduced coactions in a fixed component.
Without that bound the displayed finite-filtration argument does not apply.

Derived $I$-adic completion, where $I=\ker(A\to k)$, uses the different comparison
\[
 M\longrightarrow\varprojlim_p\bigl((A/I^p)\otimes_A^{\mathbb L}M\bigr).
\]
Identifying it with the weight product requires a comparison of the filtrations and their derived limits.
A coderived category has its own class of coacyclic complexes.
It is not the localization specified here by definition.

For a right module $R$ and a left module $M$ with lower weight bounds, the normalized two-sided bar has terms
\[
 R\widehat\otimes(s\bar A)^{\widehat\otimes r}\widehat\otimes M,
 \qquad \bar A=\prod_{u>0}A_u.
\]
Only finitely many lengths and weight decompositions occur at any total weight.
It computes their relative derived tensor product in these weighted categories.
Indeed, the preceding left contraction resolves $M$ by induced $A$-modules.
Their coefficient complexes are complexes of vector spaces, and tensoring with an acyclic right module gives acyclic associated graded components.
The finite weightwise filtration therefore justifies the derived tensor computation.

A geometric identification with a derived cotensor product needs a further two-sided twisting comparison.
That comparison must preserve the chosen completion and the collision maps.
The one-sided algebraic equivalence alone does not supply those operations.

\section{Reconstruction from an endomorphism algebra}
We now work in presentable categories, separately from the lower-bounded weighted categories above.
Let $\mathcal V$ be a presentable symmetric monoidal $\infty$-category whose tensor product preserves colimits separately.
Let $\mathcal C$ be a presentable left $\mathcal V$-module category, with the action preserving colimits separately, and choose $b\in\mathcal C$.
The functor $v\mapsto v\otimes b$ has a right adjoint
\[
 H_0(x)=\underline{\operatorname{Hom}}_{\mathcal C}(b,x)\in\mathcal V.
\]
This is an internal mapping object, not an external complex chosen without a coefficient category.

Composition makes
\[
 E=H_0(b)
\]
an associative unital algebra in $\mathcal V$.
Precomposition makes $H_0(x)$ a right $E$-module: $f\cdot e=f\circ e$.
Write $H:\mathcal C\to\operatorname{RMod}_E(\mathcal V)$ for this lift.
Evaluation makes $b$ a left $E$-module object of $\mathcal C$.
The opposite algebra is required if right modules are written as left modules.

For a right $E$-module $M$, define
\[
 P(M)=M\otimes_E b
   =\left|[r]\longmapsto M\otimes E^{\otimes r}\otimes b\right|.
\]
The bars denote geometric realization of the two-sided simplicial bar.
Its faces use the two actions and multiplication, and its degeneracies insert the unit of $E$.
The tensor--internal-Hom adjunction in every simplicial degree identifies compatible maps from this realization with right $E$-module maps into $H(x)$.
Thus $P\dashv H$.

\begin{proposition}[Relative Morita criterion]\label{cbc:modules:relative-morita}
Suppose $H_0$ is conservative and preserves geometric realizations.
Suppose also that the natural map
\[
 \theta_v:v\otimes E\longrightarrow H_0(v\otimes b)
\]
is an equivalence for every $v\in\mathcal V$.
Then $P,H$ are inverse equivalences
\[
 \operatorname{RMod}_E(\mathcal V)\simeq\mathcal C.
\]
\end{proposition}
\begin{proof}
On the free right module $v\otimes E$, the unit $M\to HP(M)$ is $\theta_v$, with its right $E$-action.
It is an equivalence by hypothesis.
Every right module is the realization of its simplicial free-module resolution.
After forgetting its action, the augmented resolution is split by insertion of the algebra unit.

The functor $P$ preserves realizations because it is a left adjoint.
The functor $H$ preserves them because $H_0$ does and the forgetful functor from right modules creates them.
Realizing the free-module unit equivalences proves the unit equivalence for every module.

For $x\in\mathcal C$, the triangle identity says that $H$ of the counit $PH(x)\to x$ is inverse to the unit at $H(x)$.
It is therefore an equivalence.
Conservativity of $H_0$, and hence of $H$, proves that the counit itself is an equivalence.
\end{proof}

This is the free-module form of relative Morita reconstruction~\cite[Proposition~4.8.5.8]{LurieHA}.
It records the projection comparison $\theta_v$ as well as generation.
No stability of a category of factorization objects is inferred merely from stability of its ambient coefficient category.

For $\mathcal V=D(k)$, let $\mathcal C$ be stable and $k$-linear, and let $b$ be a compact generator.
Here generation means that $\RHom_{\mathcal C}(b,x)=0$ implies $x=0$.
Then $H_0$ is conservative.
It is exact, and compactness makes it commute with filtered colimits and therefore with arbitrary coproducts.
Exactness and these colimit properties imply preservation of geometric realizations, by their skeletal filtrations.

Both sides of $\theta_v$ are exact colimit-preserving functors of $v$.
The map is the identity for $v=k$, and shifts and colimits of $k$ generate $D(k)$.
Hence $\theta_v$ is an equivalence for every $v$.
The proposition recovers compact-generator Morita equivalence~\cite[Theorem~7.1.2.1]{LurieHA}, with the dg version in~\cite[Theorem~3.8(c)]{Keller2006DG}.

For example, in $D(k)$ take $b=k^n$ with $n>0$.
Then $E=\operatorname{Mat}_n(k)$, $H(x)=\operatorname{Hom}_k(k^n,x)$ has the right matrix action, and $P(M)=M\otimes_E k^n$.
Evaluation and the matrix units give the inverse maps of classical Morita equivalence.

Compactness alone does not suffice.
In $D(k)\times D(k)$, the compact object $b=(k,0)$ has endomorphism algebra $k$.
Its Hom functor sends $(M,N)$ to $M$ and annihilates the nonzero object $(0,k)$.
The missing condition is generation.

\section{Factorization descent}
For a factorization category, the local reconstruction maps must also respect the maps used to assemble the category.
Choose a chart diagram with coefficient categories $\mathcal V_U$, module categories $\mathcal C_U$, and compatible objects $b_U$.
Put $E_U=\underline{\operatorname{Hom}}_{\mathcal C_U}(b_U,b_U)$.
Assume these algebras are the chart values of an algebra $A_b$ in the chosen category $\Fact_D(X)$.
Use specified descent presentations
\[
 \mathcal C\simeq\varprojlim_U\mathcal C_U,
 \qquad
 \operatorname{RMod}_{A_b}(\Fact_D(X))
   \simeq\varprojlim_U\operatorname{RMod}_{E_U}(\mathcal V_U).
\]
The limits include the full descent coherence data.

\begin{theorem}[Factorization Morita theorem]\label{cbc:modules:factorization-morita}
Suppose the relative Morita criterion holds on every chart.
Suppose the local internal-Hom, evaluation and tensor comparison maps are compatible with the chart transitions and all factorization operations, including their higher coherences.
Then the descent presentations above give an equivalence of factorization categories
\[
 \mathcal C\simeq\operatorname{RMod}_{A_b}(\Fact_D(X)).
\]
\end{theorem}
\begin{proof}
On each chart, the proposition gives mutually inverse functors $P_U,H_U$.
The stated compatibilities make them natural transformations between the two descent diagrams.
Their units and counits are equivalences on every chart and commute with the same transition and factorization maps.
Taking limits gives inverse equivalences of the displayed categories.
The preserved operation maps and their coherences give the factorization structure on this equivalence.
\end{proof}

Thus a compact generator supplies the local theorem in stable $k$-linear charts, while factorization descent supplies the assembly.
For a particular boundary object, both generation and the operation-compatible comparisons must be proved.
The tensor counit and the one-sided module equivalence do not establish either condition.

\chapter{Derived center and bulk action}
A central action must commute with multiplication together with all its homotopies.
The derived center represents these actions.
Its two compatible compositions give the little-disks structure, while a cyclic orientation supplies the additional circle action.
To recognize a physical bulk as this center, one must construct its central action and prove that the representing map is an equivalence.

\section{The action represented by the center}

Let $\mathcal V$ be a presentable stable symmetric monoidal $\infty$-category whose tensor product preserves small colimits separately.
All tensors in this section use this specified product.
Let $A$ be a unital associative algebra in $\mathcal V$, and write
\[
 A^e=A\otimes A^{\op}.
\]
Its left modules are $A$-bimodules.

The internal derived mapping object
\begin{equation}\label{eq:cbc:center:internal}
 Z_{\ord}(A)=
 \underline{\operatorname{RHom}}_{A^e}(A,A)\in\mathcal V
\end{equation}
is characterized by
\begin{equation}\label{eq:cbc:center:representing}
 \operatorname{Map}_{\mathcal V}(v,Z_{\ord}(A))
 \simeq
 \operatorname{Map}_{A\text{-}A}(v\otimes A,A).
\end{equation}
Here $v$ ranges over $\mathcal V$, and the two $A$-actions on $v\otimes A$ act on its $A$-factor.
Presentability and the colimit-preserving action give this internal Hom by the adjoint functor theorem.
Composition makes it an associative algebra.

An $E_2$-algebra is an algebra over the operad of little two-disks.
Write $\operatorname{Alg}_{E_1}(\mathcal V)$ for the category of unital associative algebras in $\mathcal V$.
Dunn additivity identifies an $E_2$-algebra with an associative algebra object in $\operatorname{Alg}_{E_1}(\mathcal V)$.
The latter category has the tensor product induced from $\mathcal V$.

For an $E_2$-algebra $R$, a coherent central action on $A$ means a left $R$-module structure on $A$ in $\operatorname{Alg}_{E_1}(\mathcal V)$.
Thus the action
\[
 \rho:R\otimes A\longrightarrow A
\]
preserves the inner algebra multiplications, while its module identities use the outer multiplication on $R$.
Units, associativity and interchange include their full higher coherences.
Write $\operatorname{Act}^{\mathrm{cen}}_R(A)$ for the space of these structures on the fixed algebra $A$.

\begin{theorem}[Higher Deligne center]\label{cbc:center:deligne}
The object~\eqref{eq:cbc:center:internal} has a canonical $E_2$-algebra structure and a coherent central action on $A$.
For every $E_2$-algebra $R$, this action induces a natural equivalence
\begin{equation}\label{eq:cbc:center:universal}
 \operatorname{Map}_{\operatorname{Alg}_{E_2}(\mathcal V)}
       (R,Z_{\ord}(A))
 \simeq \operatorname{Act}^{\mathrm{cen}}_R(A).
\end{equation}
Its underlying associative multiplication is composition.
\end{theorem}
\begin{proof}
Apply the center theorem for associative algebras in a presentably symmetric monoidal category
\cite[Corollary~5.3.1.15]{LurieHA}.
It gives a final algebra acting on $A$ in $\operatorname{Alg}_{E_1}(\mathcal V)$.
The acting algebra is associative in that category, hence is an $E_2$-algebra by Dunn additivity
\cite[Theorem~5.1.2.2]{LurieHA}.

The underlying object of this center represents morphisms from the regular $A$-module to itself in the operadic module category
\cite[Theorem~5.3.1.30]{LurieHA}.
For the associative operad these modules are $A$-bimodules
\cite[Theorem~4.4.1.28]{LurieHA}.
Its representing property is therefore exactly~\eqref{eq:cbc:center:representing}.
Uniqueness of representing objects identifies it with~\eqref{eq:cbc:center:internal}, including the evaluation map.

Restriction along a map of acting algebras pulls back the action on $A$.
Finality says that each such action has a contractible space of compatible maps to the center.
This identifies its space of actions with the left side of~\eqref{eq:cbc:center:universal}, naturally in $R$.
The associative multiplication on the represented endomorphism object is composition.
\end{proof}

The two multiplications can also be seen in the monoidal category of $A$-bimodules.
Its tensor is the relative tensor product $\otimes_A$, and its unit is the diagonal bimodule $A$.
Endomorphisms of this unit can be composed or tensored over $A$.
The interchange coherences give the two associative structures.
Their $E_2$-structure is homotopy coherent, so a cochain model need not have a strictly commutative product.

The word \emph{canonical} refers to the center with its universal action.
It does not assert uniqueness of every possible $E_2$-structure on an unmarked complex.
The tensor in~\eqref{eq:cbc:center:representing} is equally essential: changing it changes the action problem.

\section{Hochschild cochains and the first operation homotopy}

For an ordinary unital algebra $A$ over a field $k$, the two-sided bar resolution has
\[
 P_n=A\otimes_k A^{\otimes n}\otimes_k A,\qquad n\geq0,
\]
in cohomological degree $-n$.
Its boundary is the alternating sum of adjacent multiplications, starting with sign $+1$, and its augmentation is multiplication $P_0\to A$.
The map
\[
 a_0\otimes\cdots\otimes a_{n+1}
 \longmapsto1\otimes a_0\otimes\cdots\otimes a_{n+1}
\]
contracts the augmented complex as a complex of $k$-vector spaces.
On the augmented term $A$, use $a\mapsto1\otimes a$.
The first face returns the original tensor and every later face cancels with insertion after the preceding face.
Each $P_n$ is free over $A^e$, so this bounded-above resolution computes~\eqref{eq:cbc:center:internal}.

Thus the Hochschild cochain model is
\[
 C^p(A,A)=\operatorname{Hom}_k(A^{\otimes p},A),\qquad p\geq0,
\]
using the identification
$\phi\mapsto(-1)^{p(p+1)/2}\phi(1\otimes-\otimes1)$ from the ordinary Hom complex in degree $p$.
Its differential is
\begin{equation}\label{eq:cbc:center:hoch-differential}
\begin{aligned}
 (\delta f)(a_1,\ldots,a_{p+1})
 ={}&a_1f(a_2,\ldots,a_{p+1})\\
 &+\sum_{i=1}^{p}(-1)^i
     f(a_1,\ldots,a_i a_{i+1},\ldots,a_{p+1})\\
 &+(-1)^{p+1}f(a_1,\ldots,a_p)a_{p+1}.
\end{aligned}
\end{equation}
Its cohomology is denoted $HH^p(A)$.
Set $C^r(A,A)=0$ for $r<0$.
In degree zero, $\delta a(b)=ba-ab$, so $HH^0(A)$ is the ordinary center of $A$.

For $f\in C^p(A,A)$ and $g\in C^q(A,A)$, the cup product is
\[
 (f\smile g)(a_1,\ldots,a_{p+q})
   =f(a_1,\ldots,a_p)g(a_{p+1},\ldots,a_{p+q}).
\]
It represents the product of the derived center.
The bar diagonal splits the inputs at $p$.
Its Hom-tensor sign $(-1)^{pq}$ cancels the ratio of the cochain-identification signs, giving the displayed cup formula.

For $0\leq i<p$, let $f\circ_i g$ replace the block of inputs $a_{i+1},\ldots,a_{i+q}$ by $g(a_{i+1},\ldots,a_{i+q})$ and leave the other inputs in order.
The signed insertion operation~\cite[\S1]{McClureSmith} is
\begin{equation}\label{eq:cbc:center:insertion}
 f\circ g=\sum_{i=0}^{p-1}(-1)^{i(q-1)}f\circ_i g
 \in C^{p+q-1}(A,A).
\end{equation}
For $q=0$, the value of $g$ is inserted as a constant, and an empty sum is zero.

The bracket
\[
 [f,g]=f\circ g-(-1)^{(p-1)(q-1)}g\circ f
\]
induces the degree-$-1$ Gerstenhaber bracket on Hochschild cohomology.
Together with the higher operations, these give a cochain realization of the little-disks action
\cite[\S7, Corollary~1]{KontsevichSoibelmanDeligne}.

The first homotopy is visible without the full operad.
If $D,E:A\to A$ are derivations, then
\[
 \delta(D\circ E)(a,b)
   =-D(a)E(b)-E(a)D(b).
\]

Indeed, expand $D(E(ab))$ by the two Leibniz rules in~\eqref{eq:cbc:center:hoch-differential}.
The remaining terms are the negative of $D\smile E+E\smile D$.
Hence the cup products of these degree-one cocycles anticommute up to the displayed cochain homotopy.

This identity illustrates one operation of the $E_2$-structure.
The full action requires all arities and their coherences, as in the theorem.

\section{The polynomial example}

Let $k$ have characteristic zero and take $A=k[x]$.
Write $x_\ell=x\otimes1$ and $x_r=1\otimes x$ in $A^e=k[x_\ell,x_r]$.
The sequence
\begin{equation}\label{eq:cbc:center:polynomial-resolution}
 0\longrightarrow A^e e
   \xrightarrow{\ e\mapsto x_\ell-x_r\ }
   A^e\longrightarrow A\longrightarrow0
\end{equation}
is a free bimodule resolution, with $e$ in degree $-1$.
The last map sets $x_\ell=x_r=x$.
Its kernel is generated by $x_\ell-x_r$, by polynomial division, and multiplication by that element is injective.

Applying $\operatorname{Hom}_{A^e}(-,A)$ makes the differential zero.
The chain map to the bar resolution sends $e$ to $1\otimes x\otimes1$.
Consequently the derivation $D=d/dx$, which takes $x$ to $1$, represents the degree-one generator.
There are no higher cohomology groups, and therefore
\[
 HH^*(k[x])\cong k[x]\otimes_k\Lambda(\xi),
 \qquad |\xi|=1,\qquad \xi=[D].
\]

The insertion formula gives
\[
 [a(x)D,b(x)]=a(x)b'(x),\qquad
 [a(x)D,b(x)D]=(a b'-b a')D.
\]
In particular, $[\xi,x]=1$.
The ordinary commutative algebra $HH^0(A)=A$ alone does not contain this operation.

The relation $\xi^2=0$ has an explicit representative:
\[
 \delta\!\left(-\frac12D^2\right)(a,b)=D(a)D(b).
\]
Both endpoint terms of the Hochschild differential enter this identity.
For $a=x$ and $b=x^2$, they give
$\delta(D^2)(x,x^2)=-4x$, which cancels $2D(x)D(x^2)=4x$.
The cup square is thus zero in cohomology through a specified homotopy.

\section{The bulk map and a change of sector}

Let $B$ be an $E_2$-algebra in $\mathcal V$, and equip a boundary algebra $A$ with a coherent central $B$-action.
The universal property gives its classifying map
\begin{equation}\label{eq:cbc:center:bulk-map}
 \beta:B\longrightarrow Z_{\ord}(A).
\end{equation}
The evaluation of $\beta(b)$ on $A$ is the prescribed action of $b$.
The contractibility in~\eqref{eq:cbc:center:universal} specifies the map together with all action coherences.

For a field theory, this construction starts after the bulk and boundary operations have been realized in the common coefficient category with this central-action structure.
In a mixed holomorphic--topological theory, that realization must also retain the collision coefficients and their supported maps.
Constructing those operations is part of constructing the physical bulk-to-center map.
A pointwise action on states supplies only the corresponding pointwise datum.

Centrality does not require~\eqref{eq:cbc:center:bulk-map} to be an equivalence.
For example, the scalar $E_2$-algebra $B=k$ acts centrally on $k[x]$.
Its map to the derived center includes the constants and misses both $x$ and $\xi$.
This is an example of a central action, with no assertion that it describes a complete field theory.

The failure also occurs in theories satisfying surface sewing.
A finite-dimensional Frobenius algebra is a unital associative algebra over a field, with a linear functional $\varepsilon$ for which $(a,b)\mapsto\varepsilon(ab)$ is nondegenerate.
Its coproduct $\Delta$ is adjoint to multiplication.
The Frobenius algebra is symmetric when $\varepsilon(ab)=\varepsilon(ba)$.

An oriented two-dimensional open--closed topological quantum field theory assigns vector spaces to intervals and circles, and linear maps to bordisms with corners.
These assignments respect gluing and disjoint union.
Its data are a symmetric Frobenius algebra $A$ for the interval and a commutative Frobenius algebra $B$ for the circle.
They include a unital central algebra map $\iota:B\to A$.
Its adjoint $\iota^*:A\to B$ must obey
\begin{equation}\label{eq:cbc:center:cardy}
 \iota\iota^*=\mu_A\circ\sigma\circ\Delta_A,
\end{equation}
where $\mu_A$ is multiplication and $\sigma$ interchanges the two factors.
These conditions are sufficient by the open--closed classification
\cite[Definition~2.3 and Corollary~4.3]{LaudaPfeiffer2008}.

\begin{example}[Two closed theories with one boundary algebra]\label{cbc:center:two-bulks}
Take $A=k$ with $\varepsilon_A(1)=1$.
One choice of closed data is $B_1=k$, with the same trace and with $\iota=\iota^*=\id$.
A second choice is
\[
 B_2=ke_1\oplus ke_2,\qquad
 e_i e_j=\delta_{ij}e_i,\qquad
 1_{B_2}=e_1+e_2,\qquad \varepsilon_{B_2}(e_i)=1.
\]
Set $\iota(e_1)=1$, $\iota(e_2)=0$ and $\iota^*(1)=e_1$.
Both Frobenius pairings are nondegenerate.
The closed coproduct is $\Delta_{B_2}(e_i)=e_i\otimes e_i$, while $\Delta_A(1)=1\otimes1$.

The map $\iota$ is unital and multiplicative, and its image is central.
For $b=\lambda e_1+\nu e_2$ and $a\in k$,
\[
 \varepsilon_{B_2}(b\,\iota^*(a))
   =\lambda a=\varepsilon_A(\iota(b)a),
\]
so the stated maps are adjoint.
Both sides of~\eqref{eq:cbc:center:cardy} are $\id_k$.
Thus both choices define oriented open--closed theories.

They have the same open Frobenius algebra and the same image acting on it.
Their values on a circle have dimensions one and two, so the theories are not equivalent.
Since $Z_{\ord}(k)=k$, the second bulk-to-center map is the projection with kernel $ke_2$.
A single boundary algebra therefore does not determine an arbitrary open--closed surface theory.
\end{example}

Now let $L:\mathcal V\to\mathcal W$ be a specified strong symmetric monoidal functor, where $\mathcal W$ satisfies the same presentability assumptions.
Transporting the universal action gives a central action of $LZ_{\ord}^{\mathcal V}(A)$ on $LA$.
Its classifying map is
\begin{equation}\label{eq:cbc:center:sector-comparison}
 c_{L,A}:LZ_{\ord}^{\mathcal V}(A)
   \longrightarrow Z_{\ord}^{\mathcal W}(LA).
\end{equation}

The transported bulk action is classified by
\[
 \beta_L=c_{L,A}\circ L\beta:
 LB\longrightarrow Z_{\ord}^{\mathcal W}(LA).
\]
Both composites have the same evaluation, so this factorization follows from the universal property.

No equivalence of $c_{L,A}$ is assumed.
If $L$ specifies a protected sector, the boundary is \emph{open--closed complete in that sector} when $\beta_L$ is an equivalence.
This definition requires the sector functor and the central action, rather than only an equality of graded dimensions.

\section{Cyclic orientation and the circle action}

The framed little-disks operad is $E_2\rtimes SO(2)$.
Its circle rotates the disk configurations, so its action must be compatible with the $E_2$-operations.
We now state the extra datum that produces this compatibility.

Let $k$ be a field and let $\mathcal C$ be a dualizable presentable $k$-linear category.
Dualizability is taken in the symmetric monoidal category of presentable $k$-linear categories and colimit-preserving $k$-linear functors.
Write
\[
 R=HH^\bullet(\mathcal C)
   =\underline{\operatorname{RHom}}_{\operatorname{Fun}^L_k
       (\mathcal C,\mathcal C)}(\id_{\mathcal C},\id_{\mathcal C}).
\]
Composition of endofunctors makes the functor category monoidal with unit $\id_{\mathcal C}$.
The two unit-endomorphism multiplications give $R$ its $E_2$-structure
\cite[Proposition~2.1.5]{BravRozenblyum}.

For $\mathcal C=\operatorname{LMod}_A(D(k))$, this recovers $Z_{\ord}(A)$.
Indeed, an $A$-bimodule $P$ gives the colimit-preserving functor $P\otimes_A-$.
A $k$-linear colimit-preserving functor is determined by its value on the free module $A$ and its right $A$-action.
The assertion holds first on free modules $v\otimes A$ and then on every module by its simplicial free resolution.
The same argument identifies natural transformations with bimodule maps.
Under this equivalence the diagonal bimodule corresponds to the identity functor, and relative tensor corresponds to composition.

The same free-resolution argument works for $\mathcal V$-linear colimit-preserving functors on $\operatorname{LMod}_A(\mathcal V)$ under the hypotheses of the first section.

Let
\[
 T=HH_\bullet(\mathcal C)
\]
be the categorical trace of $\id_{\mathcal C}$, an object of $D(k)$ with its canonical circle action.
Explicitly, the trace is the composite of coevaluation, interchange and evaluation for the dual pair $\mathcal C,\mathcal C^\vee$.
For $\operatorname{LMod}_A(D(k))$, it is the Hochschild chain object $A\otimes_{A^e}^{\mathbb L}A$.
Its homological degree $n$ is placed in cohomological degree $-n$.

Put $H_R=\int_{S^1}R$, the factorization homology of the $E_2$-algebra $R$ around a circle.
It is an associative algebra using the remaining $E_1$-direction, and it carries the circle action.
Noncommutative calculus makes $T$ a circle-equivariant $H_R$-module
\cite[Theorem~2.3.7]{BravRozenblyum}.
An embedded interval gives $R\to H_R$.
Restricting the action along this map gives the cap action $R\otimes T\to T$.

For an integer $d$, a noncommutative orientation of dimension $d$ is a class
\[
 \theta:k\longrightarrow T^{hS^1}[-d]
\]
whose cap map is an equivalence
\begin{equation}\label{eq:cbc:center:cap-duality}
 D_\theta:R\xrightarrow{\ \sim\ }T[-d].
\end{equation}
Here $hS^1$ denotes homotopy fixed points, so the datum includes coherent circle invariance.

There is also a version using a trace.
Put $T^\vee=\underline{\operatorname{RHom}}_k(T,k)$.
A noncommutative coorientation is a circle-invariant trace
\[
 \tau:T_{hS^1}\longrightarrow k[-d]
\]
for which the pairing with the cap action gives an equivalence
\[
 D_\tau:R\xrightarrow{\ \sim\ }T^\vee[-d].
\]
The subscript $hS^1$ denotes homotopy orbits.
In particular, $\tau$ includes more data than a functional on $HH_0(\mathcal C)$.

These conditions are supplied by Calabi--Yau structures as follows.
Call $\mathcal C$ smooth, respectively proper, when its evaluation has a left, respectively right, adjoint in the category of presentable $k$-linear categories.
Put $I=\id_{\mathcal C}$.
Those adjoints evaluated on $k$ define the inverse dualizing functor $I^!$, respectively the Serre functor $S$, in $\operatorname{Fun}^L_k(\mathcal C,\mathcal C)$.
The trace identifications are
\[
 \begin{aligned}
 T&\simeq\underline{\operatorname{RHom}}(I^!,I)
   &&(\mathcal C\text{ smooth}),\\
 T^\vee&\simeq\underline{\operatorname{RHom}}(I,S)
   &&(\mathcal C\text{ proper}).
 \end{aligned}
\]

A smooth Calabi--Yau structure of dimension $d$ is a circle-invariant class $\theta$ as above whose underlying map $I^![d]\to I$ is an equivalence.
Precomposition with that map makes $D_\theta$ an equivalence.
A proper Calabi--Yau structure is a circle-invariant trace $\tau$ whose underlying map $I\to S[-d]$ is an equivalence.
Postcomposition makes $D_\tau$ an equivalence.
These are the absolute smooth and proper constructions of Brav--Rozenblyum
\cite[\S\S4.1.9 and 4.3.5]{BravRozenblyum}.

\begin{theorem}[Cyclic refinement]\label{cbc:center:cyclic}
An orientation or coorientation as above induces a framed $E_2$-algebra structure on $HH^\bullet(\mathcal C)$ refining its existing $E_2$-structure.
Thus a dualizable linear category with a smooth or proper Calabi--Yau structure has the stated cyclic refinement.
\end{theorem}
\begin{proof}
The noncommutative calculus makes $T$ a circle-equivariant $H_R$-module.
The oriented case supplies the shifted module $T[-d]$ and its circle-invariant vector $\theta$.
Its cap map from $R$ is an equivalence by~\eqref{eq:cbc:center:cap-duality}.
The framed-algebra criterion of Brav--Rozenblyum identifies precisely these data with a framed $E_2$-algebra
\cite[Theorem~3.3.9]{BravRozenblyum}.
This proves the absolute case of their cyclic Deligne theorem
\cite[Theorem~4.1.7]{BravRozenblyum}.

For a coorientation, use the dual calculus on $T^\vee[-d]$ and the invariant vector corresponding to $\tau$.
The assumed cap equivalence is the same nondegeneracy condition in that module.
The cooriented cyclic Deligne theorem supplies the framed refinement
\cite[Theorem~4.3.4]{BravRozenblyum}.
The original $E_2$-operations are retained in both constructions.
\end{proof}

\section{The circle operator for the polynomial algebra}

Return to $A=k[x]$ in characteristic zero.
Put $\bar A=A/k1$ and $C_n(A)=A\otimes_k\bar A^{\otimes n}$, placed in degree $-n$.
In this normalized Hochschild chain complex, write
$a_0[a_1|\cdots|a_n]$ for $a_0\otimes\bar a_1\otimes\cdots\otimes\bar a_n$.
The Hochschild boundary is the alternating sum of adjacent products followed by the cyclic last product.
In particular, $b(a_0[a_1])=a_0a_1-a_1a_0$.

Fix Connes' operator by the convention
\[
 B(a_0)=1[a_0].
\]
On a chain of length one it is
\[
 B(a_0[a_1])=1[a_0|a_1]-1[a_1|a_0].
\]
An inner occurrence of $1$ vanishes in the normalized complex.

The chain $\theta=1[x]$ has both $b\theta=0$ and $B\theta=0$.
The negative cyclic complex is the degreewise product $C_\bullet(A)[[u]]$ with differential $b+uB$ and $|u|=2$.
Thus the same constant series $\theta$ defines a negative cyclic class, or equivalently a coherent circle-invariant class, of dimension one.
Tensoring~\eqref{eq:cbc:center:polynomial-resolution} over $A^e$ with $A$ gives
\[
 HH_0(A)=A,\qquad HH_1(A)=A\,dx,\qquad HH_n(A)=0\quad(n>1).
\]
The class of $a_0[a_1]$ maps to $a_0a_1'(x)\,dx$.
The Leibniz rule makes this map vanish on boundaries, and it sends $1[x]$ to the generator $dx$.

Cap product with $\theta$ sends $f(x)\in HH^0(A)$ to $f(x)\,dx$ and sends $a(x)D\in HH^1(A)$ to $a(x)$.
Both maps are isomorphisms.
The category $\operatorname{LMod}_A(D(k))$ is compactly generated by $A$ and hence is dualizable.
Hence $\theta$ is an orientation in the sense of~\eqref{eq:cbc:center:cap-duality}, and the cyclic refinement applies.

With these cap and Connes conventions, the transported degree-$-1$ circle operator is
\[
 \Delta(f)=0,\qquad \Delta(aD)=a'.
\]
Indeed, $B(a)=1[a]$ represents $a'(x)\,dx$, while the next homology group is zero.
For homogeneous classes $u,v$, the bracket convention used above is
\[
 [u,v]=(-1)^{|u|-1}
 \bigl(\Delta(u\smile v)-\Delta(u)\smile v
             -(-1)^{|u|}u\smile\Delta(v)\bigr).
\]
It gives $[aD,f]=af'$ and $[aD,bD]=(ab'-ba')D$, as before.
In particular, the circle operation recovers $[\xi,x]=1$.

The cyclic theorem supplies this framed structure from the orientation data.
It does not construct a mixed collision action or identify a specified physical bulk with the center.
Those comparisons must preserve their coefficient tensor, support, completion and operation coherences.

\chapter{Cyclic closure and modular sewing}\label{chap:cbc:cyclic-sewing}
Closing an ordered interval requires a pairing that can be used in both directions.
The pairing joins two ends, and its inverse cuts an internal edge.
Independence of the cuts requires cyclic symmetry of every higher multiplication.
We first construct the resulting surface theory over a field, then state the additional maps needed for chiral sewing.

\section{The cyclic pairing}

Let $k$ be a field of characteristic zero.
Use cohomological grading and the shift $(V[r])^i=V^{i+r}$.
For a complex $V$, write $V^\vee=\operatorname{RHom}_k(V,k)$.
A complex is perfect over $k$ when it is quasi-isomorphic to a bounded complex of finite-dimensional vector spaces.
An algebra is called proper when it is perfect over $k$, and smooth when its diagonal bimodule is perfect over its derived enveloping algebra $A^e$.
Fix an integer $d$.

\begin{definition}[Finite Calabi--Yau ordered algebra]\label{cbc:cyclic:finite-model}
A finite Calabi--Yau ordered algebra of dimension $d$ consists of a finite-dimensional, strictly unital $\Ainf$-algebra $A$ and a pairing
\[
 B:A\otimes_k A\longrightarrow k[-d].
\]
The pairing is closed and graded-symmetric, and induces an isomorphism of complexes $A\cong A^\vee[-d]$.
Its higher products $m_n:A^{\otimes n}\to A$ have degree $2-n$, with $m_1=d_A$ and $m_0=0$.
They satisfy the $\Ainf$ identities and, for homogeneous inputs, the cyclic identity
\begin{equation}\label{eq:cbc:cyclic:identity}
 B(m_n(a_0,\ldots,a_{n-1}),a_n)
 =(-1)^{n+|a_0|\sum_{i=1}^n|a_i|}
 B(m_n(a_1,\ldots,a_n),a_0).
\end{equation}
The unit $1\in A^0$ is closed, is a two-sided unit for $m_2$, and makes $m_n$ vanish whenever it is an input and $n\ne2$.
\end{definition}

Here graded symmetry means $B(a,b)=(-1)^{|a||b|}B(b,a)$.
Closedness means $B(d_Aa,b)+(-1)^{|a|}B(a,d_Ab)=0$.
For $n=1$, this is precisely~\eqref{eq:cbc:cyclic:identity} after graded symmetry.
For $n=2$, the identity and the unit give
\[
 \tau(a)=B(1,a),\qquad B(a,b)=\tau(m_2(a,b)).
\]
Thus the pairing determines a trace, while the identities for all $n$ retain the compatibility with higher products.

The isomorphism to the shifted dual supplies an inverse pairing.
It is a degree-$d$ coevaluation $k\to A\otimes A$ whose two contractions with $B$ are the identity, with the graded tensor signs.
Finite dimensionality makes every contraction along a fixed finite graph a finite sum.
There is no sum over all genera in this assertion.

For a perfect derived algebra, a proper Calabi--Yau structure uses a coherent circle-invariant functional
\[
 CC_*(A)_{hS^1}\longrightarrow k[-d].
\]
Its underlying bimodule map $A\to A^\vee[-d]$ must be an equivalence.
Here $CC_*(A)$ is the Hochschild chain object, and $hS^1$ denotes homotopy orbits, as in the preceding chapter.
A functional on $HH_0(A)$ alone does not specify this coherence or a finite cyclic model.
The surface construction below uses the model of Definition~\ref{cbc:cyclic:finite-model}.

\section{Surfaces and the determinant twist}

Fix one boundary label.
An open boundary is a parametrized interval in the boundary of an oriented bordered Riemann surface.
A closed boundary is a parametrized boundary circle.
Each is designated incoming or outgoing, and the remaining boundary is called free.
The objects of the surface category are finite collections of open intervals and closed circles.
Its morphism spaces are moduli spaces of these surfaces, with composition given by gluing matching boundaries.

We use Costello's admissibility convention: every connected component has an incoming closed boundary or a free boundary.
In particular, a surface with only closed boundaries must have an incoming one on every component.
Disks and annuli with only free boundary are represented by points.
The one-input, one-output open disk is an identity, and closed cylinders are replaced by their reparametrization model, including the identity.
These conventions define the surface category used in Costello's Theorem~A \cite[\S3]{CostelloCY}.

The degree of the pairing requires a local coefficient line on each moduli space.
For a surface $\Sigma$ with $o_{\mathrm{out}}$ outgoing open intervals, put
\begin{equation}\label{eq:cbc:cyclic:determinant}
 \mathcal D(\Sigma)
 =\det\bigl(H^0(\Sigma;k)\oplus k^{o_{\mathrm{out}}}\bigr)
       \otimes\det\bigl(H^1(\Sigma;k)\bigr)^{-1},
 \qquad |\mathcal D(\Sigma)|=o_{\mathrm{out}}-\chi(\Sigma).
\end{equation}
The last equality places this line in the stated cohomological degree.
The inverse on the second determinant denotes the dual line.
This is Costello's determinant local system after changing homological to cohomological grading.

Mayer--Vietoris for a gluing identifies the determinant line of the glued surface with the tensor product of those of its two pieces.
The outgoing-open factor in~\eqref{eq:cbc:cyclic:determinant} accounts for the intervals that disappear in the gluing.
These identifications are associative and compatible with disjoint union \cite[\S3, equation defining $\det(\Sigma)$]{CostelloCY}.
For example, a disk with two incoming open intervals gives a line of degree $-1$.
Its $d$th tensor power therefore permits the pairing of degree $-d$.
The disk with two outgoing intervals gives the inverse pairing of degree $d$.

Write $\mathcal{OC}^d$ for the dg symmetric monoidal category of chains on these moduli spaces with coefficients in $\mathcal D^{\otimes d}$.
A singular $r$-chain is placed in cohomological degree $-r$, in addition to the degree of its coefficient line.
Write $\mathcal O^d$ for its full subcategory on purely open objects.
Write $\mathcal C^d$ for the subcategory with purely closed objects and no free boundaries in its morphisms.
Negative powers of $\mathcal D$ mean powers of its dual.

An open topological conformal field theory is a dg symmetric monoidal functor $\Phi:\mathcal O^d\to\operatorname{Ch}(k)$ whose tensor comparison maps are quasi-isomorphisms.
This last condition is called homological splitting.
An open--closed theory is such a functor on $\mathcal{OC}^d$.
A chain in a morphism space acts with its chain degree, its boundary acts by the Hom differential, and sewing acts by composition.
These theories retain the chains of moduli spaces, rather than just the diffeomorphism classes of surfaces.

\section{The open theory and its circle states}

\begin{theorem}[Ribbon-surface realization]\label{cbc:cyclic:surface}
A finite Calabi--Yau ordered algebra $A$ determines an open theory $\Phi_A$ on $\mathcal O^d$, up to quasi-isomorphism, with interval value quasi-isomorphic to $A$.
Let $i:\mathcal O^d\hookrightarrow\mathcal{OC}^d$ be the inclusion.
Its derived extension
\begin{equation}\label{eq:cbc:cyclic:extension}
 \Psi_A=\mathcal{OC}^d\otimes_{\mathcal O^d}^{\mathbb L}\Phi_A
\end{equation}
is homologically split and is the homotopy-universal open--closed extension of $\Phi_A$.
Its restriction to $\mathcal C^d$ has circle value quasi-isomorphic to $CC_*(A)$.
Consequently, the homology of the circle value is $HH_*(A)$.
The theorem applies in particular to a smooth proper algebra supplied with such a finite cyclic model.
\end{theorem}
\begin{proof}
We apply Costello's Theorem~A and describe the maps that enter it.
In the cellular model, disks with one outgoing interval give the higher products, and the two-input and two-output disks give the pairing and its inverse.
The inverse relations are the two contraction identities for $B$.
Rotating a disk gives~\eqref{eq:cbc:cyclic:identity}, with the cell orientation and determinant signs.
The boundary of a disk cell is a sum of two-disk gluings, and its image is the corresponding $\Ainf$ identity.
Thus the algebra defines a split functor on the cellular open category \cite[Theorem~6.2.3 and Lemma~7.3.2]{CostelloCY}.

The cellular category is quasi-isomorphic to $\mathcal O^d$.
Derived extension along this quasi-isomorphism transports the functor to $\Phi_A$ without changing its interval value \cite[Lemmas~6.1.7 and 7.3.4]{CostelloCY}.
Formula~\eqref{eq:cbc:cyclic:extension} is the derived left adjoint to restriction along $i$.
This gives its universal property and identifies its restriction to open objects with $\Phi_A$.

To compute a circle, use the cellular right module generated by annuli with ordered marked intervals on the outer boundary and one parametrized inner circle.
The unit relations remove markings carrying the unit, except at the distinguished starting point.
After tensoring with $A$, these are exactly the normalized Hochschild chains.
The annular boundary joins consecutive inputs, including the last with the first.
Its differential is therefore the Hochschild differential.

For a dg model this is the explicit identification in Costello's Proposition~7.4.3.
The same proposition and its quasi-isomorphism invariance give the statement for $\Ainf$ models.

Filter the annular module by its number of markings.
Its differential lowers this filtration, and the unit relations split on the associated graded.
This proves the flatness needed for the derived tensor and shows that disjoint unions give quasi-isomorphic tensor comparison maps \cite[Lemma~6.3.1]{CostelloCY}.
It proves homological splitting of $\Psi_A$ and completes the construction.
\end{proof}

The open category has no circle object.
The circle assertion concerns $\Psi_A$, and an arbitrary open--closed theory receiving the same open theory receives a comparison from $\Psi_A$.
That comparison need not be an equivalence.
The two surface theories in the preceding chapter illustrate why the open states alone do not identify every possible closed sector.

In a ribbon graph, a cyclic order is chosen for the flags at each vertex.
The cyclic higher product paired with the remaining input labels a vertex, and the inverse pairing contracts an internal edge.
Costello's cellular model uses nodal disks and annuli, whose boundary creates nodes.
The dual ribbon description uses edge contractions.
The comparison of these models, including their orientations, is part of the geometric theorem, not a consequence of the $\Ainf$ identities alone.

\section{A finite open algebra with unbounded circle homology}

Take $A=k[t]/(t^2)$ in degree zero, with $\tau(1)=0$ and $\tau(t)=1$.
Multiplication and $B(a,b)=\tau(ab)$ make it a finite cyclic algebra of dimension zero.
In the basis $1,t$, the pairing and inverse pairing are
\[
 [B]=\begin{pmatrix}0&1\\1&0\end{pmatrix},
 \qquad \operatorname{coev}(1)=1\otimes t+t\otimes1.
\]
For $a=\alpha+\beta t$, contraction on either side gives
\[
 1\,B(t,a)+t\,B(1,a)=\alpha+\beta t,
 \qquad B(a,1)t+B(a,t)1=\alpha+\beta t.
\]
Both terms of the inverse pairing are necessary.
Joining them by multiplication gives the handle element $2t$, whose square is zero.

The circle complex retains more information than this two-dimensional open space.
Put $R=A^e=k[x,y]/(x^2,y^2)$ and let $R\to A$ send $x,y$ to $t$.
There is a free bimodule resolution
\begin{equation}\label{eq:cbc:cyclic:dual-number-resolution}
 \cdots\xrightarrow{x-y}R\xrightarrow{x+y}R
 \xrightarrow{x-y}R\longrightarrow A\longrightarrow0.
\end{equation}
Indeed, multiplication by $x-y$ has kernel spanned by $x+y,xy$, which is the image of multiplication by $x+y$.
The reverse kernel is spanned by $x-y,xy$, and the augmentation kernel is also this latter span.
This proves exactness in every degree.

Tensoring the resolution over $R$ with $A$ gives alternating differentials $0$ and $2t$.
Hence
\begin{equation}\label{eq:cbc:cyclic:dual-number-hh}
 HH_0(A)=A,\qquad
 HH_{2r+1}(A)=A/(t),\qquad
 HH_{2r+2}(A)=kt\quad(r\ge0).
\end{equation}
The circle complex is therefore not perfect over $k$.
This algebra is proper and not smooth: a perfect bimodule resolution would make $A\otimes_{A^e}^{\mathbb L}A$ bounded, contrary to~\eqref{eq:cbc:cyclic:dual-number-hh}.

Smoothness is thus unnecessary for Theorem~\ref{cbc:cyclic:surface}, but it changes the finiteness of the closed states.
If $A$ is both smooth and proper, a finite construction of the diagonal from perfect $A^e$-modules makes its Hochschild chains perfect over $k$.
Neither conclusion identifies the bordered-surface category with the compactification by stable algebraic curves.

\section{Stable curves and modular operations}

For $2g-2+n>0$, let $\overline{\mathcal M}_{g,n}$ be the moduli stack of stable complex curves of arithmetic genus $g$ with $n$ labeled smooth marked points.
Stability means that each rational component has at least three special points and each genus-one component at least one, with nodes counted as special points.
Identifying two marked points gives the clutching maps
\begin{align}
 \overline{\mathcal M}_{g_1,n_1+1}\times
 \overline{\mathcal M}_{g_2,n_2+1}
 &\longrightarrow\overline{\mathcal M}_{g_1+g_2,n_1+n_2},
 \label{eq:cbc:cyclic:separating}\\
 \overline{\mathcal M}_{g-1,n+2}
 &\longrightarrow\overline{\mathcal M}_{g,n}.
 \label{eq:cbc:cyclic:nonseparating}
\end{align}
Their images are nodal curves, so these are maps to the compactified moduli stacks.

A stable graph is a finite connected graph with labeled external legs and a nonnegative integer $g(v)$ at each vertex, such that $2g(v)-2+n(v)>0$.
Here $n(v)$ counts all incident flags, including both flags of a loop.
Its genus is
\[
 g(G)=\sum_v g(v)+b_1(G).
\]
Contracting an edge between distinct vertices adds their vertex genera.
Contracting a loop raises its vertex genus by one.
Both operations preserve $g(G)$.

A modular operad $P$ assigns complexes $P(g,n)$ with symmetric-group actions and equivariant composition maps
\[
 \mu_G:\bigotimes_v P(g(v),n(v))\longrightarrow P(g(G),n(G))
\]
for stable graphs.
Here $n(G)$ is the number of external legs.
The corolla map is the identity, and composition is compatible with every graph contraction and substitution.
In particular, the operad includes a contraction of two flags at one vertex.
The genus-zero cyclic operations alone do not supply this map.

Clutching gives the stable-curve modular operad, and compatible chains or homology give its linear versions \cite[\S6.2]{GetzlerKapranov}.
For a finite complex $V$ with a degree-zero symmetric perfect pairing, the collection $V^{\otimes n}$ is an endomorphism modular operad.
Its compositions contract the two factors at each internal edge with the pairing.
A modular action is a morphism to this endomorphism operad \cite[\S2.25]{GetzlerKapranov}.
In the dual convention of multilinear forms, the same contraction uses the inverse pairing.
Nonzero-degree pairings require the corresponding graded determinant twists.

A chiral coefficient over pointed curves must therefore provide the appropriate pullback and pushforward maps along~\eqref{eq:cbc:cyclic:separating} and~\eqref{eq:cbc:cyclic:nonseparating}, with their support and completion conditions.
Its operadic compositions must agree under successive clutching.
The bordered-surface theorem constructs neither these coefficient maps nor their comparison with its ribbon operations.

\section{The mixed sewing equation}

Suppose a common collection of decorated graphs has edges of two kinds, transverse and chiral, and is closed under their contractions.
Assume all iterated contraction targets are coherently identified, including those involving loops.
To a graph $G$ assign a cochain complex $C(G)$ and to an edge $e$ a degree-zero cochain map $c_e^G:C(G)\to C(G/e)$.
Require equivariance under graph isomorphisms.

The coefficient equation for distinct edges is
\begin{equation}\label{eq:cbc:cyclic:mixed-square}
 c_f^{G/e}c_e^G=c_e^{G/f}c_f^G.
\end{equation}
For different edge kinds this is an additional mixed sewing condition.
For one kind it is part of the corresponding modular or ribbon composition law.

\begin{proposition}[Strict mixed contraction]\label{cbc:cyclic:mixed}
Under these assumptions, contraction of any finite set of edges is independent of its ordering.
Put $\operatorname{or}(G)=\det(kE(G))$ in degree $-|E(G)|$ and form the direct sum
\[
 K^n=\bigoplus_{[G]}
 \bigl(C^{n+|E(G)|}(G)\otimes\operatorname{or}(G)\bigr)_{\operatorname{Aut}(G)}.
\]
For homogeneous $x\in C^q(G)$, let
\[
 D(x\otimes\omega)=d_Cx\otimes\omega
       +(-1)^q\sum_{e\in E(G)}c_e^Gx\otimes\iota_e\omega,
\]
where $\iota_e$ removes $e$ from the oriented wedge with its exterior sign.
Then $D^2=0$.
The transverse and chiral edge parts anticommute.
\end{proposition}
\begin{proof}
Any two orders of a finite edge set differ by adjacent transpositions.
Equation~\eqref{eq:cbc:cyclic:mixed-square} identifies the composites for each such transposition at the intermediate graph.
The coherent target identifications then give the same total contraction.
This is the graph-contraction argument of Getzler--Kapranov's Theorem~3.7, with the mixed squares included.

For the differential, contracting two distinct edges in opposite orders gives equal coefficient maps and opposite exterior signs, since $\iota_f\iota_e=-\iota_e\iota_f$.
The internal differential anticommutes with each edge term because $c_e$ is a cochain map and the sign $(-1)^q$ changes when $d_C$ is applied.
An edge cannot be contracted twice.
These facts prove the square-zero and anticommutation identities, as in Theorem~\ref{cbc:ord:forest-d}.
The argument uses no acyclicity of $G$, so it includes loops.
Equivariance descends the differential to coinvariants, and each graph contributes finitely many terms to the direct sum.
\end{proof}

The determinant sign and the coefficient equation have distinct roles.
For two edges $e,f$, the two-edge term in $D^2$ is
\[
 \bigl(c_f^{G/e}c_e^G-c_e^{G/f}c_f^G\bigr)x.
\]
For example, take $C(G)=k^2$ at every contraction target and take the two maps to be
\[
 P=\begin{pmatrix}1&0\\0&0\end{pmatrix},\qquad
 Q=\begin{pmatrix}0&1\\0&0\end{pmatrix}.
\]
Then $QP-PQ=-Q\ne0$.
An orientation does not remove this defect.
If the coefficient equation holds only up to homotopy, the action also requires compatible higher homotopies for longer contraction sequences.

Nor does an edge coloring alone define two independent genera.
Take two genus-zero vertices joined by one transverse and one chiral edge, and give each vertex an external leg so that the graph is stable.
Contracting the transverse edge first leaves a chiral loop.
Reversing the order leaves a transverse loop.
Counting the genus increment by the loop color gives $(g_\perp,g_X)=(0,1)$ in one order and $(1,0)$ in the other.
The total genus is one in both.

Separate geometric genera must be carried by a more structured sewing object if both contraction orders are allowed.

For an ordered chiral algebra, the remaining construction is consequently a common coefficient system realizing the two geometric sewings, with~\eqref{eq:cbc:cyclic:mixed-square} or its coherent replacement.
It must preserve the coefficient tensor, support, completion, units and the chosen cyclic pairing.
Theorem~\ref{cbc:cyclic:surface} supplies the transverse surface operations when a finite cyclic model is given.
An extension to stable curves, and a comparison of the two geometric genera, require their own maps and identities.

\part{Exact noncommutative models}

\chapter{Free and square-zero pairs}\label{chap:cbc:free-square-zero}
The bar differential joins neighboring inputs by multiplication.
For a tensor algebra these joins contract every word involving more than one generator.
For a square-zero algebra all such joins vanish, so every tensor length remains.
The two calculations exhibit the dependence of the bar on multiplication and determine the resulting dual products.

\section{The tensor algebra and its primitive summand}

Let $k$ be a commutative ring and $M$ a cochain complex of $k$-modules.
All tensor products in this section are over $k$.
Put $s=[1]$, so $|sm|=|m|-1$.
The tensor algebra and its augmentation are
\[
 T(M)=\bigoplus_{n\ge0}M^{\otimes n},\qquad
 \epsilon:T(M)\longrightarrow k,
\]
where $M^{\otimes0}=k$ and $\epsilon$ projects to this summand.
Multiplication concatenates tensors, and the differential extends that of $M$ by the graded Leibniz rule.
Give each generator weight one.

For an augmented dg algebra $A=k\oplus\bar A$, use the normalized bar
\[
 BA=\bigoplus_{r\ge0}(s\bar A)^{\otimes r}
\]
with the differential and deconcatenation of Theorem~\ref{cbc:internal:normalized}.
In particular, $b_1(sa)=-s(da)$ and $b_2(sa\otimes sb)=(-1)^{|a|}s(ab)$.
This is the direct-sum bar, with no tensor-length completion.

Let $K(M)=k\oplus sM$ have its suspended internal differential and primitive reduced summand.
Thus
\[
 \Delta(sm)=1\otimes sm+sm\otimes1,
 \qquad \Delta(1)=1\otimes1.
\]
Including a generator as a one-input bar defines a dg coalgebra map
\begin{equation}\label{eq:cbc:free:inclusion}
 i:K(M)\longrightarrow B(T(M)),\qquad i(sm)=[m].
\end{equation}

\begin{theorem}[Bar of the tensor algebra]\label{cbc:free:bar}
The map $i$ is a quasi-isomorphism of coaugmented dg coalgebras.
There are natural chain maps $p:B(T(M))\to K(M)$ and a degree-$-1$ map $h$ satisfying
\[
 pi=\id,\qquad bh+hb=\id-ip,\qquad h^2=hi=ph=0.
\]
No flatness hypothesis on $M$ is needed for these chain identities.
\end{theorem}
\begin{proof}
Expand each bar input into its finite sum of generator tensors.
A summand of positive generator weight $n$ is specified by a subset $S\subseteq\{1,\ldots,n-1\}$ of cuts between adjacent generators.
For homogeneous $m_1,\ldots,m_n$, denote its tensor by $w_S$.
Its degree is
\[
 |w_S|=\sum_{a=1}^n|m_a|-|S|-1.
\]
This notation describes multilinear tensor maps and does not require a basis of $M$.

For $j\in S$, let $r_S(j)=\#\{a\in S:a\le j\}$.
The multiplication part $b_\mu$ removes the cut $j$ with coefficient
\begin{equation}\label{eq:cbc:free:cut-sign}
 (-1)^{\sum_{a\le j}|m_a|-r_S(j)+1}.
\end{equation}
Indeed, $r_S(j)-1$ suspended blocks precede the block ending at $j$.
Passing those blocks and applying $b_2$ gives exactly this sign.

For $n\ge2$, define
\begin{equation}\label{eq:cbc:free:homotopy}
 h(w_S)=
 \begin{cases}
  (-1)^{|m_1|}w_{S\cup\{1\}},&1\notin S,\\
  0,&1\in S.
 \end{cases}
\end{equation}
Set $h=0$ in weights zero and one.
Let $p$ be the identity on those two weights and zero on every larger weight.
Both $b_\mu$ and the internal differential preserve generator weight, so $p$ and $i$ are chain maps.

If $1\notin S$, removing the inserted first cut contributes $w_S$ to $b_\mu h(w_S)$.
Every other cut has its rank increased by one after insertion, so the two orders of removing that cut and inserting the first cancel.
If $1\in S$, only removal of that first cut can contribute to $hb_\mu(w_S)$, and its coefficient is again one.
This proves $b_\mu h+hb_\mu=\id$ in weights at least two.

The internal differential replaces $m_j$ by $d_Mm_j$ with coefficient
\[
 (-1)^{\sum_{a<j}|m_a|-\#\{a\in S:a<j\}+1}.
\]
For $j>1$, inserting the first cut reverses this sign.
For $j=1$, applying $d_M$ reverses the sign $(-1)^{|m_1|}$ in $h$.
Thus the internal differential anticommutes with $h$ in every case.
Combining the two parts proves $bh+hb=\id-ip$.

Inserting the first cut twice is impossible, giving $h^2=0$.
The remaining side identities follow from the weight projections.
Finally, a one-input bar is primitive, so $i$ preserves deconcatenation, the counit and the coaugmentation.
The chain contraction makes it a quasi-isomorphism.
\end{proof}

The projection $p$ need not preserve the coproduct.
Over a field, take $M=kx$ in degree zero and consider $[x|x]$.
Its image under $p$ is zero, while its reduced deconcatenation gives
\[
 (p\otimes p)\overline\Delta[x|x]=sx\otimes sx\ne0.
\]
Thus the displayed contraction is a contraction of complexes, and $i$ is the specified coalgebra quasi-isomorphism.
The distinction matters when dualizing multiplication.

\section{The square-zero algebra}

The square-zero extension of $M$ is the dg algebra
\[
 S(M)=k\oplus M,\qquad
 (a,m)(a',m')=(aa',am'+a'm).
\]
Its unit is $(1,0)$, its augmentation projects to $k$, and its differential restricts to $d_M$.

\begin{proposition}[Square-zero bar]\label{cbc:free:square-zero}
There is an equality of dg coalgebras
\[
 B(S(M))=T^c(sM)=\bigoplus_{r\ge0}(sM)^{\otimes r}
\]
with deconcatenation and the internal tensor differential.
The multiplication part of the bar differential is zero.
\end{proposition}
\begin{proof}
Every adjacent multiplication in the augmentation ideal is a product of two elements of $M$.
Such products vanish by the definition of $S(M)$.
The unary bar differential remains $b_1(sm)=-s(d_Mm)$, and its tensor extensions give precisely the stated differential.
The coproduct is unchanged.
\end{proof}

If $M$ is concentrated in degree zero with zero differential, every reduced bar tensor is a cycle.
For a $d$-dimensional vector space over a field, the homological bar degrees therefore give
\[
 \dim_k\operatorname{Tor}^{S(M)}_r(k,k)=d^r,
 \qquad
 \sum_{r\ge0}\dim_k\operatorname{Tor}^{S(M)}_r(k,k)\,z^r
 =\frac1{1-dz}.
\]
The tensor algebra instead has $\operatorname{Tor}_0=k$, $\operatorname{Tor}_1=M$ and zero higher groups, by Theorem~\ref{cbc:free:bar}.
These are statements about the two specified differentials, not a comparison of their chain dimensions.

\section{Dual multiplication and completion}

Now assume that $M$ is a bounded complex of finitely generated projective $k$-modules.
It represents a dualizable object of $D(k)$.
Write $M^\vee=\operatorname{Hom}_k(M,k)$, with its Hom differential.
For a dg coalgebra $C$, its full linear dual is $C^\vee=\operatorname{Hom}_k(C,k)$.
Its convolution product is
\[
 (f\star g)(c)=\sum_{(c)}(-1)^{|g||c_{(1)}|}
                    f(c_{(1)})g(c_{(2)}).
\]
Every deconcatenation here is finite on an individual input.

For an augmented algebra $A$, write $A^!=\operatorname{RHom}_A(k,k)$ with its derived composition product.
The completion of the tensor algebra on a complex $N$ means
\[
 \widehat T(N)=\prod_{r\ge0}N^{\otimes r}.
\]
Products are degreewise, and the length filtration has neighborhoods $\prod_{r\ge p}N^{\otimes r}$ of zero.
Concatenation is continuous, because an output of length $r$ receives only the finitely many products with input lengths summing to $r$.

\begin{theorem}[Free--square-zero duality]\label{cbc:free:duality}
Under these hypotheses there are dg algebra quasi-isomorphisms representing
\[
 T(M)^!\simeq k\oplus M^\vee[-1],\qquad
 S(M)^!\simeq\widehat T\bigl(M^\vee[-1]\bigr).
\]
The first algebra has square-zero reduced summand and its internal differential.
The second has the completed concatenation product and its internal tensor differential.
Its completion comes from dualizing a direct sum of bar tensors.
\end{theorem}
\begin{proof}
For either algebra, the left bar resolution $P=A\otimes BA\to k$ computes derived Hom.
For homogeneous inputs its differential is
\[
 \begin{aligned}
 d_P(a\otimes[a_1|\cdots|a_r])
 &=d_Aa\otimes[a_1|\cdots|a_r]
   +(-1)^{|a|}a\otimes b[a_1|\cdots|a_r]\\
 &\quad+(-1)^{|a|+1}aa_1\otimes[a_2|\cdots|a_r],
 \end{aligned}
\]
with the last term omitted when $r=0$.

Writing $\bar a=a-\epsilon(a)1$, its augmented contraction inserts $\bar a$ as a first bar input, with coefficient $-1$.
The left multiplication term recovers $\bar a$.
In positive bar length, first multiplication followed by insertion also recovers the scalar part of $a$.
The other internal and merging terms cancel in pairs.
This gives $d_Ph_P+h_Pd_P=\id-\eta\epsilon$, where $\eta:k\to P$ inserts the algebra unit in bar length zero.

Here is the projectivity assertion needed for that statement.
Filter $P$ by bar length.
Its associated graded pieces are $A\otimes(s\bar A)^{\otimes r}$ with their internal differential.
The coefficient complexes are direct sums of bounded complexes of projective $k$-modules, since $M$ has this property and tensor powers preserve it.
Consequently their Hom complexes into any acyclic complex are acyclic.

Let $Q$ be an acyclic dg $A$-module and let $f:P\to Q$ be a closed homogeneous map.
Construct a primitive for its restriction to each finite bar filtration, extending the previous primitive.
The filtration is split as graded $A$-modules, so a graded extension of that primitive exists.
Its differential defect vanishes on the preceding filtration and is a closed map from the associated graded piece to $Q$.
The Hom acyclicity just proved supplies a correction.

Induction gives compatible primitives, and their union is a primitive for $f$ on $P$.
Thus $\operatorname{Hom}_A(P,Q)$ is acyclic, which is the required homotopical projectivity.
The augmented bar contraction proves $P\to k$ is a quasi-isomorphism.

The right $BA$-coaction on $P$ lifts convolution to composition:
\[
 L_f=(\id_P\otimes f)(\id_A\otimes\Delta):P\longrightarrow P.
\]
Coassociativity gives $L_fL_g=L_{f\star g}$, with the graded tensor sign displayed above.
The chain-coaction identity gives $[d_P,L_f]=L_{\partial f}$.
Postcomposition with $P\to k$ is a quasi-isomorphism from $\operatorname{End}_A(P)$ to $\operatorname{Hom}_A(P,k)$ by homotopical projectivity.
Its composite with $L$ is the identity under $\operatorname{Hom}_A(P,k)=\operatorname{Hom}_k(BA,k)$.
Therefore $(BA)^\vee$ is a dg algebra model for $A^!$.

For $A=T(M)$, dualize the coalgebra map~\eqref{eq:cbc:free:inclusion}.
It gives a dg algebra map
\[
 i^\vee:(B(T(M)))^\vee\longrightarrow K(M)^\vee
       =k\oplus M^\vee[-1].
\]
Dualizing the explicit chain contraction preserves its homotopy identities, so this map is a quasi-isomorphism.
The reduced coproduct of $K(M)$ is zero, hence the reduced convolution product of its dual is zero.
This proves the first formula without treating $p$ as a coalgebra map.

For $A=S(M)$, the preceding proposition gives
\[
 (BA)^\vee
 =\prod_{r\ge0}\operatorname{Hom}_k((sM)^{\otimes r},k)
 \cong\prod_{r\ge0}(M^\vee[-1])^{\otimes r}.
\]
The last comparison uses finite projectivity and the graded tensor evaluation in each length.
Under it, deconcatenation dualizes to concatenation with exactly the convolution sign.
The differentials agree by the Hom and tensor conventions.
This proves the second formula, including its multiplication and completion.
\end{proof}

Over a field, take $M=kx$ in degree zero, so its dual generator has degree one.
The square-zero algebra $S(M)$ has a nonzero dual class in every length.
In particular, the square of that odd generator is nonzero.
The completed tensor algebra is associative and has no imposed graded-commutativity relation.

The full dual in the proof is a product because the coalgebra is a direct sum.
It is not the continuous dual of an unspecified completed coalgebra.
Changing either topology requires a separate comparison.

Finite projectivity cannot simply be omitted from the tensor-dual identification.
Over a field, let $M=\bigoplus_{j\ge1}ke_j$.
The bilinear form $\beta(e_i,e_j)=\delta_{ij}$ belongs to $(M\otimes M)^\vee$.
It does not belong to the image of $M^\vee\otimes M^\vee$.
Indeed, a sum of $r$ elementary tensor functionals restricts to a matrix of rank at most $r$ on any finite coordinate subspace.
The restriction of $\beta$ to the first $r+1$ coordinates has rank $r+1$.

For a general complex, the actual dual $\prod_r\operatorname{Hom}_k((sM)^{\otimes r},k)$ must therefore be retained unless its tensor comparison is proved.

\section{Realization in the chiral coefficient category}

The preceding computations take place in the ordinary coefficient tensor.
Their formal tensor identities also apply in a specified symmetric monoidal dg coefficient model with the sums, suspensions and normalization used above.
For example, an exact strong monoidal dg functor preserving those sums carries the free-bar inclusion and its contraction to the corresponding ambient coefficients.
Preserving the dual algebra formulas additionally requires the tensor-dual and product comparisons occurring in their proof.

These requirements do not make a componentwise tensor algebra a factorization object.
The normalized summands in Theorem~\ref{cbc:internal:normalized} are ambient coefficients, whereas factorization compares the entire object on separated configurations.
Already $k^2\otimes\omega_{\Ran}$ has rank two at a singleton and at a separated pair, while factorization would require rank four at the pair.
Thus tensoring a factorizing unit by an arbitrary algebra does not supply the missing factorization equivalences.

To interpret $T_X(M)$ or $\mathbbm1\oplus M$ as ordered chiral algebras, construct their factorization structures and compatible multiplication and augmentation maps.
Apply the realization hypotheses of Theorem~\ref{cbc:internal:normalized} to their full simplicial bars.
Then specify the comparison with the free or square-zero coefficient model and prove compatibility of~\eqref{eq:cbc:free:inclusion} with the retained structures.
The actual factorizing independence examples of Theorem~\ref{cbc:two:independence} remain available, without this generic tensoring assertion.

A nonlocal vertex algebra supplies neither that factorization structure nor a free geometric associative chiral algebra by itself.
The conditional recognition of Theorem~\ref{thm:formal-disc} starts from a specified geometric realization.
Constructing its reverse, and comparing the two associative bars on the resulting object, remain separate mathematical problems.

\chapter{The quantum plane}\label{chap:cbc:quantum-plane}

The quantum plane has the same graded dimensions as a polynomial algebra in two variables.
Its augmentation module also has the same Betti numbers.
The parameter remains in the bar coproduct and the Yoneda multiplication, which we compute explicitly.

\section{Normal forms and multiplication}

Let $k$ be a field, choose $q\in k^\times$, and set
\[
 A_q=k\langle x,y\rangle/(yx-qxy).
\]
The augmentation $\epsilon:A_q\to k$ sends $x$ and $y$ to zero.
Give both generators polynomial weight one and cohomological degree zero.
The differential is zero, and all tensor products below are over $k$.

\begin{theorem}[PBW basis of the quantum plane]\label{cbc:qplane:pbw}
The monomials $x^ay^b$, for $a,b\ge0$, form a basis of $A_q$.
Their product and the Hilbert series are
\[
 (x^ay^b)(x^cy^d)=q^{bc}x^{a+c}y^{b+d},\qquad
 H_{A_q}(z):=\sum_{w\ge0}\dim_k(A_q)_wz^w=\frac1{(1-z)^2}.
\]
\end{theorem}
\begin{proof}
Moving each $y$ past a later $x$ by $yx=qxy$ expresses every word in the stated form.
To prove independence, take a vector space with basis $e_{a,b}$ and multiplication
\[
 e_{a,b}e_{c,d}=q^{bc}e_{a+c,b+d}.
\]

This multiplication is associative because
\[
 bc+(b+d)e=de+b(c+e).
\]
Its unit is $e_{0,0}$, and $e_{0,1}e_{1,0}=q e_{1,0}e_{0,1}$.
The resulting homomorphism from $A_q$ sends $x^ay^b$ to $e_{a,b}$, proving independence.
Commuting $y^b$ past $x^c$ gives the product formula.
There are $w+1$ basis monomials of weight $w$, giving the Hilbert series.
\end{proof}

For $q\ne1$, the relation has a nonzero commutator $(q-1)xy$.
Thus the common Hilbert series does not identify $A_q$ with the commutative plane.

\section{The linear resolution}

Consider a nonnegatively graded algebra with weight-zero part $k$.
It is \emph{Koszul} if its augmentation module has a graded free resolution whose term in homological degree $r$ is generated in polynomial weight $r$.
Write $e_x,e_y$ for free module generators of weight one and $e_r$ for one of weight two.
Place them in cohomological degrees $-1,-1,-2$, respectively.

\begin{theorem}[Koszulness]\label{cbc:qplane:resolution}
The following complex of free left $A_q$-modules is an augmented resolution:
\[
 0\longrightarrow A_qe_r
 \xrightarrow{d_2}A_qe_x\oplus A_qe_y
 \xrightarrow{d_1}A_q\xrightarrow{\epsilon}k\longrightarrow0,
\]
where
\[
 d_1(ue_x+ve_y)=ux+vy,\qquad
 d_2(ce_r)=cy e_x-qcx e_y.
\]
In particular, $A_q$ is Koszul.
\end{theorem}
\begin{proof}
The composite is $d_1d_2(ce_r)=c(yx-qxy)=0$.
The image of $d_1$ is the augmentation ideal, so exactness remains to be proved at the two left terms.

Write the finite sums
\[
 u=\sum_{i,j\ge0}u_{ij}x^iy^j,\qquad
 v=\sum_{i,j\ge0}v_{ij}x^iy^j.
\]
The coefficient of $x^ay^b$ in $ux+vy$ is
\[
 q^b u_{a-1,b}+v_{a,b-1},
\]
with negative indices interpreted as zero.

If $ux+vy=0$, then $u_{i,0}=0$, $v_{0,j}=0$, and
\[
 v_{a,b}=-q^{b+1}u_{a-1,b+1}.
\]
Set $c=\sum_{i,j\ge0}u_{i,j+1}x^iy^j$.
These identities give $u=cy$ and $v=-qcx$, proving $\ker d_1=\operatorname{im}d_2$.
Finally, $cy=0$ forces $c=0$ by the ordered basis, so $d_2$ is injective.
The stated generator weights make this a linear resolution.
\end{proof}

\section{The finite bar coalgebra}

Use the normalized bar $BA_q$ with $s=[1]$, so $|sx|=|sy|=-1$.
Its differential merges adjacent inputs with alternating signs, beginning with a plus sign.
The element
\[
 \omega_q=[y|x]-q[x|y]
\]
is a cycle because its boundary is $yx-qxy$.

Let $C_q\subset BA_q$ be the subspace spanned by $1$, $sx=[x]$, $sy=[y]$ and $\omega_q$.
Its differential is zero.
It is a coaugmented subcoalgebra: $sx,sy$ are primitive and
\begin{equation}\label{eq:cbc:qplane:coproduct}
 \overline\Delta(\omega_q)=sy\otimes sx-q\,sx\otimes sy.
\end{equation}
Here $\omega_q$ has cohomological degree $-2$ and polynomial weight two.

\begin{proposition}\label{cbc:qplane:coalgebra}
The inclusion $j:C_q\longrightarrow BA_q$ is a dg coalgebra quasi-isomorphism.
Consequently the augmentation Tor groups have dimensions $1,2,1$ in homological degrees $0,1,2$ and vanish in higher degrees.
Their coalgebra structure is the one displayed above.
\end{proposition}
\begin{proof}
Let $P=A_q\otimes BA_q\to k$ be the left bar resolution, with the differential used in Theorem~\ref{cbc:free:duality}.
The linear resolution $K$ of Theorem~\ref{cbc:qplane:resolution} maps to $P$ by
\[
 1\longmapsto1\otimes1,\qquad
 e_x\longmapsto-1\otimes[x],\qquad
 e_y\longmapsto-1\otimes[y],\qquad
 e_r\longmapsto1\otimes\omega_q.
\]
The minus signs compensate for the left action term in that bar differential.
On $e_r$ the chain-map equation is
\[
 d_P(1\otimes\omega_q)=-y\otimes[x]+qx\otimes[y].
\]

Both complexes are free resolutions, and the map lifts the identity of $k$.
Lift an inverse successively through the exact kernels in each degree, using freeness.
The same lifting applied to the differences from the identities constructs chain homotopies.
Thus this comparison is a homotopy equivalence.

Tensoring these homotopies with the augmentation right module preserves their identities.
The resulting map $k\otimes_{A_q}K\to BA_q$ identifies its source with $C_q$, with a minus sign on each degree-one generator.
It therefore proves that $j$ is a quasi-isomorphism.
The coproduct formula proves the claimed coalgebra structure, and the generator degrees give the Tor groups.
\end{proof}

\section{The Yoneda product}

Let $V=kx\oplus ky$ and $R=k(y\otimes x-qx\otimes y)\subset V^{\otimes2}$.
The \emph{quadratic dual} is the algebra $Q_q=T(V^\vee)/(R^\perp)$, where the annihilator uses the ordinary tensor evaluation.
For the basis $\xi,\eta$ dual to $x,y$, this gives
\[
 Q_q=k\langle\xi,\eta\rangle/
       (\xi^2,\eta^2,\xi\eta+q\eta\xi).
\]

Source weight records the polynomial weight on which a cochain evaluates.
Assign these dual generators cohomological degree one and source weight one.
Its graded Hom degree is the negative of that weight.
Use the Yoneda convention that $\alpha\beta$ means first $\beta$, then $\alpha$.

\begin{theorem}[Yoneda algebra of the quantum plane]\label{cbc:qplane:yoneda}
There is an isomorphism of bigraded algebras
\[
 \operatorname{Ext}_{A_q}^*(k,k)\cong Q_q.
\]
The full linear dual of $BA_q$ is a dg algebra model for $\operatorname{RHom}_{A_q}(k,k)$.
Restriction along $j$ gives a dg algebra quasi-isomorphism from that model to $C_q^\vee\cong Q_q$, with zero differential on the latter.
\end{theorem}
\begin{proof}
Applying $\operatorname{Hom}_{A_q}(-,k)$ to $K$ kills its differential, because all matrix entries lie in the augmentation ideal.
Let $\xi,\eta,\zeta$ be the resulting cochains dual to $e_x,e_y,e_r$.
Together with the unit they give the full Ext vector space.

Define degree-one $A_q$-linear endomorphisms $F_\xi,F_\eta$ of $K$ by
\[
 \begin{array}{c|ccc}
       &e_x&e_y&e_r\\ \hline
 F_\xi&1&0&-e_y\\
 F_\eta&0&1&q e_x
 \end{array}
\]
and set them to zero on $K^0=A_q$.
On $e_r$, the equations $dF_\xi+F_\xi d=0$ and $dF_\eta+F_\eta d=0$ reduce to $-y+y=0$ and $qx-qx=0$.
On the other generators both equations vanish termwise.
These are closed lifts of the two degree-one cochains.
Their compositions give
\begin{equation}\label{eq:cbc:qplane:yoneda-products}
 \xi^2=\eta^2=0,\qquad
 \xi\eta=q\zeta,\qquad \eta\xi=-\zeta.
\end{equation}
Postcomposition with $K\to k$ computes the derived composition product, since the finite free resolution is homotopically projective.

The relations of $Q_q$ hold, and $\eta\xi=-\zeta$ shows that the two degree-one classes generate Ext.
Those relations reduce every word of length at least three to zero and leave at most four basis elements.
Ext has dimension four, so this surjection from $Q_q$ is an isomorphism.
All maps preserve the source weights.

The left bar resolution also computes derived Hom: its bar-length filtration has free coefficient modules, and the successive lifting argument proves homotopical projectivity.
Its right bar coaction identifies dual convolution with derived composition, as in the proof of Theorem~\ref{cbc:free:duality}.
The homotopy equivalence proving the preceding proposition dualizes to a homotopy equivalence of complexes.
Because $j$ is a coalgebra map, its dual is also an algebra map.

Finally, the convolution sign for two degree-one functionals on a two-input bar is minus one.
Dualizing~\eqref{eq:cbc:qplane:coproduct} therefore gives the same product table~\eqref{eq:cbc:qplane:yoneda-products}.
The two primitive duals pull back to $-\xi,-\eta$ under $K\to P$.
Their simultaneous sign change leaves that table unchanged.
This identifies $C_q^\vee$ with $Q_q$ as asserted, including the multiplication.
\end{proof}

When $q=1$, this Yoneda algebra is the exterior algebra on two generators.
For $q\ne1$, its odd anticommutator is
\[
 \xi\eta+\eta\xi=(q-1)\zeta\ne0.
\]
Thus it is not graded commutative, although its Betti numbers are unchanged.
Interchanging $x$ and $y$ replaces the presentation parameter by $q^{-1}$.
The displayed formula keeps their order fixed.

\section{Geometric realization}

These computations concern the ordinary augmented algebra $A_q$ and its coefficient tensor.
A geometric realization must supply a factorization object with a compatible transverse multiplication and augmentation, together with the relation $yx=qxy$ in that tensor.
It must also compare the full geometric bar with the coefficient bar and preserve~\eqref{eq:cbc:qplane:coproduct}.
Tensoring a constant algebra by a factorizing coefficient does not supply those maps, as the rank example in Chapter~\ref{chap:cbc:free-square-zero} shows.

The marked-point examples of Theorem~\ref{cbc:two:independence} retain their actual factorization structures.
To equip a quantum-plane realization with the second associative chiral operation, construct that operation and its mixed coherences as well.
The coefficient presentation and its Yoneda algebra do not by themselves construct this geometric extension.

\chapter{Quiver path algebras}\label{chap:cbc:quivers}
A quiver records several objects and the maps between them.
Its bar construction must retain the source and target of every map.
The vertex idempotents provide the base over which tensors compose, and a quadratic relation becomes a two-input cycle in the relative bar.

\section{Vertices, paths and left modules}

Let $k$ be a field and $Q$ a finite quiver, with vertex set $Q_0$ and arrow set $Q_1$.
Put
\[
 S=\bigoplus_{i\in Q_0}ke_i,\qquad
 e_ie_j=\delta_{ij}e_i,\qquad \sum_i e_i=1.
\]
An arrow $a:i\to j$ belongs to $e_jMe_i$, where $M$ is the $S$-bimodule spanned by the arrows.
Give each arrow weight one and each vertex idempotent weight zero.

The tensor algebra
\[
 A=T_S(M)=\bigoplus_{n\ge0}M^{\otimes_S n},\qquad M^{\otimes_S0}=S,
\]
has multiplication by concatenation and augmentation $\epsilon:A\to S$ given by projection to weight zero.
Thus paths multiply from right to left: for $a:i\to j$ and $b:j\to\ell$, the composite is $ba$.
The balancing over $S$ makes a tensor of noncomposable arrows zero.

Write $L_i=Se_i$ for the left simple $A$-module obtained through $\epsilon$, and $P_i=Ae_i$ for the corresponding left projective module.
For $Q=(1\xrightarrow{a}2)$, the algebra has basis $e_1,e_2,a$ and $a=e_2ae_1$.
The sequence
\[
 0\longrightarrow P_2\xrightarrow{\,u\mapsto ua\,}P_1
 \longrightarrow L_1\longrightarrow0
\]
is exact: $P_2=ke_2$, and the kernel of $P_1\to L_1$ is $ka$.
Also $P_2=L_2$.
Applying $\operatorname{Hom}_A(-,L_j)$, and using
$\operatorname{Hom}_A(P_i,L_j)=e_iL_j$, gives
\[
 \operatorname{Ext}^1_A(L_1,L_2)=k,\qquad
 \operatorname{Ext}^1_A(L_i,L_i)=0,\qquad
 \operatorname{Ext}^1_A(L_2,L_1)=0.
\]
There are no higher Ext groups.

These endpoints also fix the dual-arrow convention.
Let $E=\operatorname{Ext}^*_A(S,S)$, with product given by composition, the right-hand factor acting first.
Its vertex corner is
\[
 e_jE^re_i=\operatorname{Ext}^r_A(L_i,L_j).
\]
Hence the degree-one Ext arrow in this example runs from $1$ to $2$.
The linear dual $M^\vee=\operatorname{Hom}_k(M,k)$ has bimodule action
\[
 (u f v)(m)=f(vmu),\qquad u,v\in S,
\]
so the functional dual to $a$ lies in $e_1M^\vee e_2$.
This bimodule dual reverses endpoints.

For any finite quiver with no relations, multiplication identifies $A\otimes_S M$ with the positive-weight ideal of $A$.
Thus $0\to A\otimes_S M\to A\to S\to0$ is a projective resolution.
Applying Hom to the simples shows that
\[
 e_jE^1e_i\cong(e_jMe_i)^\vee,\qquad E^{\ge2}=0.
\]
With $u\cdot_{\mathrm{op}}v=vu$ defining the opposite product, this yields
\[
 E^{\mathrm{op}}\cong S\oplus M^\vee[-1],
\]
where the reduced product is zero and $[-1]$ places the dual arrows in cohomological degree one.
The ordinary and graded opposite products agree here because every product of two positive-degree elements vanishes.

\section{The relative free bar}

Allow now $M$ to be a cochain complex of $S$-bimodules, with differential preserving its vertex corners.
For an augmented dg $S$-algebra $A=S\oplus\bar A$, define
\[
 B_SA=\bigoplus_{r\ge0}(s\bar A)^{\otimes_S r},\qquad s=[1].
\]
Thus $|sm|=|m|-1$.
Use $b_1(sa)=-s(da)$ and $b_2(sa\otimes sb)=(-1)^{|a|}s(ab)$, extended with the graded tensor signs.

The coproduct is deconcatenation over $S$.
This makes $B_SA$ an $S$-coring: an $S$-bimodule with coassociative coproduct to $B_SA\otimes_S B_SA$ and counit to $S$.
The length-zero summand supplies its coaugmentation.
All sums are direct sums, without a completion.

Give $S\oplus sM$ its suspended internal differential and zero reduced coproduct.
For $m\in e_jMe_i$, its full coproduct is
\[
 \Delta(sm)=e_j\otimes sm+sm\otimes e_i,
 \qquad \Delta(e_i)=e_i\otimes e_i.
\]

\begin{theorem}[Relative path-algebra bar]\label{cbc:quiver:free-bar}
The map
\[
 \iota:S\oplus sM\longrightarrow B_S(T_S(M)),\qquad
 \iota(sm)=[m],\quad \iota(e_i)=e_i,
\]
is a quasi-isomorphism of coaugmented dg $S$-corings.
The underlying complex of the target represents $S\otimes^{\mathbb L}_{T_S(M)}S$, with the first $S$ a right module and the second a left module.
\end{theorem}
\begin{proof}
Expand a bar input into generator tensors.
A tensor of $n$ homogeneous generators $m_1,\ldots,m_n$, with a set $C\subseteq\{1,\ldots,n-1\}$ of cuts, defines a bar tensor $w_C$.
Every cut and concatenation is balanced over $S$, so this notation retains the two endpoint idempotents.
Let $p$ project to generator weights zero and one.
For $n\ge2$, set
\[
 h(w_C)=
 \begin{cases}
 (-1)^{|m_1|}w_{C\cup\{1\}},&1\notin C,\\
 0,&1\in C,
 \end{cases}
\]
and put $h=0$ in weights zero and one.
These are $S$-bimodule maps.

Removing the inserted first cut contributes $w_C$ with coefficient one.
Insertion increases the rank of every other cut by one, so its two orders with any other cut removal have opposite signs.
For an internal differential on $m_j$ with $j>1$, the inserted suspension reverses the tensor sign.
For $j=1$, the change in $|m_1|$ reverses the coefficient of $h$ instead.
These are the cancellations in~\eqref{eq:cbc:free:cut-sign}, now in each vertex corner.

They give
\[
 bh+hb=\id-\iota p,\qquad p\iota=\id.
\]
The inclusion $\iota$ preserves deconcatenation, since each $[m]$ is primitive relative to its endpoints.
It also preserves the internal differential, counit and coaugmentation.
The contraction proves the asserted quasi-isomorphism.

For the derived interpretation, use the left relative bar resolution $A\otimes_S B_SA\to S$.
Its augmented contraction inserts the reduced part of the leading coefficient, as in the proof of Theorem~\ref{cbc:free:duality}.
The same formula is $S$-balanced and proves exactness.

Its filtration by bar length has quotients induced from complexes of $S$-modules.
Since $S$ is semisimple, acyclic complexes of $S$-modules contract, so these coefficient complexes are homotopically projective.
The filtration argument of that proof therefore makes the relative bar a homotopically projective resolution as well.
Tensoring it with the right augmentation module $S$ gives $B_SA$.
\end{proof}

The projection $p$ is generally not a coring map.
For composable arrows $a:1\to2$ and $b:2\to3$,
\[
 p[b|a]=0,\qquad
 (p\otimes_Sp)\overline\Delta[b|a]=sb\otimes_Ssa\ne0.
\]
The coring quasi-isomorphism in the theorem is the inclusion $\iota$.

\section{Relations and the quadratic complex}

Return to an ordinary finite-dimensional arrow bimodule $M$, placed in cohomological degree zero and weight one.
Let $R\subseteq M\otimes_S M$ be an $S$-subbimodule and put $A=T_S(M)/(R)$.
For $n\ge2$, define
\[
 W_n=\bigcap_{j=0}^{n-2}
 M^{\otimes_S j}\otimes_S R\otimes_S M^{\otimes_S(n-j-2)},
 \qquad W_0=S,\quad W_1=M,
\]
where the intersection is inside $M^{\otimes_S n}$.

The relative quadratic complex has $K_n=A\otimes_S W_n$ in homological degree $n$, augmented to $S$, with differential
\[
 d(a\otimes m_1\otimes\cdots\otimes m_n)
 =am_1\otimes m_2\otimes\cdots\otimes m_n.
\]
The containment $W_n\subseteq M\otimes_S W_{n-1}$ makes this map well-defined.
For $n\ge2$, the containment $W_n\subseteq R\otimes_S M^{\otimes_S(n-2)}$ gives $d^2=0$.
Each $K_n$ is projective and generated in weight $n$.
Call $A$ Koszul if $S$ admits a graded projective resolution whose term in homological degree $n$ is generated in weight $n$.

\begin{theorem}[Relative quadratic criterion]\label{cbc:quiver:koszul}
The algebra $A$ is Koszul if and only if $K_\bullet\to S$ is exact.
When it is exact, the subspace
\[
 C_R=\bigoplus_{n\ge0}s^nW_n\subseteq B_SA
\]
is a dg $S$-coring with zero differential, and its inclusion is a quasi-isomorphism.
Here $s^nW_n$ denotes the image of $W_n$ in $(sM)^{\otimes_S n}$, in cohomological degree $-n$.
\end{theorem}
\begin{proof}
Exactness of $K_\bullet\to S$ gives the required linear resolution.
For the converse, write $A_+$ for the positive-weight ideal.
A bounded-below graded syzygy $Z$ has a projective surjection obtained by choosing an $S$-linear section of $Z/A_+Z$ and mapping
$A\otimes_S(Z/A_+Z)$ onto $Z$.
Surjectivity follows by induction on weight.
Choosing sections successively constructs a minimal graded resolution, whose differentials have image in $A_+$ times the next term.
The choices can preserve the right vertex action, since $S\otimes_k S^{\mathrm{op}}$ is also semisimple.

After tensoring with the right augmentation module $S$, all minimal differentials vanish.
Its generator spaces are therefore $\operatorname{Tor}^A_n(S,S)$.
A linear resolution has the same property and puts this Tor group entirely in weight $n$.
Consequently, if $A$ is Koszul, the $n$th minimal generator space is concentrated in weight $n$.

Identify the minimal resolution inductively with $K_\bullet$.
In degree one, its generators are $M$, and its differential is multiplication.
If the identification holds through degree $n-1$, its next syzygy is generated in weight $n$.
That weight lies in $M\otimes_S W_{n-1}$, and multiplication of the first two factors has kernel
\[
 (M\otimes_S W_{n-1})\cap
 (R\otimes_S M^{\otimes_S(n-2)})=W_n.
\]
Thus $A\otimes_S W_n$ maps onto the entire syzygy, with precisely the displayed differential.

This proves exactness in every homological degree.
It is the relative quadratic criterion of Beilinson--Ginzburg--Soergel~\cite[Theorem~2.6.1]{BGS1996}.

Each adjacent pair in $W_n$ belongs to the corresponding copy of $R$, so every bar merger vanishes on $s^nW_n$.
Deconcatenation at a cut $n=p+q$ lands in $s^pW_p\otimes_S s^qW_q$.
Indeed, the defining intersection imposes all relations internal to either block, and tensoring over the semisimple algebra $S$ preserves these intersections.
Hence $C_R$ is the stated subcoring.

Finally, bar length $n$ and weight $n$ force every input to lie in $M$.
In that bidegree the cycle space is exactly $s^nW_n$, because the different possible mergers have distinct input-weight patterns.
There are no incoming boundaries of that weight.
The linear resolution proves that all Tor groups off this diagonal vanish.
The inclusion of $C_R$ therefore induces an isomorphism on every homology group.
\end{proof}

The relation quiver $1\xrightarrow{a}2\xrightarrow{b}3$, with $ba=0$, illustrates the last assertion in Chapter~\ref{chap:quiver}.
Its class $[b|a]$ has nonzero reduced coproduct $sb\otimes_Ssa$.
That coproduct retains the composable pair even though its product in $A$ is zero.

\section{Geometric realization}

The vertex algebra and arrow bimodules above belong to ordinary coefficient complexes.
For a boundary realization, the vertex labels must act on boundary objects, and the arrows must represent fields changing one boundary condition into another.
Their composition must realize the path product.

A realization over $X$ requires factorization coefficients carrying these vertex actions, a compatible transverse product and augmentation, and the full relative bar comparison.
The signed contraction applies in a coefficient model admitting the displayed balanced tensors and direct sums.
Its geometric use requires those constructions and their comparison maps.

A second associative chiral operation requires further maps on composable arrow coefficients and their mixed coherences with transverse composition.
An operation on each arrow coefficient alone does not provide the maps between different vertex corners.
These are the additional structures needed to turn the relative algebra calculation into a doubly associative quiver construction.

\chapter{Weyl and Clifford deformations}\label{cbc:weyl:chapter}

A constant in a quadratic relation can preserve the ordered monomial basis while preventing the generators from acting by zero on a one-dimensional module.
The relation $yx-xy=\hbar$ exhibits both phenomena.
Its constant term survives on the dual side as curvature.
We first prove the basis theorem, then construct the curved coalgebra and the cobar map that recover the deformed algebra.

\section{Polynomial families and ordered bases}
Fix a field $k$, a finite-dimensional vector space $V$, and the polynomial ring $R_\hbar=k[\hbar]$.
All generators have cohomological degree zero.
For an alternating bilinear form $\omega$ on $V$, define
\[
 W_\hbar(V,\omega)
 =T_{R_\hbar}(V\otimes_k R_\hbar)/
   (uv-vu-\hbar\omega(u,v)).
\]
Alternating means $\omega(v,v)=0$, including in characteristic two.
No nondegeneracy assumption is imposed.

For the Clifford family, choose a quadratic form $q:V\to k$.
Thus $q(av)=a^2q(v)$ and
$\beta(u,v)=q(u+v)-q(u)-q(v)$ is bilinear.
Define
\[
 \operatorname{Cl}_\hbar(V,q)
 =T_{R_\hbar}(V\otimes_k R_\hbar)/(v^2-\hbar q(v):v\in V).
\]
Its relations imply $uv+vu=\hbar\beta(u,v)$.
If $2$ is invertible, the latter relations recover the former with $q(v)=\beta(v,v)/2$.

The Clifford generators also have odd internal parity.
That parity is distinct from the cohomological grading used in the bar below.

Choose an ordered basis $x_1,\ldots,x_d$ of $V$.
Write $q_i=q(x_i)$, $\beta_{ji}=\beta(x_j,x_i)$ and $\omega_{ji}=\omega(x_j,x_i)$.
Filter each algebra increasingly by the number of generators, with $\hbar$ in filtration zero.

\begin{theorem}[Filtered PBW theorem]\label{cbc:weyl:pbw}
The ordered monomials
\[
 x_1^{a_1}\cdots x_d^{a_d},\qquad
 \begin{cases}
 a_i\geq0&\text{for }W_\hbar(V,\omega),\\
 a_i\in\{0,1\}&\text{for }\operatorname{Cl}_\hbar(V,q),
 \end{cases}
\]
are $R_\hbar$-bases.
Consequently their associated graded algebras for the word filtration are
$\operatorname{Sym}_k(V)\otimes_k R_\hbar$ and
$\Lambda_k(V)\otimes_k R_\hbar$, respectively.
Both families are flat over $R_\hbar$ in every characteristic.
\end{theorem}
\begin{proof}
Move a larger index to the right of a smaller one, using the defining relation.
Each correction has two fewer generators.
The reordered term has fewer inversions.
In the Clifford case, replace $x_i^2$ by $\hbar q_i$.
Induction on length and then inversions proves spanning.

For independence in the Weyl case, act on $R_\hbar[t_1,\ldots,t_d]$ by
\[
 X_i=t_i+\hbar\sum_{r<i}\omega_{ir}\frac{\partial}{\partial t_r}.
\]
The multiplication and differentiation operators satisfy
$[X_i,X_j]=\hbar\omega(x_i,x_j)$.
An ordered product applied to $1$ gives
\[
 X_1^{a_1}\cdots X_d^{a_d}(1)=t_1^{a_1}\cdots t_d^{a_d}.
\]
Indeed, when $X_i$ acts, only variables with index at least $i$ are present.
Its derivative terms therefore vanish.
The resulting polynomial monomials are independent over $R_\hbar$.

For the Clifford case, let $B$ be the bilinear form defined by
\[
 B(x_i,x_i)=q_i,\qquad B(x_j,x_i)=\beta_{ji}\ (j>i),\qquad
 B(x_i,x_j)=0\ (i<j).
\]
Then $B(v,v)=q(v)$.

On $\Lambda_k(V)\otimes_k R_\hbar$, let $L_v$ denote exterior multiplication by $v$ and let $\iota_v$ be contraction by $B(v,-)$.
The operators $T_v=L_v+\hbar\iota_v$ satisfy
\[
 T_v^2=\hbar q(v),\qquad T_uT_v+T_vT_u=\hbar\beta(u,v).
\]
These identities follow from $L_v^2=\iota_v^2=0$ and
$\iota_uL_v+L_v\iota_u=B(u,v)\id$.

For $i_1<\cdots<i_r$, the triangular choice of $B$ gives
\[
 T_{x_{i_1}}\cdots T_{x_{i_r}}(1)
       =x_{i_1}\wedge\cdots\wedge x_{i_r}.
\]
These vectors form the exterior basis, proving independence.
The same ordered bases identify every quotient of the word filtration and prove the associated graded assertions.
\end{proof}

In characteristic two, the square relation is indispensable.
For one generator, $2x^2=0$ imposes no relation, whereas $x^2=\hbar$ gives a free rank-two Clifford family.
If one instead prescribes $2x^2=\hbar$, the resulting algebra has $\hbar=0$ and is not flat over $k[\hbar]$.
Thus a symmetric anticommutator without the quadratic data does not state the characteristic-two PBW theorem.

\section{Specialization, Rees construction and completion}
Write $A_\hbar$ for either polynomial family and $A_a=A_\hbar/(\hbar-a)$ for $a\in k$.
The fiber $A_0$ is the symmetric or exterior algebra, and the same ordered monomials form a basis of every $A_a$.
For the word filtration $F$ on $A_1$, the Rees algebra is
\[
 \mathcal R_F(A_1)=\bigoplus_{m\geq0}F_mA_1\,z^m\subseteq A_1[z].
\]
It is generated by $z$ and $X_i=zx_i$.
Its constant terms are quadratic in $z$: in the first Weyl example,
$YX-XY=z^2$.
The ordered bases prove the isomorphism
\[
 A_\hbar\otimes_{k[\hbar],\,\hbar\mapsto z^2}k[z]
       \ \xrightarrow{\ \sim\ }\ \mathcal R_F(A_1),
 \qquad x_i\longmapsto zx_i.
\]

The formal family is instead the inverse limit
\[
 \widehat A_\hbar=\varprojlim_N A_\hbar/\hbar^N A_\hbar.
\]
The ordered bases identify its underlying module with
$\operatorname{Sym}(V)[[\hbar]]$ in the Weyl case and
$\Lambda(V)[[\hbar]]$ in the Clifford case.
In the first expression, every coefficient of $\hbar^n$ is a polynomial, with no uniform bound on its degree.

For example, $\sum_{n\geq0}\hbar^n x_1^n$ belongs to the completion but to no bounded word-filtration piece.
Thus that bounded-word filtration is not exhaustive on a nontrivial completed Weyl family.

Nor is there a unital specialization $k[[\hbar]]\to k$ sending $\hbar$ to $1$, since $1-\hbar$ is invertible.
The fiber at one belongs to the polynomial family.

\section{The curved quadratic dual}
Let $A_0=T_k(V)/(\mathcal R)$ denote the homogeneous fiber.
Its quadratic coalgebra has
\[
 C_0=k,\quad C_1=sV,\quad
 C_n=s^n\!\bigcap_{i+j=n-2}
       V^{\otimes i}\otimes\mathcal R\otimes V^{\otimes j}
       \quad(n\geq2).
\]
Here $|sv|=-1$.
The notation sends a tensor $v_1\otimes\cdots\otimes v_n$ to $[v_1|\cdots|v_n]$ without an additional sign.

The coalgebra $C=\bigoplus_n C_n$ has degree $-n$ and weight $n$ on $C_n$, zero differential and deconcatenation coproduct.
The intersection argument preceding Theorem~\ref{cbc:quadratic:criterion} proves that deconcatenation restricts to $C$.
Extend scalars to $R_\hbar$ for the remainder of the construction.

Write a deformed relation as $r-\hbar\kappa(r)$ for $r\in\mathcal R$.
Define a degree-two functional $\theta:C\otimes_k R_\hbar\to R_\hbar$ by
\[
 \theta(s^2r)=\hbar\kappa(r),\qquad \theta(C_n)=0\quad(n\ne2).
\]
Let $E=\bigoplus_n\operatorname{Hom}_k(C_n,k)\otimes_k R_\hbar$, graded by $n$.
Its multiplication is the signed convolution
\[
 (fg)(u)=\sum_{\Delta u=\sum u'\otimes u''}
          (-1)^{|g||u'|}f(u')g(u'').
\]
The functional $\theta$ corresponds to an element $c\in E^2$.

Equivalently, it determines the zero-input operation $m_0:R_\hbar\to E$, $m_0(1)=c$.
This operation and the functional $\theta:C\otimes R_\hbar\to R_\hbar$ have different sources and targets.

Let $\xi_i$ be dual to $[x_i]$.
The two curved duals are
\begin{align}
 E_W&=\Lambda_{R_\hbar}(\xi_1,\ldots,\xi_d),&
 c_W&=\hbar\sum_{i<j}\omega_{ji}\xi_i\xi_j,\label{eq:cbc:weyl:curvature}\\
 E_{\mathrm{Cl}}&=R_\hbar[\xi_1,\ldots,\xi_d],&
 c_{\mathrm{Cl}}&=-\hbar\left(\sum_i q_i\xi_i^2+
                        \sum_{i<j}\beta_{ji}\xi_i\xi_j\right).
 \label{eq:cbc:clifford:curvature}
\end{align}
Every $\xi_i$ has cohomological degree one.
The polynomial algebra in the second line has ordinary commuting generators.
It is not asserted to be graded commutative.
Both dual differentials are zero.

To check the signs, set
$\rho_{ji}=[x_j|x_i]-[x_i|x_j]$ in the Weyl case.
The convolution gives $(\xi_i\xi_j)(\rho_{ji})=1$ for $i<j$.
For Clifford relations, use $\rho_{ii}=[x_i|x_i]$ and
$\rho_{ji}=[x_j|x_i]+[x_i|x_j]$.
Then $(\xi_i^2)(\rho_{ii})=-1$ and $(\xi_i\xi_j)(\rho_{ji})=-1$.
These evaluations give exactly $\theta$.

The algebra presentations follow by orthogonal quadratic relations and finite-dimensional duality in each weight.
The coalgebra $C_W$ consists of alternating tensors.
For $C_{\mathrm{Cl}}$, sum the distinct permutations of each multiset to obtain a basis.
The latter sums carry no factorial denominators.

In each case $c$ is central, so $(E,0,c)$ satisfies the curved dg identities
\[
 d^2=[c,-]=0,\qquad dc=0.
\]
This is nonhomogeneous quadratic duality in the sense of~\cite[\S\S2.1--2.2]{PositselskiCurvature2014}, with the suspension signs fixed above.
The underlying complex $(E,0)$ has ordinary cohomology $E$.
Forgetting $c$ discards the constant terms of the deformation.

\section{A cobar resolution with a scalar term}
Put $t=s^{-1}$, so $|tu|=|u|+1$, and $\bar C=\bigoplus_{n>0}C_n$.
On the unital tensor algebra $T_{R_\hbar}(t\bar C\otimes_k R_\hbar)$, define the degree-one derivation
\begin{equation}\label{eq:cbc:weyl:curved-cobar}
 D(tu)=-\sum_{\bar\Delta u=\sum u'\otimes u''}
                  (-1)^{|u'|}(tu')(tu'')-\theta(u).
\end{equation}
Denote this dg algebra, once the square-zero identity is proved, by $\Omega_\theta C$.
It has no augmentation sending all $t\bar C$ to zero when $\theta\ne0$.

\begin{theorem}\label{cbc:weyl:cobar}
The derivation in~\eqref{eq:cbc:weyl:curved-cobar} satisfies $D^2=0$.
The assignments
\[
 p:\Omega_\theta C\longrightarrow A_\hbar,
 \qquad p(t[v])=v,\qquad p(tC_n)=0\quad(n>1),
\]
define a quasi-isomorphism of unital dg $R_\hbar$-algebras.
The degreewise $\hbar$-adic completion of $p$ is also a quasi-isomorphism.
\end{theorem}
\begin{proof}
The quadratic part of $D$ squares to zero by coassociativity, as in Proposition~\ref{cbc:cobar:square}.
The remaining terms on a generator are
\[
 D^2(tu)=t\bigl((\theta\otimes\id-\id\otimes\theta)\Delta u\bigr).
\]
Centrality of $c$ makes the expression in parentheses zero.
Indeed, pairing it with any homogeneous element of $E$ gives the commutator with $c$.
The pairings in each weight are perfect.
Thus $D^2=0$ on generators and hence everywhere.

For $u=s^2r$, the image of $D(tu)$ is $r-\hbar\kappa(r)$, which vanishes in $A_\hbar$.
For $u\in C_n$ with $n>2$, every reduced cut has a factor of weight greater than one, and $\theta(u)=0$.
Its image also vanishes.
The weight-one generators are closed, so $p$ is a dg algebra map.

We next verify the homogeneous Koszul comparison.
For the symmetric algebra, tensor the one-variable resolutions
$0\to k[x_i]\xrightarrow{x_i}k[x_i]\to k\to0$ over $k$.
They give a linear free resolution, and hence the diagonal criterion applies.

For the exterior algebra, use the complex
\[
 K_n=\Lambda(V)\otimes_k\Gamma^n(V),\qquad
 d(a\otimes u^{[\alpha]})=
       \sum_{i:\alpha_i>0}a x_i\otimes u^{[\alpha-e_i]}.
\]
Here $u^{[\alpha]}$ are divided-power basis symbols with $|\alpha|=n$.
No division is used in the differential.
The exterior relations give $d^2=0$.

For an ordered exterior monomial $x_I$, fix the total multidegree
$\nu=\mathbf1_I+\alpha$.
If $\nu\ne0$, choose $j$ with $\nu_j>0$.

The map
\[
 H(x_I\otimes u^{[\alpha]})=
 \begin{cases}
 (-1)^{\#\{i\in I:i>j\}}x_{I\setminus\{j\}}
           \otimes u^{[\alpha+e_j]},&j\in I,\\
 0,&j\notin I
 \end{cases}
\]
satisfies $dH+Hd=\id$ on that multidegree.
The $j$-term gives the identity.
Every other term cancels because the two exterior insertions have opposite signs.
Multidegree zero supplies the augmentation to $k$.
This proves a linear free resolution in every characteristic.

Theorem~\ref{cbc:quadratic:criterion} now gives
$\Omega C\simeq A_0$ for both homogeneous algebras.
Tensoring over the field $k$ extends this comparison to $R_\hbar$.
Filter $\Omega_\theta C$ by total $C$-weight, with $\hbar$ in filtration zero, and filter $A_\hbar$ by word length.
The quadratic term preserves this filtration and the scalar term lowers it by two.
By the PBW theorem, the associated graded map is precisely
$\Omega C\otimes_k R_\hbar\to A_0\otimes_k R_\hbar$.

Its cone is acyclic in every weight.
For a cycle in the filtered cone, remove its highest-weight term by a boundary and continue in lower weights.
Every element has finite weight and the filtration starts at zero, so this procedure terminates.
The cone of $p$ is therefore acyclic.

The same argument works modulo $\hbar^N$ for every $N$.
The transition kernels of these cones are copies of the cone modulo $\hbar$, and are acyclic.
Compatible primitives can thus be lifted recursively from level $N$ to level $N+1$.
Their inverse limit proves acyclicity of the completed cone.
\end{proof}

For $V=kx\oplus ky$ with $\omega(y,x)=1$, put
$\rho=[y|x]-[x|y]$.
The formulas give
\[
 c=\hbar\xi\eta,\qquad D(t\rho)=yx-xy-\hbar.
\]
The scalar term is required for $p$ to respect the differential.
Curvature in this calculation belongs to a specified ordinary quadratic dual.
A geometric anomaly curvature belongs to its own coefficient complex.
Comparing them requires a map carrying the curvature, differential and products, together with the relevant factorization and collision operations.

\chapter{Associative chiral algebras on the formal disc}\label{cbc:disc:chapter}

The value of a field at the origin recovers a multiplication when that field has no pole.
Translation then recovers the entire field from this multiplication.
We prove this algebraic correspondence, construct a complete braided example, and state the geometric realization problem that remains for the resulting fields.

\section{Regular fields and differential algebras}
Work over $\C$ in a fixed formal coordinate $z$.
A nonlocal vertex algebra $(V,Y,\mathbf1)$, in the sense of Definition~\ref{cbc:two:nonlocal}, is \emph{regular} if
\[
 Y(a,z)b\in V[[z]]\qquad(a,b\in V).
\]
Define its product and translation operator by
\[
 ab=Y(a,z)b\big|_{z=0},\qquad
 Ta=\partial_zY(a,z)\mathbf1\big|_{z=0}.
\]
A morphism of nonlocal vertex algebras is a linear map preserving the vacuum and satisfying
$f(Y(a,z)b)=Y'(f(a),z)f(b)$.

\begin{theorem}[Reconstruction of regular fields]\label{cbc:disc:regular}
The product above makes $V$ a unital associative algebra, $T$ is a derivation, and
\begin{equation}\label{eq:cbc:disc:regular}
 Y(a,z)b=(e^{zT}a)b.
\end{equation}
Conversely, this formula gives a regular nonlocal vertex algebra for every unital associative $\C$-algebra with derivation $T$.
These assignments are inverse functors between regular nonlocal vertex algebras and unital associative differential algebras.
Morphisms on the latter side preserve multiplication and commute with $T$.
\end{theorem}
\begin{proof}
Regularity puts both sides of weak associativity in $V[[z,w]]$.
Multiplication by $(z+w)^N$ is injective there: the change of variables $u=z+w$ reduces it to multiplication by $u^N$.
Thus weak associativity becomes the exact identity
\begin{equation}\label{eq:cbc:disc:regular-assoc}
 Y(a,z+w)Y(b,w)c=Y(Y(a,z)b,w)c.
\end{equation}
No choice of Laurent expansion is needed for this power-series identity.
Setting $z=w=0$ proves $a(bc)=(ab)c$.
Vacuum and creation give $\mathbf1a=a\mathbf1=a$.

Put $F_z(a)=Y(a,z)\mathbf1$, so $F_0=\id$.
In~\eqref{eq:cbc:disc:regular-assoc}, first set $w=0$ and $b=\mathbf1$.
Then set $b=c=\mathbf1$ without specializing the variables.
These substitutions give
\[
 Y(a,z)c=F_z(a)c,\qquad F_{z+w}(a)=F_w(F_z(a)).
\]
Differentiating the second identity in $w$ at zero gives
$\partial_zF_z=TF_z$.
The initial value $F_0=\id$ determines every coefficient, hence
$F_z=e^{zT}$.

Finally set $z=0$ and $c=\mathbf1$ in~\eqref{eq:cbc:disc:regular-assoc}.
It gives $F_w(ab)=F_w(a)F_w(b)$.
The coefficient of $w$ is the derivation identity
$T(ab)=(Ta)b+a(Tb)$, while $F_w(\mathbf1)=\mathbf1$ gives $T\mathbf1=0$.
This proves~\eqref{eq:cbc:disc:regular}.

Conversely, a derivation satisfies the iterated Leibniz rule, so $e^{wT}$ preserves products.
For the field map~\eqref{eq:cbc:disc:regular}, both sides of~\eqref{eq:cbc:disc:regular-assoc} equal
\[
 (e^{(z+w)T}a)(e^{wT}b)c.
\]
Vacuum and creation follow from $T\mathbf1=0$.
The field is regular, and its constant term and vacuum derivative recover the given product and $T$.

A field morphism preserves those two coefficients.
A unital algebra map commuting with $T$ preserves the exponential formula.
Thus the assignments are inverse on objects and morphisms.
\end{proof}

For $V=M_2(\C[t])$ and $T=\partial_t$, the fields are
\[
 Y(a,z)b=a(t+z)b(t).
\]
Let $E_{ij}$ be the matrix units and let $L_a$ denote left multiplication by $a$.
The constant matrix fields need not commute:
\[
 [Y(E_{12},z),Y(E_{21},w)]=L_{E_{11}-E_{22}},
\]
Regularity therefore retains ordinary associative algebras with noncommuting multiplication.

The same correspondence has a complete form.
Let $R=\C[[\hbar]]$ and $V=V_0[[\hbar]]$, where $V_0$ is a complex vector space.
Use the completed tensor product
$V\widehat\otimes_R V=(V_0\otimes_{\C}V_0)[[\hbar]]$.
A regular topological field map is a continuous $R$-linear map
$Y:V\widehat\otimes_R V\to V[[z]]$.
Require vacuum, creation and weak associativity modulo every $\hbar^m$.

\begin{corollary}\label{cbc:disc:complete-regular}
These field systems are equivalent, by~\eqref{eq:cbc:disc:regular}, to continuous unital associative products on $V$ with continuous $R$-linear derivation $T$.
The morphisms are continuous and $R$-linear on both sides.
\end{corollary}
\begin{proof}
Apply the preceding coefficient argument to $V/\hbar^mV$ for every $m$.
These quotients are complex vector spaces, so cancellation and division by factorials remain valid.
The recovered products, operators and identities commute with reduction.
Their inverse limits give the claimed correspondence, and the exponential is defined coefficientwise in $z$.
\end{proof}

\section{A braiding defined on every state}
The complete quantum torus gives an example where the multiplication and the exchange operator can both be written explicitly.
Keep $R=\C[[\hbar]]$ and put $q=e^\hbar$ as a formal power series.
Let $\Gamma=\mathbb Z^2$ and let $e_g$, for $g=(g_1,g_2)\in\Gamma$, be the basis of $\C[\Gamma]$.
On $V=\C[\Gamma][[\hbar]]$, define
\begin{equation}\label{eq:cbc:disc:torus-product}
 e_g e_j=q^{g_2j_1}e_{g+j},\qquad \mathbf1=e_0,
 \qquad T=0,\qquad Y(a,z)b=ab.
\end{equation}
Every coefficient of $\hbar$ has finite support in $\Gamma$.
The formula therefore extends continuously to $V\widehat\otimes_R V$.

Define the alternating integer form
$\omega(g,j)=g_1j_2-g_2j_1$ and the operator
\begin{equation}\label{eq:cbc:disc:torus-braiding}
 S(e_g\otimes e_j)=q^{\omega(g,j)}e_g\otimes e_j.
\end{equation}
It is continuous on the completed tensor product and independent of the spectral variable.
In particular, $S=1+O(\hbar)$.

\begin{proposition}\label{cbc:disc:torus}
The data~\eqref{eq:cbc:disc:torus-product}--\eqref{eq:cbc:disc:torus-braiding} form a topological quantum vertex algebra in the convention of~\cite[Definitions 3.5 and 3.12]{DSGKQVA}.
The braiding satisfies exact unitarity, the quantum Yang--Baxter equation and both hexagon identities with the multiplication.
The algebra $(V,\cdot)$ has no unital $R$-algebra homomorphism to $R$.
\end{proposition}
\begin{proof}
Associativity of the product is the integer identity
\[
 g_2j_1+(g_2+j_2)k_1
       =j_2k_1+g_2(j_1+k_1).
\]
The vector $e_0$ is a two-sided unit.
Corollary~\ref{cbc:disc:complete-regular} gives the vacuum, translation and weak associativity axioms for $Y$.

The locality condition~\eqref{eq:cbc:two:top-local} holds with exponent zero because
\[
 q^{\omega(g,j)}(e_g e_j)e_k=(e_j e_g)e_k.
\]
Indeed, $\omega(g,j)+g_2j_1=j_2g_1$.
This is an identity on every triple of basis states.
Linearity modulo $\hbar^m$ and passage to the inverse limit prove it for all states.
These checks give the asserted quantum vertex algebra.

Write $S^{ij}$ for the operator on tensor factors $i,j$ and $S^{21}$ for its conjugate by interchange.
Alternation gives $S^{21}S=\id$.
The operators $S^{12},S^{13},S^{23}$ are diagonal on the same basis, so
\[
 S^{12}S^{13}S^{23}=S^{23}S^{13}S^{12}.
\]

Let $m$ denote the multiplication~\eqref{eq:cbc:disc:torus-product}.
The two bicharacter identities
\[
 \omega(g+j,k)=\omega(g,k)+\omega(j,k),\qquad
 \omega(g,j+k)=\omega(g,j)+\omega(g,k)
\]
give the exact hexagons
\[
 S(m\otimes1)=(m\otimes1)S^{23}S^{13},\qquad
 S(1\otimes m)=(1\otimes m)S^{12}S^{13}.
\]
The unit normalization follows from $\omega(0,g)=\omega(g,0)=0$.
All translation identities are zero since $T=0$ and $S$ is constant.

Put $X=e_{(1,0)}$ and $Y_0=e_{(0,1)}$.
They are invertible, with inverses $e_{(-1,0)}$ and $e_{(0,-1)}$, and
$Y_0X=qXY_0$.
A unital $R$-algebra map $\varepsilon:V\to R$ would therefore satisfy
\[
 (1-q)\varepsilon(X)\varepsilon(Y_0)=0.
\]
The two images are units, forcing $q=1$, contrary to $q=e^\hbar$ in $R$.
\end{proof}

The sign in the braiding is tied to the fixed field multiplication.
Replacing $S$ by $S^{-1}$ preserves its unit normalization, unitarity, Yang--Baxter equation and both hexagons.
On the same states $X,Y_0,\mathbf1$, the difference in the locality equation becomes
\[
 (q^{-1}-q)e_{(1,1)}=-2\hbar e_{(1,1)}\pmod{\hbar^2}.
\]
Multiplication by any power of $z-w$ cannot kill this constant nonzero distribution modulo $\hbar^2$.
Thus $S^{-1}$ is incompatible with locality for these fields.
The vacuum $e_0$ coexists with the absence of a scalar algebra augmentation.

\section{From generating fields to an exchange operator}
For singular fields, the Laurent expansions and their substitution maps are part of the data.
The field space is $\operatorname{Hom}(V,V((z)))$ in the ordinary convention of Definition~\ref{cbc:two:nonlocal}, with a pole bound for each input vector.
The completed convention instead uses $V_\hbar((z))$ as defined before~\eqref{eq:cbc:two:top-local}.
The regular correspondence proved above uses the smaller space $V[[z]]$.

Return to an ordinary nonlocal vertex algebra $V$ over $\C$.
Let $U\subseteq V$ be a finite-dimensional subspace generating $V$ from the vacuum by coefficients of its fields, and choose a basis $a_i$ of $U$.
A spectral operator $R(u)\in\operatorname{End}(U\otimes U)\otimes\C((u))$ may prescribe an exchange relation
\[
 Y(a_i,z_1)Y(a_j,z_2)
 =\sum_{k,l}R_{ij}^{kl}(z_1-z_2)
          Y(a_l,z_2)Y(a_k,z_1).
\]
Expand the Laurent coefficients in the specified regions and clear the required poles by a power of $z_1-z_2$.
This is $S$-locality as in~\eqref{eq:cbc:two:s-local}, with the coefficient functions evaluated at the negative of its difference variable.

To state the hypothesis that recovers an operator on all states, put
$K_3=\C((z_1))((z_2))((z_3))$ and define
\begin{align*}
 \pi_3:V^{\otimes3}\otimes_{\C}K_3
 &\longrightarrow
 \operatorname{Hom}_{\C}\bigl(V,V((z_1))((z_2))((z_3))\bigr),\\
 \pi_3(a\otimes b\otimes c\otimes f)(d)
 &=fY(a,z_1)Y(b,z_2)Y(c,z_3)d.
\end{align*}

\begin{proposition}[Generating exchange relations]\label{cbc:disc:generators}
If the fields of $U$ are $S$-local, then $V$ is a weak quantum vertex algebra.
If, in addition, $\pi_3$ is injective, the exchange relations determine a unique operator
\[
 \mathcal S(u):V\otimes V\longrightarrow
                    V\otimes V\otimes\C((u)).
\]
It satisfies unitarity $\mathcal S^{21}(-u)\mathcal S(u)=\id$ and
\[
 \mathcal S^{12}(u)\mathcal S^{13}(u+v)\mathcal S^{23}(v)
   =\mathcal S^{23}(v)\mathcal S^{13}(u+v)\mathcal S^{12}(u),
\]
with the Laurent expansions in the indicated order.
With that operator, $V$ is a quantum vertex algebra in the convention of~\cite[Definition 5.10 and Theorem 5.11]{LiNonlocal}.
\end{proposition}
\begin{proof}
Write $\mathcal E(V)=\operatorname{Hom}_{\C}(V,V((z)))$.
Call fields $A(z),B(z)$ compatible if some $N\geq0$ satisfies
\[
 (z_1-z)^N A(z_1)B(z)
       \in\operatorname{Hom}_{\C}(V,V[[z_1,z]][z_1^{-1},z^{-1}]).
\]
For such an $N$, their field products are the coefficients of
\[
 Y_{\mathcal E}(A(z),x)B(z)
   =x^{-N}\left((z_1-z)^NA(z_1)B(z)\right)\big|_{z_1=z+x},
\]
where substitution expands in nonnegative powers of $x$ relative to $z$.
The expression is independent of a larger clearing exponent.
Li's construction identifies the subspace $H\subseteq\mathcal E(V)$ generated by $\id_V$ and the fields of $U$ under these coefficients as a weak quantum vertex algebra~\cite[Definition 2.11 and Theorem 5.8]{LiNonlocal}.

The state--field map $\Phi:V\to\mathcal E(V)$, $\Phi(a)=Y(a,z)$, is injective by creation.
The self-module identity gives
\[
 \Phi(Y(a,x)b)=Y_{\mathcal E}(\Phi(a),x)\Phi(b)
\]
by~\cite[Proposition 2.22]{LiNonlocal}.
Thus $\Phi$ preserves field products and identifies $V$ with $H$, because $U$ generates $V$.
This proves the first assertion.

For the second assertion, apply~\cite[Theorems 4.8 and 5.11]{LiNonlocal} to the faithful self-module $V$.
Injectivity of $\pi_3$ also gives injectivity for two inputs by inserting the vacuum in the third slot.
Consequently the coefficients of an exchange are unique.
Reversing twice gives unitarity, and comparing the two reversals of three fields gives the Yang--Baxter equation.
The cited theorem controls the Laurent expansions in these comparisons and gives the stated operator.
\end{proof}

The tensor injectivity hypothesis is stronger than injectivity of $a\mapsto Y(a,z)$.
For the constant fields of $M_2(\C)$, the nonzero tensor $E_{11}\otimes E_{22}$ acts by the zero product, whereas the state--field map is injective.
Injectivity of $\pi_3$ is a sufficient reconstruction hypothesis.
An explicit operator can also be supplied directly, as for the quantum torus above.
The ordinary convention in this proposition does not impose a formal deformation or the normalization $S=1+O(\hbar)$.

\section{The geometric realization problem}

\begin{problem}[Formal-disc realization of fields]\label{cbc:disc:realization}
Fix a nonlocal vertex algebra with covariance, topology and completion.
Construct coefficient objects on every finite power of the formal disc.
Construct their collision maps, relabeling maps, factorization comparisons and associative chiral operations with all substitution coherences.
Give a comparison whose field expansions recover $Y$, its vacuum and its weak associativity identity~\eqref{eq:cbc:two:weak-assoc}.
For a quantum vertex algebra, also realize the full exchange operator and its compatibility with these operations.

Determine the coefficient, support and continuity hypotheses under which these constructions are functorial, and prove both inverse comparisons between geometric objects and fields.
Specify any scalar augmentation separately from the vacuum.
The full field class remains to be treated, including singular fields outside the regular correspondence.
\end{problem}

For a given geometric realization satisfying the stated comparison hypotheses, Proposition~\ref{cbc:two:field-realization} proves the forward field implication.
The converse requires the coefficient objects and collision-compatible maps in the problem.
Neither the regular algebraic correspondence nor the explicit quantum-torus braiding supplies those geometric constructions.

\chapter{Zamolodchikov--Faddeev exchange}
\label{cbc:exchange:chapter}
An exchange relation prescribes how two generators change when their order is
reversed.  For three generators there are two ways to reverse the order.
Their agreement is the Yang--Baxter equation.  To deduce a basis from this
agreement, one must also specify which exchanges are allowed and what
happens when a spectral label is repeated.

\section{Ordered colors and normal forms}
Let $B$ be a commutative ring and $1<\cdots<m$ a finite set of colors.
For each color~$a$, let $V_a$ be a finite free $B$-module.
All tensor products in this section are over~$B$.

For each $a>b$, choose a $B$-linear map
\[
 c_{ab}:V_a\otimes V_b\longrightarrow V_b\otimes V_a.
\]
Write $T_B(V)$ for the unital tensor algebra of a $B$-module~$V$.
Define the \emph{oriented exchange algebra}
\begin{equation}\label{cbc:exchange:algebra}
 A(c)=T_B\Bigl(\bigoplus_{a=1}^m V_a\Bigr)
       \big/\bigl\langle x\otimes y-c_{ab}(x\otimes y):
                         x\in V_a,\ y\in V_b,\ a>b\bigr\rangle.
\end{equation}

Thus an adjacent descending pair may be exchanged.  No relation is imposed
between two generators of the same color.  Choose a basis of each~$V_a$.
A word is \emph{ordered} if its colors are nondecreasing.  Its generators
within any one color retain their order.

\begin{theorem}[The ordered-word criterion]\label{cbc:exchange:normal}
Multiplication induces an isomorphism of $B$-modules
\begin{equation}\label{cbc:exchange:ordered-module}
 T_B(V_1)\otimes\cdots\otimes T_B(V_m)\xrightarrow{\ \sim\ }A(c)
\end{equation}
if and only if the following identity holds for every $a>b>c$:

\begin{equation}\label{cbc:exchange:colored-braid}
 (c_{bc}\otimes1)(1\otimes c_{ac})(c_{ab}\otimes1)
 = (1\otimes c_{ab})(c_{ac}\otimes1)(1\otimes c_{bc}).
\end{equation}
Both sides map $V_a\otimes V_b\otimes V_c$ to $V_c\otimes V_b\otimes V_a$.
Under this condition every sequence of descending adjacent exchanges
terminates at the same linear combination of ordered words.
The maps $c_{ab}$ need not be invertible.
\end{theorem}
\begin{proof}
For a word with color sequence $(a_1,\ldots,a_n)$, set
$\operatorname{inv}(a_1,\ldots,a_n)
 =\#\{(i,j):i<j,\ a_i>a_j\}$.
Every term produced by a descending adjacent exchange has inversion number
one smaller.  Consequently every reduction terminates.

We prove uniqueness of its result by induction on the inversion number.
An ordered word has itself as its normal form.  Otherwise choose a first
exchange and apply the already defined normal form linearly to its output.
Two choices of first exchange either coincide, occupy disjoint pairs, or
meet in one position.

Disjoint exchanges commute as maps on the tensor
word, so their two outputs have a common next reduction.  Overlapping
exchanges require three strictly descending colors.  Performing the next
two exchanges on each side gives the two maps in
\eqref{cbc:exchange:colored-braid}.  They have equal outputs.  In both
cases every subsequent word has smaller inversion number, so induction
identifies its normal form.  This defines a $B$-linear map $N$ from the
free tensor algebra to the free module on ordered words, independent of
all choices.

The two sides of each defining relation, in every word context, have the
same image under~$N$.  Hence $N$ descends to $A(c)$ and is a left inverse
to ordered multiplication.  Conversely, reducing a word changes it only by a defining relation.
Its normal form therefore multiplies back to the original class.  These maps are inverse, proving
\eqref{cbc:exchange:ordered-module}.

For necessity, both reductions of a descending three-generator tensor
represent the same element of $A(c)$.  Their outputs already have increasing
colors.  Injectivity in \eqref{cbc:exchange:ordered-module} identifies
them as tensors, which is \eqref{cbc:exchange:colored-braid}.
\end{proof}

Repeated colors cause no additional overlap: two adjacent reducible pairs
can meet only in a strictly decreasing triple.  In particular, the theorem
concerns all words, including repeated labels.

If reverse maps $c_{ba}:V_b\otimes V_a\to V_a\otimes V_b$ are also
prescribed, preserving the ordered basis forces $c_{ab}c_{ba}=1$.
Indeed, reversing an ordered pair and then sorting it must return that pair.
Both modules are finite free of the same rank, so a right inverse is also
a left inverse.  Equivalently,
\begin{equation}\label{cbc:exchange:reverse}
 c_{ba}c_{ab}=1_{V_a\otimes V_b}\qquad(a>b).
\end{equation}
Conversely, taking $c_{ba}=c_{ab}^{-1}$ makes the reverse relations
consequences of the oriented ones.  The oriented presentation does not require reversal.

\section{Spectral matrices and the rational solution}
Suppose now that $V_a$ is a labeled copy of one finite free module~$V$.
Let $t_1,\ldots,t_m$ be central spectral parameters and suppose the
coefficients of $R(t_a-t_b)\in\operatorname{End}_B(V\otimes V)$ lie in~$B$.
If $P(x\otimes y)=y\otimes x$, put
$c_{ab}=P R(t_a-t_b)$, with the target labeled $V_b\otimes V_a$.
In a basis $e_i^{(a)}$ this means
\[
 e_i^{(a)}e_j^{(b)}
 =\sum_{k,l}R_{ij}^{kl}(t_a-t_b)e_l^{(b)}e_k^{(a)}
 \qquad(a>b).
\]

For operators on $V^{\otimes3}$, subscripts specify the original tensor
legs.  For example $R_{13}=P_{23}(R\otimes1)P_{23}$.
Removing the common reversal of the three tensor factors from
\eqref{cbc:exchange:colored-braid} gives
\begin{equation}\label{cbc:exchange:spectral-ybe}
 \begin{split}
 R_{12}(t_a-t_b)R_{13}(t_a-t_c)R_{23}(t_b-t_c)\\
 =R_{23}(t_b-t_c)R_{13}(t_a-t_c)R_{12}(t_a-t_b).
 \end{split}
\end{equation}
This is the spectral Yang--Baxter equation on the chosen labels.
It is a rational identity when the labels are independent variables.
A single numerical specialization proves only that specialization.

Writing $R_{21}(u)=P R(u)P$, the reverse condition becomes
$R_{21}(-u)R(u)=1$.

Take a field $k$ of characteristic zero, $V=k^d$, and
\[
 B_0=k[\hbar,t_1,\ldots,t_m,\Delta^{-1}],\qquad
 \Delta=\prod_{a>b}(t_a-t_b).
\]
Extend $V$ to~$B_0$ and define
\begin{equation}\label{cbc:exchange:rational}
 R(u)=1-\frac{\hbar P}{u}.
\end{equation}
\begin{proposition}\label{cbc:exchange:rational-result}
The matrix \eqref{cbc:exchange:rational} satisfies
\eqref{cbc:exchange:spectral-ybe}.  Its oriented exchange algebra has the
ordered basis of Theorem~\ref{cbc:exchange:normal}.  Its reverse product is
\begin{equation}\label{cbc:exchange:rational-reversal}
 R_{21}(-u)R(u)=1-\frac{\hbar^2}{u^2}.
\end{equation}
After adjoining the inverses of $t_a-t_b\pm\hbar$, the matrix
\begin{equation}\label{cbc:exchange:normalized}
 \widehat R(u)=\frac{u-\hbar P}{u-\hbar}
\end{equation}
satisfies both Yang--Baxter and the reverse condition.
\end{proposition}
\begin{proof}
Put $u=t_a-t_b$ and $v=t_b-t_c$, so $t_a-t_c=u+v$.
Expand the difference between the two sides of Yang--Baxter in~$\hbar$.
The constant and linear coefficients cancel.  If
$C=P_{12}P_{13}$ and $D=P_{13}P_{12}$, the permutation identities
\[
 P_{13}P_{23}=P_{23}P_{12}=C,\qquad
 P_{12}P_{23}=P_{23}P_{13}=D
\]
make the quadratic coefficient equal to
\[
 (C-D)\left(\frac1{u(u+v)}-\frac1{uv}
                         +\frac1{v(u+v)}\right)=0.
\]
The cubic coefficient vanishes because
$P_{12}P_{13}P_{23}=P_{23}P_{13}P_{12}$.
This proves the identity over~$B_0$ in every dimension.

Since $P^2=1$, direct multiplication gives
\eqref{cbc:exchange:rational-reversal}.
Multiplying $R(u)$ by $f(u)=u/(u-\hbar)$ preserves Yang--Baxter:
both sides acquire the same product of three central factors.
Finally,
$f(u)f(-u)=u^2/(u^2-\hbar^2)$ cancels the reverse product.
\end{proof}

The unnormalized matrix can therefore be used for oriented exchange,
but not for a reversible presentation with the same free basis.
Indeed, imposing both unnormalized orientations forces
$\hbar^2 e_i^{(a)}e_j^{(b)}=0$ for $a\ne b$.
Over the generic rational field these mixed quadratic products vanish.
Over $k[t_1,\ldots,t_m,\Delta^{-1}][[\hbar]]$, with every word-degree
component completed in~$\hbar$, their reductions modulo~$\hbar$ are
nonzero ordered words.  They are thus nonzero $\hbar^2$-torsion classes.

The normalized construction is also available over this complete ring,
where $u\pm\hbar$ are units.

\section{Scalar normalization and the product}
There is an explicit comparison between oriented normalizations, but it
changes the multiplication.  Suppose the maps $c_{ab}$ satisfy
\eqref{cbc:exchange:colored-braid}, and choose central units $f_{ab}\in B$
for $a>b$.  The maps $c'_{ab}=f_{ab}c_{ab}$ satisfy the same equation.
Both algebras are graded by the color count $\alpha\in\mathbb N^m$.

For two color counts define
\[
 \sigma(\alpha,\beta)=\prod_{a>b}f_{ab}^{\alpha_a\beta_b}.
\]
It satisfies
\[
 \sigma(\alpha,\beta)\sigma(\alpha+\beta,\gamma)
 =\sigma(\beta,\gamma)\sigma(\alpha,\beta+\gamma),
 \qquad \sigma(0,\alpha)=\sigma(\alpha,0)=1.
\]
Consequently $x\star y=\sigma(\alpha,\beta)xy$ for homogeneous $x,y$
of degrees $\alpha,\beta$ defines an associative unital algebra on the
module $A(c)$.  Denote it by $A(c)^\sigma$.
Sending each generator to itself gives an algebra isomorphism
\begin{equation}\label{cbc:exchange:twist}
 A(c')\xrightarrow{\ \sim\ }A(c)^\sigma.
\end{equation}

To verify this, a descending pair acquires the factor $f_{ab}$,
whereas its increasing counterpart acquires none.  Thus the new relations
hold.  An ordered product acquires no factor at any stage, so the map
identifies the two ordered bases and is invertible.

More generally, for a word with colors $(a_1,\ldots,a_n)$ its normal
form for $c'$ is its normal form for $c$ multiplied by
\[
 \prod_{i<j,\ a_i>a_j}f_{a_i a_j}.
\]
Each inverted pair crosses exactly once during sorting.
In particular, \eqref{cbc:exchange:normalized} is a scalar twist of the
oriented rational algebra.  This comparison does not identify their
untwisted products.

\section{The obstruction and the realization problem}
Failure of Yang--Baxter obstructs the stated free ordered basis.
For example, on $k^2$ let $E_{ij}$ be the matrix units and take the
invertible constant matrix $R=1+E_{12}\otimes E_{21}$, whose inverse is
$1-E_{12}\otimes E_{21}$.  Direct multiplication gives
\[
 R_{12}R_{13}R_{23}-R_{23}R_{13}R_{12}
 =E_{12}\otimes(E_{22}-E_{11})\otimes E_{21}.
\]
For $x=e_1$ and $y=e_2$ this sends $y\otimes x\otimes x$ to
$-x\otimes x\otimes y$.  Thus three colors with these exchange maps
cannot have the ordered basis of Theorem~\ref{cbc:exchange:normal}.

Their associative quotient nevertheless exists and has a nonzero unit:
its relations are homogeneous of degree two and leave the scalar summand
unchanged.  With constant fields $Y(a,z)b=ab$, it is also a regular
nonlocal vertex algebra as constructed in the preceding chapter.
The obstruction concerns the specified independent ordered products,
not the existence of every associative algebra or field system.

The finite-color algebra retains spectral coefficients and word order.
An associative chiral realization requires coefficient objects over
configurations and their extensions across collisions.  One must also
construct the substitution maps with all required coherences.  A realization
by quantum fields must additionally give the state space, vacuum,
translation, Laurent expansions and full-state exchange operator.
Neither construction follows from an ordered basis alone.

\chapter{The RTT Yangian of \texorpdfstring{$\mathfrak{gl}_2$}{gl2}}

The rational exchange matrix also defines an algebra whose representations
can be tensored.  Its generators are the coefficients of a matrix series.
Matrix multiplication supplies the coproduct.  Keeping the deformation
parameter in the coefficient ring reveals two different associated graded
algebras and determines the differential in the bar construction.

\section{Coefficients and evaluation}

Work over $B=\mathbb C[\hbar]$.  The presentation and Hopf formulas hold
for every finite $n\geq1$.  The determinant calculation below concerns
$n=2$.  Let $V=B^n$, let $E_{ij}$ be its matrix units, and let
$P(x\otimes y)=y\otimes x$ on $V\otimes_B V$.

Write
\[
 T^{(r)}=(t_{ij}^{(r)})_{i,j=1}^n,\qquad
 T(u)=I+\sum_{r\geq1}T^{(r)}u^{-r},\qquad
 t_{ij}^{(0)}=\delta_{ij}.
\]
Here $u$ is central, $T_1(u)=T(u)\otimes I$ and
$T_2(v)=I\otimes T(v)$ refer to the two auxiliary copies of $V$.
Their matrix entries belong to the same coefficient algebra.

Define $Y_\hbar(\mathfrak{gl}_n)$ to be the unital associative
$B$-algebra on the $t_{ij}^{(r)}$, $r\geq1$, with relations
\begin{equation}\label{cbc:rtt:cleared}
 \bigl((u-v)I-\hbar P\bigr)T_1(u)T_2(v)
 =T_2(v)T_1(u)\bigl((u-v)I-\hbar P\bigr).
\end{equation}
The equality is coefficientwise in simultaneous formal Laurent series
with bounded positive powers of $u$ and $v$.  Thus the presentation
requires no choice of expansion for $(u-v)^{-1}$.
The customary rational notation uses $R(w)=I-\hbar P/w$.

\begin{proposition}[Coefficient relations]\label{cbc:rtt:coefficients}
The defining relations are equivalent to
\begin{equation}\label{cbc:rtt:recurrence}
 [t_{ij}^{(r+1)},t_{kl}^{(s)}]-[t_{ij}^{(r)},t_{kl}^{(s+1)}]
 =\hbar\bigl(t_{kj}^{(r)}t_{il}^{(s)}-t_{kj}^{(s)}t_{il}^{(r)}\bigr),
 \qquad r,s\geq0.
\end{equation}
For $r,s\geq1$ this gives
\begin{equation}\label{cbc:rtt:commutator}
 [t_{ij}^{(r)},t_{kl}^{(s)}]
 =\hbar\sum_{a=0}^{\min(r,s)-1}
 \bigl(t_{kj}^{(a)}t_{il}^{(r+s-1-a)}
       -t_{kj}^{(r+s-1-a)}t_{il}^{(a)}\bigr).
\end{equation}
\end{proposition}
\begin{proof}
The $(ij,kl)$ matrix entry of \eqref{cbc:rtt:cleared} is
\[
 (u-v)[t_{ij}(u),t_{kl}(v)]
 =\hbar\bigl(t_{kj}(u)t_{il}(v)-t_{kj}(v)t_{il}(u)\bigr).
\]
Its coefficient of $u^{-r}v^{-s}$ is
\eqref{cbc:rtt:recurrence}, including the boundary cases $r=0$ or
$s=0$.  All entries and coefficients have been retained, so the
implication reverses.  Iterating the recurrence down to first index
zero gives the sum in \eqref{cbc:rtt:commutator} with $a=0,\ldots,r-1$.
If $r>s$, its terms with $s\leq a\leq r-1$ cancel in pairs under
$a\mapsto r+s-1-a$.  This gives the displayed upper limit.
\end{proof}

To identify the first modes, define $U_\hbar(\mathfrak{gl}_n)$ as the
$B$-algebra generated by $X_{ij}$ with relations
\[
 [X_{ij},X_{kl}]
 =\hbar(\delta_{jk}X_{il}-\delta_{il}X_{kj}).
\]
For each $a\in\mathbb C$ there is an evaluation homomorphism
\begin{equation}\label{cbc:rtt:evaluation}
 \operatorname{ev}_a:Y_\hbar(\mathfrak{gl}_n)\longrightarrow
 U_\hbar(\mathfrak{gl}_n),\qquad
 t_{ij}^{(r)}\longmapsto a^{r-1}X_{ij}.
\end{equation}
For positive $r,s$, substituting these coefficients in
\eqref{cbc:rtt:recurrence} makes both sides zero.  When $r=0$ or $s=0$,
the same substitution gives precisely the displayed relation for the
$X_{ij}$.  This proves that evaluation is well defined.

The relation with $r=0,s=1$ also defines a homomorphism
$X_{ij}\mapsto t_{ij}^{(1)}$, and $\operatorname{ev}_0$ is its left
inverse.  The first modes therefore contain an actual copy of
$U_\hbar(\mathfrak{gl}_n)$.

The map $X_{ij}\mapsto\hbar e_{ij}$, where $e_{ij}$ are the standard
generators of $U(\mathfrak{gl}_n)$, gives the familiar expression
\[
 T(u)\longmapsto I+\frac{\hbar(e_{ij})}{u-a}.
\]
This is a generating series for an algebra map into
$U(\mathfrak{gl}_n)\otimes_{\mathbb C}B$.  Its expansion is at $u=\infty$.
The original generators $t_{ij}^{(r)}$ have not been divided by $\hbar$.

\section{The coefficient Hopf algebra}

\begin{theorem}[Hopf structure]\label{cbc:rtt:hopf}
The Yangian is a Hopf algebra over $B$, with
\begin{align*}
 \Delta(t_{ij}^{(r)})
 &=\sum_{k=1}^n\sum_{p=0}^r
       t_{ik}^{(p)}\otimes t_{kj}^{(r-p)},&
 \varepsilon(t_{ij}^{(r)})&=0\quad(r\geq1),\\
 S(T^{(r)})
 &=\sum_{m=1}^r(-1)^m
   \sum_{\substack{r_1+\cdots+r_m=r\\r_\nu\geq1}}
       T^{(r_1)}\cdots T^{(r_m)}.
\end{align*}
The coproduct takes values in the ordinary algebraic tensor product
$Y_\hbar\otimes_B Y_\hbar$.  The antipode formula specifies the matrix
of its values on the generators.
\end{theorem}
\begin{proof}
Let $T'$ and $T''$ be RTT matrices whose coefficient algebras commute,
and put $\mathcal R=(u-v)I-\hbar P$.  Their product satisfies RTT:
\begin{align*}
 \mathcal R T'_1T''_1T'_2T''_2
 &=\mathcal R T'_1T'_2T''_1T''_2
  =T'_2T'_1\mathcal R T''_1T''_2\\
 &=T'_2T'_1T''_2T''_1\mathcal R
  =T'_2T''_2T'_1T''_1\mathcal R.
\end{align*}
In the first and last equalities the crossed factors have different
auxiliary spaces and commuting coefficients.  Taking $T'$ and $T''$
in the two tensor factors defines $\Delta$.  Its coefficients are the
finite sums above.  Threefold matrix multiplication proves
coassociativity, and substitution $T=I$ gives the counit identities.

Give $t_{ij}^{(r)}$ weight $r$ and $B$ weight zero.  The resulting
filtration has degree-zero part $B$, and the reduced coproduct of a
positive-weight element has smaller positive weights in both factors.
For $J=\operatorname{id}-\eta\varepsilon$, the convolution series
$\sum_{m\geq0}(-J)^{*m}$ is consequently finite on each filtered piece.
Multiplying this series by $\eta\varepsilon+J$ on either side gives
$\eta\varepsilon$.  It is the antipode.

On the generating matrix the
two identities say $S(T(u))T(u)=T(u)S(T(u))=I$.
Expanding the inverse series gives the formula in the theorem.
\end{proof}

Formal series in $u^{-1}$ thus package finite algebraic formulas.
Completion in $\hbar$ is a further operation, discussed after the
coefficient basis has been established.

\section{Two filtrations}

Fix a total order on the triples $(r,i,j)$.  An ordered monomial means
a finite product of generators in that order, including the empty
product.  Let $G_d$ denote the filtration by total weight
$\sum r$, and let $F_p$ denote the filtration by total loop degree
$\sum(r-1)$.  Both give $\hbar$ degree zero.
For the current Lie algebra $L=\mathfrak{gl}_n(\mathbb C)[z]$, write
$L_B=L\otimes_{\mathbb C}B$ and equip it with bracket
$[x,y]_\hbar=\hbar[x,y]$.  Its enveloping algebra over $B$ is denoted
$U_B(L_\hbar)$.

\begin{theorem}[PBW and the current algebra]\label{cbc:rtt:pbw}
The ordered monomials form a $B$-basis of $Y_\hbar(\mathfrak{gl}_n)$.
Moreover,
\[
 \operatorname{gr}_G Y_\hbar
   =B[\,\overline t_{ij}^{(r)}:r\geq1\,],\qquad
 \operatorname{gr}_F Y_\hbar\cong U_B(L_\hbar),\qquad
 e_{ij}z^{r-1}\longmapsto\overline t_{ij}^{(r)}.
\]
If $Y_1$ is the specialization at $\hbar=1$ and $G_dY_1$ has weights
$\deg t_{ij}^{(r)}=r$, then
\[
 Y_\hbar\cong\bigoplus_{d\geq0}\hbar^dG_dY_1
 \ \subset\ Y_1[\hbar].
\]
In this last expression $\hbar$ is the Rees variable.
\end{theorem}
\begin{proof}
At $\hbar=1$, the ordered-monomial theorem is
Molev--Nazarov--Olshanskii's PBW theorem
\cite[Theorem~1.22]{MNOYangians}.  Denote its generators by
$\tau_{ij}^{(r)}$.  The assignment
$t_{ij}^{(r)}\mapsto\hbar^r\tau_{ij}^{(r)}$ respects
\eqref{cbc:rtt:recurrence} and hence defines a homomorphism to
$Y_1[\hbar]$.  Equation~\eqref{cbc:rtt:commutator} lowers total weight
by at least one when two generators are commuted.  Induction on weight,
followed by sorting at fixed weight, therefore expresses every word as
a $B$-linear combination of ordered monomials.

The image of an ordered monomial of weight $d$ is $\hbar^d$ times its
corresponding PBW monomial in $Y_1$.  These images are $B$-linearly
independent, since the PBW monomials are independent and $B$ has no
$\hbar$-torsion.  This proves the basis claim and identifies the image
with the displayed Rees algebra.  It also proves the polynomial
description of $\operatorname{gr}_G$.

For the loop filtration, the $a=0$ term of
\eqref{cbc:rtt:commutator} has degree $r+s-2$.  Every other term has
degree at most $r+s-3$.  The leading relation is consequently
\[
 [\overline t_{ij}^{(r)},\overline t_{kl}^{(s)}]
 =\hbar\bigl(\delta_{kj}\overline t_{il}^{(r+s-1)}
             -\delta_{il}\overline t_{kj}^{(r+s-1)}\bigr).
\]
It defines the stated surjection from $U_B(L_\hbar)$.
PBW for the free $B$-module Lie algebra $L_\hbar$ gives an ordered basis
of its enveloping algebra.  Its image is the basis just proved, so the
map is injective.
\end{proof}

In particular, the loop-associated graded retains its scaled bracket.
Specializing $\hbar=1$ gives $U(L)$, whereas specializing $\hbar=0$
gives the symmetric algebra on $L$.  If $W$ is the complex vector space
with basis the Yangian's ordered monomials, its $\hbar$-adic completion
has underlying module $W[[\hbar]]$.  A coefficient at each power of
$\hbar$ is a finite linear combination of PBW monomials.
A single finite set of monomials need not suffice for the whole series.

\section{The determinant in rank two}

For $n=2$ the antisymmetric square of the auxiliary space has rank one.
This produces a central scalar series.  Put
$A_{12}=(I-P)/2$ and define
\begin{equation}\label{cbc:rtt:determinant}
 D(u)=t_{11}(u)t_{22}(u-\hbar)
       -t_{21}(u)t_{12}(u-\hbar).
\end{equation}
Each shifted series is expanded at infinity.

\begin{proposition}[Quantum determinant]\label{cbc:rtt:qdet}
The coefficients of $D(u)=1+\sum_{r\geq1}d_ru^{-r}$ are central, and
\[
 A_{12}T_1(u)T_2(u-\hbar)=D(u)A_{12},\qquad
 \Delta D(u)=D(u)\otimes D(u),\qquad \varepsilon(D(u))=1.
\]
There is a unique central series $c(u)=1+\sum_{r\geq1}c_ru^{-r}$
such that $c(u)c(u-\hbar)=D(u)$.  Writing
$Y_\hbar^0=Y_\hbar(\mathfrak{gl}_2)/(d_1,d_2,\ldots)$, one has
an algebra isomorphism
\[
 Y_\hbar(\mathfrak{gl}_2)
 \cong B[c_1,c_2,\ldots]\otimes_B Y_\hbar^0.
\]
The determinant-one quotient is a Hopf algebra.  At $\hbar=1$ it is
the usual $Y(\mathfrak{sl}_2)$.
\end{proposition}
\begin{proof}
Let $\mathcal T(v)$ be the standard Yangian matrix at $\hbar=1$.
After inverting $\hbar$, the preceding PBW proof identifies
$T(u)$ with $\mathcal T(u/\hbar)$.
The antisymmetrizer identity and centrality are then the standard
quantum-determinant identity and theorem
\cite[\S2.1 and Theorem~2.10]{MNOYangians}.
Their coefficient identities descend to $B$ because $Y_\hbar$ is
$B$-free.  In particular, the order of the cross term in
\eqref{cbc:rtt:determinant} is fixed by the column formula.

For two RTT matrices $T',T''$ with commuting coefficient algebras,
the antisymmetrizer identity gives
\begin{align*}
 A_{12}(T'T'')_1(u)(T'T'')_2(u-\hbar)
 &=A_{12}T'_1(u)T'_2(u-\hbar)T''_1(u)T''_2(u-\hbar)\\
 &=D'(u)D''(u)A_{12}.
\end{align*}
This proves the coproduct formula.  The counit follows by setting $T=I$.
The coefficient of $u^{-r}$ in $c(u)c(u-\hbar)$ is $2c_r$ plus a
polynomial in earlier $c_s$ and $\hbar$.  Recursion therefore constructs
the unique $c_r$ as central polynomials in the $d_s$.
The same uniqueness applied in $Y_\hbar\otimes_B Y_\hbar$ proves
$\Delta c(u)=c(u)\otimes c(u)$.

The normalized matrix $\widetilde T(u)=c(u)^{-1}T(u)$ satisfies RTT
and has determinant one.  Conversely, multiply a determinant-one RTT
matrix by an independent central series $q(u)=1+\sum_{r\geq1}q_ru^{-r}$.
Its determinant is $q(u)q(u-\hbar)$, so the recursion recovers $q$ and
normalization recovers the original matrix.  These two coefficientwise
algebra maps are inverse.  This proves the tensor factorization and the
algebraic independence of the $c_r$.

Group-likeness of $D$ makes
$(d_1,d_2,\ldots)$ a Hopf ideal, since $S(D)=D^{-1}$.
The identification at $\hbar=1$ is the determinant-one presentation
of the special-linear Yangian \cite[Theorem~2.21]{MolevYangianApplications}.
\end{proof}

The cross-term order can already be distinguished on a two-dimensional
evaluation module.  At $\hbar=1$, $a=0$ and $u=3$, let
$t_{ij}(u)=\delta_{ij}I+E_{ij}/u$.  Then $D(3)=\frac43 I$, whereas
\[
 t_{11}(3)t_{22}(2)-t_{12}(3)t_{21}(2)
 =\begin{pmatrix}7/6&0\\0&3/2\end{pmatrix}.
\]
The commutator of this last matrix with $E_{12}$ is
$-\frac13E_{12}$.  Thus the cross-term indices cannot be interchanged.

\section{Exchange on evaluation modules}

An exchange operator must specify the representations it exchanges.
Let $V_a$ be the rank-$n$ evaluation module obtained from
\eqref{cbc:rtt:evaluation} and the dual fundamental action
$X_{ij}\mapsto-\hbar E_{ji}$.  Its auxiliary matrix is
$R_{0V}(u-a)$, where $0$ denotes the auxiliary space.

For central parameters $a,b$, extend scalars to
\[
 B_{a,b}=\mathbb C[\hbar,a,b,(a-b)^{-1},
                          (a-b-\hbar)^{-1},(a-b+\hbar)^{-1}].
\]
With the coproduct action, there is an invertible intertwiner
\begin{equation}\label{cbc:rtt:intertwiner}
 C_{a,b}=P_{12}R_{12}(b-a):
       V_a\otimes_{B_{a,b}}V_b\longrightarrow
       V_b\otimes_{B_{a,b}}V_a.
\end{equation}

Indeed, put $M=R_{01}(u-a)R_{02}(u-b)$ and
$M^{\mathrm{op}}=R_{02}(u-b)R_{01}(u-a)$.
Yang--Baxter gives
\[
 M R_{12}(a-b)=R_{12}(a-b)M^{\mathrm{op}}.
\]
Since $R(w)R(-w)=(1-\hbar^2/w^2)I$, this implies
$R_{12}(b-a)M=M^{\mathrm{op}}R_{12}(b-a)$.
Conjugating the right-hand auxiliary product by $P_{12}$ puts the
parameter $b$ in the first module and $a$ in the second, proving
\eqref{cbc:rtt:intertwiner}.  The same scalar identity proves
invertibility over $B_{a,b}$.

We have obtained an ordinary coefficient algebra, a coproduct and
exchange maps on specified modules.  The series $T(u)$ by itself
does not supply a state--field map or coefficient objects on configuration
spaces.  An associative chiral realization must construct those objects,
their collision and substitution maps, and the comparison preserving
these operations.  A realization by line operators must also identify
the tensor product and exchange above with the corresponding physical
operations.  These are additional reconstruction problems.

\chapter{The Hall algebra of the \texorpdfstring{$A_2$}{A2} quiver}\label{cbc:hall:main}
Counting extensions gives an algebra in which the quantum Serre relation is a
cancellation among flags.  For the quiver with one arrow, rank classifies every
representation.  This makes it possible to prove both the relation and the
basis theorem that identifies the entire Hall algebra.  The calculation is an
elementary instance of the Hall-algebra approach to quantum groups
\cite{Ringel,Green}.

\section{Representations and multiplication}
Fix a finite field $k=\mathbb F_q$, put $v=\sqrt q>0$, and use the coefficient
field $K_q=\mathbb Q(v)$.  A representation of $Q:1\to2$ is a linear map
$f:V_1\to V_2$ between finite-dimensional $k$-spaces.  Its dimension vector is
$\dim f=(\dim V_1,\dim V_2)$.

The three indecomposable representations are
\[
 S_1=(k\to0),\qquad S_2=(0\to k),\qquad M=(k\xrightarrow{1}k).
\]
Indeed, bases adapted to the kernel and image of $f$ give
\begin{equation}\label{cbc:hall:rank}
 f\simeq S_2^a\oplus M^b\oplus S_1^c,
 \qquad b=\operatorname{rank}f,\quad
 a=\dim V_2-b,\quad c=\dim V_1-b.
\end{equation}
These integers determine the isomorphism class.  The arrow in $M$ also gives
the nonsplit exact sequence $0\to S_2\to M\to S_1\to0$.

For representations $A,B,C$, let $g_{A,B}^C$ be the number of subrepresentations
$W\subset C$ with $W\simeq B$ and $C/W\simeq A$.  Thus the right factor is the
subobject and the left factor is the quotient.  The Hall space and its
untwisted multiplication are
\[
 H_q=\bigoplus_{[C]}K_q[C],\qquad
 [A]*[B]=\sum_{[C]}g_{A,B}^C[C].
\]
The sum is finite: its dimension vector is $\dim A+\dim B$, and
\eqref{cbc:hall:rank} leaves only finitely many ranks.  The zero representation
gives the unit $[0]$.

Both bracketings of $[A]*[B]*[D]$ count the same flags
\[
 0\subset W\subset U\subset C,
 \qquad W\simeq D,\quad U/W\simeq B,\quad C/U\simeq A.
\]
This proves associativity, with the dimension vector as an additive grading.

The Euler form for this orientation is
\[
 \langle(a,b),(c,d)\rangle=ac+bd-ad.
\]
Define the twisted multiplication by
\begin{equation}\label{cbc:hall:twist}
 [A]\diamond[B]=v^{\langle\dim A,\dim B\rangle}[A]*[B].
\end{equation}
It is associative because bilinearity gives
\[
 \langle\alpha,\beta\rangle+\langle\alpha+\beta,\gamma\rangle
 =\langle\beta,\gamma\rangle+\langle\alpha,\beta+\gamma\rangle.
\]
The unit remains $[0]$.  Write $e_i=[S_i]$ and, unless a star is displayed,
use $\diamond$ for products in $H_q$.

\section{The root vector and the Serre relation}
A representation of dimension $(1,1)$ has rank zero or one.  Each has exactly
one subobject $S_2$ with quotient $S_1$.  Only the rank-zero representation has
a subobject $S_1$.  Consequently
\[
 e_1e_2=v^{-1}\bigl([S_1\oplus S_2]+[M]\bigr),
 \qquad e_2e_1=[S_1\oplus S_2].
\]
The root vector has the precise normalization
\begin{equation}\label{cbc:hall:root}
 e_{12}:=e_1e_2-v^{-1}e_2e_1=v^{-1}[M].
\end{equation}

\begin{theorem}[Quantum Serre relations]\label{cbc:hall:serre}
The generators $e_1,e_2$ satisfy
\begin{align*}
 e_1^2e_2-(v+v^{-1})e_1e_2e_1+e_2e_1^2&=0,\\
 e_2^2e_1-(v+v^{-1})e_2e_1e_2+e_1e_2^2&=0.
\end{align*}
\end{theorem}
\begin{proof}
Read a word from its top quotient to its bottom quotient.  For a map
$f:k^2\to k$, the three untwisted flag counts are
\[
\begin{array}{c|ccc}
 \operatorname{rank}f&112&121&211\\ \hline
 0&q+1&q+1&q+1\\
 1&q+1&1&0.
\end{array}
\]
If $f=0$, every line in $k^2$ is allowed.  For rank one, $112$ allows every line and $121$ forces the kernel line.
The word $211$ would require the whole space $k^2$ to map to zero.

The total Euler exponents are $-1,0,1$.
The coefficients of the first Serre expression on the two rank classes are
therefore
\[
 (q+1)\bigl(v^{-1}-(v+v^{-1})+v\bigr)=0,
 \qquad v^{-1}(q+1)-(v+v^{-1})=0.
\]
Here the second equality uses $q=v^2$.

Let $D$ be the contravariant involution obtained by dualizing a map and
exchanging the two vertices.  It exchanges $S_1,S_2$, fixes
$M$, and reverses a flag.  Moreover
$\langle\dim DB,\dim DA\rangle=\langle\dim A,\dim B\rangle$.  It therefore induces an anti-automorphism of the twisted Hall algebra.
Applying it to the first relation proves the second.
\end{proof}

\section{A basis and the positive quantum group}
To distinguish a parameter from its numerical specialization, let
$\mathcal A=\mathbb Z[v,v^{-1}]$ in this section, with $v$ indeterminate.
Define $U_{\mathcal A}^+$ to be the associative $\mathcal A$-algebra generated
by $E_1,E_2$ subject to the two relations of
Theorem~\ref{cbc:hall:serre}, with $E_i$ in place of $e_i$.  Its extension of
scalars to $\mathbb Q(v)$ is the usual positive quantum group
$U_v^+(\mathfrak{sl}_3)$.  Put
\[
 X=E_2,\qquad Z=E_1,\qquad Y=ZX-v^{-1}XZ.
\]
The degrees of $E_1,E_2$ are $(1,0),(0,1)$, so $Y$ has degree $(1,1)$.

\begin{proposition}[Normal forms]\label{cbc:hall:normal}
The algebra $U_{\mathcal A}^+$ is free over $\mathcal A$ with basis
$X^aY^bZ^c$, for $a,b,c\geq0$.  Equivalently, it has the presentation
\begin{equation}\label{cbc:hall:three-relations}
 YX=vXY,\qquad ZY=vYZ,\qquad ZX=v^{-1}XZ+Y.
\end{equation}
\end{proposition}
\begin{proof}
In the three-generator presentation, reduce each left side of
\eqref{cbc:hall:three-relations} to its right side.  Order words first by
length and then by their number of inversions for $X<Y<Z$.  Each replacement
strictly decreases this pair.  The only overlapping pair of reductions is
$ZYX$, and both choices give
\[
 ZYX\longmapsto vXYZ+vY^2.
\]
Reductions at disjoint positions commute.

Induction on the decreasing pair
now shows that every word has a unique linear combination of ordered words
as its normal form.  The resulting normal-form map kills every defining
relation in every word context and fixes the ordered words.  Thus these words
are both spanning and linearly independent in the quotient, over
$\mathcal A$ itself.

Eliminating $Y$ from $YX=vXY$ gives
$ZX^2-(v+v^{-1})XZX+X^2Z=0$.  Eliminating it from $ZY=vYZ$ gives
$Z^2X-(v+v^{-1})ZXZ+XZ^2=0$.  These are exactly the two defining relations
of $U_{\mathcal A}^+$.  Conversely, those relations imply the first two
identities in \eqref{cbc:hall:three-relations} after setting
$Y=ZX-v^{-1}XZ$.  The two substitutions are inverse algebra homomorphisms.
\end{proof}

Specialize along $\mathcal A\to K_q$, $v\mapsto\sqrt q$.  In the following
formula $v$ again denotes this positive square root.  Define
\[
 [m]_v=\sum_{j=0}^{m-1}v^{m-1-2j},\qquad
 [m]_v!=\prod_{j=1}^m[j]_v,\qquad [0]_v!=1.
\]
All these factorials are nonzero in $K_q$.

\begin{theorem}[Hall realization in type $A_2$]\label{cbc:hall:realization}
The map
\[
 K_q\otimes_{\mathcal A}U_{\mathcal A}^+\longrightarrow H_q,
 \qquad E_i\longmapsto e_i,
\]
is an isomorphism of dimension-graded associative algebras.  More precisely,
\begin{equation}\label{cbc:hall:monomial}
\begin{aligned}
 e_2^a e_{12}^b e_1^c
 &=v^{a(a-1)+b(b-1)+c(c-1)-b+ab+bc}\\
 &\quad{}\cdot[a]_v![b]_v![c]_v!\,[S_2^a\oplus M^b\oplus S_1^c].
\end{aligned}
\end{equation}
\end{theorem}
\begin{proof}
For $S=S_1,S_2$ or $M$, an iterated extension with all factors $S$ is
isomorphic to $S^m$.  For $M$ this follows because the arrow is invertible
on every successive quotient, hence invertible on the whole representation.
Its flags are complete flags in $k^m$.  The same count applies to the two
simples.

Their number is
\[
 [m]_q!:=\prod_{j=1}^m(1+q+\cdots+q^{j-1})
 =v^{m(m-1)/2}[m]_v!.
\]
Since $\langle\dim S,\dim S\rangle=1$, the twist contributes another
$v^{m(m-1)/2}$.  Hence $[S]^m=v^{m(m-1)}[m]_v![S^m]$.
The power of $e_{12}=v^{-1}[M]$ contributes the additional factor $v^{-b}$.

In $[S_2^a]*[M^b]$, the subobject must have the full source space and its
image at the target.  Thus the middle representation is $S_2^a\oplus M^b$,
and that subobject is unique.  In
$[S_2^a\oplus M^b]*[S_1^c]$, the subobject must be the full kernel, again
unique.  The corresponding Euler exponents are $ab$ and $bc$.  Multiplying
these factors proves \eqref{cbc:hall:monomial}.

The Serre relations define the displayed homomorphism.  By
Proposition~\ref{cbc:hall:normal}, its source has the ordered monomial basis.
Equation~\eqref{cbc:hall:monomial} sends that basis bijectively to nonzero
multiples of all rank classes \eqref{cbc:hall:rank}.  It proves both
injectivity and surjectivity.
\end{proof}

The specialization uses the Laurent-polynomial form $U_{\mathcal A}^+$.
It does not require evaluating arbitrary elements of $\mathbb Q(v)$ at an
algebraic number.  A geometric Hall realization requires, in addition, a
coefficient category and extension correspondences compatible with its
factorization operations.  Chapter~\ref{chap:hall} describes these data and
the flag correspondences that the finite-field calculation supplies.

\chapter{Matrix traces and the necklace bracket}\label{cbc:neck:main}
A trace remembers a word only up to cyclic rotation.  Differentiating that
trace cuts the word at every occurrence of the varied matrix.  These cuts
give a Lie bracket on cyclic words and Hamiltonian vector fields on every
matrix space.  Passing to commuting matrices leads to multisymmetric
functions, where the trace of the identity retains the matrix rank.

\section{Cyclic words and matrix Hamiltonians}
Work over $\mathbb C$ and let $F=\mathbb C\langle x,y\rangle$ be the free
associative algebra.  Write $[F,F]$ for the linear span of commutators and put
$F_{\mathrm{cyc}}=F/[F,F]$.  This is a vector-space quotient, with a basis
consisting of cyclic rotation classes of words, including the empty word.
Indeed, a commutator of words is a difference of cyclic rotations.  Moving
the initial generator of a word to the end changes it by a commutator.

For a word $w=a_1\cdots a_n$, where each $a_i$ is $x$ or $y$, define
\begin{equation}\label{cbc:neck:derivative}
 D_x[w]_{\mathrm{cyc}}
 =\sum_{a_i=x}a_{i+1}\cdots a_na_1\cdots a_{i-1}\in F,
\end{equation}
and define $D_y$ in the same way.  Both derivatives annihilate the empty
word.  Rotating $w$ only permutes its cuts, so these maps are well-defined
on $F_{\mathrm{cyc}}$.  Define the bilinear operation
\begin{equation}\label{cbc:neck:bracket}
 \{f,g\}_{\mathrm{neck}}
 =[D_xf\,D_yg-D_yf\,D_xg]_{\mathrm{cyc}}.
\end{equation}
We will prove its Jacobi identity by evaluating actual matrix traces.  These
constructions are the two-generator case of noncommutative symplectic
geometry \cite[\S\S6.2,8.4]{CBEGNecklace}.

Fix $N\geq1$.  The coordinate algebra of a pair of $N\times N$ matrices is
\[
 O_N=\mathbb C[X_{ij},Y_{ij}:1\leq i,j\leq N].
\]
Equip it with the Poisson bracket determined by
\[
 \{X_{ij},Y_{kl}\}=\delta_{il}\delta_{jk},\qquad
 \{X_{ij},X_{kl}\}=\{Y_{ij},Y_{kl}\}=0.
\]
This convention fixes the sign of \eqref{cbc:neck:bracket}.  Cyclicity of
trace defines a linear map
\[
 J_N:F_{\mathrm{cyc}}\longrightarrow O_N^{\mathrm{GL}_N},
 \qquad J_N(f)=\operatorname{Tr}f(X,Y),
\]
where $\mathrm{GL}_N$ acts by simultaneous conjugation.  In particular,
$J_N([1]_{\mathrm{cyc}})=N$.

\begin{theorem}[Trace realization]\label{cbc:neck:trace}
For $f,g\in F_{\mathrm{cyc}}$,
\[
 \{J_N(f),J_N(g)\}=J_N(\{f,g\}_{\mathrm{neck}}).
\]
\end{theorem}
\begin{proof}
Moving a varied matrix entry to the front of the trace gives
\[
 \frac{\partial J_N(f)}{\partial X_{ij}}
 =(D_xf(X,Y))_{ji},\qquad
 \frac{\partial J_N(f)}{\partial Y_{ij}}
 =(D_yf(X,Y))_{ji}.
\]
The entrywise Poisson contraction is therefore
\[
 \sum_{i,j}(D_xf)_{ji}(D_yg)_{ij}
 -(D_yf)_{ji}(D_xg)_{ij}
 =\operatorname{Tr}(D_xf\,D_yg-D_yf\,D_xg).
\]
This is the required identity.
\end{proof}

For example, $\{[x],[y]\}_{\mathrm{neck}}=[1]$ gives
$\{\operatorname{Tr}X,\operatorname{Tr}Y\}=N$.  Cutting the two occurrences
of $x$ in $x^2y$ gives $D_x[x^2y]=xy+yx$, so
\[
 \{\operatorname{Tr}(X^2Y),\operatorname{Tr}Y\}
 =2\operatorname{Tr}(XY).
\]

\begin{lemma}[Separation by traces]\label{cbc:neck:separation}
If $J_N(f)=0$ for every $N\geq1$, then $f=0$ in $F_{\mathrm{cyc}}$.
\end{lemma}
\begin{proof}
Fix a nonempty word $w=a_1\cdots a_n$.  On an oriented cycle with $n$
vertices, assign the entry $t_iE_{i,i+1}$ to $X$ or $Y$ according as $a_i=x$
or $a_i=y$, with indices read modulo $n$.  The variables $t_i$ are independent.
The coefficient of $t_1\cdots t_n$ in the trace of a word is zero unless
that word belongs to the cyclic class of $w$.

For that class, the coefficient is its positive number of rotation
symmetries.  Words of other lengths have a different total degree in the
$t_i$.  Since the characteristic is zero, this coefficient detects the
coefficient of $[w]_{\mathrm{cyc}}$ in $f$.  Evaluating at $N=1$, $X=Y=0$
detects the constant class.  These tests detect every basis coefficient.
\end{proof}

\begin{proposition}[The necklace Lie algebra]\label{cbc:neck:lie}
Equation~\eqref{cbc:neck:bracket} makes $F_{\mathrm{cyc}}$ a Lie algebra.
The constant class is central, so $F_{\mathrm{cyc}}/\mathbb C[1]$ is also a
Lie algebra.
\end{proposition}
\begin{proof}
Skew symmetry follows by moving the two factors cyclically.  By
Theorem~\ref{cbc:neck:trace}, every matrix trace sends the Jacobi expression
to the Jacobi expression of an ordinary Poisson bracket, hence to zero.
Lemma~\ref{cbc:neck:separation} proves the claimed identity in
$F_{\mathrm{cyc}}$.  The constant class has zero cyclic derivatives.
\end{proof}

\section{Hamiltonian derivations and the moment element}
For $f\in F_{\mathrm{cyc}}$, define a derivation $\delta_f$ of $F$ by
\begin{equation}\label{cbc:neck:hamiltonian}
 \delta_f(x)=-D_yf,\qquad \delta_f(y)=D_xf.
\end{equation}
Its matrix evaluation is the Hamiltonian vector field
$\{J_N(f),-\}$, as the entrywise derivative formula shows.  Consequently
\begin{equation}\label{cbc:neck:commutator}
 [\delta_f,\delta_g]=\delta_{\{f,g\}_{\mathrm{neck}}}.
\end{equation}

To justify the implication back in $F$, evaluate both sides on $x$ and $y$.
Their difference vanishes on matrices of every size.  A nonzero polynomial
of degree at most $d$ acts nontrivially on $1$ by left multiplication in
$F/(\text{words of length}>d)$.  This finite-dimensional representation
separates free polynomials and proves \eqref{cbc:neck:commutator}.

The moment element is $\mu=[x,y]$.  Telescoping the cyclic cuts gives
\begin{equation}\label{cbc:neck:euler}
 [x,D_xf]+[y,D_yf]=0,
 \qquad [xD_xf+yD_yf]_{\mathrm{cyc}}=nf
\end{equation}
for $f$ homogeneous of word length $n$.  The first equality implies
$\delta_f(\mu)=0$.  The second identifies the kernel of $f\mapsto\delta_f$
with the constant classes.

\begin{proposition}[Recovering the cyclic Hamiltonian]\label{cbc:neck:potential}
Let $\operatorname{Der}_{\mu}(F)$ be the Lie algebra of derivations fixing
$\mu$.  Then $f\mapsto\delta_f$ is an isomorphism
\[
 F_{\mathrm{cyc}}/\mathbb C[1]\xrightarrow{\ \sim\ }
 \operatorname{Der}_{\mu}(F).
\]
If $D(x)=a$ and $D(y)=b$ are homogeneous of degree $m$, its inverse is
\[
 f=\frac{[xb-ya]_{\mathrm{cyc}}}{m+1}\pmod{\mathbb C[1]}.
\]
\end{proposition}
\begin{proof}
Put $n=m+1$.  In word length $n$, the map
\[
 \rho:F_m\oplus F_m\longrightarrow[F,F]_n,
 \qquad (u,v)\longmapsto[x,u]+[y,v]
\]
is onto, because successive rotations by one position generate every
commutator of words.  Since $2\dim F_m=\dim F_n$, its kernel has dimension
$\dim(F_{\mathrm{cyc}})_n$.  By \eqref{cbc:neck:euler}, the cyclic gradient
$(D_x,D_y)$ embeds that space in $\ker\rho$, so it is an isomorphism onto
the kernel.

The condition $D\mu=0$ says $\rho(b,-a)=0$.  Its unique cyclic potential
satisfies $D_xf=b$, $D_yf=-a$, and the second Euler identity gives the
stated inverse.  Apply this argument to every homogeneous component of a
derivation.  Equation~\eqref{cbc:neck:commutator} proves compatibility with
the Lie brackets.
\end{proof}

\section{Power sums and the matrix rank}
On $N$ particle pairs $(x_i,y_i)$, set $\{x_i,y_j\}=\delta_{ij}$ and let
the symmetric group $S_N$ permute the pairs.  Define
\[
 R_N=\mathbb C[x_1,y_1,\ldots,x_N,y_N]^{S_N},\qquad
 p_{ab}^{(N)}=\sum_{i=1}^N x_i^ay_i^b\quad(a,b\geq0).
\]
Thus $p_{00}^{(N)}=N$.  Direct differentiation gives
\begin{equation}\label{cbc:neck:particle-bracket}
 \{p_{ab}^{(N)},p_{cd}^{(N)}\}
 =(ad-bc)p_{a+c-1,b+d-1}^{(N)}.
\end{equation}
When the coefficient is zero, the right side is zero without requiring a
power sum with a negative exponent.  In particular,
$\{p_{10}^{(N)},p_{01}^{(N)}\}=N$.

A uniform Poisson algebra retaining this constant is
\begin{equation}\label{cbc:neck:universal-rank}
 \mathcal P=\mathbb C[\nu,p_{ab}:a,b\geq0,\ a+b>0],
 \qquad p_{00}:=\nu,
\end{equation}
with $\nu$ central and bracket given by \eqref{cbc:neck:particle-bracket}.
To construct it, regard $\mathbb C[x,y]$ as a Lie algebra under its usual
Poisson bracket and take its symmetric algebra.  The generator $\nu$ is the
constant polynomial viewed as a Lie-algebra element.  The multiplicative
unit of the symmetric algebra is a separate element.

The Leibniz extension of this Lie bracket makes $\mathcal P$ a Poisson
algebra.  For each $N$, evaluation gives a Poisson map
\begin{equation}\label{cbc:neck:rank-evaluation}
 \pi_N:\mathcal P\longrightarrow R_N,
 \qquad \nu\longmapsto N,\quad p_{ab}\longmapsto p_{ab}^{(N)}.
\end{equation}
Chapter~\ref{chap:trace} proves its surjectivity and describes the stable
commutative algebra by an explicit orbit basis.

The substitution setting an additional particle pair to zero preserves all
positive-degree power sums.  It is not a Poisson map for these particle
brackets: the bracket of the two linear power sums changes from $N+1$ to
$N$.  The central parameter in \eqref{cbc:neck:universal-rank} records the
rank needed for each Poisson specialization.

\section{The relative two-particle surface}
For $N=2$, put $r=x_1-x_2$, $s=y_1-y_2$.  Then $\{r,s\}=2$, and the
transposition negates both coordinates.  The invariant functions
\[
 A=\frac{r^2}{2},\qquad B=\frac{rs}{2},\qquad C=\frac{s^2}{2}
\]
satisfy
\[
 B^2=AC,\qquad \{A,B\}=2A,\qquad
 \{B,C\}=2C,\qquad \{A,C\}=4B.
\]

The change of variables $e=A/2$, $f=-C/2$, $h=-B$ gives
\[
 \{h,e\}=2e,\qquad \{h,f\}=-2f,\qquad
 \{e,f\}=h,\qquad h^2+4ef=0.
\]
This is the Kirillov--Kostant bracket on the nilpotent cone of
$\mathfrak{sl}_2$.  The full invariant presentation and the Poisson splitting
of the center-of-mass plane are proved in Chapter~\ref{chap:trace}.
The derived matrix action constructed there retains more information than
this degree-zero surface.  Comparing that action with a full closed
observable algebra requires a further chain-level construction.

\part{The nonlocal ladder}

\chapter{Spectral exchange over configuration spaces}\label{cbc:spectral:chapter}
An exchange matrix varies with the positions of its inputs.  When Yang--Baxter
holds, the ordered-word theorem gives a holomorphic family of algebras
away from the poles.  Reversible exchange compares different orders and determines the
product between two groups of generators.  We construct these maps before
asking for differential structure or extension across collisions.

\section{The coefficient space and its algebra}
Let $V$ be a finite-dimensional complex vector space and let
$R(u)\in\operatorname{End}(V\otimes V)$ be a generically invertible
meromorphic matrix.  Choose a negation-invariant open set
$\Omega\subset\mathbb C\setminus\{0\}$ on which $R$ and $R^{-1}$ are
holomorphic.  For a finite set $I$, define
\[
 X_I=\{(t_a)_{a\in I}\in\mathbb C^I:
                  t_a-t_b\in\Omega\text{ whenever }a\ne b\}.
\]
This is an open subset of the configuration space of distinct labeled
points.  A total order $\prec$ on $I$ orders the labels, independently of
their complex coordinates.  All tensor products below are ordinary tensor
products over $\mathcal O_{X_I}$, unless another base is displayed.

Put $E_a=\mathcal O_{X_I}\otimes_{\mathbb C}V$, with its label $a$, and let
$P(v\otimes w)=w\otimes v$.  The exchange map is
\[
 c_{ab}=P R(t_a-t_b):E_a\otimes E_b\longrightarrow E_b\otimes E_a.
\]
Define $\mathcal A_I^{\prec}$ as the unital tensor algebra on
$\bigoplus_{a\in I}E_a$, modulo the relations
$x y=c_{ab}(x\otimes y)$ for $a\succ b$.
There is no relation within a single $E_a$.
We use the direct sum of word-degree components, with no completion.

Choose a basis $v_1,\ldots,v_d$ of $V$, and write $a_i^{(a)}$ for its
corresponding generator in $E_a$.  With
$R(u)(v_i\otimes v_j)=\sum_{k,l}R_{ij}^{kl}(u)v_k\otimes v_l$, the relation is
\begin{equation}\label{eq:vol1-zf}
 a_i^{(a)}a_j^{(b)}
 =\sum_{k,l}R_{ij}^{kl}(t_a-t_b)a_l^{(b)}a_k^{(a)}
 \qquad(a\succ b).
\end{equation}
An ordered word has nondecreasing colors and retains the order of its
generators within each color.

For operators on $V^{\otimes3}$, subscripts refer to the original tensor
factors.  The spectral Yang--Baxter equation is
\begin{equation}\label{eq:vol1-spectral-ybe}
 R_{12}(u)R_{13}(u+v)R_{23}(v)
 =R_{23}(v)R_{13}(u+v)R_{12}(u).
\end{equation}
It is evaluated where $u,v,u+v\in\Omega$.

\begin{proposition}[The holomorphic ordered-word family]\label{cbc:spectral:family}
Multiplication defines the map of $\mathcal O_{X_I}$-modules
\begin{equation}\label{cbc:spectral:ordered}
 \bigotimes_{a\in I}^{\prec}T_{\mathcal O_{X_I}}(E_a)
 \longrightarrow\mathcal A_I^{\prec}.
\end{equation}
It is an isomorphism if and only if \eqref{eq:vol1-spectral-ybe} holds on every triple
$a\succ b\succ c$, with $u=t_a-t_b$ and $v=t_b-t_c$.
Under this condition every word-degree component is finite free.
The presentation and its normal-form inverse commute with base change.
\end{proposition}
\begin{proof}
At every stalk, apply Theorem~\ref{cbc:exchange:normal} to the free modules
$E_a$ and the descending maps $c_{ab}$.  Its three-color equation becomes
\eqref{eq:vol1-spectral-ybe} after removing the common reversal permutation.
This proves both directions of the criterion.

The normal-form map consists of finitely many adjacent exchanges on each
word.  Its coefficients are holomorphic, so the stalkwise formulas define
a sheaf map inverse to multiplication.  In any fixed degree there are
only finitely many ordered words in the chosen bases.  They give the
asserted free basis, including words with repeated colors.

Both the defining relations and the normal-form formulas can be extended
to any new coefficient algebra.  Their two inverse identities remain
identities after this extension.  This proves the base-change assertion,
without requiring the new coefficient algebra to be flat.
\end{proof}

Assume henceforth that \eqref{eq:vol1-spectral-ybe} is a meromorphic
identity.  For $J\subseteq I$, forgetting coordinates gives
$p_J:X_I\to X_J$.  The generator inclusions define
\[
 p_J^*\mathcal A_J^{\prec|_J}\longrightarrow\mathcal A_I^{\prec}.
\]
The ordered bases show that this map is injective.  These inclusions
commute with further forgetting and with base change.

\section{Changing the order of tensor factors}
Write $R_{21}(u)=P R(u)P$.  A reversible exchange must also satisfy
\begin{equation}\label{cbc:spectral:reverse}
 R_{21}(-u)R(u)=1.
\end{equation}
Indeed, $c_{ba}c_{ab}$ is the left side of this equation on
$E_a\otimes E_b$.  Under this hypothesis the reverse relations are
consequences of the descending ones.  Imposing both orientations therefore
defines one algebra $\mathcal A_I$, independent of the chosen total order.

For an ordering $\sigma=(a_1,\ldots,a_m)$ of $I$, put
$W_\sigma=E_{a_1}\otimes\cdots\otimes E_{a_m}$.
Fix a reference order $\sigma_0$.  Let
$N_\sigma:W_\sigma\to W_{\sigma_0}$ be its normal-form map.
It is invertible because each elementary exchange is invertible.

\begin{proposition}[Transport between orders]\label{cbc:spectral:transport}
Assume \eqref{cbc:spectral:reverse}.  Every sequence of adjacent exchanges
from $W_\sigma$ to $W_\tau$ has the same composite,
\[
 C_{\sigma,\tau}=N_\tau^{-1}N_\sigma.
\]
Consequently $C_{\tau,\upsilon}C_{\sigma,\tau}=C_{\sigma,\upsilon}$ and
$C_{\sigma,\sigma}=1$.
\end{proposition}
\begin{proof}
For an adjacent exchange $c:W_\sigma\to W_\tau$ descending relative to
$\sigma_0$, confluence gives $N_\tau c=N_\sigma$.
For an ascending exchange, the same equality follows by applying this
argument to $c^{-1}$.  Multiplying these identities along any path gives
the displayed formula.  Its composition and identity laws follow directly.
\end{proof}

These maps compare orders of tensor factors.  They do not specify parallel
transport along paths in $X_I$: that would require a connection.

For the raw rational matrix $R(u)=1-\hbar P/u$,
Proposition~\ref{cbc:exchange:rational-result} gives
$R_{21}(-u)R(u)=1-\hbar^2/u^2$.
It supplies the oriented family, but fails \eqref{cbc:spectral:reverse}.
The normalized matrix
\[
 \widehat R(u)=\frac{u-\hbar P}{u-\hbar}
\]
is reversible over
\[
 \mathbb C[\hbar,(t_a)_{a\in I},
       (t_a-t_b)^{-1},(t_a-t_b-\hbar)^{-1},(t_a-t_b+\hbar)^{-1}:a\ne b].
\]
Its algebra is the explicit scalar twist of the oriented raw algebra in
\eqref{cbc:exchange:twist}.  This comparison changes the product.

\section{Multiplication between disjoint groups of generators}
Suppose exchange is reversible, and write $I=J\sqcup K$.
Pull the two smaller algebras to $X_I$ and abbreviate
\[
 B=p_J^*\mathcal A_J,\qquad C=p_K^*\mathcal A_K.
\]
Multiplication in $\mathcal A_I$ gives a map
\[
 \Phi_{J,K}:B\otimes C\longrightarrow\mathcal A_I,
 \qquad b\otimes c\longmapsto bc.
\]
It is an isomorphism of modules: choose the total order with all labels
of $J$ before those of $K$ and apply \eqref{cbc:spectral:ordered}.
Its inverse is the actual normal-form map for this order.

Define
\[
 \tau_{K,J}:C\otimes B\longrightarrow B\otimes C,
 \qquad \tau_{K,J}(c\otimes b)=\Phi_{J,K}^{-1}(cb).
\]
Thus $\tau_{K,J}$ sorts the $K$-word past the $J$-word.
If $m_B,m_C$ are the two internal multiplications, put
\begin{equation}\label{cbc:spectral:crossed}
 (b\otimes c)\star(b'\otimes c')
 =(m_B\otimes m_C)\bigl(b\otimes\tau_{K,J}(c\otimes b')\otimes c'\bigr).
\end{equation}

\begin{proposition}[The crossed tensor product]\label{cbc:spectral:blocks}
The product \eqref{cbc:spectral:crossed} is associative with unit
$1\otimes1$, and $\Phi_{J,K}$ is an algebra isomorphism for this product.
The multiplication maps for any finite partition of $I$ agree under
regrouping and insertion of empty subsets.
\end{proposition}
\begin{proof}
Write $\tau_{K,J}(c\otimes b')=\sum b_s\otimes c_s$.
Its definition gives $cb'=\sum b_sc_s$ in $\mathcal A_I$.  Hence
\[
 \Phi_{J,K}\bigl((b\otimes c)\star(b'\otimes c')\bigr)
 =\sum bb_sc_sc'=bc b'c'.
\]
The displayed product is therefore the multiplication of $\mathcal A_I$
transported by the module isomorphism $\Phi_{J,K}$.  It inherits
associativity and the stated unit.

For three subsets, both regroupings send $b\otimes c\otimes d$ to $bcd$.
For any number of subsets, every regrouping gives the same product in
$\mathcal A_I$.  Their normal-form inverses agree as well.
The empty subset contributes the scalar algebra and its unit.
These equalities prove the claimed coherence.
\end{proof}

The ordinary tensor-product multiplication generally differs from
\eqref{cbc:spectral:crossed}.  For $B=\mathbb C[x]$, $C=\mathbb C[y]$ and
$yx=qxy$, the crossing sends $y\otimes x$ to $q x\otimes y$.
Ordinary tensor multiplication sends
$(1\otimes y)(x\otimes1)$ to $x\otimes y$, whose image under $\Phi$ is $xy$.
The product of their images is instead $yx=qxy$.
For $q\ne1$, independence of the ordered monomials prevents these maps
from being algebra homomorphisms for the ordinary tensor product.

\section{Diagonal exchange and its three-particle coefficient}
Choose meromorphic units $q_{ij}(u)$ satisfying
\[
 q_{ij}(u)q_{ji}(-u)=1,\qquad q_{ii}(u)=1,
\]
and set $R(u)(v_i\otimes v_j)=q_{ij}(u)v_i\otimes v_j$.
The three fixed-leg operators in \eqref{eq:vol1-spectral-ybe} are diagonal
and commute.  The family is therefore reversible and has the ordered basis.
Its relations are
\[
 a_i^{(a)}a_j^{(b)}
 =q_{ij}(t_a-t_b)a_j^{(b)}a_i^{(a)}\qquad(a\succ b).
\]

For $a\succ b\succ c$, both complete reductions give
\begin{equation}\label{cbc:spectral:triple}
 \begin{split}
 a_i^{(a)}a_j^{(b)}a_k^{(c)}
 &=q_{ij}(t_a-t_b)q_{ik}(t_a-t_c)q_{jk}(t_b-t_c)\\
 &\hspace{2em}\cdot a_k^{(c)}a_j^{(b)}a_i^{(a)}.
 \end{split}
\end{equation}
Each inverted pair is exchanged exactly once.  Scalar commutativity makes
the product independent of the chosen sequence.  For a general word,
the coefficient is the product over all initially inverted color pairs,
with the component indices attached to their respective occurrences.
Repeated colors retain their internal order.

For constant $q_{12}=q$, $q_{21}=q^{-1}$, select two distinct colors
$a\succ b$ and put $x=a_1^{(a)}$, $y=a_2^{(b)}$.
Their subalgebra is the quantum plane $\mathbb C\langle x,y\rangle/(xy-qyx)$:
the independent monomials $y^r x^s$ prove that there are no further relations.
This specialization uses distinct colors and imposes no same-color relation.

\section{A coefficient connection is additional structure}
A varying exchange coefficient need not admit a connection preserving its
generator lines.  Set $K=\mathbb C(u)$ and consider
\[
 A=K\langle x,y\rangle/(yx-q(u)xy),\qquad
 q(u)=\frac{u+1}{u-1}.
\]
Here $q(u)q(-u)=1$.  This scalar example occurs in the diagonal family with
$q_{12}=q_{21}=q$ and $q_{11}=q_{22}=1$.  Take $x=a_1^{(b)}$,
$y=a_2^{(a)}$ with $a\succ b$, and put $u=t_a-t_b$.

A connection over $\partial_u$ is a $\mathbb C$-linear operator $\nabla$ with
$\nabla(fa)=f'a+f\nabla(a)$ for $f\in K$.
Suppose it is also an algebra derivation and preserves $Kx$ and $Ky$.
Then $\nabla x=\alpha x$, $\nabla y=\beta y$ for some $\alpha,\beta\in K$.

Applying $\nabla$ to the defining relation gives
\[
 0=(\alpha+\beta)(yx-qxy)-q'(u)xy
   =\frac{2}{(u-1)^2}xy.
\]
The last element is nonzero by the ordered basis.  No connection with
these stated properties exists.  This includes the prescription making
both generators horizontal.

For constant $q$, coefficientwise differentiation with horizontal
generators does preserve the relation and defines such a connection.
The obstruction concerns this varying algebra and the prescribed generator
lines.  A different carrier or differential structure requires its own analysis.

The family $\mathcal A_I$ and its crossed products are defined on the
distinct-point opens $X_I$.  An associative chiral realization must specify
the coefficient category and tensor, differential structure, and extensions
at collision diagonals and the excluded poles.  It must construct
collision and substitution maps, compatible relabelings, factorization
comparisons and all their required coherences.

The comparison must identify its chiral operation with the exchange maps
constructed here.  A field realization additionally needs the state space,
vacuum, translation, Laurent expansions and an exchange operator on all
states.  Constructing these structures and their comparison maps remains
the realization problem.

\chapter{RTT and reflection algebras}\label{cbc:reflection:chapter}
Matrix multiplication and reflection impose different compatibility laws.
The product of two independent RTT matrices is again an RTT matrix, which
constructs the coproduct.  The reflection equation instead gives operators
for braids with a boundary.  We construct their covariance and then ask
when the matrix entries admit an ordered basis.

\section{Matrix composition and the two coproduct conventions}
Let $B$ be a commutative ring and $V=B^n$ with a chosen basis.
All algebras are unital $B$-algebras, and all tensor products are ordinary
tensor products over $B$.
For $R\in\operatorname{End}_B(V\otimes V)$, let
$\mathcal T_R$ be the free algebra on $t_{ij}$, modulo the entries of
\begin{equation}\label{cbc:reflection:constant-rtt}
 R T_1T_2=T_2T_1R,\qquad T=(t_{ij}),\quad
 T_1=T\otimes I_V,\quad T_2=I_V\otimes T.
\end{equation}
The subscripts here denote auxiliary copies of $V$.
The matrix entries multiply in the coefficient algebra.

There is also a coefficientwise current version.  Fix a polynomial matrix
$Q(u,v)\in\operatorname{End}_B(V\otimes V)[u,v]$ and put
\[
 T(u)=I_V+\sum_{r\geq1}t^{(r)}u^{-r}.
\]
Define $\mathcal T_Q^{\mathrm{cur}}$ by the entries of every Laurent
coefficient in
\begin{equation}\label{eq:vol1-rtt-general}
 Q(u,v)T_1(u)T_2(v)=T_2(v)T_1(u)Q(u,v).
\end{equation}

Each coefficient is a finite expression because $Q$ is polynomial.
The algebra is generated by the $t_{ij}^{(r)}$ without completion.
The displayed current is a formal series with coefficients in that algebra.
The choice of the cleared kernel $Q$ is part of this presentation.

\begin{lemma}[Multiplication of independent matrices]\label{cbc:reflection:product}
Suppose $X$ and $Y$ satisfy \eqref{cbc:reflection:constant-rtt} with
coefficients in $B$-algebras $D$ and $E$, respectively.
Their product $XY$, with coefficients in $D\otimes E$, satisfies the same
equation.  The analogous assertion holds for
\eqref{eq:vol1-rtt-general}.
\end{lemma}
\begin{proof}
Entries from the two coefficient factors commute.  Operators on different
auxiliary factors also commute, so
\[
 \begin{aligned}
 R X_1Y_1X_2Y_2
 &=R X_1X_2Y_1Y_2
  =X_2X_1R Y_1Y_2\\
 &=X_2X_1Y_2Y_1R
  =X_2Y_2X_1Y_1R.
 \end{aligned}
\]
This is the required equation for $XY$.
Replacing $R$ by $Q(u,v)$ gives the current calculation coefficientwise.
No Yang--Baxter identity is needed for this argument.
\end{proof}

Use superscripts $[1]$ and $[2]$ for the two coefficient factors.
Thus $(T^{[1]})_{ij}=t_{ij}\otimes1$ and
$(T^{[2]})_{ij}=1\otimes t_{ij}$.

\begin{theorem}[RTT coproduct]\label{cbc:reflection:coproduct}
The constant RTT algebra has the bialgebra structure
\[
 \Delta_+(t_{ij})=\sum_k t_{ik}\otimes t_{kj},\qquad
 \epsilon(t_{ij})=\delta_{ij}.
\]
Its opposite coproduct is another bialgebra structure:
\[
 \Delta_-(t_{ij})=\sum_k t_{kj}\otimes t_{ik}.
\]
In matrix notation these are
$\Delta_+(T)=T^{[1]}T^{[2]}$ and
$\Delta_-(T)=T^{[2]}T^{[1]}$.
For currents, set $t_{ij}^{(0)}=\delta_{ij}$.  Then
\begin{equation}\label{cbc:reflection:mode-coproduct}
 \Delta_+(t_{ij}^{(r)})
 =\sum_{a+b=r}\sum_k t_{ik}^{(a)}\otimes t_{kj}^{(b)},\qquad
 \epsilon(t_{ij}^{(r)})=0\quad(r\geq1),
\end{equation}
and $\Delta_-$ is the flip of this formula.
Both coproducts take values in the ordinary algebraic tensor product.
\end{theorem}
\begin{proof}
Lemma~\ref{cbc:reflection:product} shows that both matrix products preserve
the defining relations.  Evaluation at the identity matrix defines the
counit.  In three coefficient factors, the two iterates of $\Delta_+$
are $T^{[1]}T^{[2]}T^{[3]}$, while those of $\Delta_-$ are
$T^{[3]}T^{[2]}T^{[1]}$.
Matrix associativity proves coassociativity.
Replacing either coefficient factor by the identity proves the counit
identities.  The same calculation applies to currents, and every mode sum
in \eqref{cbc:reflection:mode-coproduct} is finite.
\end{proof}

If the auxiliary factor is numbered first, the two matrix formulas read
$\Delta_+(T)=T_{12}T_{13}$ and $\Delta_-(T)=T_{13}T_{12}$.
These conventions describe opposite coproducts on the same algebra.

For the Yangian, take $B=\mathbb C[\hbar]$ and
$Q(u,v)=(u-v)I-\hbar P$, where $P(x\otimes y)=y\otimes x$.
The coefficient calculation in Proposition~\ref{cbc:rtt:coefficients}
gives
\[
 [t_{ij}^{(1)},t_{kl}^{(1)}]
 =\hbar\bigl(\delta_{kj}t_{il}^{(1)}-\delta_{il}t_{kj}^{(1)}\bigr).
\]
The map $X_{ij}\mapsto t_{ij}^{(1)}$ identifies the first-mode subalgebra
with $U_\hbar(\mathfrak{gl}_n)$.
Its inverse on that subalgebra is the restriction of the evaluation map
\eqref{cbc:rtt:evaluation} at zero.  It sends $t^{(1)}$ to $X$ and all
higher modes to zero.  This uses the unscaled current convention
$T(u)=I+\sum t^{(r)}u^{-r}$ over $\mathbb C[\hbar]$.

\section{The reflection equation and boundary braids}
Assume now that $R$ is invertible and satisfies
\[
 R_{12}R_{13}R_{23}=R_{23}R_{13}R_{12}
 \quad\text{on }V^{\otimes3}.
\]
Write $R_{21}=PRP$.
The reflection algebra $\mathcal K_R$ is the free algebra on entries
$k_{ij}$ of a matrix $K$, modulo the entries of
\begin{equation}\label{eq:vol1-reflection}
 R_{21}K_1R_{12}K_2=K_2R_{21}K_1R_{12}.
\end{equation}
The matrix $K$ is not assumed invertible.
Set $C=PR$.
Since $PK_1=K_2P$, the reflection equation is exactly
\begin{equation}\label{cbc:reflection:checked}
 C K_2 C K_2=K_2 C K_2 C.
\end{equation}

For $m\geq2$, define the positive Artin monoid of type $B_m$ by generators
$\sigma_1,\ldots,\sigma_{m-1},\kappa$ and relations
\begin{equation}\label{cbc:reflection:artin}
 \begin{aligned}
 \sigma_i\sigma_{i+1}\sigma_i
  &=\sigma_{i+1}\sigma_i\sigma_{i+1},\qquad 1\leq i<m-1,\\
 \sigma_i\sigma_j&=\sigma_j\sigma_i,\qquad |i-j|>1,\\
 \kappa\sigma_i&=\sigma_i\kappa,\qquad i<m-1,\\
 \sigma_{m-1}\kappa\sigma_{m-1}\kappa
  &=\kappa\sigma_{m-1}\kappa\sigma_{m-1}.
 \end{aligned}
\end{equation}
Adjoining inverses to all generators gives the Artin group of type $B_m$.
The generator $\kappa$ represents a boundary operation on the last factor.

\begin{theorem}[Boundary braid coherence]\label{cbc:reflection:braid}
On the free right $\mathcal K_R$-module
$M_m=V^{\otimes m}\otimes\mathcal K_R$, the operators
\[
 \sigma_i\longmapsto C_{i,i+1},\qquad
 \kappa\longmapsto K_m
\]
define a representation of the monoid \eqref{cbc:reflection:artin}.
Here $K_m$ acts on the last auxiliary factor, with its entries acting by
left multiplication on $\mathcal K_R$.
After an algebra map $\mathcal K_R\to D$, these operators give a group
representation if the resulting matrix $K\in\operatorname{Mat}_n(D)$
is invertible.
\end{theorem}
\begin{proof}
Moving the flips to the left gives
\[
 \begin{aligned}
 C_{12}C_{23}C_{12}&=P_{12}P_{23}P_{12}R_{23}R_{13}R_{12},\\
 C_{23}C_{12}C_{23}&=P_{23}P_{12}P_{23}R_{12}R_{13}R_{23}.
 \end{aligned}
\]
The permutation products agree, and Yang--Baxter identifies the remaining
factors.  Operators on disjoint auxiliary factors commute.
This proves the first three families of relations in
\eqref{cbc:reflection:artin}.
The final relation is \eqref{cbc:reflection:checked} on the last two
factors.  Thus every defining relation is respected.

When $K$ is invertible, all operators are invertible and the representation
extends to the group.
\end{proof}

This theorem compares boundary braid words acting on $M_m$.
A normal form for words in the entries $k_{ij}$ is a separate algebraic
question.

\section{Covariance under an RTT Hopf algebra}
Suppose $H$ is a Hopf $B$-algebra with a matrix $T=(t_{ij})$ satisfying
\eqref{cbc:reflection:constant-rtt},
$\Delta(T)=T^{[1]}T^{[2]}$ and $\epsilon(T)=I$.
If $S_H$ is its antipode, the matrix $U=(S_H(t_{ij}))$ is the inverse of
$T$.  Place the entries of $K$ and $T$ in the separate factors of
$\mathcal K_R\otimes H$.

\begin{proposition}[Right covariance]\label{cbc:reflection:coaction}
The formula
\begin{equation}\label{cbc:reflection:delta}
 \delta:\mathcal K_R\longrightarrow\mathcal K_R\otimes H,
 \qquad
 \delta(k_{ij})=\sum_{a,b}k_{ab}\otimes u_{ia}t_{bj},
\end{equation}
defines a right $H$-comodule algebra.  In matrix notation,
$\delta(K)=UKT$.
\end{proposition}
\begin{proof}
The RTT equation, its conjugate by $P$, and the inverse identities give
\[
 T_1R U_2=U_2R T_1,\qquad
 R_{21}U_1U_2=U_2U_1R_{21},\qquad
 U_1R_{21}T_2=T_2R_{21}U_1.
\]
For example, multiplying $RT_1T_2=T_2T_1R$ on the left and right by
$U_2$ gives the first identity.
Put $K'=UKT$.  Using these identities in order, and commuting only
entries from separate coefficient and auxiliary factors, gives
\[
 \begin{aligned}
 R_{21}K'_1R K'_2
 &=R_{21}U_1U_2K_1R K_2T_1T_2\\
 &=U_2U_1R_{21}K_1R K_2T_1T_2\\
 &=U_2U_1K_2R_{21}K_1R T_1T_2\\
 &=U_2K_2U_1R_{21}T_2K_1T_1R\\
 &=K'_2R_{21}K'_1R.
 \end{aligned}
\]
Thus \eqref{cbc:reflection:delta} preserves every defining relation and
extends to an algebra homomorphism.

In $\mathcal K_R\otimes H\otimes H$, both
$(\delta\otimes1)\delta(K)$ and $(1\otimes\Delta)\delta(K)$ equal
\[
 (T^{[2]})^{-1}(T^{[1]})^{-1}K T^{[1]}T^{[2]}.
\]
The counit reduces this matrix to $K$.
These are the right coaction identities.
\end{proof}

The side of the coaction and the inverse placement are fixed by the
chosen reflection equation.  Other transpose and inverse placements have
different covariance formulas
\cite[\S6.2]{KulishSklyaninReflection1992}.

\section{Ordered bases for reflection entries}
Two invertible Yang--Baxter matrices already distinguish braid coherence
from the usual polynomial-sized basis.
For $R=I$, equation~\eqref{eq:vol1-reflection} is $K_1K_2=K_2K_1$.
Its entries are exactly the pairwise commutators of the $n^2$ generators,
so $\mathcal K_I$ is a polynomial algebra.

For $R=P$, the identity $PK_1P=K_2$ makes both sides equal to $K_2^2$
before any quotient.  Thus $\mathcal K_P$ is the free algebra on the same
$n^2$ generators.  Over a field, their degree-two dimensions are
\[
 \dim(\mathcal K_I)_2=\frac{n^2(n^2+1)}2,\qquad
 \dim(\mathcal K_P)_2=n^4.
\]

For $n\geq2$ they differ.  Invertibility and Yang--Baxter therefore do not
force a basis indexed by nondecreasing words in the matrix entries.
Theorem~\ref{cbc:reflection:braid} still applies in both cases.

A nontrivial positive example is the standard two-dimensional matrix
\cite[(5)]{KulishSklyaninReflection1992}.
Work over $A=\mathbb C[q,q^{-1}]$, with $V=A^2$ and basis $v_0,v_1$.
Write $e_{i_1\cdots i_m}=v_{i_1}\otimes\cdots\otimes v_{i_m}$ and put
$\omega=q-q^{-1}$.  In the ordered basis
$(e_{00},e_{01},e_{10},e_{11})$, take
\[
 R=\begin{pmatrix}
 q&0&0&0\\
 0&1&0&0\\
 0&\omega&1&0\\
 0&0&0&q
 \end{pmatrix},\qquad
 K=\begin{pmatrix}a&b\\c&d\end{pmatrix}.
\]

The checked matrix acts by
\[
 C(e_{00})=qe_{00},\quad C(e_{01})=\omega e_{01}+e_{10},\quad
 C(e_{10})=e_{01},\quad C(e_{11})=qe_{11}.
\]

Both three-factor braid products act by $q^3$ on $e_{000}$ and $e_{111}$.
In each ordered basis $(e_{001},e_{010},e_{100})$ and
$(e_{011},e_{101},e_{110})$ their common
matrix is
\[
 \begin{pmatrix}
 q^2\omega&q\omega&q\\
 q\omega&q&0\\
 q&0&0
 \end{pmatrix}.
\]
These eight basis vectors prove the braid identity, using
$q\omega+1=q^2$.  The preceding flip calculation gives quantum
Yang--Baxter for $R$, which is invertible because its determinant is $q^2$.

Put $\rho=q^{-2}$ and $\lambda=1-\rho$.
Multiplying the matrices in \eqref{eq:vol1-reflection} gives the six
relations
\begin{equation}\label{cbc:reflection:rules}
 \begin{aligned}
 bd&=\rho db,& cd&=q^2dc,& ad&=da,\\
 cb&=bc+\lambda da-\lambda d^2,&
 ab&=ba-\lambda\rho db,&
 ac&=ca+\lambda dc.
 \end{aligned}
\end{equation}

For completeness, the intermediate matrix is
\[
 R_{21}K_1R=
 \begin{pmatrix}
 q^2a&q\omega b&qb&0\\
 q\omega c&a+\omega^2d&\omega d&qb\\
 qc&\omega d&d&0\\
 0&qc&0&q^2d
 \end{pmatrix}.
\]

Let $F=R_{21}K_1RK_2-K_2R_{21}K_1R$, with rows and columns numbered
from zero.  The first five relations in the order
$bd,cd,ad,cb,ac$ follow respectively from
$F_{23},F_{31},F_{22},F_{12},F_{10}$ by multiplying by units in $A$.
The remaining relation follows from $F_{02}$ and the first one.
Conversely, substituting \eqref{cbc:reflection:rules} makes every entry
of $F$ vanish.  Thus these are equivalent presentations.

\begin{theorem}[A complete reflection basis]\label{cbc:reflection:pbw}
For this matrix $R$, the monomials
\[
 d^i b^j c^k a^\ell,\qquad i,j,k,\ell\geq0,
\]
form an $A$-basis of $\mathcal K_R$.
They remain a basis after every specialization $q\mapsto q_0\ne0$.
\end{theorem}
\begin{proof}
Order the generators by $d<b<c<a$ and orient
\eqref{cbc:reflection:rules} from left to right.
Each replacement strictly decreases the length-lexicographic order, so
reductions terminate.  There are no inclusions between leading words.

The only overlapping pairs of reductions occur in $cbd$, $abd$, $acd$
and $acb$.  Reducing first either the left or the right pair gives the
same normal expression:
\[
 \begin{array}{c|l}
 \text{word}&\text{normal expression from either first reduction}\\\hline
 cbd&dbc+\lambda d^2a-\lambda d^3\\
 abd&\rho dba-\lambda\rho^2d^2b\\
 acd&q^2dca+\lambda q^2d^2c\\
 acb&bca+\lambda da^2-\lambda d^2a.
 \end{array}
\]
Reductions on disjoint pairs commute.

Induct on the word order.  Two distinct first reductions either act on
disjoint pairs or form one of the four overlaps.  In either case they reduce
to the same expression.  Every remaining word is smaller, so the induction
hypothesis gives a unique complete reduction.

Extend the normal map $A$-linearly to the free algebra.
It annihilates every defining relation in every word context, because one
may perform that particular replacement first.
It therefore descends to $\mathcal K_R$ and is the identity on ordered
monomials.  Multiplication of ordered monomials into the quotient is its
inverse, proving both spanning and independence.
All rules and inverse identities are defined over $A$, so base change to
any nonzero $q_0$ preserves the basis.
\end{proof}

The RTT coproduct, reflection coaction and boundary braid action are
explicit algebraic operations on different carriers.
A physical line realization must specify its bulk line objects, boundary
modules and their compositions, and construct maps comparing those
operations with the ones above.
An interpretation in terms of charged fields additionally requires the
state space and field maps of the exchange family.
These representation and comparison constructions are not consequences
of imposing the RTT or reflection equations.

\chapter{Quantum vertex realization}
\section{Two expansion regions}
Let $V$ be a complete topological vector space with vacuum $\mathbf1$, translation $T$, state-field map $Y$, and a unitary rational operator
\[
 S(z):V\tensor V\longrightarrow V\tensor V((z)).
\]
The two formal regions $|z_1|>|z_2|$ and $|z_2|>|z_1|$ are related by $S(z_1-z_2)$.

\begin{definition}[Rational quantum vertex datum]
A rational quantum vertex datum satisfies vacuum and translation covariance together with
\begin{equation}\label{eq:vol1-s-locality}
 (z_1-z_2)^N Y(a,z_1)Y(b,z_2)
 =(z_1-z_2)^N
 Y^{21}\!\left(S(z_1-z_2)(a\tensor b);z_2,z_1\right)
\end{equation}
for a sufficiently large $N$, and weak associativity in every triple of states.
\end{definition}

\begin{theorem}[Exchange coherence]
Suppose $S$ is unitary,
\[
 S^{21}(-z)S(z)=1,
\]
and satisfies the spectral Yang--Baxter equation.  Then the pairwise continuation maps obtained from \eqref{eq:vol1-s-locality} define a braid-group action on every $n$-point field product.  Weak associativity makes the action compatible with iterated operator products.
\end{theorem}
\begin{proof}
Unitarity identifies a crossing followed by its reverse with the identity.  Yang--Baxter identifies the two continuations around every codimension-two triple-collision stratum.  These are the defining relations of the braid group.  Weak associativity identifies the collision limit in each chamber with the corresponding iterated field.
\end{proof}

An exchange datum describes relations among field products.
Its realization by an associative chiral algebra requires the geometric construction and operation comparison of Problem~\ref{cbc:two:reverse}.
The exchange normalization must specify how $S$ combines with the reversal denoted by $Y^{21}$ in~\eqref{eq:vol1-s-locality}.
It must recover ordinary locality when their combined action is the ordinary interchange of the two fields.
A comparison with the $\hbar$-adic convention also requires a formal deformation with $S=1+O(\hbar)$ and the stated completed tensor products.
The field theories of~\cite{LiNonlocal,LiQVA} do not by themselves supply this geometric comparison.

\section{The Yangian field}
Let $T(u)$ act on a family of evaluation modules.  The RTT relation becomes $S$-locality for the matrix coefficients of $T(u)$.  The normalized Yang operator
\[
 S(u)=f(u)\left(1-\frac{\hbar P}{u}\right),
 \qquad
 f(u)f(-u)=\left(1-\frac{\hbar^2}{u^2}\right)^{-1},
\]
is unitary.  Its expansion at infinity gives the quantum vertex exchange of Yangian currents.  The scalar factor changes normalization and preserves the Yang--Baxter equation.

\chapter{Shuffle realization and the positive current algebra}
A shuffle product combines spectral variables by assigning
a scalar factor to each cross-block pair.
Associativity follows from the combinatorics of these pairs.
The coefficient lattice and the tensor used by a coproduct require
further choices.  We shall make those choices explicit, beginning
with the product of two currents.

\section{The product and its normalization}
Let $k$ be a field of characteristic zero and let $K(u)\in k(u)$.
Set
\[
 \operatorname{Sh}_n=k(z_1,\ldots,z_n)^{\mathfrak S_n},
 \qquad \operatorname{Sh}=\bigoplus_{n\geq0}\operatorname{Sh}_n.
\]
The integer $n$ is the \emph{charge}.  All cochain degrees and parities
are zero.  The direct sum is algebraic, so each element has finite
charge support.

For $F\in\operatorname{Sh}_m$ and
$G\in\operatorname{Sh}_n$, define
\begin{equation}\label{eq:vol1-shuffle}
 (F*_KG)(Z)=
 \sum_{\substack{I\sqcup J=[m+n]\\|I|=m}}
 F(Z_I)G(Z_J)\prod_{i\in I,\,j\in J}K(z_i-z_j).
\end{equation}
Here $[r]=\{1,\ldots,r\}$, and $Z_I$ is the subtuple indexed by $I$.
The inputs are symmetric, so their internal order is immaterial.
There is one summand for each subset $I$ and no factorial divisor.

The algebra unit is $1_0=1\in\operatorname{Sh}_0=k$.
The constant $1_n$ in positive charge is a different element.

\begin{theorem}[Associativity of the shuffle product]
Formula~\eqref{eq:vol1-shuffle} defines an associative unital
charge-graded $k$-algebra.  Every graded subspace containing $1_0$
and closed under this product is a unital subalgebra.
\end{theorem}
\begin{proof}
For inputs of charges $m,n,p$, both parenthesizations sum over
ordered partitions $I\sqcup J\sqcup L=[m+n+p]$ of those sizes.
The summand for this partition is
\[
 F(Z_I)G(Z_J)H(Z_L)
 \prod_{i\in I,j\in J}K(z_i-z_j)
 \prod_{i\in I,l\in L}K(z_i-z_l)
 \prod_{j\in J,l\in L}K(z_j-z_l).
\]
Each factor occurs once in either parenthesization.  Permuting the
variables permutes these summands, proving symmetry of the output.
An empty block has one choice and an empty product, which proves
both unit equations.  The same argument for any number of ordered
blocks identifies every iterated product.
\end{proof}

Averaging the displayed summand over all $(m+n)!$ permutations gives
$m!n!/(m+n)!$ times~\eqref{eq:vol1-shuffle}.
Multiplication by $n!$ in charge $n$ is an algebra isomorphism from
that averaged convention to ours.  Summing over all permutations
without either normalization is different: for $K=1$ and constant
inputs of charges $1,1,2$, the two bracketings give $48$ and $144$.

\section{The cubic kernel and two currents}
To retain parameter specialization, work over
\[
 R=\mathbb Q[h_1,h_2],\qquad h_3=-h_1-h_2,
 \qquad a=\sum_{i<j}h_ih_j,\qquad b=h_1h_2h_3.
\]
Thus $a=-(h_1^2+h_1h_2+h_2^2)$ and $b=-h_1h_2(h_1+h_2)$.
Write
\[
 N(u)=\prod_{i=1}^3(u+h_i)=u^3+au+b,
 \qquad K_j(u)=u^{2j}\frac{N(u)}u\quad(j=-1,0,1).
\]

In particular, $K_{-1}(u)=\prod_i(u+h_i)/u^3$ is the cubic
rational kernel.  Its coefficients lie in $R$.  No parameter has been inverted.  

For this kernel take the carrier
\[
 \mathcal A^{(-1)}_n
 =R[z_1,\ldots,z_n,D_n^{-1}]^{\mathfrak S_n},
 \qquad D_n=\prod_{i<l}(z_i-z_l)^2,\quad D_0=D_1=1.
\]
Set $\mathcal A^{(-1)}=\bigoplus_n\mathcal A^{(-1)}_n$.
Subset shuffle preserves these modules.  We use $*_j$ for the
product with kernel $K_j$.

For $f,g\in R[z]$ and $u=x-y$, direct multiplication gives
\begin{align*}
 f*_jg&=f(x)g(y)K_j(u)+f(y)g(x)K_j(-u),\\
 [f,g]_{*_j}
 &=2b\,u^{2j-1}\bigl(f(x)g(y)-g(x)f(y)\bigr).
\end{align*}
The bracket is the associative commutator $F*_jG-G*_jF$.
If $e_r(z)=z^r$ is the charge-one current, then
\begin{equation}\label{cbc:l027:binary-currents}
 [e_0,e_1]_{*_{-1}}=-\frac{2b}{u^2},
 \qquad e_0*_{-1}e_0=2\left(1+\frac a{u^2}\right).
\end{equation}
The alternating numerator of the commutator cancels one diagonal
factor.  It does not create poles at shifted differences.

At generic parameters the negative shifts $u=-h_i$ are zeros of
$K_j$, and the only other zero or pole is at $u=0$.
In divisor notation, with $D_\lambda=\{u=\lambda\}$,
\[
 \operatorname{div}(K_j)=\sum_iD_{-h_i}+(2j-1)D_0.
\]

The reverse-to-forward exchange ratio is
\begin{equation}\label{cbc:l027:ratio}
 r(u)=\frac{K_j(-u)}{K_j(u)}
 =\frac{u^3+au-b}{u^3+au+b}
 =\frac{\prod_i(u-h_i)}{\prod_i(u+h_i)}.
\end{equation}
It is independent of $j$.  The negative shifts are generic poles
of this ratio, although they are zeros of the product kernel.
When parameters specialize, divisors can coincide and factors can
cancel.  The rational formulas determine the resulting orders.

\section{Polynomial coefficients and the generated algebra}
For $j=0,1$ define the full polynomial carrier
\[
 \mathcal A^{(j)}=
 \bigoplus_{n\geq0}R[z_1,\ldots,z_n]^{\mathfrak S_n}.
\]
Both carriers are closed under their respective shuffle products.
For $K_1$ this follows because the kernel is polynomial.

For $K_0$, every possible denominator in a summand is a simple
cross-block difference.  Pair that summand with the one obtained by
exchanging the two colliding labels.  On their diagonal all other
factors agree, and the residues of $K_0(u)$ and $K_0(-u)$ cancel.
Cancellation at every possible prime denominator in the polynomial
unique factorization domain proves polynomiality.  The
ordered-partition proof gives associativity over $R$.

The \emph{current-generated algebra} $\mathcal C^{(j)}$ is the
unital $R$-subalgebra generated by the whole charge-one module
$R[z]\subset\mathcal A^{(j)}$, for $j=-1,0,1$.
Its charge-two part can already be smaller than the ambient carrier.

Put
\[
 c=(x+y)/2,\qquad u=x-y,\qquad t=u^2,\qquad B=R[c,t].
\]
Then $R[x,y]=B\oplus uB$, and $B=R[x,y]^{\mathfrak S_2}$.
For an arbitrary input polynomial $P=A+uC$, its linearized shuffle
product is
\begin{equation}\label{cbc:l027:binary-row}
 \mu_j(P)=K_j(u)(A+uC)+K_j(-u)(A-uC)
 =2t^j\bigl((t+a)A+bC\bigr).
\end{equation}
This computes the entire binary map, since pure tensors span
$R[x,y]$.  Consequently
\[
 \mathcal C^{(0)}_2=(t+a,b)B,
 \qquad \mathcal C^{(1)}_2=t(t+a,b)B.
\]

The first cokernel is $R[c]/(b)$, with $R$-annihilator $(b)$.
The second cokernel maps onto $R[c]$ by setting $t=0$, so its
$R$-annihilator is zero.  A charge-preserving $R$-linear algebra
isomorphism of the two full polynomial carriers would induce an
isomorphism of these multiplication cokernels.  Thus no such
ambient isomorphism exists.

There is nevertheless an actual comparison on the generated
algebras.  For every split into subsets $I,J$,
\begin{equation}\label{cbc:l027:discriminant}
 D_{m+n}(Z)=D_m(Z_I)D_n(Z_J)
 \prod_{i\in I,l\in J}(z_i-z_l)^2.
\end{equation}
Multiplying each shuffle summand by this identity proves
\[
 D_{m+n}(F*_jG)=(D_mF)*_{j+1}(D_nG),\qquad j=-1,0.
\]

Starting with charge-one polynomials, the right side is polynomial
for $j=-1$, by the polynomiality just proved for $K_0$.
Since multiplication by $D_n$ is injective and fixes charge one,
its image on the generated algebra is precisely the algebra
generated by those same currents.  We obtain algebra isomorphisms
\begin{equation}\label{cbc:l027:generated-comparison}
 \mathcal C^{(-1)}\xrightarrow{\;D\;}
 \mathcal C^{(0)}\xrightarrow{\;D\;}
 \mathcal C^{(1)}.
\end{equation}

The inverses divide by $D_n$ on the specified images.
The map on the full polynomial carriers is not surjective in charge
$n\geq2$, because its image vanishes at $z_1=z_2$.
After diagonal localization it becomes an isomorphism of the full
localized carriers.  Factorization~\eqref{cbc:l027:discriminant}
for any ordered partition proves compatibility with all iterated
products.  These maps also commute with simultaneous translation
of the spectral variables, since $D_n$ depends only on differences.

A wheel condition is a vanishing condition on a triple of shifted
roots.  Every polynomial in $\mathcal C^{(0)}_n$ vanishes when three
distinct labels satisfy
\[
 z_i-z_l=h_\alpha,\qquad z_l-z_m=h_\beta,
 \qquad z_m-z_i=h_\gamma,
 \qquad \{\alpha,\beta,\gamma\}=\{1,2,3\}.
\]

Indeed, express a product of currents as a sum over their total
orders.  No total order respects all three edges of this cycle.
A reversed edge supplies a zero factor $K_0(-h_\alpha)$.
Over the fraction field of $R$, the three roots are distinct and
this evaluation avoids their diagonal poles.  The resulting
polynomial vanishing identity holds over $R$ as well.

Such conditions do not determine the integral generated algebra:
in charge two there is no triple wheel, whereas
$(t+a,b)B$ is a proper submodule of $B$.

\section{The generic positive affine Yangian}
Now choose complex parameters $h_1,h_2,h_3$ with sum zero such that
\[
 m_1h_1+m_2h_2+m_3h_3=0,\quad m_i\in\mathbb Z
 \quad\Longrightarrow\quad m_1=m_2=m_3.
\]
Form the shuffle algebra at these numerical parameters, and let
$\mathcal C^{(-1)}_h$ be the subalgebra generated by $\mathbb C[z]$.
This means specialization of the kernel followed by generation
inside the rational carrier.

The positive affine Yangian $Y_h^+$ is the unital complex algebra
on $E_r$, $r\geq0$, with relations
\begin{align}
 &[E_{r+3},E_s]-3[E_{r+2},E_{s+1}]
   +3[E_{r+1},E_{s+2}]-[E_r,E_{s+3}]
 \notag\\[-2pt]
 &\hspace{1cm}
   +a\bigl([E_{r+1},E_s]-[E_r,E_{s+1}]\bigr)
   =b(E_rE_s+E_sE_r),\label{cbc:l027:yangian-quadratic}\\
 &\sum_{\sigma\in\mathfrak S_3}
 [E_{i_{\sigma(1)}},[E_{i_{\sigma(2)}},E_{i_{\sigma(3)}+1}]]=0,
 \qquad i_1,i_2,i_3\geq0.\label{cbc:l027:yangian-serre}
\end{align}
Both relations hold for all nonnegative indices.
This presentation of the positive half and its shuffle realization
are the results of~\cite[Proposition~1.4(b), Theorem~7.6]{Tsymbaliuk}:
\begin{equation}\label{cbc:l027:yangian-map}
 Y_h^+\xrightarrow{\;\sim\;}\mathcal C^{(-1)}_h,
 \qquad E_r\longmapsto z^r.
\end{equation}

The kernel there is written
$\omega(x,y)=\prod_i(x-y-h_i)/(x-y)^3$ and inserted as
$\omega(z_j,z_i)$ for a cross-block pair.  Thus it equals our
$K_{-1}(z_i-z_j)$.  The chargewise factorial conversion above
accounts for the averaged symmetrization convention.
Composition with~\eqref{cbc:l027:generated-comparison} gives the
same generic positive algebra in each of the two polynomial
realizations.  This conclusion concerns the generated subalgebras.
Their ambient coefficient lattices retain the distinction detected
by~\eqref{cbc:l027:binary-row}.

\section{The tensor required by cooperation}
The product alone does not specify a multiplicative coproduct.
Here a relation among four currents decides whether the ordinary
tensor product can serve as its recipient.

\begin{theorem}[Obstruction to ordinary counital cooperation]
Fix $j\in\{-1,0,1\}$.  There is no $R$-linear,
charge-preserving, unital algebra map
\[
 \Delta:\mathcal C^{(j)}\longrightarrow
 \mathcal C^{(j)}\tensor_R\mathcal C^{(j)}
\]
with ordinary tensor multiplication and two-sided counit the
projection onto charge zero.  The same assertion holds for each
full ambient carrier $\mathcal A^{(j)}$.  Coassociativity is not
needed for this obstruction.
\end{theorem}
\begin{proof}
The binary formula gives
\[
 [e_0,e_3]_{*_j}=-2bt^j(x^2+xy+y^2),\quad
 [e_1,e_2]_{*_j}=-2bt^jxy,\quad
 [e_0,e_1]_{*_j}=-2bt^j.
\]
Together with $e_0*_je_0=2t^j(t+a)$, these give
\begin{equation}\label{cbc:l027:obstruction-relation}
 [e_0,e_3]_{*_j}-3[e_1,e_2]_{*_j}+a[e_0,e_1]_{*_j}
 =-b(e_0*_je_0).
\end{equation}
Charge preservation and the two counits force
$\Delta(e_r)=e_r\tensor1_0+1_0\tensor e_r$.
Commutators of these even primitive elements have no component
of bidegree $(1,1)$ in the ordinary tensor algebra.
The left side of~\eqref{cbc:l027:obstruction-relation} therefore
has zero such component, whereas the right side has
\[
 -2b\,e_0\tensor e_0\ne0
 \quad\text{in }R[x]\tensor_RR[y]=R[x,y].
\]
This is a contradiction over the original parameter ring.
\end{proof}

At $(h_1,h_2,h_3)=(1,1,-2)$ the detector is $4e_0\tensor e_0$.
The proof also applies to a derived tensor recipient with ordinary
Koszul multiplication: the detecting $(1,1)$ component is the
ordinary tensor of two free modules.  It makes no claim of a
nonzero detector on the divisor $b=0$.
A coalgebra for a cut convolution has a different tensor, whose
coefficient comparison must be specified before applying this
obstruction.

\section{Rational cooperation with exchange}
There is a positive construction after localizing the spectral
factors and changing the middle interchange.  Fix $j$ and retain
$R$.  Take ordered blocks $X,Y$ of sizes $m,n$.
Localize $R[X,Y]$ at differences of distinct variables and their
translates by $\pm h_1,\pm h_2,\pm h_3$.
Let $\mathcal B_{m,n}$ be its separately symmetric submodule.

The recipient is the algebraic direct sum over $(m,n)$.
For three or more blocks use the same localization, with invariance
under permutations within each block.  These localizations invert
no parameter by definition.  For
$F\in\mathcal A^{(j)}_{m+n}$, put
\[
 Q(X,Y)=\prod_{x\in X,y\in Y}K_j(x-y),
 \qquad
 \Delta^{\mathrm{rat}}_{m,n}F=\frac{F(X,Y)}{Q(X,Y)}.
\]

The numerator is the same symmetric function evaluated on the
concatenated variables.  The denominator is invertible in the
specified recipient.  Each charge has only finitely many splits.
Further splits use this substitution and division formula.

For $F\in\mathcal B_{m,n}$ and $G\in\mathcal B_{p,q}$, define
$F\circledast G$ by summing over $X=X_1\sqcup X_2$,
$Y=Y_1\sqcup Y_2$, with $|X_1|=m$, $|Y_1|=n$, and summand
\begin{equation}\label{cbc:l027:recipient-product}
 \begin{split}
 F(X_1,Y_1)G(X_2,Y_2)&
 \prod_{x\in X_1,x'\in X_2}K_j(x-x')\\
 {}\cdot&\prod_{y\in Y_1,y'\in Y_2}K_j(y-y')
 \prod_{y\in Y_1,x'\in X_2}r(x'-y).
 \end{split}
\end{equation}
The last product is the exchange of an earlier right input with
a later left input.  Its orientation is fixed by
$r(u)=K_j(-u)/K_j(u)$.

\begin{proposition}[Rational cooperation]
The product $\circledast$ is associative, with unit $1$ in
bidegree $(0,0)$.  The maps $\Delta^{\mathrm{rat}}$ preserve
shuffle multiplication into this recipient, are coassociative,
and have the two charge-zero counits.  Coassociativity is an
identity in the stated three-block localization.
\end{proposition}
\begin{proof}
A summand in $\Delta^{\mathrm{rat}}(F*_jG)$ partitions the first
input into $X_1,Y_1$ and the second into $X_2,Y_2$.
Its cross factors before division are
\[
 K_j(X_1-X_2)K_j(X_1-Y_2)
 K_j(Y_1-X_2)K_j(Y_1-Y_2),
\]
where each expression denotes the product over all pairs in the
indicated blocks.  Factor the total denominator as
\[
 Q(X,Y)=Q(X_1,Y_1)Q(X_1,Y_2)
 Q(X_2,Y_1)Q(X_2,Y_2).
\]

The factor $K_j(X_1-Y_2)$ cancels.  The remaining middle ratio is
$K_j(Y_1-X_2)/K_j(X_2-Y_1)=r(X_2-Y_1)$.
The two internal denominators give the corresponding summands of
$\Delta^{\mathrm{rat}}F$ and $\Delta^{\mathrm{rat}}G$.
This proves multiplicativity with~\eqref{cbc:l027:recipient-product}.

For three inputs in the recipient, either parenthesization sums
over the same partitions of $X$ and of $Y$.  Each earlier-later
pair contributes its internal kernel, and each earlier-right,
later-left pair contributes its exchange ratio exactly once.
Thus the products are associative.  Splitting into three blocks
$X,Y,Z$ divides by $Q(X,Y)Q(X,Z)Q(Y,Z)$ in either order, proving
coassociativity.  Empty blocks give the unit and both counits.
All these are identities of rational functions with finite charge
support.
\end{proof}

For example,
\[
 \Delta^{\mathrm{rat}}_{1,1}(1_2)=K_j(u)^{-1}.
\]
This function is generally not polynomial: the construction has
removed the shifted zero loci of $K_j$.
After extending coefficients to $\operatorname{Frac}(R)$, an
iterated Laurent field, such as
$\operatorname{Frac}(R)((z_1))\cdots((z_n))$, gives a specified
chamber expansion of these rational identities.
That expansion retains the exchanged multiplication in the
recipient.  It does not restore the removed collision loci.

\section{Geometric realization and the double}
A geometric Hall comparison must identify the coefficient objects
and the pull--push maps along extension flags with the product
above.  The preprojective Hall setting of~\cite{SVYangians} fixes
such geometric data separately from the rational kernel.
For the critical theory of the tripled Jordan quiver with potential
$\operatorname{Tr}X[Y,Z]$, a comparison must also specify its
vanishing-cycle coefficients, orientation, parameter ring and
completion.  An identification of generic positive algebras alone
does not determine these integral coefficient objects.

The supported problem is to realize the shuffle multiplication,
its charge-zero unit and its iterated flag comparisons.
The realization must retain the true and shifted collision loci
and their derived intersections.  Any cooperation must come with a coefficient
tensor whose interchange respects
\eqref{cbc:l027:obstruction-relation}.
The rational construction proves that an exchanged recipient can
satisfy this relation.  Extending it across the removed loci is a
further construction.

A completed Drinfeld double requires negative and neutral currents
and Hopf operations in the chosen tensor recipient.
It also requires a continuous Hopf pairing between the appropriate halves.
To obtain a quasitriangular structure, the canonical tensor of
that pairing must exist in the declared completion.
These data must be constructed and compared with the geometric
operations.  The positive shuffle product supplies their starting
algebra.

\chapter{The ladder in one diagram}
The nonlocal structures form the chain
\[
\text{ordered coefficient}
\longrightarrow \mathcal Z_R
\longrightarrow \mathcal T_R
\longrightarrow (V,Y,S),
\]
\[
(V,Y,S)
\longrightarrow \operatorname{Sh}^{+}
\longrightarrow D\!\left(\operatorname{Sh}^{+}\right).
\]
The arrows retain different data:
\begin{center}
\small
\begin{tabularx}{0.96\textwidth}{@{}>{\raggedright\arraybackslash}p{0.18\textwidth}>{\raggedright\arraybackslash}X>{\raggedright\arraybackslash}p{0.29\textwidth}@{}}
\toprule
\textbf{Object}&\textbf{Defining operation}&\textbf{Coherence}\\
\midrule
$\mathcal Z_R$&exchange of charged fields&Yang--Baxter equation\\
$\mathcal T_R$&matrix-current product and coproduct&RTT and coassociativity\\
$\mathcal K_R$&boundary exchange&reflection equation\\
$(V,Y,S)$&formal field product&$S$-locality and weak associativity\\
$\operatorname{Sh}^{+}$&extension shuffle&pairwise-kernel associativity\\
$D(\operatorname{Sh}^{+})$&positive--negative pairing&quasitriangularity\\
\bottomrule
\end{tabularx}
\end{center}
This ladder gives a concrete generator-and-relation theory of doubly noncommutative chiral observables.  The first noncommutativity is ordered composition.  The second is exchange between distinct order chambers.  Chiral singularity remains the diagonal coefficient throughout.


\part{Field theory and master equations}

\chapter{A finite cotangent BF model}\label{cbc:finite-bv}
The classical cotangent action exists for every finite dg Lie algebra.
Its quantization depends on a further datum: the divergence of the
Chevalley vector field.  The obstruction is a cohomology class, and its
vanishing depends on the algebra of functions in which a correction is
allowed.  We first prove this assertion for a finite homological vector
field, then specialize it to Lie algebras.

\section{Coordinates and the cotangent action}
Let $k$ be a field of characteristic zero and $V$ a finite-dimensional
graded vector space.  Choose homogeneous linear coordinates $q_i$ on
$V$, of degrees $d_i$.  On $E=T^*[-1]V$ their conjugate coordinates $p_i$
have degrees $-1-d_i$.  Functions are graded commutative: odd coordinates
anticommute and have square zero.  We use either the polynomial pair
\[
 A=\operatorname{Sym}(V^*),\qquad B=\mathcal O(E),
\]
or their completions in the ideals generated by all base coordinates
and all cotangent coordinates, respectively.  A statement about one
pair does not assert the same result for the other.

All derivatives below are left derivatives.  Fix the constant
Berezinian, the graded analogue of a translation-invariant volume
element, and set
\begin{equation}\label{cbc:bv:laplacian}
 \Delta=\sum_i(-1)^{d_i}\partial_{q_i}\partial_{p_i},\qquad
 \{f,g\}=(-1)^{|f|}
 \bigl(\Delta(fg)-(\Delta f)g-(-1)^{|f|}f\Delta g\bigr).
\end{equation}
The bracket and $\Delta$ have degree one.  These conventions give
$\{q_i,p_j\}=\delta_{ij}$ and $\{p_i,q_j\}=-\delta_{ij}$.
They fix the signs in the geometric BV formalism of
\cite{SchwarzBV1993,MnevBV2016}.

Exactly one coordinate in each pair $(q_i,p_i)$ is odd.  Thus each
summand of $\Delta$ squares to zero, and distinct summands anticommute.
It follows that $\Delta^2=0$.

Expanding \eqref{cbc:bv:laplacian} shows
that the bracket is a biderivation.  Its Jacobi expression is a
triderivation which vanishes on coordinate generators, hence vanishes
on $B$.  The same calculation gives
\begin{equation}\label{cbc:bv:compatibility}
 \Delta\{f,g\}=\{\Delta f,g\}+(-1)^{|f|+1}\{f,\Delta g\}.
\end{equation}
In the completed case all these identities extend by continuity.
$\Delta$ lowers the field-ideal filtration by at most two.

Let $Q:A\to A$ be a degree-one derivation with $Q^2=0$, and assume
$Q(q_i)$ has no constant term.  Such a derivation is a
\emph{homological vector field}.  In the completed case it is
continuous.  Define its cotangent Hamiltonian, with the indicated
order of factors, by
\begin{equation}\label{cbc:bv:action}
 S_Q=-\sum_i Q(q_i)p_i\in B^0,\qquad D_Q=\{S_Q,-\}.
\end{equation}
The action has momentum degree one and field order at least two.

\begin{theorem}[Cotangent master equation]\label{cbc:bv:cme}
The derivation $D_Q$ restricts to $Q$ on $A$, and
\[
 \frac12\{S_Q,S_Q\}=0,\qquad D_Q^2=0.
\]
\end{theorem}
\begin{proof}
For any even Hamiltonian $S$ of degree zero, direct differentiation gives
\[
 \{S,q_i\}=(-1)^{d_i}\partial_{p_i}S,\qquad
 \{S,p_i\}=(-1)^{d_i}\partial_{q_i}S.
\]
In \eqref{cbc:bv:action}, differentiating past $Q(q_i)$ contributes
$(-1)^{d_i+1}$.  Therefore $\{S_Q,q_i\}=Q(q_i)$.

The bracket $\{S_Q,S_Q\}/2$ is linear in the momenta.  By Jacobi its
Hamiltonian vector field restricts to $Q^2=0$ on the base.
Its derivatives in every $p_i$ consequently vanish.  A function linear
in the momenta with this property is zero.  This proves the master
equation without leaving an undetermined scalar term.
Jacobi now gives $D_Q^2=\{\{S_Q,S_Q\}/2,-\}=0$ on all of $B$.
\end{proof}

\section{The divergence class}
For the chosen base Berezinian $\rho$, the divergence of an odd vector
field is
\[
 \chi_Q=\operatorname{div}_{\rho}Q
       =\sum_i\partial_{q_i}Q(q_i)\in A^1.
\]
The order and signs in \eqref{cbc:bv:action} give
\begin{equation}\label{cbc:bv:divergence}
 \Delta S_Q=\chi_Q,\qquad Q\chi_Q=0.
\end{equation}
Indeed, the two factors $(-1)^{d_i}$ in
$\Delta S_Q$ cancel.  Applying
\eqref{cbc:bv:compatibility} to the classical master equation gives
$\{S_Q,\Delta S_Q\}=0$, which is the second assertion.
The class $[\chi_Q]\in H^1(A,Q)$ is the \emph{modular class}
for this function algebra and density.

The base and cotangent divergences have different normalizations.
If $\mu$ is the constant Berezinian on $E$, the coordinate formula for
$D_Q$ gives
\begin{equation}\label{cbc:bv:half-divergence}
 \operatorname{div}_{\mu}D_Q=2\Delta S_Q=2\chi_Q.
\end{equation}
The $q$ and $p$ derivatives each contribute $\chi_Q$.
Thus $\Delta S_Q$ is the full divergence of the base vector field and
half the divergence of its Hamiltonian lift.

Give $\hbar$ degree zero.  A \emph{quantum correction} in the chosen
function algebra is a series
$S_\hbar=S_Q+\hbar S_1+\hbar^2S_2+\cdots$ in $B^0[[\hbar]]$
satisfying the quantum master equation
\begin{equation}\label{cbc:bv:qme}
 \frac12\{S_\hbar,S_\hbar\}+\hbar\Delta S_\hbar=0.
\end{equation}
The fields, coefficient algebra and Laplacian are fixed in this
definition.  No negative powers of $\hbar$ are allowed.

\begin{theorem}[Cotangent quantization criterion]\label{cbc:bv:criterion}
A quantum correction exists if and only if
$[\chi_Q]=0$ in $H^1(A,Q)$.  More precisely, if $F\in A^0$ satisfies
$QF=\chi_Q$, then
\begin{equation}\label{cbc:bv:correction}
 S_\hbar=S_Q-\hbar F
\end{equation}
solves the quantum master equation to every order.
\end{theorem}
\begin{proof}
Set all momenta equal to zero to obtain $\pi:B\to A$.
The $q$ component of $D_Q$ is $Q$, and its $p$ component is linear
in the momenta.  Hence $\pi D_Q=Q\pi$.
The coefficient of $\hbar$ in \eqref{cbc:bv:qme} is
$D_QS_1+\chi_Q=0$.  Applying $\pi$ gives
$Q(\pi S_1)=-\chi_Q$, so a correction requires the stated vanishing.
This argument permits arbitrary momentum dependence in $S_1$.

Conversely, a base function $F$ satisfies
$\Delta F=0$ and $\{F,F\}=0$.  Substituting
\eqref{cbc:bv:correction} into the left side of \eqref{cbc:bv:qme}
therefore gives
\[
 \frac12\{S_Q,S_Q\}
 +\hbar(\Delta S_Q-\{S_Q,F\})
 +\frac{\hbar^2}{2}\{F,F\}-\hbar^2\Delta F
 =\hbar(\chi_Q-QF)=0.
\]
There are no higher terms.
\end{proof}

For any degree-zero $S_\hbar$, put
$\delta_{S_\hbar}=\{S_\hbar,-\}+\hbar\Delta$.
Jacobi and \eqref{cbc:bv:compatibility} give the exact identity
\begin{equation}\label{cbc:bv:quantum-square}
 \delta_{S_\hbar}^2
 =\left\{\frac12\{S_\hbar,S_\hbar\}
                   +\hbar\Delta S_\hbar,-\right\}.
\end{equation}
A quantum master action thus gives a differential on $B[[\hbar]]$.
Its compatibility with multiplication is
\begin{equation}\label{cbc:bv:product-defect}
 \delta_{S_\hbar}(fg)
 =(\delta_{S_\hbar}f)g+(-1)^{|f|}f\delta_{S_\hbar}g
                    +\hbar(-1)^{|f|}\{f,g\}.
\end{equation}
Consequently the ordinary commutative product is generally not a
product of complexes.  Applying a construction that requires a
derivation demands further algebraic data.

The exponential form of the master equation requires another
completion.  Let $\mathfrak m$ be the ideal of all field coordinates
in $B$, and define
\[
 \mathcal W=\varprojlim_N(B/\mathfrak m^N)((\hbar)).
\]
Here the lower bound on powers of $\hbar$ may depend on $N$.

Since $S_Q\in\mathfrak m^2$, its exponential $e^{S_Q/\hbar}$
exists in $\mathcal W$.  In \eqref{cbc:bv:correction} one may subtract
the constant of $F$ and assume $F(0)=0$.  Then $e^{S_\hbar/\hbar}$
also exists and is invertible.  Differentiation is defined by taking
two extra orders of field precision before applying $\Delta$.

The product identity for $\Delta$ gives
\begin{equation}\label{cbc:bv:exponential}
 \hbar^2e^{-S_\hbar/\hbar}\Delta(e^{S_\hbar/\hbar})
 =\frac12\{S_\hbar,S_\hbar\}+\hbar\Delta S_\hbar
\end{equation}
in $\mathcal W$.  It follows first for powers of an even action by
induction and then for the exponential at each field precision.
Thus the quantum master equation is equivalent there to
$\hbar\Delta(e^{S_\hbar/\hbar})=0$.  This formal identity asserts no
convergence of a numerical integral.

In the formal coefficient algebra, the same primitive has a density
interpretation.  Set
\[
 \rho'=e^{-F}\rho,\qquad \mu'=e^{-2F}\mu .
\]
The divergence rule gives
$\operatorname{div}_{\rho'}Q=\chi_Q-QF=0$.
For the cotangent density the corresponding Laplacian is
$\Delta'=\Delta-\{F,-\}$, so $\Delta'S_Q=0$ as well.

The factor two in $\mu'$ agrees with
\eqref{cbc:bv:half-divergence}.  The exponentials use the formal
completion.  A polynomial primitive need not have a polynomial
exponential.

\section{Finite dg Lie algebras}
Let $L$ be a finite-dimensional dg Lie algebra and take $V=L[1]$.
Then
\[
 E=L[1]\oplus L^*[-2].
\]
For a homogeneous basis $e_i$ of $L$ the coordinates have degrees
$|q_i|=1-|e_i|$ and $|p_i|=|e_i|-2$.

Let $Q$ be the Chevalley--Eilenberg derivation on
$A=\operatorname{Sym}(L^*[-1])$, or its stated completion.
It is the signed suspension of $d_L$ and the bracket.
The equations $d_L^2=0$, the derivation identity for $d_L$ and Jacobi
are respectively the linear, quadratic and cubic parts of $Q^2=0$.
Equation \eqref{cbc:bv:action} defines the cotangent BF action with
these signs.

For a degree-zero endomorphism $T$ of $L$, write
$\operatorname{str}_L(T)=\sum_r(-1)^r\operatorname{Tr}(T|_{L^r})$.
Define the degree-zero linear functional on $L$
\[
 \chi_L(x)=\operatorname{str}_L(\operatorname{ad}_x)
       \quad(x\in L^0),\qquad \chi_L|_{L^r}=0\quad(r\ne0).
\]
Its suspension is a degree-one linear function on $L[1]$.
We call $L$ \emph{unimodular} if $\chi_L=0$.

\begin{proposition}[The modular term]\label{cbc:bv:modular}
The divergence $\chi_Q$ is the suspended functional $\chi_L$.
In particular, a unimodular finite dg Lie algebra has
$\Delta S_Q=0$, and its uncorrected cotangent action satisfies the
quantum master equation.
\end{proposition}
\begin{proof}
The linear part of $Q$ has degree one, so it has no diagonal
coefficient in the homogeneous coordinates.  It contributes zero
to $\sum_i\partial_{q_i}Q(q_i)$.

For $x\in L^0$, the terms involving its odd coordinate $q^x$
in the quadratic part act on the other coordinates as
$-q^x(\operatorname{ad}_x)^*$.
Left differentiation in $q_i$ contributes
$-(-1)^{|q_i|}(\operatorname{ad}_x)_{ii}q^x
=(-1)^{|e_i|}(\operatorname{ad}_x)_{ii}q^x$.
Summing gives the stated supertrace.
Only coordinates of degree one can occur in the resulting linear
function, so these coefficients determine it.

Its closedness also follows from the vanishing of supertraces of
graded commutators, including
$\operatorname{ad}_{d_Ly}=[d_L,\operatorname{ad}_y]$.
The last assertion follows from Theorem~\ref{cbc:bv:cme}.
\end{proof}

For an ordinary Lie algebra $\mathfrak g$ with
$[e_a,e_b]=f^c_{ab}e_c$, write $c^a$ for its degree-one coordinates
and $b_a$ for the degree $-2$ conjugates.  Our convention becomes
\begin{equation}\label{cbc:bv:ordinary}
 Qc^a=-\frac12f^a_{ij}c^ic^j,\qquad
 S_Q=\frac12b_a f^a_{ij}c^ic^j,\qquad
 \Delta=-\sum_a\partial_{c^a}\partial_{b_a}.
\end{equation}
Differentiation gives
\[
 \Delta S_Q=-\sum_{a,j}f^a_{aj}c^j
          =\sum_j\operatorname{Tr}(\operatorname{ad}_{e_j})c^j.
\]

Here $A^0=k$ and $Q(A^0)=0$.
Thus an ordinary finite-dimensional $\mathfrak g$ admits a quantum
correction on this cotangent space if and only if it is unimodular.
Allowing momenta in the correction does not remove the obstruction.
For a general dg Lie algebra, exactness of $\chi_Q$ is weaker than
its vanishing as a function.

For the affine Lie algebra $[H,E]=E$, one has
\[
 S_Q=b_Ec^Hc^E,\qquad \Delta S_Q=c^H,\qquad
 \delta_{S_Q}^2(b_H)=\hbar.
\]
The nonzero character prevents a quantum correction.
For the Heisenberg algebra $[X,Y]=Z$, instead,
$S_Q=b_Zc^Xc^Y$ and $\Delta S_Q=0$.
More generally, every ordinary finite-dimensional nilpotent Lie
algebra is unimodular because each adjoint endomorphism is nilpotent.

Nilpotence is sufficient but unnecessary.  In $\mathfrak{sl}_2$,
with $[H,E]=2E$, $[H,F]=-2F$ and $[E,F]=H$, the action is
\[
 S_Q=2b_Ec^Hc^E-2b_Fc^Hc^F+b_Hc^Ec^F.
\]
The first two terms have opposite divergences, and the third has
zero divergence.  Jacobi supplies the classical master equation.
The vanishing divergence supplies the quantum one.

\section{A correction that requires completion}
Take $A_{\mathrm{pol}}=k[t]\otimes\Lambda(c)$ with $|t|=0$,
$|c|=1$, and define
\[
 Qt=c(1+t),\qquad Qc=0.
\]
The derivation is homological and has $\chi_Q=c$.
The cotangent coordinates have degrees $|p_t|=-1$, $|p_c|=-2$, and
\[
 S_Q=-c(1+t)p_t,\qquad
 \Delta=\partial_t\partial_{p_t}-\partial_c\partial_{p_c}.
\]
Polynomial differentiation is surjective on $k[t]$, so
\[
 H^1(A_{\mathrm{pol}},Q)
 =c\,k[t]\big/c(1+t)k[t]\cong k\,c .
\]
Evaluation at $t=-1$ detects the nonzero class of $c$.
Theorem~\ref{cbc:bv:criterion} therefore excludes a polynomial
quantum correction, including corrections depending on the momenta.

In $A_{\mathrm{for}}=k[[t]]\otimes\Lambda(c)$, however,
\[
 F=\log(1+t)=\sum_{m\ge1}\frac{(-1)^{m+1}}m t^m,
 \qquad QF=c .
\]
The derivative is again surjective and $1+t$ is now a unit, so
$H^1(A_{\mathrm{for}},Q)=0$.
The completed cotangent action
\[
 S_\hbar=-c(1+t)p_t-\hbar\log(1+t)
\]
satisfies the quantum master equation exactly.

For the truncation $F_N=\sum_{m=1}^N(-1)^{m+1}t^m/m$, its residual is
\[
 \frac12\{S_Q-\hbar F_N,S_Q-\hbar F_N\}
       +\hbar\Delta(S_Q-\hbar F_N)
 =(-1)^N\hbar c\,t^N .
\]
It vanishes through total base-field degree $N$.
The infinite identity follows from $Q\log(1+t)=c$.  The displayed
finite identities specify its output precision.

\section{Local fields and interval gluing}
Let $X$ be a complex curve and $\mathfrak g$ an ordinary
finite-dimensional nilpotent Lie algebra.  The local dg Lie algebra
\[
 \mathcal L_X
 =\Omega^{0,*}(X)\widehat\otimes\Omega^*(\mathbb R)\otimes\mathfrak g
\]
has differential $\bar\partial+d_{\mathbb R}$ and the bracket obtained
from wedge product and the bracket of $\mathfrak g$.
The completed tensor product here is that of smooth form spaces.
Its density dual pairs with compactly supported fields by integration.
These data define the classical cotangent action: its master equation
uses the differential and Jacobi identities, before a Laplacian at
coincident points is introduced.

A finite cochain deformation retract alone does not identify the
transferred interaction or its divergence.  A finite model with an
actual transferred homological vector field and a specified
Berezinian falls under Theorem~\ref{cbc:bv:criterion}.  The vector field
may include higher Taylor terms.
To compare its quantization with the local theory requires a
compatible BV pushforward or another construction preserving the
quantum master equation.
An infinite-dimensional realization further requires gauge fixing,
renormalized local interactions, the renormalized master equation and
compatible operations on opens.  The finite trace calculation
supplies none of these constructions by itself.

There is nevertheless a definite algebraic gluing model.
Give $U=U(\mathfrak g)$ its scalar augmentation
$\epsilon:U\to k$ and let $\overline U=\ker\epsilon$.
The two boundary modules in this model are the right and left
augmentation modules $k$.
The normalized bar complex has
$\mathrm B_q(U)=\overline U^{\otimes q}$ and differential
\[
 b[a_1|\cdots|a_q]
 =\sum_{i=1}^{q-1}(-1)^{i-1}
       [a_1|\cdots|a_ia_{i+1}|\cdots|a_q].
\]

On $C_q(\mathfrak g)=\Lambda^q\mathfrak g$ use
\[
 \partial(x_1\wedge\cdots\wedge x_q)
 =\sum_{i<j}(-1)^{i+j+1}[x_i,x_j]\wedge
 x_1\wedge\cdots\widehat{x_i}\cdots\widehat{x_j}\cdots\wedge x_q .
\]
The remaining entries keep their order.  In particular,
$\partial(x\wedge y)=[x,y]$.

\begin{proposition}[Augmentation gluing]\label{cbc:bv:augmentation}
The antisymmetrizer
\[
 \operatorname{Alt}(x_1\wedge\cdots\wedge x_q)
 =\sum_{\sigma\in S_q}\operatorname{sgn}(\sigma)
       [x_{\sigma(1)}|\cdots|x_{\sigma(q)}]
\]
gives a quasi-isomorphism
$C_*(\mathfrak g)\to\mathrm B(U)$.
Consequently $k\otimes_U^{\mathbb L}k\simeq C_*(\mathfrak g)$
as complexes.
\end{proposition}
\begin{proof}
Adjacent permutations in $b\operatorname{Alt}$ pair to give
$xy-yx=[x,y]$, with the displayed Chevalley sign.
Thus $\operatorname{Alt}$ is a chain map.

The PBW filtration on $U$ induces total word weight on its bar.
Give $\Lambda^q\mathfrak g$ weight $q$.
PBW identifies the associated graded algebra with
$\operatorname{Sym}(\mathfrak g)$.

Order a basis and rewrite $yx$ as $xy+[y,x]$.
The only overlaps involve words of length three, and their reductions
differ by Jacobi.  Disjoint reductions commute.
Reduction lowers length or the lexicographic order at fixed length,
so the ordered words form a basis.

The associated graded comparison is the antisymmetrizer from the
exterior complex with zero differential to the polynomial bar
complex.  It is the Koszul comparison, a quasi-isomorphism: the
polynomial Koszul resolution is exact by successive contraction in
each basis variable.  With negative first face in the left bar
resolution the comparison has the degree sign $(-1)^q$.
After augmentation the exterior differential is zero, so this sign
does not change the quasi-isomorphism assertion.

The mapping cone of $\operatorname{Alt}$ has acyclic associated
graded.  Every cone cycle has finite largest weight.  Subtracting a
boundary in that weight lowers it and proves exactness by finite
descent.  Finally, the normalized left bar resolution is free over
$U$, so its augmentation computes the stated derived tensor product.
\end{proof}

To obtain this gluing from the field theory requires an actual
boundary quantization.  Its algebra must be identified with
$U(\mathfrak g)$ and its two modules with the augmentation modules.
Its composition must correspond to the derived tensor product.
An interval polarization alone
does not give these identifications.
The scalar augmentation in this algebraic model also supplies no
identification with a native transverse vacuum or a prepared state.

\chapter{BRST and ordered reduction}
Reduction by a differential and the bar construction interact through
the multiplication.  If the differential is a derivation, it extends
to tensor words and anticommutes with the operation that joins
adjacent entries.  This gives a bar complex before taking cohomology.
Replacing its entries by cohomology classes requires a comparison
of complexes, and replacing it by a chiral construction requires
a comparison of operations.

\section{The signed bar complex}
Fix a field $k$ of characteristic zero.
Let $(A,Q,\epsilon)$ be a unital augmented differential graded
associative $k$-algebra.  Thus $Q$ has cohomological degree one,
$Q^2=0$, $Q(ab)=Q(a)b+(-1)^{|a|}aQ(b)$, and
$\epsilon:A\to k$ is a unital algebra map satisfying $\epsilon Q=0$.
Write $I=\ker\epsilon$.

The suspension $I[1]$ places $sa$ in degree $|a|-1$.
The \emph{algebraic reduced bar} is
\[
 B(A)=\bigoplus_{r\geq0}I[1]^{\otimes_k r},
 \qquad
 [a_1|\cdots|a_r]=sa_1\otimes\cdots\otimes sa_r.
\]
Every element has finite support in word length.
The empty word $[\,]$ spans the length-zero copy of $k$.

For homogeneous entries put
$\sigma_i=\sum_{j<i}(|a_j|-1)$ and define
\begin{align*}
 Q_B[a_1|\cdots|a_r]
 &=-\sum_i(-1)^{\sigma_i}
 [a_1|\cdots|Qa_i|\cdots|a_r],\\
 b[a_1|\cdots|a_r]
 &=\sum_{i<r}(-1)^{\sigma_i+|a_i|}
 [a_1|\cdots|a_i a_{i+1}|\cdots|a_r].
\end{align*}
Both maps have total degree one and kill the empty word.
Their sum is denoted by $D=Q_B+b$.
Deconcatenation is the coproduct
\[
 \Delta[a_1|\cdots|a_r]
 =\sum_{i=0}^r
 [a_1|\cdots|a_i]\otimes[a_{i+1}|\cdots|a_r].
\]
The inclusion of the empty word is its coaugmentation.
Projection onto that word is its counit.

\begin{proposition}[Bar differential and coproduct]
\label{prop:brst-signed-bar}
The identities
\[
 Q_B^2=b^2=0,
 \qquad Q_Bb+bQ_B=0
\]
hold on $B(A)$.
Both $Q_B$ and $b$ are coderivations of $\Delta$.
Consequently $(B(A),D,\Delta)$ is a coaugmented differential graded
coalgebra.
\end{proposition}
\begin{proof}
The unary differential on $I[1]$ is $s a\mapsto-sQa$.
In its tensor extension, terms acting in distinct positions cancel
in opposite orders, and terms acting twice in one position vanish
by $Q^2=0$.
For $b^2$, disjoint multiplications cancel by the same degree-one
sign rule.  On three adjacent entries the remaining expression is
\[
 (-1)^{|a_2|}
 [(a_1a_2)a_3-a_1(a_2a_3)],
\]
which vanishes by associativity.

For the mixed identity, operations in disjoint positions again
cancel.  On two adjacent entries the expression is
\[
 (-1)^{|a|}[Q(a)b]+[aQ(b)]-(-1)^{|a|}[Q(ab)]=0.
\]
This is precisely the derivation identity.

For the coproduct, a unary operation lies on one side of each cut.
A multiplication lies on one side of every cut that survives the
joining of its two entries.
These are exactly the terms in
$(D\otimes1+1\otimes D)\Delta$, with
$(1\otimes D)(u\otimes v)=(-1)^{|u|}u\otimes Dv$.
The same argument applies separately to $Q_B$ and $b$.
\end{proof}

The suspension signs can be seen without a general word.
In $k[x]\otimes\Lambda(\theta)$, with $|x|=0$, $|\theta|=1$,
$Qx=\theta$ and $Q\theta=0$, they give
\[
 Q_B[x|x]=-[\theta|x]+[x|\theta],
 \qquad
 bQ_B[x|x]=2[x\theta],
 \qquad Q_Bb[x|x]=-2[x\theta].
\]
An overall opposite sign for $Q_B$, with the same $b$, gives
an isomorphic complex: conjugate by
$[a_1|\cdots|a_r]\mapsto(-1)^{\sum_i|a_i|}[a_1|\cdots|a_r]$.
The prefix signs within $Q_B$ are still required.

\section{Length and cohomology}
Let $F_pB(A)$ be the subcomplex of words of length at most $p$,
with $F_p=0$ for $p<0$.
The filtration is increasing and exhaustive.
The operator $Q_B$ preserves length and $b$ lowers it by one.
Use an index $q$ for the sum of the unsuspended entry degrees,
so that a word in bidegree $(p,q)$ has total degree $q-p$.

\begin{proposition}[The length spectral sequence]
\label{prop:brst-length-sequence}
The first page and the bidegrees of the differentials are
\[
 E_1^{p,q}=\bigl(H(I,Q)^{\otimes_k p}\bigr)^q,
 \qquad
 d_r:E_r^{p,q}\longrightarrow E_r^{p-r,q-r+1}.
\]
The differential $d_1$ is induced by $b$.
Thus $(E_1,d_1)$ is the bar complex of the augmented graded algebra
$H(A,Q)$, and $E_2$ is its cohomology.
For each filtration index, the eventual quotient is
$\operatorname{gr}_F H(B(A),D)$.
This assertion does not require a common finite page on which
all incoming differentials vanish.
\end{proposition}
\begin{proof}
The associated graded differential is $Q_B$.
Tensor products of complexes over a field obey the K\"unneth
formula, including for unbounded complexes.
This identifies the first page.  The part of $D$ lowering length
by one is $b$, which gives $d_1$ and its stated signs.
Total degree one gives the displayed bidegree of $d_r$.

For the limiting statement, a representative in $F_p$ of a cycle
on page $r$ has differential in $F_{p-r}$.
Once $r>p$, such representatives are actual cycles.
The incoming boundaries form an increasing union.
Every primitive in the algebraic bar has a finite largest word
length, so this union contains exactly the boundaries in $F_p$.
The resulting quotient is
\[
 \frac{\ker D\cap F_p}
 { (\ker D\cap F_{p-1})+(\operatorname{im}D\cap F_p)}
 =\operatorname{gr}_p^F H(B(A),D).
\]

The filtration on cohomology is exhaustive because every cycle
has finite word support.
The finite support argument supplies no uniform bound on the
lengths of primitives for different classes.
\end{proof}

A completion needs its own convergence argument.
Suppose, in addition, that
$A=\bigoplus_{N\geq0}A^{(N)}$ has an integral weight grading,
$A^{(0)}=k$, $I=\bigoplus_{N>0}A^{(N)}$,
$Q$ preserves weight and multiplication adds weights.
Let $B(A)^{(N)}$ denote total weight $N$.
Every word in it has length at most $N$.

Define the \emph{weight-completed bar} and its tensor recipient by
\begin{align*}
 \widehat B_{\mathrm{wt}}(A)&=\prod_{N\geq0}B(A)^{(N)},\\
 \widehat B_{\mathrm{wt}}(A)\widehat\otimes
 \widehat B_{\mathrm{wt}}(A)
 &=\prod_{u,v\geq0}B(A)^{(u)}\otimes_k B(A)^{(v)}.
\end{align*}
All products here are taken degree by degree as products of
complexes.  The differential and coproduct extend componentwise.

\begin{corollary}[Positive-weight convergence]
\label{cor:brst-weight-completion}
At each weight $N$, the length spectral sequence has a finite
filtration and computes $H(B(A)^{(N)},D)$.
Moreover
\[
 H(\widehat B_{\mathrm{wt}}(A),D)
 =\prod_{N\geq0}H(B(A)^{(N)},D).
\]
\end{corollary}
\begin{proof}
The bound $p\leq N$ bounds both outgoing and incoming
differentials at fixed weight.
Products of vector spaces are exact, so taking cohomology commutes
with the displayed product of complexes.
\end{proof}
This completion is specified by weight.
A product over arbitrary word lengths is a different complex,
for which the finite support proof of
\cref{prop:brst-length-sequence} does not apply.

\section{Comparison before and after reduction}
An augmented differential graded algebra map $f:A\to A'$ induces
a coalgebra map $B(f)$ by applying $f$ to every entry.
It commutes with $Q_B$ by the chain-map identity and with $b$
by multiplicativity.
This functorial map gives a comparison stronger than the
identification of a spectral-sequence page.

\begin{theorem}[Homotopy invariance of the algebraic bar]
\label{thm:brst-bar-invariance}
If $f:A\to A'$ is an augmented differential graded algebra
quasi-isomorphism over $k$, then $B(f)$ is a quasi-isomorphism
of differential graded coalgebras.
No bound on the cohomological degrees is required.
For positive-weight algebras and weight-preserving maps, the same
holds for $\widehat B_{\mathrm{wt}}(f)$ when $f$ is a
quasi-isomorphism at every weight.
\end{theorem}
\begin{proof}
The splittings $A=k\oplus I$ and $A'=k\oplus I'$ imply that
$I\to I'$ is a quasi-isomorphism.
Over a field an acyclic complex is contractible, so tensoring
with any complex preserves quasi-isomorphisms.
Every finite tensor power of $I\to I'$ is therefore a
quasi-isomorphism.

Filter the mapping cone of $B(f)$ by word length.
Its associated graded complexes are acyclic by this observation.

For a cone cycle of largest length $p$, solve for a primitive
of its length-$p$ component in the associated graded.
Subtract its differential.  The remaining cycle has smaller
largest length.  Repetition ends at length zero and expresses
the original cycle as a boundary.
This proves the algebraic assertion.

At fixed positive weight the same argument has at most $N+1$
steps, and exactness of products proves the completed assertion.
\end{proof}

Suppose now that $H^i(A,Q)=0$ for $i\ne0$.
The truncation $A_{\leq0}$ is defined by
\[
 A_{\leq0}^i=
 \begin{cases}
 A^i,&i<0,\\
 \ker(Q:A^0\to A^1),&i=0,\\
 0,&i>0.
 \end{cases}
\]
Its multiplication and augmentation are inherited from $A$.
Products of degree-zero cycles are cycles, so it is an augmented
differential graded subalgebra.
Projection of its degree-zero cycles onto cohomology, and of its
negative degrees onto zero, defines an algebra map
$A_{\leq0}\to H^0(A,Q)$.

\begin{corollary}[Concentrated reduction]
\label{cor:brst-concentrated-reduction}
Under the concentration hypothesis there is a canonical zigzag
of differential graded coalgebra quasi-isomorphisms
\[
 B(A)\ \longleftarrow\ B(A_{\leq0})
 \ \longrightarrow\ B(H^0(A,Q)).
\]
The positive-weight completed version holds with the preceding
weight and tensor conventions.
\end{corollary}
\begin{proof}
The inclusion $A_{\leq0}\to A$ is a quasi-isomorphism because
the positive cohomology of $A$ vanishes.
The projection $A_{\leq0}\to H^0(A,Q)$ is a quasi-isomorphism
because its negative cohomology vanishes.
Both preserve the augmentation.
Apply \cref{thm:brst-bar-invariance}, weight by weight when
the completion is present.
\end{proof}
The zigzag is canonical.  A strict reverse map
$B(H^0(A,Q))\to B(A)$ needs an additional construction.
Even a linear section of cycles need not preserve multiplication.
\Cref{chap:brst} constructs a bar map with a nonzero quadratic
correction and an example with nonzero $d_2$.
Over a general coefficient ring, the tensor and K\"unneth
arguments above require replacement hypotheses.
In particular, they must hold for the intermediate truncation
as well as for the original algebra.

\section{The vertex-algebra application}
For a finite-dimensional complex simple Lie algebra $\mathfrak g$,
let $V^\kappa(\mathfrak g)$ be its universal affine vertex algebra
at level $\kappa$.  The superscript distinguishes it from a
simple quotient.
For a principal nilpotent element, let
$C_{\mathrm{DS}}^\kappa(\mathfrak g)$ be the charged-ghost
Drinfeld--Sokolov BRST complex defined in \cref{chap:brst}.
Its cohomology is concentrated in degree zero for every
$\kappa\in\mathbb C$, by the universal-complex theorem of
Kac--Wakimoto \cite[Theorem 4.1]{KacWakimoto2004}.
Its degree-zero cohomology is the universal principal
$W$-algebra.  A noncritical level is needed for the Sugawara
conformal vector, rather than for this concentration theorem.

This vertex-algebra theorem does not yet supply the augmented
associative algebra in \cref{cor:brst-concentrated-reduction}.
The normal product of vertex fields can fail associativity.
In the principal $\mathfrak{sl}_2$ complex, the ordinary vacuum
coefficient functional fails to commute with $Q$.
Both failures are computed in \cref{chap:brst}.

An application of the bar comparison therefore requires an
associative realization of the coefficient operations with a
compatible differential.
It needs an augmentation, or specified relative boundary modules
in place of that augmentation.
It also requires the tensor product, completion and comparison
with the full factorization operations.
Neither word deconcatenation nor the empty word constructs
those operations or identifies a physical state.

\chapter{Stokes and algebraic identities}\label{ch:stokes-identities}
The two bracketings of a product give two points of an interval.
To turn that interval into an associativity homotopy, one needs
an operation on its interior whose differential is the difference
of the endpoint operations.
The boundary of a pentagon then asks whether the chosen homotopy
is compatible with a fourth input.
Oriented boundaries determine these equations.
Coefficient-valued operations must solve them.

We use cohomological grading over a field $k$.
For cochain complexes $E,F$, the complex
$\operatorname{Hom}^\bullet(E,F)$ consists of homogeneous linear maps,
with differential
\[
 \delta f=d_Ff-(-1)^{|f|}fd_E.
\]
Tensor products use the Koszul rule
$(f\otimes g)(a\otimes b)=(-1)^{|g||a|}f(a)\otimes g(b)$.
In particular, if $V$ is a cochain complex and
$\mu:V\otimes V\to V$ is a degree-zero chain map, put
\[
 L_3=\mu(\mu\otimes1),\qquad
 R_3=\mu(1\otimes\mu).
\]
An associativity homotopy is a degree-$-1$ map
$H:V^{\otimes3}\to V$ satisfying $\delta H=L_3-R_3$.
Its existence is a condition on the product and the coefficient complex.

\section{Cellular coefficients and the full square}
Let $K$ be a finite oriented regular CW complex.
Write $\tau<\sigma$ when $\tau$ is a proper face of $\sigma$,
and write $\epsilon_{\tau\sigma}$ for the cellular incidence number
when $\dim\sigma-\dim\tau=1$.
Choose a cochain complex $(E_\sigma,d_\sigma)$ for every cell.
Give the orientation generator $[\sigma]$ degree $-\dim\sigma$ and set
\[
 \begin{aligned}
 C^m(K;E)&=\bigoplus_\sigma
 [\sigma]\otimes E_\sigma^{m+\dim\sigma},\\
 D_0([\sigma]\otimes e)&=
 (-1)^{\dim\sigma}[\sigma]\otimes d_\sigma e.
 \end{aligned}
\]
Thus $D_0$ has degree one.

For each face pair of codimension $r$, choose a map
\[
 t_{\tau\sigma}:E_\sigma\longrightarrow E_\tau,
 \qquad |t_{\tau\sigma}|=1-r.
\]
The chosen orientation coefficients are included in $t_{\tau\sigma}$.
Define a degree-one endomorphism of the finite cellular direct sum by
\[
 D([\sigma]\otimes e)=D_0([\sigma]\otimes e)
       +\sum_{\tau<\sigma}[\tau]\otimes t_{\tau\sigma}(e).
\]

\begin{theorem}[Cellular curvature]\label{thm:cellular-curvature}
The component of $D^2$ from $[\sigma]\otimes E_\sigma$
to $[\tau]\otimes E_\tau$, for $\tau<\sigma$, is
\begin{equation}\label{eq:cellular-curvature}
 (-1)^{\dim\tau}\delta t_{\tau\sigma}
 +\sum_{\tau<\nu<\sigma}t_{\tau\nu}t_{\nu\sigma}.
\end{equation}
Consequently $D$ is a differential if and only if all these expressions vanish.
For covers, the condition says that $t_{\tau\sigma}$ is a chain map.
\end{theorem}
\begin{proof}
The diagonal component is $d_\sigma^2=0$.
For a face pair of codimension $r$, the two terms involving an internal
differential are
\[
 \begin{aligned}
 &(-1)^{\dim\tau}d_\tau t_{\tau\sigma}
       +(-1)^{\dim\sigma}t_{\tau\sigma}d_\sigma\\
 &\qquad =(-1)^{\dim\tau}
       \bigl(d_\tau t_{\tau\sigma}+(-1)^r
                        t_{\tau\sigma}d_\sigma\bigr).
 \end{aligned}
\]
Since $|t_{\tau\sigma}|=1-r$, the last parenthesis is
$\delta t_{\tau\sigma}$.
Every remaining term passes through exactly one intermediate cell $\nu$,
which gives the sum in~\eqref{eq:cellular-curvature}.
These are all components of the square.
\end{proof}

Suppose only cover maps occur, with
$t_{\tau\sigma}=\epsilon_{\tau\sigma}F_{\tau\sigma}$
for degree-zero chain maps $F_{\tau\sigma}$.
The theorem reduces to the codimension-two identities
\begin{equation}\label{eq:strict-corner-identity}
 \sum_{\tau<\nu<\sigma}
 \epsilon_{\tau\nu}\epsilon_{\nu\sigma}
 F_{\tau\nu}F_{\nu\sigma}=0,
 \qquad \dim\sigma-\dim\tau=2.
\end{equation}
For constant coefficients and identity maps this is the cellular equation
$\partial^2=0$.
For general coefficients, opposite incidence signs cancel equal composites.
Equality of the composites is a separate condition.

There is also a mixed term to retain when a face map is written as
an identity plus a correction.
For constant $E$, let $d_K$ be the signed cellular boundary with identity
coefficients, and write $D=D_0+d_K+\rho$.
Even when $D_0\rho+\rho D_0=0$, its square contains
\[
 d_K\rho+\rho d_K+\rho^2.
\]

For example, orient the unit square $Q$ by $dx\wedge dy$,
and orient its bottom edge $e_b$ by increasing $x$.
With scalar coefficients in degree zero, set $\rho(Q)=e_b$
and let $\rho$ vanish on every other cell.
Writing $v_{ab}$ for the vertex $(a,b)$, one obtains
\[
 \rho^2=0,\qquad
 (d_K+\rho)^2(Q)=\partial e_b=v_{10}-v_{00}\ne0.
\]
The nonzero term is $d_K\rho(Q)$.

\section{Choosing higher homotopies}
Maps of codimension two have coefficient degree $-1$.
They can therefore fill a failure of~\eqref{eq:strict-corner-identity}
when that failure is a boundary in the appropriate Hom complex.
Subsequent codimensions impose further conditions on these choices.

\begin{proposition}[Extension obstruction]\label{prop:cellular-obstruction}
Fix the maps $t_{\tau\sigma}$ for all face pairs of codimension less than
$r$, where $r\ge2$, and suppose their curvature vanishes in those
codimensions.
For a face pair of codimension $r$, put
\[
 O_{\tau\sigma}=
       \sum_{\tau<\nu<\sigma}t_{\tau\nu}t_{\nu\sigma}.
\]
This is a closed element of
$\operatorname{Hom}^{2-r}(E_\sigma,E_\tau)$.
The chosen lower maps extend through codimension $r$ if and only if
every class $[O_{\tau\sigma}]$ vanishes in
$H^{2-r}\operatorname{Hom}(E_\sigma,E_\tau)$.
An extension is a choice of primitives satisfying
\[
 \delta t_{\tau\sigma}
       =(-1)^{\dim\tau+1}O_{\tau\sigma}.
\]
\end{proposition}
\begin{proof}
Set all as yet unchosen maps equal to zero and form the resulting $D$.
Filter its matrix entries by the drop in cell dimension.
The first possible nonzero drop in $D^2$ is $r$, and its component there
is $O_{\tau\sigma}$.
The identity $DD^2-D^2D=0$ in drop $r$ involves only $D_0$ and $O$.
Its $(\tau,\sigma)$ component is
\[
 (-1)^{\dim\tau}
 \bigl(d_\tau O_{\tau\sigma}-(-1)^r
                     O_{\tau\sigma}d_\sigma\bigr)=0.
\]
This is $(-1)^{\dim\tau}\delta O_{\tau\sigma}=0$,
because $|O_{\tau\sigma}|=2-r$.

Adding maps of codimension $r$ changes the curvature in that codimension
only by $(-1)^{\dim\tau}\delta t_{\tau\sigma}$.
Thus the displayed primitive equation is necessary and sufficient.
The equations for different face pairs at this stage are independent.
There are only finitely many stages for a finite-dimensional $K$.
\end{proof}

A cube shows why chosen face homotopies need not extend.
Orient $Q^3=[0,1]^3$ by $dx\wedge dy\wedge dz$,
and orient each face by its increasing free coordinates.
Take constant coefficients
$E=k\alpha\oplus k\beta$, with $|\alpha|=0$, $|\beta|=-1$ and $d_E=0$.
All cover maps are incidence numbers times the identity.

Let $F=\{z=0\}$ and $v=(0,0,0)$.
Among the codimension-two maps choose only
\[
 t_{vF}=h,\qquad h(\alpha)=\beta,\qquad h(\beta)=0
\]
to be nonzero.
All equations through codimension two hold: the cover composites cancel,
and $\delta h=0$.

The incidence of $F$ in $Q^3$ is $-1$, so
\[
 O_{v,Q^3}=t_{vF}t_{F,Q^3}=-h\ne0.
\]
The coefficient differential is zero and
$\operatorname{Hom}^{-2}(E,E)=0$.
Hence this nonzero class has no primitive.
The obstruction is to extending the specified choice of $h$.
The identity coefficient system also admits the different choice
in which every higher map is zero.

\section{Integration constructs an operation homotopy}
Now take $k=\mathbb R$ or $\mathbb C$ and a finite-dimensional cochain
complex $E$.
Suppose the closed cells of $K$ are compact smooth oriented manifolds
with corners, with compatible smooth face inclusions.
A form on $K$ means a family of forms on the closed cells whose
restrictions agree on faces.

Let
\[
 \omega=\sum_{p\ge0}\omega_p,
 \qquad \omega_p\in\Omega^p(K)\otimes E^{-p},
 \qquad d_{\mathrm{tot}}=d_{\mathrm{form}}+(-1)^p d_E
\]
on the summand of form degree $p$.
Thus $\omega$ has total degree zero.
For a cell $\sigma$ of dimension $r$, define
\begin{equation}\label{eq:signed-cell-integration}
 \eta_r=(-1)^{r(r-1)/2},\qquad
 I_\omega([\sigma])=\eta_r\int_\sigma\omega_r\in E^{-r}.
\end{equation}
The integral is taken coefficient by coefficient.

\begin{proposition}[Signed Stokes map]\label{prop:signed-stokes-map}
If $d_{\mathrm{tot}}\omega=0$, then $I_\omega$ is a chain map
from cellular chains placed in degree minus cell dimension to $E$:
\[
 d_E I_\omega([\sigma])=I_\omega(\partial[\sigma]).
\]
\end{proposition}
\begin{proof}
The component of $d_{\mathrm{tot}}\omega=0$ in form degree $r$ is
$d_{\mathrm{form}}\omega_{r-1}+(-1)^r d_E\omega_r=0$.
For $r\ge1$, Stokes' theorem gives
\[
 d_E I_\omega([\sigma])
 =\eta_r(-1)^{r+1}\int_{\partial\sigma}\omega_{r-1}
 =\eta_{r-1}\int_{\partial\sigma}\omega_{r-1}.
\]
The last equality uses
$\eta_r(-1)^{r+1}=\eta_{r-1}$.
The assertion at a vertex is the form-degree-zero equation
$d_E\omega_0=0$.
\end{proof}

For operation coefficients, take
$E=\operatorname{End}(V)(n):=\operatorname{Hom}(V^{\otimes n},V)$,
where $V$ is a finite-dimensional cochain complex and $E$ has differential
$\delta$.
On an interval with coordinate $t$, a total-degree-zero form is
\[
 \omega=m_t+dt\,H_t,
 \qquad |m_t|=0,\quad |H_t|=-1.
\]
It is closed precisely when
$\delta m_t=0$ and $\delta H_t=\partial_t m_t$.
Consequently
\[
 \delta\left(\int_0^1 H_t\,dt\right)=m_1-m_0.
\]
This formula constructs a homotopy from a closed coefficient-valued form.

The integration sign also respects products of cells.
Let $\alpha\otimes f$ and $\beta\otimes g$ have total degree zero,
with form degrees $p,q$ and coefficient degrees $-p,-q$.
For a composition of coefficient operations their product is
\[
 (\alpha\otimes f)(\beta\otimes g)
       =(-1)^{pq}(\alpha\wedge\beta)\otimes(f\circ g).
\]
The factor $(-1)^{pq}$ moves $f$ past $\beta$.

Since $\eta_{p+q}(-1)^{pq}=\eta_p\eta_q$, Fubini's theorem implies
that~\eqref{eq:signed-cell-integration} preserves this composition
on product cells with product orientation.
Suppose closed forms on a cellular operad restrict on every composition
face to the prescribed products of forms.
Their signed integrals then give an operation-valued chain map
preserving substitution.
Smooth extension to the closed cells and these face equations
are hypotheses of this construction.

\section{The associahedral equations}
Write $K_n$ for the associahedron of dimension $n-2$, for $n\ge2$.
Its cells are indexed by planar rooted trees with $n$ inputs.
The cell of a tree is a product of associahedra, one for each vertex.
Grafting trees gives the nonsymmetric cellular operad.
Its identification with the operad for associative algebras up to homotopy
is proved in \cite[\S9.2, Proposition~9.2.8]{LV}.
In that source the polytope is denoted $K^{n-2}$.

Place its cellular chains in negative degrees and change the orientation
of every $d$-cell by $(-1)^d$ relative to that convention.
This change respects product orientations and reverses each boundary
coefficient.
If $\theta_n$ is the oriented top cell, its boundary is
\begin{equation}\label{eq:associahedral-boundary}
 \partial\theta_n=
 -\sum_{\substack{r+s+t=n\,;\ s\ge2\,;\ r+1+t\ge2}}
 (-1)^{r+st}\theta_{r+1+t}\circ_{r+1}\theta_s.
\end{equation}
Here $\circ_i$ means substitution in the $i$th input.
The unary operadic identity acts by $1_V$.
No nullary operation is assumed.

\begin{proposition}[Operations on all associahedra]
\label{prop:associahedral-operations}
An operation-valued chain map from this cellular operad to
$\operatorname{End}(V)$ preserving substitution is equivalent to maps
\[
 m_n:V^{\otimes n}\longrightarrow V,
 \qquad |m_n|=2-n\quad(n\ge2),
\]
which, together with $m_1=d_V$, satisfy
\begin{equation}\label{eq:stokes-stasheff}
 \sum_{\substack{r+s+t=n\,;\ s\ge1}}
 (-1)^{r+st}
 m_{r+1+t}(1^{\otimes r}\otimes m_s\otimes1^{\otimes t})=0
 \qquad(n\ge1).
\end{equation}
These are the cohomological $A_\infty$ equations.
On homogeneous inputs the indicated tensor map has the additional
Koszul sign $(-1)^{(2-s)(|a_1|+\cdots+|a_r|)}$.
\end{proposition}
\begin{proof}
The graded cellular operad is freely generated by the corollas
$\theta_n$, $n\ge2$: every other cell is the unique substitution
specified by its planar tree.
Hence a substitution-preserving map is determined by
$m_n=I(\theta_n)$.
It commutes with differentials if and only if
$\delta m_n=I(\partial\theta_n)$ for each generator,
since the cellular boundary is a derivation for substitution.
Using~\eqref{eq:associahedral-boundary} gives~\eqref{eq:stokes-stasheff}
after adding the terms involving $m_1$.
Those terms are exactly
$d_Vm_n-(-1)^{2-n}m_n d_{V^{\otimes n}}$.

The arity-one equation is $d_V^2=0$.
This also proves the converse on every tree cell.
The resulting convention agrees with \cite[\S3.1]{KellerAInfinity2001}.
\end{proof}

\section{A nonassociative product with a coherent extension}
The interval operation and all subsequent equations can be made explicit.
A strict $A_\infty$ morphism will mean a chain map preserving each $m_n$.
It is a quasi-isomorphism when it induces an isomorphism on cohomology.
Over $\mathbb Q$, take
\[
 V^0=\mathbb Qx\oplus\mathbb Qy,\qquad
 V^{-1}=\mathbb Qz,\qquad dz=y.
\]
All other degrees and differentials vanish.

Define chain maps $\phi:V\to\mathbb Q$ and $i:\mathbb Q\to V$ by
\[
 \phi(x)=1,\quad\phi(y)=\phi(z)=0,\qquad i(1)=x-y,
\]
and put $P=1-i\phi$.
Let $h:V\to V$ have degree $-1$ with
$h(x)=h(y)=z$ and $h(z)=0$.
Directly on this basis,
\[
 P^2=P,\qquad dh+hd=P,\qquad
 h^2=0,\quad\phi h=0,\quad hi=0,\quad Ph=hP=h.
\]

Set $R=1+P$ and define
\[
 \mu(a,b)=R(a)\phi(b),\qquad
 H(a,b,c)=2h(a)\phi(b)\phi(c).
\]
The product is a chain map, because $R$ and $\phi$ are chain maps.
Since $\phi R=\phi$ and $R^2-R=2P$, its associator is
\[
 L_3(a,b,c)-R_3(a,b,c)=2P(a)\phi(b)\phi(c)=\delta H(a,b,c).
\]
In particular, $\mu(x,x)=x+y$ and the associator at $(x,x,x)$ is $2y$.

After extension to real coefficients, the form
\[
 R_3+t(L_3-R_3)+dt\,H
\]
is closed and has signed integral $H$ over the interval oriented from
$R_3$ to $L_3$.
The orientation in~\eqref{eq:associahedral-boundary} is the opposite one,
so the corresponding ternary operation is $m_3=-H$.

\begin{proposition}\label{prop:explicit-coherent-product}
The operations
\[
 m_1=d,\qquad m_2=\mu,\qquad m_3=-H,\qquad
 m_n=0\quad(n\ge4)
\]
satisfy~\eqref{eq:stokes-stasheff} in every arity.
Moreover $\phi$ and $i$ are strict $A_\infty$ quasi-isomorphisms between
this algebra and $\mathbb Q$, with $\phi i=1$.
\end{proposition}
\begin{proof}
The first three equations are $d^2=0$, the chain-map identity for $\mu$,
and $\delta m_3+L_3-R_3=0$, already proved.
In arity four, the terms involving $m_4$ vanish.
The term $\mu(1\otimes m_3)$ vanishes because $\phi h=0$.
In the order
\[
 \begin{gathered}
 -\mu(m_3\otimes1),\quad
 m_3(\mu\otimes1\otimes1),\\
 -m_3(1\otimes\mu\otimes1),\quad
 m_3(1\otimes1\otimes\mu),
 \end{gathered}
\]
the remaining terms are respectively $4,-4,2,-2$ times
$h(a)\phi(b)\phi(c)\phi(d)$.
Here $Rh=hR=2h$ and $\phi R=\phi$ give the coefficients.
They sum to zero.

In arity five, each possible composite inserts $m_3$ into $m_3$.
Insertion in the first input gives $h^2=0$.
Insertion in either other input gives $\phi h=0$.
For $n>5$, two nonzero operations cannot have combined arity $n$.
This proves all equations, including their graded tensor signs.

For strictness, the identities are
\[
 \begin{aligned}
 \phi\mu(a,b)&=\phi(a)\phi(b),&\qquad \phi m_3&=0,\\
 \mu(i(1),i(1))&=i(1),& m_3(i(1),i(1),i(1))&=0.
 \end{aligned}
\]
All higher operations vanish, and the unary identities are the chain-map
conditions for $i$ and $\phi$.
Finally $dh+hd=1-i\phi$ and $\phi i=1$ prove that both maps
induce inverse isomorphisms on cohomology.
\end{proof}

The product $\mu$ has no strict two-sided unit.
A right unit $e$ would give $\phi(e)=1$ on the input $x-y$
and $2\phi(e)=1$ on the input $y$, which is impossible.
The induced cohomology algebra is nevertheless $\mathbb Q$.
The example gives an operation on every associahedral cell by the
substitution construction of Proposition~\ref{prop:associahedral-operations}.

A fixed nonassociative product need not admit even the first homotopy.
Take $W=\mathbb Qx\oplus\mathbb Qy$ in degree zero with zero differential,
and set
\[
 \mu(x,x)=\mu(y,x)=y,\qquad
 \mu(x,y)=\mu(y,y)=0.
\]
Its associator at $(x,x,x)$ is $y$.
But $\operatorname{Hom}^{-1}(W^{\otimes3},W)=0$,
so no $H$ can satisfy $\delta H=L_3-R_3$.
There is no $A_\infty$ extension with this fixed binary operation.

For Fulton--MacPherson collision spaces or stable-graph parameter spaces,
the corresponding construction requires an actual coefficient-valued map
on the compactification.
This map must respect every face restriction, composition and gluing.
It also requires a tensor product and, where infinite sums occur,
a completion in which the operations and integrals are defined.
Singular forms require an extension or renormalization compatible with
those equations before Stokes' theorem can be applied.

The associahedral construction proves what such data yield in its stated
carrier.
The chiral and modular coefficient maps remain to be constructed.
Residue relations alone leave these operation homotopies and their
higher compatibility equations to be supplied.

\chapter{Two edge colors and two loop operators}\label{ch:two-edge-operators}

A loop is determined by its endpoints at the moment of contraction.
Its color alone does not determine when it becomes a loop.
Consider two genus-zero vertices joined by a transverse edge $e$ and a chiral edge $f$.
Give each vertex one labeled external leg, so that both vertices are stable.

Contracting $e$ first leaves a chiral loop.
Contracting $f$ first leaves a transverse loop.
Both orders produce the same genus-one corolla.
This graph will determine which coefficients may be assigned to the two kinds of contraction.

\section{Weights compatible with contraction}

Let $k$ be a field of characteristic zero and let $R$ be a commutative $k$-algebra concentrated in degree zero.
We use the stable graphs and coefficient maps of Proposition~\ref{cbc:cyclic:mixed}.
Thus a graph $G$ has a complex $C(G)$, and each edge $e$ has an equivariant cochain map
\[
 c_e^G:C(G)\longrightarrow C(G/e).
\]
The two composites for every pair of distinct edges agree after the specified identification of their targets.
These equations include pairs of different colors and pairs containing loops.

Write $r$ for the transverse color and $b$ for the chiral color.
Let $E_a(G)$ be the set of edges of color $a$.
Choose scalars $u_r,u_b,v_r,v_b\in R$.

A contraction of color $a\in\{r,b\}$ has weight $u_a$ when its endpoints are distinct, and weight $v_a$ when it is a loop.
The first condition allows the edge to lie on a cycle.
Denote this weight by $w_a(e,G)$.

Put the edge orientation line $\operatorname{or}(G)=\det(kE(G))$ in degree $-|E(G)|$.
Let $d_C$ be the coefficient differential and let $\iota_e$ remove $e$ from the oriented wedge with its exterior sign.
For $x\in C^q(G)$ and $\omega\in\operatorname{or}(G)$, define
\begin{equation}\label{eq:cbc:two-edge:weighted-d}
 D_{u,v}(x\otimes\omega)
 =d_Cx\otimes\omega
 +(-1)^q\sum_{e\in E(G)}w_{\operatorname{col}(e)}(e,G)
       c_e^Gx\otimes\iota_e\omega.
\end{equation}
The carrier is the direct sum over graph isomorphism classes, with automorphism coinvariants, as in that proposition.
Each summand has finitely many outgoing terms.

\begin{proposition}[The weighting equation]\label{prop:cbc:two-edge:weights}
The identity $D_{u,v}^2=0$ holds for every such coefficient system on the full category of stable graphs if and only if
\begin{equation}\label{eq:cbc:two-edge:weight-relation}
 u_rv_b=u_bv_r.
\end{equation}
In particular, $v_a=\hbar u_a$ gives a differential over $k[u_r,u_b,\hbar]$.
If $u_r=u_b=1$, independent indeterminates $v_r,v_b$ do not give a differential on this graph category.
\end{proposition}
\begin{proof}
The internal differential anticommutes with the edge terms because every $c_e^G$ is a cochain map.
For two distinct edges, the two contraction orders have opposite exterior signs and the same coefficient composite.
Equal scalar products therefore cancel the pair.

Contraction can change another edge from a nonloop to a loop only when the two edges join the same pair of vertices.
All other pairs retain their loop types, so their weights agree by commutativity of $R$.
For parallel edges of the same color, both products are $u_av_a$.
For differently colored parallel edges, they are $u_rv_b$ and $u_bv_r$.
Equation~\eqref{eq:cbc:two-edge:weight-relation} consequently cancels every two-edge term.

For necessity, use the graph in the opening paragraph with $C(G)=k$ and all coefficient maps equal to the identity.
Orient its edges by $e\wedge f$, with $e$ transverse.
The coefficient of the final corolla in $D_{u,v}^2$ is exactly
\[
 u_rv_b-u_bv_r.
\]
The two labeled external legs and the different edge colors prevent an automorphism from reversing this orientation.
Thus this coefficient must vanish.
\end{proof}

For a particular coefficient system, a nonzero scalar difference can annihilate the relevant composite.
The proposition gives the condition valid for every coefficient system.
It does not require either $u_a$ to be invertible.
For example, $u_r=u_b=0$ permits arbitrary loop weights.

The common parameter $\hbar$ counts the total number of independent graph cycles.
If a connected graph is contracted to a corolla, then exactly
\[
 b_1(G)=|E(G)|-|V(G)|+1
\]
of its contractions are loop contractions, in every ordering.
The other contractions reduce the number of vertices by one.
Consequently, when $v_a=\hbar u_a$, the scalar weight of the whole contraction is
\begin{equation}\label{eq:cbc:two-edge:whole-weight}
 u_r^{|E_r(G)|}u_b^{|E_b(G)|}\hbar^{b_1(G)}.
\end{equation}
Both edge colors remain visible in this formula.
The exponent of $\hbar$ is the ordinary graph cycle rank.

\section{A domain with two invariant genera}

There is also a contraction domain on which the two loop weights can remain independent.
Let $G_a$ be the spanning subgraph containing all vertices of $G$ and just the edges of color $a$.
Circuits include loops and pairs of parallel edges.
Suppose every simple circuit of $G$ is monochromatic.
Equivalently,
\[
 H_1(G;k)=H_1(G_r;k)\oplus H_1(G_b;k),
\]
where the summands have disjoint edge supports.

Indeed, simple circuits generate the cycle space.
A circuit using both colors cannot split in this way, since neither of its proper color subgraphs contains a cycle.

Give each vertex two nonnegative genus labels $g_r(v),g_b(v)$, whose sum is its ordinary genus.
A nonloop contraction adds the two labels separately.
A loop contraction of color $a$ raises $g_a(v)$ by one.
Define
\[
 \gamma_a(G)=\sum_{v\in V(G)}g_a(v)+b_1(G_a).
\]

\begin{proposition}[Monochromatic circuits]\label{prop:cbc:two-edge:two-genera}
Graphs whose simple circuits are monochromatic are closed under contraction.
On this category, each $\gamma_a$ is invariant and the independently weighted operator $D_{u,v}$ squares to zero.
For contraction of a connected graph to a corolla, the scalar weight is
\[
 \prod_{a\in\{r,b\}}
 u_a^{|E_a(G)|-b_1(G_a)}v_a^{b_1(G_a)}.
\]
\end{proposition}
\begin{proof}
A mixed circuit after contraction would lift to a mixed circuit before contraction, possibly by inserting the contracted edge.
Thus the class is closed under contraction.

Contract a nonloop edge of color $r$.
Its endpoints cannot be joined by a path consisting of blue edges, since that path together with the edge would be a mixed circuit.
Identifying the endpoints therefore creates no blue cycle.
The red cycle rank is unchanged by a red nonloop contraction.
The same argument applies with the colors exchanged.

A loop contraction lowers the cycle rank of its color by one and raises the corresponding vertex genus by one.
These observations prove invariance of both $\gamma_a$.
They also show that a complete contraction has exactly $b_1(G_a)$ loop contractions of color $a$, which gives the displayed weight.
There are no differently colored parallel edges in this class.
Every two-edge square in the proof of Proposition~\ref{prop:cbc:two-edge:weights} therefore cancels for arbitrary weights.
\end{proof}

This contraction category is not closed under arbitrary graph substitution.
Start with a red loop on one vertex having two external legs.
Replace that vertex by a blue edge joining two vertices, placing one flag of the red loop and one external leg at each vertex.
Both graphs used in the substitution satisfy the circuit condition.
The substituted graph is the mixed parallel-edge graph from the opening paragraph.
An operadic sewing theory on the restricted domain must therefore specify which substitutions it admits.

\section{The dual graph construction}

The graph differential above contracts an existing edge.
The Feynman transform uses the transpose operation: its differential expands a vertex by one edge.
We describe the grading and the generator maps before formulating its master equation.

Use cohomological complexes throughout.
For a finite edge set, let
\[
 \mathcal K(G)=\det(kE(G))
 \quad\text{in degree }|E(G)|.
\]
This is the degree-reversed version of the determinant cocycle in the homological convention of Getzler--Kapranov \cite[\S\S4.5--5.2]{GetzlerKapranov}.
It is distinct from the negatively graded orientation line in~\eqref{eq:cbc:two-edge:weighted-d}.

A graph contraction partitions its edge set into quotient edges and edges of the contracted subgraphs.
This partition identifies the determinant of all edges with the tensor product of the corresponding determinants.
These identifications include their Koszul signs.

A cocycle $\mathcal D$ assigns graded lines $\mathcal D(G)$ with such coherent tensor identifications under contraction and substitution.
Its value on a corolla is $k$.
Put $\mathcal D^\vee=\mathcal K\otimes\mathcal D^{-1}$.

A modular $\mathcal D$-operad with two edge kinds is a stable collection $P(g,n)$, with symmetric-group actions and degree-zero cochain maps
\[
 \mu_G^P:\mathcal D(G)\otimes
       \bigotimes_{v\in V(G)}P(g(v),n(v))
       \longrightarrow P(g(G),n(G)).
\]
They are equivariant and compatible with every graph substitution.
In particular, they obey all mixed contraction squares with the cocycle identifications.
The corolla composition is the identity.
The colors here distinguish edge operations on a common collection.
They are not two independent genus labels.

Assume each $P(g,n)$ is a bounded complex of finite-dimensional vector spaces and vanishes unless $2g-2+n>0$.
Write $P(g,n)^\vee$ for its cochain dual.
The graded Feynman transform is the free modular $\mathcal D^\vee$-operad generated by this dual collection:
\begin{equation}\label{eq:cbc:two-edge:free-transform}
 (\mathsf F_{\mathcal D}P)(g,n)
 =\bigoplus_{[G]}
 \left(\mathcal D^\vee(G)\otimes
       \bigotimes_vP(g(v),n(v))^\vee\right)_{\operatorname{Aut}(G)}.
\end{equation}
The sum is over stable graphs of type $(g,n)$ with labeled external legs and either color on each edge.
Graph substitution gives its free operadic composition.

For a one-edge graph $H$ of type $(g,n)$, transpose $\mu_H^P$ and tensor with its edge generator in $\mathcal K(H)$.
This defines a degree-one expansion
\[
 \delta_H:P(g,n)^\vee\longrightarrow
 \mathcal D^\vee(H)\otimes
       \bigotimes_vP(g(v),n(v))^\vee.
\]
In the free operad, apply this map at each vertex with the tensor-product differential sign.
Weight a two-vertex expansion of color $a$ by $u_a$ and a one-vertex loop expansion by $v_a$.
Together with the internal dual differential, these maps define a derivation $d_{\mathsf F,u,v}$ of degree one.

If~\eqref{eq:cbc:two-edge:weight-relation} holds, then
\[
 d_{\mathsf F,u,v}^{\,2}=0.
\]

To see this, expand a vertex by two edges.
The two expansion orders are the transposes of the two contraction orders.
Operadic compatibility identifies their coefficient maps, the determinant signs are opposite, and the scalar products agree by Proposition~\ref{prop:cbc:two-edge:weights}.
The internal differential anticommutes with the expansion terms because the original compositions are cochain maps.
This proves the assertion on generators and hence on the free operad.
It is the two-kind, weighted version of the construction and square-zero argument in Getzler--Kapranov's Theorem~5.2.

\section{Generator amplitudes and the master equation}

Let $Q$ be a modular $\mathcal D^\vee$-operad on the same graph category, whose components are also bounded complexes of finite-dimensional vector spaces.
For example, if $\mathcal D=\mathcal K$, a bounded finite-dimensional complex $V$ with a symmetric degree-zero perfect cochain pairing gives the untwisted target
\[
 Q(g,n)=V^{\otimes n}.
\]
Its graph compositions contract the factors at the two flags of each edge with the pairing.
Other pairing degrees require the corresponding complementary twist.

The generator carrier is
\begin{equation}\label{eq:cbc:two-edge:amplitudes}
 \mathcal A=
 \prod_{2g-2+n>0}
 \operatorname{Hom}_{S_n}\bigl(P(g,n)^\vee,Q(g,n)\bigr).
\end{equation}
For homogeneous $\phi\in\mathcal A$, its internal differential is
\[
 (d\phi)_{g,n}=d_Q\phi_{g,n}-(-1)^{|\phi|}\phi_{g,n}d_{P^\vee}.
\]
A degree-zero element $\Theta$ extends uniquely to a graded operad morphism
\[
 \widehat\Theta:\mathsf F_{\mathcal D}P\longrightarrow Q.
\]
On a decorated graph this extension applies $\Theta$ at every vertex and then $\mu_G^Q$.
Thus the carrier in~\eqref{eq:cbc:two-edge:amplitudes} records maps on the generators, rather than arbitrary linear maps on the whole free operad.

For a degree-zero generator map $\phi$, define its loop term by
\[
 (\Delta_a\phi)_{g,n}
 =-\sum_{\substack{H\text{ one-vertex, one-edge}\\
                   \operatorname{col}(H)=a}}
       \mu_H^Q(\operatorname{id}\otimes\phi)\delta_H.
\]
The internal-edge flags in this sum are unlabeled, and the external legs retain their labels.
Equivalently, use the automorphism coinvariants in~\eqref{eq:cbc:two-edge:free-transform}.

For the two-vertex graphs of color $a$, set
\[
 q_a(\phi)_{g,n}
 =-\sum_{\substack{H\text{ two-vertex, one-edge}\\
                   \operatorname{col}(H)=a}}
       \mu_H^Q(\operatorname{id}\otimes\phi\otimes\phi)\delta_H.
\]
This sum uses unordered vertices.
Let $B_a$ be its polarization, so that $q_a(\phi)=\tfrac12 B_a(\phi,\phi)$.

For homogeneous inputs, introduce formal coefficients of opposite degrees to make the generator maps degree zero.
Move those coefficients to the left by the Koszul rule in the preceding formulas.
This defines the operations in every degree and gives
\[
 |\Delta_a|=|B_a|=1,
 \qquad B_a(x,y)=(-1)^{|x||y|}B_a(y,x).
\]
The factor one-half removes the double counting introduced by polarization.

The product in~\eqref{eq:cbc:two-edge:amplitudes} has a concrete filtration.
Give its $(g,n)$ component the positive stability weight $w=2g-2+n$ and let $F^p\mathcal A$ consist of components of weight at least $p$.
There are finitely many stable pairs of each weight.
A separating term adds its two input weights, and a loop term preserves the input weight.
The one-edge sums at a fixed output have finitely many genus splittings and partitions of its external legs.
Hence all the displayed operations are continuous on this completion.

Formal scalar parameters may be completed separately in their adic topology.
Write $\mathcal A\widehat\otimes R$ for the product of the scalar-extended weight components.
The stability condition excludes the unstable vacuum components from this carrier.

\begin{theorem}[The two-kind master equation]\label{thm:cbc:two-edge:master}
Suppose~\eqref{eq:cbc:two-edge:weight-relation} holds.
Degree-zero amplitudes $\Theta\in\mathcal A\widehat\otimes R$ satisfying
\begin{equation}\label{eq:cbc:two-edge:master}
 d\Theta
 +\frac{u_r}{2}B_r(\Theta,\Theta)
 +\frac{u_b}{2}B_b(\Theta,\Theta)
 +v_r\Delta_r\Theta+v_b\Delta_b\Theta=0
\end{equation}
are in bijection with differential graded operad morphisms
\[
 (\mathsf F_{\mathcal D}P\otimes R,d_{\mathsf F,u,v})\longrightarrow Q\otimes R.
\]
For adically complete $R$, interpret the target and morphism componentwise in that completion.
\end{theorem}
\begin{proof}
Freeness already gives the unique graded morphism $\widehat\Theta$.
Both source and target differentials are operadic derivations.
Their compatibility can therefore be checked on $P(g,n)^\vee$.

On a generator, the difference
\[
 d_Q\widehat\Theta-\widehat\Theta d_{\mathsf F,u,v}
\]
has one internal term, the two one-vertex loop terms, and the two two-vertex terms.
Their definitions give exactly the left side of~\eqref{eq:cbc:two-edge:master}.
Thus the equation is necessary and sufficient.
The positive-weight filtration justifies the operations component by component.
This is the generator-map argument of Barannikov, with the edge kinds and their weights retained \cite[\S\S4--5]{BarannikovBV2007}.
\end{proof}

The same construction specifies the Maurer--Cartan convention.
On the graded space of generator maps, the difference of the two differentials defines an odd polynomial vector field $\mathcal Q$.
Its linear term is $d+v_r\Delta_r+v_b\Delta_b$, and its quadratic term is $\tfrac12(u_rB_r+u_bB_b)$.
The identities $d_Q^2=d_{\mathsf F,u,v}^2=0$ imply $\mathcal Q^2=0$.
Indeed, the difference is tangent to the space of graded operad maps, and applying it twice cancels the two mixed compositions.
Its linear, quadratic and cubic coefficients are respectively the differential, derivation and Jacobi identities.

To express these as an ordinary differential graded Lie algebra, put $L^i=\mathcal A^{i-1}$ and write $sx\in L^{|x|+1}$ for $x\in\mathcal A$.
Set
\[
 d_L(sx)=s(d+v_r\Delta_r+v_b\Delta_b)x,
 \qquad
 [sx,sy]_a=(-1)^{|x|}sB_a(x,y).
\]
Then $u_r[-,-]_r+u_b[-,-]_b$ is a degree-zero Lie bracket and $d_L$ is its differential.
For $\theta=s\Theta\in L^1$, equation~\eqref{eq:cbc:two-edge:master} becomes
\[
 d_L\theta+\frac12\bigl(u_r[\theta,\theta]_r+u_b[\theta,\theta]_b\bigr)=0.
\]
The amplitude has degree zero before this shift and degree one afterwards.

Taking all four weights to be one gives the ordinary two-kind Feynman transform.
Taking $v_a=\hbar u_a$ gives a family with a common loop parameter.

One may instead mark a genus-$g$ amplitude by $\hbar^g$.
A loop composition raises the genus by one, while a separating composition adds the input genera.
Thus the marking produces one factor of $\hbar$ in the loop term.
It does not justify multiplying twice for the same genus increment.

\section{The crossed derivation identity}

There is an algebraic form of the parallel-edge obstruction.
Use the operations on $L$ just constructed, and denote the shifted loop operators again by $\Delta_a$.
For homogeneous $X,Y\in L$, define
\[
 T_{a,b}(X,Y)=\Delta_a[X,Y]_b
       -[\Delta_aX,Y]_b-(-1)^{|X|}[X,\Delta_aY]_b.
\]
The unweighted graph identities give
\begin{gather*}
 \Delta_a^2=0,\qquad
 \Delta_r\Delta_b+\Delta_b\Delta_r=0,\\
 T_{r,r}=T_{b,b}=0,\qquad T_{r,b}+T_{b,r}=0.
\end{gather*}

For example, regard $u_r,u_b,\hbar$ as independent indeterminates in the valid family $v_a=\hbar u_a$.
Expand its identities and compare the coefficients of $u_r^2$, $u_b^2$ and $u_ru_b$.
This gives the two individual derivation identities and their crossed sum.
The separating brackets likewise satisfy the individual and mixed Jacobi identities.

Consequently, the failure of $v_r\Delta_r+v_b\Delta_b$ to be a derivation of $u_r[-,-]_r+u_b[-,-]_b$ is
\begin{equation}\label{eq:cbc:two-edge:crossed-defect}
 (v_ru_b-v_bu_r)T_{r,b}.
\end{equation}
With unweighted separating brackets and independent formal loop parameters, this vanishes identically precisely when $T_{r,b}=0$.
Then $T_{b,r}=0$ as well.
Thus independent loop parameters can occur on a particular amplitude carrier, but they require this additional compatibility.
Anticommutation of the two loop operators supplies only the linear part of that requirement.

\section{An explicit polynomial model}

The crossed defect need not vanish even for constant-coefficient operators.
Consider the graded commutative algebra
\[
 A=k[x,y]\otimes\Lambda(\xi,\eta,c),
 \qquad |x|=|y|=0,\quad |\xi|=|\eta|=-1,\quad |c|=1.
\]
Give $A$ the zero internal differential.
Use left derivatives in the exterior variables and define
\[
 \Delta_r=\partial_x\partial_\xi,
 \qquad \Delta_b=\partial_y\partial_\eta.
\]
Both operators have degree one, square to zero and annihilate $1$.
They anticommute because the two exterior derivatives anticommute.

For homogeneous $f,g\in A$, their second-order product defects define brackets
\begin{equation}\label{eq:cbc:two-edge:BV-bracket}
 \{f,g\}_a=(-1)^{|f|}
 \bigl(\Delta_a(fg)-(\Delta_af)g-(-1)^{|f|}f\Delta_ag\bigr).
\end{equation}
These brackets have degree one and satisfy
\[
 \{f,g\}_a=-(-1)^{(|f|+1)(|g|+1)}\{g,f\}_a.
\]
For example,
\[
 \{f,g\}_r=(\partial_xf)(\partial_\xi g)
              +(-1)^{|f|}(\partial_\xi f)(\partial_xg).
\]
The blue formula replaces $x,\xi$ by $y,\eta$.

These formulas prove the product derivation law
\[
 \{f,gh\}_a=\{f,g\}_ah+(-1)^{(|f|+1)|g|}g\{f,h\}_a.
\]
Applying the product derivation law twice shows that the Jacobiator is a derivation in each of its three inputs.
On the five algebra generators it vanishes, since every bracket of generators is constant.
It therefore vanishes on $A$.
This proof applies to every scalar linear combination of the two brackets.
Thus they are compatible odd Poisson brackets.

Finally, expand $\Delta_a^2(fg)=0$ using~\eqref{eq:cbc:two-edge:BV-bracket}.
The result is the derivation identity for $\Delta_a$ and its own bracket.
Applying the same argument to $\Delta_r+\Delta_b$ gives the sum of the two crossed derivation identities.

For $f=xy\xi$ and $g=x\eta$, one obtains
\[
 \{f,g\}_b=x^2\xi,\qquad
 \Delta_r\{f,g\}_b=2x,\qquad
 \{\Delta_rf,g\}_b=x,\qquad \Delta_rg=0.
\]
Hence the crossed derivation defect is $x$.
The defect with the colors exchanged is $-x$.
In the odd bracket convention the last term of a derivation identity has sign $(-1)^{|f|+1}$.
It follows that $v_r\Delta_r+v_b\Delta_b$ fails to be a derivation of the unweighted bracket sum by $(v_r-v_b)x$ on these two inputs.
This gives an actual coefficient realization of~\eqref{eq:cbc:two-edge:crossed-defect}.

There are also nonzero amplitudes for which both colors contribute to a valid quantum master equation.
Over $k[u_r,u_b,\hbar]$, put
\[
 \Theta=(xy)^2+c(u_bx\xi-u_ry\eta).
\]
It has degree zero.

The four color contributions are
\begin{align*}
 \Delta_r\Theta&=-u_bc,&
 \Delta_b\Theta&=u_rc,\\
 \tfrac12\{\Theta,\Theta\}_r&=-2u_b(xy)^2c,&
 \tfrac12\{\Theta,\Theta\}_b&=2u_r(xy)^2c.
\end{align*}
For the bracket calculation, the terms containing two factors of $c$ vanish and the two cross terms agree.
The displayed terms therefore prove
\begin{equation}\label{eq:cbc:two-edge:polynomial-solution}
 \frac12\bigl(u_r\{\Theta,\Theta\}_r+u_b\{\Theta,\Theta\}_b\bigr)
 +\hbar(u_r\Delta_r+u_b\Delta_b)\Theta=0.
\end{equation}
Both colors are needed for these cancellations.
Replacing the two loop coefficients by arbitrary $v_r,v_b$ leaves the residual $(u_rv_b-u_bv_r)c$.

The factor one-half can also be seen without graph notation.
Let $\Delta$ be either Laplacian or any scalar linear combination, and let $\{-,-\}$ be its bracket from~\eqref{eq:cbc:two-edge:BV-bracket}.
For an even element $S$, induction using the second-order product formula gives
\[
 \Delta(S^n)=nS^{n-1}\Delta S
       +\frac{n(n-1)}2S^{n-2}\{S,S\}.
\]
With a degree-zero adic variable $t$ and invertible $\hbar$, this implies the formal identity
\[
 \hbar^2e^{-tS/\hbar}\Delta(e^{tS/\hbar})
 =t\hbar\Delta S+\frac{t^2}{2}\{S,S\}.
\]
Every coefficient is a finite algebraic expression.
No analytic convergence is used.

The polynomial construction supplies a master equation on its specified graded algebra.
A geometric realization for an ordered chiral algebra requires a common coefficient object and the relevant cyclic pairings.
It also requires maps along both surface and curve clutchings.
Those maps must preserve the tensor products, support conditions and completions, and must satisfy the mixed contraction identities or their coherent replacements.

If two geometric genera are retained, their behavior under all admitted sewings must also be defined.
Theorem~\ref{thm:cbc:two-edge:master} applies once these data give its source and target operads and its generator amplitudes.
It does not construct those data from an edge coloring or a scalar partition function.

\chapter{Typed anomalies}\label{cbc:anomaly:chapter}

To cancel an anomaly one must specify the equation to be solved and the trivialization required.
On $X=\mathbb C^\times$, consider the trivial holomorphic line with connection
\[
 \nabla=d+\frac12\frac{dz}{z}.
\]
Its curvature $\nabla^2$ vanishes.
A horizontal section satisfies $df+\tfrac12f\,dz/z=0$, so it is locally a multiple of $z^{-1/2}$.
Continuation around a positively oriented circle changes its sign.
Thus its holonomy, the parallel transport around that circle, is $-1$, and no nowhere-zero global horizontal section exists.

The same distinction between an equation and the data it retains occurs in deformation theory.
A leading obstruction can be a cohomology class, while the curvature of an arbitrary partial construction need not belong to a complex at all.
We first construct the obstruction class and then identify its different carriers.
All algebraic sections use cohomological degrees over a field $k$ of characteristic zero.

\section{The differential at the classical solution}

Let $(L,d,[-,-])$ be a differential graded Lie algebra, with bracket of degree zero.
Its degree-one differential satisfies $d^2=0$ and is a graded derivation of the bracket.
Fix a classical Maurer--Cartan element $a_0\in L^1$, meaning
\[
 da_0+\frac12[a_0,a_0]=0,
 \qquad \partial=d+[a_0,-].
\]
The graded Jacobi identity gives $\partial^2=0$.
This is the differential at the chosen classical solution.

For a degree-zero formal variable $t$, write $L[[t]]=\prod_{j\geq0}t^jL$.
The product denotes completion in powers of $t$.
Each coefficient of a bracket is a finite sum.
For $a\in tL^1[[t]]$, define its curvature and the corresponding operator by
\[
 \mathcal R(a)=\partial a+\frac12[a,a],
 \qquad \partial_a=\partial+[a,-].
\]

Direct expansion gives
\begin{equation}\label{eq:cbc:anomaly:bianchi}
 \partial_a^2=[\mathcal R(a),-],
 \qquad \partial_a\mathcal R(a)=0.
\end{equation}
Indeed, the square is the sum of $[\partial a,-]$ and $\tfrac12[[a,a],-]$.
For the second identity, $\partial\mathcal R(a)=-[a,\partial a]$, while $[a,[a,a]]=0$ by graded Jacobi.
The second equation is the Bianchi identity.
The first shows why $\partial_a$ need not define an ordinary cochain complex.

\begin{proposition}[The leading obstruction]\label{cbc:anomaly:extension}
Let $n\geq2$ and let $a^{<n}=\sum_{j=1}^{n-1}t^ja_j$ satisfy $\mathcal R(a^{<n})=0$ modulo $t^n$.
The coefficient
\[
 \Omega_n=[t^n]\mathcal R(a^{<n})\in L^2
\]
is $\partial$-closed.
An extension $a^{<n}+t^nb$ solves the Maurer--Cartan equation modulo $t^{n+1}$ precisely when
\[
 \partial b=-\Omega_n.
\]
Consequently the extension exists if and only if $[\Omega_n]=0$ in $H^2(L,\partial)$.
\end{proposition}
\begin{proof}
In the Bianchi identity, the term $[a^{<n},\mathcal R(a^{<n})]$ is divisible by $t^{n+1}$.
Its coefficient of $t^n$ therefore gives $\partial\Omega_n=0$.
Adding $t^nb$ changes the curvature, modulo $t^{n+1}$, by $t^n\partial b$.
The bracket with $a^{<n}$ and the quadratic term in $b$ have higher order.
This proves both implications of the criterion.
\end{proof}

If $H^2(L,\partial)=0$, every $\partial$-closed first coefficient $a_1$ extends recursively to a solution in $tL^1[[t]]$.
At each step choose a primitive of $-\Omega_n$.
The resulting coefficients define a formal series, so no additional convergence assertion is needed.

For a nontrivial extension, take $L^1=kx\oplus ky$, $L^2=kz$, with
\[
 dx=dz=0,\qquad dy=-z,\qquad [x,x]=2z,
\]
and all other brackets zero.
These formulas define a differential graded Lie algebra because $z$ is central.
At $a_0=0$, the first term $tx$ has curvature $t^2z$, and the exact solution is
\[
 a=tx+t^2y.
\]
Removing $y$ instead gives a two-term graded Lie algebra with zero differential.
For that algebra the same fixed first term $tx$ cannot extend, since its second obstruction is the nonzero class of $z$.

Let $f:L\to M$ be a strict morphism of differential graded Lie algebras carrying $a_0$ to the chosen classical solution in $M$.
It intertwines the classical differentials and satisfies
\[
 f\bigl(\mathcal R(a)\bigr)=\mathcal R\bigl(f(a)\bigr).
\]
It therefore transports the leading obstruction classes.
Both differential and bracket are used in this equation.

On $kx\oplus kz$, the identity map to the same complex with zero bracket is an isomorphism of complexes.
It sends the source curvature $t^2z$ to $t^2z$, whereas the target curvature of $tx$ is zero.
An isomorphism of complexes alone does not compare these nonlinear equations.

\section{Central extensions and nilpotency}

For an ordinary Lie algebra $\mathfrak g$, use the Chevalley--Eilenberg complex
\[
 C^p(\mathfrak g;k)=\operatorname{Hom}_k(\Lambda^p\mathfrak g,k)
\]
with trivial coefficient action.
Our differential convention is
\[
 (\delta f)(x_0,\ldots,x_p)
 =\sum_{i<j}(-1)^{i+j}
 f\bigl([x_i,x_j],x_0,\ldots,\widehat x_i,\ldots,
                    \widehat x_j,\ldots,x_p\bigr).
\]
The remaining arguments retain their order.
In $\delta^2$, disjoint brackets cancel in pairs and the terms on three arguments cancel by Jacobi.
Thus this is a cochain complex.

An alternating form $\omega\in C^2(\mathfrak g;k)$ defines a bracket on $\mathfrak g\oplus kK$ by
\[
 [(x,aK),(y,bK)]=([x,y],\omega(x,y)K).
\]
The element $K$ is central.
The central component of Jacobi vanishes exactly when
\[
 -\omega([x,y],z)+\omega([x,z],y)-\omega([y,z],x)=0,
\]
which is $\delta\omega=0$.

Every central extension with fixed kernel $kK$ has this form after choosing a linear section of its projection to $\mathfrak g$.
Replacing the section $\sigma$ by $\sigma+\beta K$ changes its form to
\[
 \omega'=\omega-\beta\circ[-,-]=\omega+\delta\beta.
\]
Conversely, if $E_\omega$ denotes the displayed extension, this equation gives an isomorphism
\[
 E_{\omega'}\longrightarrow E_\omega,
 \qquad (x,aK)\longmapsto(x,(a+\beta(x))K).
\]
It follows that $H^2(\mathfrak g;k)$ classifies these extensions with identity maps on the quotient and kernel.

For instance, let $\mathfrak g=kh\oplus kx\oplus kz$ with $[h,x]=x$ and $z$ central.
The form with $\omega(h,z)=1$ and other basis values zero is closed and is not a coboundary, since every $\delta\beta$ vanishes on $(h,z)$.
The form with $\omega(x,z)=1$ instead has
\[
 (\delta\omega)(h,x,z)=-1
\]
and does not define a central extension.
A local central extension uses its specified local cochain complex.
The choice of admissible cochains is part of that problem.

For BRST nilpotency, let $(V,Q_0)$ be a cochain complex and set
\[
 L^j=\prod_i\operatorname{Hom}_k(V^i,V^{i+j}),
 \qquad [S,T]=ST-(-1)^{|S||T|}TS.
\]
Composition is the multiplication, and the differential on $L$ is $[Q_0,-]$.
A formal degree-one operator $Q(t)=Q_0+a(t)$ satisfies
\[
 Q(t)^2=[Q_0,a(t)]+a(t)^2
       =[Q_0,a(t)]+\frac12[a(t),a(t)].
\]
Once the lower coefficients vanish, its leading failure of nilpotency defines the degree-two class of Proposition~\ref{cbc:anomaly:extension}.
If the classical operator is changed, the endomorphism differential must be changed with it.

\begin{example}\label{cbc:anomaly:brst-example}
Let $V=ke_0\oplus ke_1\oplus ke_2$, with $|e_i|=i$, and let $Q_0=0$.
Define $q(e_0)=e_1$, $q(e_1)=e_2$ and $q(e_2)=0$.
The first term $tq$ solves nilpotency modulo $t^2$, but
\[
 [t^2]Q(t)^2=q^2,\qquad q^2(e_0)=e_2.
\]
The endomorphism differential is zero, so this nonzero obstruction cannot be removed by any coefficient of order $t^2$.
\end{example}

\section{A master anomaly invisible to operator nilpotency}

Let $A$ be a unital graded-commutative BV algebra.
Thus its degree-one operator $\Delta$ satisfies $\Delta1=0$, $\Delta^2=0$, and is of order at most two.
Use the derived bracket
\begin{equation}\label{eq:cbc:anomaly:BV-bracket}
 \{f,g\}=(-1)^{|f|}
 \bigl(\Delta(fg)-(\Delta f)g-(-1)^{|f|}f\Delta g\bigr).
\end{equation}
This bracket has degree one, is odd Poisson, and $\Delta$ is a graded derivation of it.
These are the BV identities used below.
Let $Q$ be a degree-one derivation of the product and bracket, with $Q^2=0$ and $Q\Delta+\Delta Q=0$.

Fix $S_0\in A^0$ satisfying $QS_0+\tfrac12\{S_0,S_0\}=0$.
For a degree-zero formal variable $\hbar$ and $S\in S_0+\hbar A^0[[\hbar]]$, put
\[
 C(S)=QS+\frac12\{S,S\}+\hbar\Delta S,
 \qquad D_S=Q+\{S,-\}+\hbar\Delta.
\]
The master equation is $C(S)=0$, with curvature of degree one.

The BV identities give
\begin{equation}\label{eq:cbc:anomaly:quantum-square}
 D_S^2=\{C(S),-\},\qquad D_SC(S)=0.
\end{equation}
To verify the first formula, the mixed anticommutators with $\{S,-\}$ are $\{QS,-\}$ and $\hbar\{\Delta S,-\}$, while its square is $\tfrac12\{\{S,S\},-\}$.
The second formula follows by expanding the curvature and using the same derivation identities and odd Jacobi.
In particular, $D_S^2=0$ detects the master curvature only through its Hamiltonian action.

\begin{proposition}[The quantum obstruction]\label{cbc:anomaly:BV-obstruction}
Write $\delta_0=Q+\{S_0,-\}$, so $\delta_0^2=0$.
Suppose $S^{<n}=S_0+\sum_{i=1}^{n-1}\hbar^iS_i$ solves the master equation modulo $\hbar^n$, for $n\geq1$.
Then
\[
 \Omega_n=\Delta S_{n-1}
   +\frac12\sum_{\substack{i+j=n\\i,j\geq1}}\{S_i,S_j\}
 \in A^1
\]
is $\delta_0$-closed.
An extension $S^{<n}+\hbar^nb$ solves the equation modulo $\hbar^{n+1}$ if and only if $\delta_0b=-\Omega_n$.
Its obstruction therefore lies in $H^1(A,\delta_0)$.
\end{proposition}
\begin{proof}
The displayed expression is the coefficient of $\hbar^n$ in $C(S^{<n})$.
Taking that coefficient in $D_{S^{<n}}C(S^{<n})=0$ gives $\delta_0\Omega_n=0$, since all lower curvature coefficients vanish.
Adding $\hbar^nb$ changes this coefficient by $Qb+\{S_0,b\}$.
Its Laplacian contributes only at order $n+1$.
This proves the extension criterion.
\end{proof}

The kernel of the Hamiltonian action can contain a nonzero obstruction class.
Take
\[
 A=k[x]\otimes\Lambda(\xi,c),\qquad
 |x|=0,\quad |\xi|=-1,\quad |c|=1,
 \qquad Q=0,\quad \Delta=\partial_x\partial_\xi,
\]
where $\partial_\xi$ is the left exterior derivative and monomials place $\xi$ before $c$.
The variable $c$ has zero bracket with every element.
For $S_0=x\xi c$, direct calculation gives
\[
 \{S_0,S_0\}=0,\qquad \Delta S_0=c,
 \qquad \delta_0x=xc,\quad \delta_0\xi=\xi c,\quad \delta_0c=0.
\]

Every degree-zero element is $f(x)+g(x)\xi c$, and
\[
 \delta_0\bigl(f(x)+g(x)\xi c\bigr)=xf'(x)c.
\]
All degree-one elements $k[x]c$ are closed.
Since $xk[x]c$ is the image from degree zero, there is an exact identification
\begin{equation}\label{eq:cbc:anomaly:central-class}
 H^1(A,\delta_0)=k\cdot[c],\qquad [c]\ne0.
\end{equation}
Thus no polynomial counterterm extends this fixed $S_0$ even to first quantum order.
Nevertheless $C(S_0)=\hbar c$ and $D_{S_0}^2=0$ on all of $A[[\hbar]]$.

Changing the counterterm algebra can change the answer.
In
\[
 \widetilde A=k[x,x^{-1},u]\otimes\Lambda(\xi,c),
 \qquad |u|=0,\quad \partial_xu=x^{-1},
\]
extend $\partial_x$ as a derivation and retain $\Delta=\partial_x\partial_\xi$.
The symbol $u$ is a formal logarithmic coordinate.
Now $\{S_0,u\}=c$, $\Delta u=0$ and $\{u,u\}=0$, so
\[
 S=S_0-\hbar u,\qquad C(S)=\hbar c-\hbar c=0.
\]

This is an exact solution in the enlarged algebra.
The polynomial obstruction remains valid on its original carrier.
For a local field theory, the admissible local functionals and renormalized operations must be specified before either obstruction calculation is applied.

\section{Curvature as algebraic structure}

A curved differential graded associative algebra consists of a graded unital algebra $B$, a degree-one derivation $d_B$, and an element $\kappa\in B^2$ such that
\[
 d_B^2(b)=[\kappa,b],\qquad d_B\kappa=0.
\]
The commutator is graded.
The zero-input operation is the degree-two map $m_0:k\to B$ with $m_0(1)=\kappa$.
If $[\kappa,-]\ne0$, the pair $(B,d_B)$ is not a cochain complex, so $\kappa$ is not a class in its ordinary cohomology.

For the graded vector space and operator $q$ of Example~\ref{cbc:anomaly:brst-example}, take
\[
 B=\operatorname{End}_k(V),\qquad d_B=[q,-],\qquad \kappa=q^2.
\]
Associativity gives $d_B^2=[q^2,-]$ and $d_B(q^2)=0$.
If $p_0$ is the projection onto $ke_0$, then
\[
 d_B^2(p_0)(e_0)=[q^2,p_0](e_0)=e_2.
\]
The failure to be a complex is already visible on this one operator.

There is a precise way to change curved data.
For $b\in B^1$, define
\begin{equation}\label{eq:cbc:anomaly:curved-change}
 d_{B,b}=d_B+[b,-],\qquad
 \kappa_b=\kappa+d_Bb+b^2.
\end{equation}
Then $d_{B,b}^2=[\kappa_b,-]$ and $d_{B,b}\kappa_b=0$.

The square follows by expanding the two derivations.
For the second identity, $d_B^2b=[\kappa,b]$ cancels $[b,\kappa]$, $d_B(b^2)=(d_Bb)b-b(d_Bb)$ cancels $[b,d_Bb]$, and $[b,b^2]=0$.
In the matrix example, $b=-q$ makes both the new differential and curvature zero.
The existence of such a $b$ is additional data in a general curved algebra.

If $\kappa$ is central, $d_B^2=0$ and a class $[\kappa]\in H^2(B,d_B)$ is defined.
Its meaning still depends on the construction.
For example, the curvature in \eqref{eq:cbc:weyl:curvature} records the scalar term of a Weyl relation in its quadratic dual.
Identifying it with a quantum master obstruction requires a map between the stated carriers that preserves the differential, curvature and relevant operations.

\section{The Atiyah extension, curvature and holonomy}

Let $L$ be a holomorphic line bundle on a complex manifold $X$.
The sheaf $\operatorname{At}(L)$ consists of first-order $\mathbb C$-linear operators on its sections whose symbol is a vector field: an operator with symbol $v$ obeys
\[
 P(fs)=fP(s)+v(f)s.
\]
Its symbol sequence is
\begin{equation}\label{eq:cbc:anomaly:Atiyah}
 0\longrightarrow\mathcal O_X
  \longrightarrow\operatorname{At}(L)
  \longrightarrow T_X\longrightarrow0.
\end{equation}

A holomorphic connection is an $\mathcal O_X$-linear splitting of this sequence.
It preserves Lie brackets precisely when its curvature vanishes.
The kernel consists of multiplication operators, with $[P,f]=v(f)$ for an operator of symbol $v$.
These are the connection and extension constructions of Atiyah (1958), Sections~1--3 \cite{AtiyahConnections1958}.

Choose local frames with $e_j=g_{ij}e_i$.
Writing $\nabla e_i=A_i e_i$, the Leibniz rule gives
\begin{equation}\label{eq:cbc:anomaly:connection-change}
 A_j=A_i+d\log g_{ij},\qquad d\log g_{ij}=g_{ij}^{-1}dg_{ij}.
\end{equation}
The forms $d\log g_{ij}$ are a \v{C}ech cocycle in $\Omega_X^1$.
Put $\alpha(L)=[\{d\log g_{ij}\}]\in H^1(X,\Omega_X^1)$.
For the connection $\sigma_i$ with zero one-form in the frame $e_i$, we have $\sigma_j-\sigma_i=-d\log g_{ij}$.
Thus the extension class of \eqref{eq:cbc:anomaly:Atiyah}, represented by those differences, is $-\alpha(L)$.

\begin{proposition}[The holomorphic connection obstruction]\label{cbc:anomaly:line-obstruction}
The class $\alpha(L)$ vanishes if and only if $L$ admits a holomorphic connection.
When a connection exists, its curvature is a closed holomorphic two-form $F_\nabla$.
Changing the connection by a global holomorphic one-form $\eta$ changes the curvature by $d\eta$.
\end{proposition}
\begin{proof}
Changing frames to $e_i'=h_i e_i$ replaces $d\log g_{ij}$ by
\[
 d\log g_{ij}+d\log h_j-d\log h_i,
\]
so its class is unchanged.
The class is zero precisely when, after refinement if necessary, there are holomorphic forms $A_i$ satisfying \eqref{eq:cbc:anomaly:connection-change}.
The corresponding local connections then glue, and every global connection supplies such forms.
Their curvatures $dA_i$ agree on overlaps because $d(d\log g_{ij})=0$.
Thus $F_\nabla=dA_i$ is global and closed.
Two connections differ by a global one-form, which proves the last assertion.
\end{proof}

On $\mathbb P^1$, take coordinates $z$ and $w=z^{-1}$ and the line with transition $e_\infty=z^m e_0$.
This is our frame convention for $\mathcal O(m)$.
Its transition one-form is $m\,dz/z$.

A holomorphic form on the $z$-chart has only nonnegative powers of $z$ in its Laurent expansion.
A holomorphic form $h(w)\,dw$ on the other chart becomes $-h(z^{-1})z^{-2}dz$, with powers at most $-2$.
Their difference has no $z^{-1}dz$ coefficient.
Therefore \eqref{eq:cbc:anomaly:connection-change} has a holomorphic solution exactly when $m=0$.
For $m=0$, take both connection forms to be zero.

For a chosen connection, its curvature form, the de Rham class of that form, and the integral first Chern class $c_1(L)\in H^2(X,\mathbb Z)$ are different data.
For local logarithms $\ell_{ij}$ of the transition functions, the latter is represented by the integral cocycle
\[
 n_{ijk}=\frac{\ell_{ij}+\ell_{jk}-\ell_{ik}}{2\pi i}.
\]

A flat connection additionally retains its holonomy.
For $\lambda\in\mathbb C$, the trivial line on $\mathbb C^\times$ with $\nabla=d+\lambda\,dz/z$ has parallel transport around the positive generator of $\pi_1(\mathbb C^\times,1)$ equal to
\[
 \exp(-2\pi i\lambda).
\]
This follows from the local horizontal solution $z^{-\lambda}$.
Both $\alpha(L)$ and $c_1(L)$ vanish because the underlying holomorphic line is trivial.
Its curvature is zero for every $\lambda$, whereas a horizontal trivialization exists exactly when $\lambda\in\mathbb Z$.

The determinant gives an actual comparison of curvatures with a fixed coefficient.
Let $E$ be an ordinary rank-$r$ complex vector bundle with connection $d+A$ in a local frame.
The induced connection on $\det E$ has one-form $\operatorname{tr}A$, so
\begin{equation}\label{eq:cbc:anomaly:determinant}
 F_{\det E}=d\operatorname{tr}A
           =\operatorname{tr}(dA+A\wedge A)
           =\operatorname{tr}F_E.
\end{equation}

Here $\operatorname{tr}(A\wedge A)=0$, since its off-diagonal terms cancel in pairs and each diagonal one-form squares to zero.
In particular, if $F_E=\omega\,\mathrm{id}_E$, then $F_{\det E}=r\omega$.
The scalar coefficient is supplied by this determinant construction.
Another family or normalization requires its own comparison.

\section{Comparison maps}

The preceding constructions locate their obstruction data as follows.
\begin{center}
\begin{tabularx}{\textwidth}{@{}p{0.30\textwidth}X@{}}
\toprule
datum&carrier and condition\\ \midrule
central extension&$H^2(\mathfrak g;k)$ in the declared cochain category.\\
leading BRST obstruction&$H^2(\operatorname{End}(V),[Q_0,-])$ after the lower nilpotency equations.\\
leading BV obstruction&$H^1(A,Q+\{S_0,-\})$ after fixing the classical solution.\\
holomorphic connection obstruction&$\alpha(L)\in H^1(X,\Omega_X^1)$.\\
connection curvature&$F_\nabla\in\Omega_X^2$, determined by a connection.\\
curved associative datum&$\kappa\in B^2$ with $d_B^2=[\kappa,-]$, without an automatic cohomology class.\\
\bottomrule
\end{tabularx}
\end{center}

To compare cohomology classes of different degrees, specify a degree-$r$ map $f:C^\bullet\to D^{\bullet+r}$ satisfying $d_Df=(-1)^rfd_C$ and prove the equality of the transported classes.
To compare nonlinear deformation problems, also preserve the brackets and curvature equations, or supply the higher comparison maps and their identities.
Numerical evaluations retain only the information detected by the specified cycle, trace, or background.

For a specified renormalized mixed bulk--boundary theory, start with its local BV deformation cochains.
A chiral comparison must map these to the vacuum-normalized supported insertion complex, with a degree shift and an anomaly identity.
The source, target differential, admissible counterterms, support, completion and compatibility with factorization and gluing must be constructed.
The finite algebraic comparisons above do not supply that map.

\part{Calabi--Yau and automorphic interfaces}

\chapter{Calabi--Yau categories}\label{cbc:cy:chapter}

The trace pairing on matrices is a first example of the geometry in this chapter.
On $\mathfrak{gl}_n(k)$, the form $(a,b)\mapsto\operatorname{tr}(ab)$ is symmetric and nondegenerate: the matrix units $E_{ij}$ and $E_{ji}$ are dual.
The classifying stack $BGL_n$ parametrizes rank-$n$ vector bundles, and its tangent complex at a bundle is its endomorphism complex shifted by $[1]$.
The shift changes the symmetric trace pairing into an alternating pairing with values in $k[2]$.

A Calabi--Yau structure extends this construction to categories whose endomorphism complexes have duality in dimension $d$.
The resulting moduli form has degree $2-d$.
Its closedness requires the cyclic part of that structure, in addition to the nondegenerate pairing.

\section{The two Calabi--Yau variances}

Let $k$ be a field of characteristic zero.
Use cohomological grading, so $V[r]^i=V^{i+r}$ and the shift map $s:V\to V[1]$ has degree $-1$.
Let $\mathcal C$ be a small stable idempotent-complete $k$-linear category, and put $\widehat{\mathcal C}=\operatorname{Ind}(\mathcal C)$.
Thus $\mathcal C$ is the subcategory of compact objects of $\widehat{\mathcal C}$.
Call $\mathcal C$ \emph{smooth} when its diagonal bimodule is perfect, and \emph{proper} when every morphism complex between its objects is perfect over $k$.
These conditions concern different duals.

Write $I$ for the identity endofunctor of $\widehat{\mathcal C}$.
For smooth $\mathcal C$, the left adjoint of categorical evaluation defines the inverse dualizing endofunctor $I^!$.
For proper $\mathcal C$, the right adjoint of evaluation defines the Serre endofunctor $S$.
The trace identifications are
\[
 T=CC_*(\mathcal C)\simeq
 \underline{\operatorname{RHom}}(I^!,I),\qquad
 T^\vee\simeq\underline{\operatorname{RHom}}(I,S),
\]
under the respective hypotheses, with $T^\vee=\operatorname{RHom}_k(T,k)$.
Here $CC_*(\mathcal C)$ is the Hochschild chain complex, with homological degree $j$ placed in cohomological degree $-j$.
The Hom objects on the right are taken in the category of colimit-preserving $k$-linear endofunctors.

The canonical circle action on $T$ is represented by its mixed-complex structure.
Its operators $b$ and $B$ have cohomological degrees $1$ and $-1$, with
\[
 b^2=B^2=bB+Bb=0.
\]
The negative cyclic complex is the completed product
\begin{equation}\label{cbc:cy:negative-cyclic}
 CC^-_*(\mathcal C)=(T[[u]],b+uB),\qquad |u|=2.
\end{equation}
It models the homotopy fixed points $T^{hS^1}$.
A class $\theta:k[d]\to CC^-_*(\mathcal C)$ has a representative
$\theta_0+u\theta_1+u^2\theta_2+\cdots$, where $|\theta_j|=-d-2j$ and
\begin{equation}\label{cbc:cy:lift-equations}
 b\theta_0=0,\qquad b\theta_{j+1}+B\theta_j=0\quad(j\geq0).
\end{equation}
The product in~\eqref{cbc:cy:negative-cyclic} retains all these components.

A \emph{smooth Calabi--Yau structure of dimension $d$} is such a class whose underlying Hochschild class induces an equivalence
\begin{equation}\label{cbc:cy:smooth-duality}
 \alpha_\theta:I^![d]\xrightarrow{\sim}I.
\end{equation}

For proper $\mathcal C$, a \emph{proper Calabi--Yau structure of dimension $d$} is a map
\[
 \tau:T_{hS^1}\longrightarrow k[-d]
\]
whose induced map $I\to S[-d]$ is an equivalence.
The subscript $hS^1$ denotes homotopy orbits.
The first structure is a class in a chain complex.
The second is a cyclic trace functional.
Both include coherent circle data.
Neither is specified merely by choosing a nondegenerate bilinear form on a space of morphisms.

Put $R=CC^*(\mathcal C)=\underline{\operatorname{RHom}}(I,I)$.
Composition of endofunctors gives $R$ its Hochschild $E_2$-structure.
Under~\eqref{cbc:cy:smooth-duality}, precomposition gives the cap equivalence
\[
 D_\theta:R\xrightarrow{\sim}T[-d].
\]
The proper structure instead gives $D_\tau:R\xrightarrow{\sim}T^\vee[-d]$ by postcomposition.

These are the orientation and coorientation of Theorem~\ref{cbc:center:cyclic}.
The circle-equivariant Hochschild calculus and either cap equivalence give a framed $E_2$-structure on $R$, retaining its original $E_2$-operations
\cite[Theorems~4.1.7 and 4.3.4]{BravRozenblyum}.
In this statement, \emph{framed} means an action of $E_2\rtimes SO(2)$, including the compatibility of rotation with every disk operation.
The dimension $d$ changes the cap shift while the operad remains $E_2$.

\section{The moduli functor and its tangent complex}

For a connective commutative differential graded $k$-algebra $A$, let $\operatorname{Perf}(A)$ denote the category of perfect $A$-complexes.
Fix the covariant convention
\begin{equation}\label{cbc:cy:moduli}
 \mathcal M_{\mathcal C}(A)=
 \operatorname{Fun}^{\mathrm{ex}}_k
       (\mathcal C,\operatorname{Perf}(A))^{\simeq}.
\end{equation}
The superscript denotes the space of exact functors and their equivalences.
These are pseudo-perfect functors: every object of $\mathcal C$ has perfect image over $A$.
Using contravariant modules would replace the source by $\mathcal C^{\mathrm{op}}$ and require the corresponding opposite Calabi--Yau class.

The functors in~\eqref{cbc:cy:moduli} assemble to a universal evaluation
\[
 F_{\mathcal C}:\mathcal C\longrightarrow
                  \operatorname{Perf}(\mathcal M_{\mathcal C}).
\]

If $\mathcal C$ is smooth, an $A$-point $F$ is corepresentable after scalar extension.
Writing $\mathcal C_A=\widehat{\mathcal C}\otimes D(A)$, there is a compact object $E_A\in\mathcal C_A$ with
\[
 F_A(-)\simeq\underline{\operatorname{RHom}}_{\mathcal C_A}(E_A,-).
\]
Here $F_A$ is the $A$-linear extension of $F$, and its values on compact objects are perfect over $A$.
This is the relative corepresentability theorem
\cite[Corollary~2.6 and Example~3.7]{BravDyckerhoff}.
The requirement on those values is the properness condition on the individual moduli object.
The whole category need not be proper.

The tangent complex is
\begin{equation}\label{cbc:cy:tangent}
 T_F\mathcal M_{\mathcal C}\simeq
 \underline{\operatorname{RHom}}_{\mathcal C_A}(E_A,E_A)[1].
\end{equation}
To see the shift, consider the loop space of $\mathcal M_{\mathcal C}$ at $F$.
It parametrizes automorphisms of $F$, or equivalently automorphisms of its corepresenter, using inversion for the opposite Yoneda variance.

Over a square-zero extension $A\oplus J$, an automorphism reducing to the identity has the form $1+f$, with
\[
 f:E_A\longrightarrow E_A\otimes_A J.
\]
Compactness and the properness condition identify this deformation functor with the one represented by the endomorphism complex of $E_A$.
The tangent of the loop space is therefore that complex.
Delooping gives~\eqref{cbc:cy:tangent}, compatibly with scalar extension
\cite[Proposition~3.3]{BravDyckerhoff}.

For the symplectic assertion below, assume $\mathcal M_{\mathcal C}$ is locally geometric and locally of finite presentation, with perfect cotangent complex $L_{\mathcal M_{\mathcal C}}$.
Thus it is covered by open derived Artin substacks on which the usual descent theory of differential forms applies.
A sufficient hypothesis is that $\mathcal C$ be of finite type, meaning compact in the Morita category of small stable idempotent-complete $k$-linear categories
\cite[Remark~3.9]{BravDyckerhoff}.
More generally, the construction applies whenever these geometric and perfectness hypotheses hold.

\section{From a cyclic class to a closed form}

Closedness of a derived form includes homotopies.
For an affine $U=\operatorname{Spec}(A)$, the derived de Rham complex is the graded mixed complex
\[
 \operatorname{DR}(A/k)(p)=\bigwedge_A^p L_{A/k}[p].
\]
Its internal differential has degree $1$.
The de Rham differential raises the weight $p$ by one and has degree $-1$ in this shifted complex.

For any such graded mixed complex $V$, put
\[
 \operatorname{NC}^{w}(V)(p)
   =\left(\prod_{j\geq0}u^j V(p+j),\ d+u\epsilon\right),
 \qquad |u|=2,
\]
where $\epsilon$ is its mixed differential.

The space of closed $p$-forms of degree $n$ is
\[
 \mathcal A^{p,\mathrm{cl}}(U,n)=
 \operatorname{Map}_{D(k)}
 \left(k,\operatorname{NC}^{w}(\operatorname{DR}(A/k))(p)[n-p]\right).
\]
Projection to $j=0$ gives the underlying $p$-form, a map
$k\to\bigwedge_A^pL_{A/k}[n]$.

For a derived stack, take the limit of these spaces over the opposite of the category of its affine charts.
Every limit over charts below uses this variance.
These definitions retain the higher components that specify closedness
\cite[Definitions~1.7--1.12]{PTVV}.

If $Y$ has perfect cotangent complex, its tangent complex is $T_Y=L_Y^\vee$.
A closed two-form $\omega$ of degree $n$ is \emph{$n$-shifted symplectic} when the adjoint of its underlying pairing is an equivalence
\begin{equation}\label{cbc:cy:nondegenerate}
 \omega^\flat:T_Y\xrightarrow{\sim}L_Y[n].
\end{equation}
Closedness and~\eqref{cbc:cy:nondegenerate} are separate conditions.

Apply Hochschild chains to the universal functor $F_{\mathcal C}$, restrict to each affine chart of $\mathcal M_{\mathcal C}$, and use the circle-equivariant derived Hochschild--Kostant--Rosenberg map.
Taking homotopy fixed points and projecting to weight two gives a map of complexes
\begin{equation}\label{cbc:cy:closed-map}
 \begin{aligned}
 \kappa_2:CC^-_*(\mathcal C)
 &\longrightarrow
 \lim_{U\to\mathcal M_{\mathcal C}}
                 CC^-_*(\operatorname{Perf}(U))\\
 &\longrightarrow
 \lim_{U\to\mathcal M_{\mathcal C}}
       \operatorname{NC}^{w}(\operatorname{DR}(U/k))(2).
 \end{aligned}
\end{equation}
All maps in the limit are pullbacks along affine charts.
The circle-equivariance and the weight projection make this a map of the full negative cyclic complexes, rather than only a map on Hochschild homology
\cite[Theorem~4.6 and Proposition~5.2]{BravDyckerhoff}.

We use the trace normalization of the last arrow: the induced pairing on shifted endomorphisms is composition followed by trace, with the suspension sign below.
Let $\operatorname{At}_E:E\to E\otimes L[1]$ be the Atiyah class.
With the Chern-character convention
$\operatorname{Ch}_2(E)=\tfrac12\operatorname{Tr}(\operatorname{At}_E^2)$, this normalization uses twice the weight-two projection.
The factor $\tfrac12$ in the Chern-character convention is explicit in
\cite[proof of Theorem~2.12]{PTVV}.
Scaling the whole closed-form map by $2$ scales its homotopies as well as its underlying pairing.

\begin{theorem}[Cyclic and shifted outputs]\label{cbc:cy:outputs}
Let $\mathcal C$ be smooth over the characteristic-zero field $k$, with a smooth Calabi--Yau structure $\theta$ of dimension $d$.
Its Hochschild cochains carry a framed $E_2$-algebra structure refining their usual $E_2$-structure.
If $\mathcal M_{\mathcal C}$ satisfies the geometric and perfectness hypotheses above, the composite
\[
 k[d]\xrightarrow{\theta}CC^-_*(\mathcal C)
       \xrightarrow{\kappa_2}
 \lim_{U\to\mathcal M_{\mathcal C}}
       \operatorname{NC}^{w}(\operatorname{DR}(U/k))(2)
\]
defines a $(2-d)$-shifted symplectic form.
At a field-valued point with corepresenter $E$, its pairing on homogeneous shifted endomorphisms is
\begin{equation}\label{cbc:cy:shifted-trace}
 \omega_\theta(sa,sb)=(-1)^{|a|}\operatorname{tr}_\theta(ab),
 \qquad
 \operatorname{tr}_\theta:\operatorname{RHom}(E,E)\to k[-d].
\end{equation}
\end{theorem}
\begin{proof}
The cap equivalence $D_\theta$ is induced by~\eqref{cbc:cy:smooth-duality}.
The circle-equivariant Hochschild calculus makes $T[-d]$ a module for the circle factorization homology of $R$.
The vector $\theta$ is circle-invariant, and its cap map identifies this module with $R$ on the distinguished interval.
The framed-algebra criterion used in Theorem~\ref{cbc:center:cyclic} supplies the compatible disk and rotation operations.
This proves the first assertion with its existing $E_2$-multiplication intact.

For the second, the composite in the statement is a cocycle of degree $-d$ in the weight-two negative cyclic complex.
By the definition of $\mathcal A^{2,\mathrm{cl}}$, it is a closed two-form of degree $2-d$.
Thus closedness is constructed before nondegeneracy is tested.

On a chart $\operatorname{Spec}(A)\to\mathcal M_{\mathcal C}$, put
$K_A=\underline{\operatorname{RHom}}_{\mathcal C_A}(E_A,E_A)$.
Relative Serre duality, applied to $\alpha_\theta$, gives an equivalence
\[
 K_A[d]\xrightarrow{\sim}K_A^\vee.
\]
Its adjoint is the pairing
\[
 K_A\otimes_A K_A\xrightarrow{\ \circ\ }K_A
             \xrightarrow{\operatorname{tr}_\theta}A[-d].
\]

The map on Hochschild chains induced by a corepresentable functor is computed by its unit, this duality map and evaluation.
The weight-two Hochschild--Kostant--Rosenberg component symmetrizes the two endomorphism inputs.
Cyclicity identifies the two terms, giving the displayed trace pairing with the chosen normalization.
This is the functorial trace calculation of
\cite[Corollary~2.5 and Propositions~4.5 and 5.3]{BravDyckerhoff}.
It identifies the pairing with the underlying form of~\eqref{cbc:cy:closed-map} on every affine chart, not only at individual points.

Shifting both inputs by $[1]$ gives~\eqref{cbc:cy:shifted-trace}.
Indeed, the interchange of a homogeneous endomorphism with the second suspension contributes $(-1)^{|a|}$.
Cyclicity then yields
\[
 \omega_\theta(sa,sb)
 =-(-1)^{(|a|-1)(|b|-1)}\omega_\theta(sb,sa).
\]
The shift differential is $d(sa)=-s(da)$.
On a strict fiber model, the differential compatibility reduces to
\[
 (-1)^{|a|}\operatorname{tr}_\theta((da)b)
       +\operatorname{tr}_\theta(a\,db)=0,
\]
which is the closed-trace identity applied to $d(ab)$.
The same suspension construction is a map of complexes over every chart.

Finally, $T_{\mathcal M_{\mathcal C}}|_A\simeq K_A[1]$ and
$L_{\mathcal M_{\mathcal C}}|_A\simeq K_A^\vee[-1]$.
The duality $K_A[d]\simeq K_A^\vee$ therefore identifies the adjoint of the form with
\[
 K_A[1]\xrightarrow{\sim}K_A^\vee[1-d]
      =L_{\mathcal M_{\mathcal C}}|_A[2-d].
\]
These equivalences are compatible with pullback.
Descent proves~\eqref{cbc:cy:nondegenerate} globally, completing the symplectic construction
\cite[Theorem~5.5(1)]{BravDyckerhoff}.
\end{proof}

For $\mathcal C=\operatorname{Perf}(k)$ and its unit orientation, the theorem recovers the degree-two form on the stack of perfect complexes.
On $BGL_n$ it is the trace pairing from the opening example.

For a graded complex $V$, the same calculation uses the supertrace
$\operatorname{str}(a)=\sum_i(-1)^i\operatorname{tr}(a|_{V^i})$ and differential $[Q,-]$, where $Q$ is the differential of $V$.
For example, let $|e_0|=0$, $|e_1|=1$ and $Qe_0=e_1$.
Taking $a=E_{00}$ and $b=E_{01}$ in the shifted chain identity gives zero with~\eqref{cbc:cy:shifted-trace}.
Omitting its suspension sign gives $2$.
The sign is required even though the underlying ungraded trace pairing is familiar.

\section{A smooth category with nonproper total space}

The category need not be proper for the moduli theorem.
Let $A=k[t]$ and $\mathcal C=\operatorname{Perf}(A)$.
An exact functor from this category is specified by an object and one endomorphism.
These finite data commute with filtered colimits of the recipient category, so $\mathcal C$ is of finite type.

The diagonal has the free bimodule resolution
\[
 0\longrightarrow A^e
   \xrightarrow{\ t\otimes1-1\otimes t\ }A^e
   \longrightarrow A\longrightarrow0,
 \qquad A^e=k[t_1,t_2],
\]
in degrees $-1,0$ before augmentation.
Multiplication identifies the cokernel with $k[t]$, and $t_1-t_2$ is a non-zero-divisor.
Thus $\mathcal C$ is smooth, with inverse dualizing bimodule $A^!\simeq A[-1]$.
It is not proper, since $\operatorname{RHom}_A(A,A)=k[t]$ is not perfect over $k$.

The degree-$-1$ generator of the diagonal resolution maps to $1[t]1$ in the normalized bimodule bar resolution, whose boundary is $t_1-t_2$.
After tensoring over $A^e$ with $A$, its Hochschild class maps to $dt$ under the Hochschild--Kostant--Rosenberg identification.
This class has Hochschild degree one, is closed under the de Rham operator, and generates the dualizing line.
It therefore gives a smooth one-dimensional Calabi--Yau structure.

For the covariant moduli convention, the corepresenters are the perfect complexes $E$ for which $\operatorname{RHom}_A(E,A)$ is perfect over $k$.
For $k[t]$, these are precisely the perfect complexes with finite-dimensional total cohomology, equivalently proper support on the affine line.
Decompose their finitely generated cohomology modules over the principal ideal domain $k[t]$.
Free summands persist under duality, while torsion summands have finite-dimensional dual cohomology.

At the residue module $k_a=A/(t-a)$, the two-term resolution
\[
 A\xrightarrow{\ t-a\ }A\longrightarrow k_a
\]
gives
\[
 \operatorname{Ext}_A^*(k_a,k_a)
       =k\cdot1\oplus k\cdot\varepsilon,
 \qquad |\varepsilon|=1,\quad\varepsilon^2=0.
\]
The differential becomes zero after applying $\operatorname{Hom}_A(-,k_a)$, and there is no degree-two term.

Choose the generator so that $\operatorname{tr}_{dt}(\varepsilon)=1$.
The unshifted pairing has matrix
$\left(\begin{smallmatrix}0&1\\1&0\end{smallmatrix}\right)$ in the basis $(1,\varepsilon)$.
After shifting, it pairs the degree-$-1$ automorphism direction with the degree-zero deformation direction and has values in $k[1]$.
This is the $2-d=1$ shift of the general theorem.
Discarding either direction destroys the duality.

\section{Cyclicity and closedness}

Nondegeneracy alone does not supply cyclicity.
For $D=\operatorname{diag}(1,2)$ on $k^2$, the functional
$\tau_D(a)=\operatorname{tr}(Da)$ gives a nondegenerate pairing $\tau_D(ab)$ on $\operatorname{Mat}_2(k)$, since left multiplication by $D$ is invertible.
But
\[
 \tau_D(E_{12}E_{21})=1,\qquad
 \tau_D(E_{21}E_{12})=2.
\]
It is not a cyclic trace and does not define the proper Calabi--Yau datum used above.

Nor does an everywhere nondegenerate geometric pairing have to be closed.
On
$U=\operatorname{Spec}k[x_1,x_2,x_3,x_4,(1+x_1)^{-1}]$, put
\[
 \eta=dx_1\wedge dx_2+(1+x_1)dx_3\wedge dx_4.
\]
Its determinant is $(1+x_1)^2$, while
\[
 d_{\mathrm{dR}}\eta=dx_1\wedge dx_3\wedge dx_4\ne0.
\]
Since $U$ is an ordinary smooth affine scheme, a closed degree-zero two-form is an ordinary closed differential form.
Thus $\eta$ has no closing structure, despite being nondegenerate at every point.

On a derived complex, higher components can instead supply the required homotopy.
For instance, let a mixed complex have basis $v_0,a,v_1$ in degrees $0,-1,-2$, with
\[
 Bv_0=a,\qquad bv_1=-a,
\]
and all other values of $b,B$ zero.
Then $v_0+uv_1$ is a negative cyclic cycle, although $Bv_0\ne0$.
Its second component implements $bv_1+Bv_0=0$ in~\eqref{cbc:cy:lift-equations}.
This example explains why a cyclic lift includes data beyond strict invariance of a chosen Hochschild representative.

\section{The ordered chiral comparison}

The categorical outputs above have a specified coefficient category and cyclic class.
To obtain an ordered chiral algebra on a curve $X$, one needs geometric factorization data on $X^{\mathrm{an}}\times\mathbb R$.

Concretely, let $\mathcal F$ be a unital factorization theory whose interval coefficients
$B_J(U)=\mathcal F(U\times J)$ lie in a symmetric monoidal category $(\mathcal H,\odot,\mathbf1)$ of holomorphic factorization objects.
Assume interval inclusions act by equivalences and that these coefficients have a compatible augmentation.

Supply a strong symmetric monoidal comparison
\[
 \Phi:(\mathcal H,\odot,\mathbf1)
       \longrightarrow
       (\mathcal C_D,\tensor^!,\omega_{\operatorname{Ran}}),
 \qquad \mathcal C_D\subseteq\operatorname{Fact}_D(X),
\]
where $\mathcal C_D$ specifies the admitted algebraic coefficients and $\Phi$ preserves their coherent collision and separated factorization maps.
Theorem~\ref{cbc:ord:recognition} then constructs the ordered chiral algebra $\Phi(B_{J_0})$.
If $\phi_2$ is the tensorator of $\Phi$, its multiplication is
\[
 \Phi(B_{J_0})\tensor^!\Phi(B_{J_0})
   \xrightarrow{\phi_2}\Phi(B_{J_0}\odot B_{J_0})
   \xrightarrow{\Phi(\mu)}\Phi(B_{J_0}),
\]
with unit and augmentation transported by the same monoidal comparison.
Its locality and higher compatibilities follow from that comparison and the interval operations.

For a given Calabi--Yau category, constructing $\mathcal F$ and $\Phi$, and identifying their cyclic pairing with $\theta$, are additional geometric problems.
The framed Hochschild operations do not themselves specify collision supports, algebraic coefficients or vacuum insertions on $X$.
A further bulk-to-center equivalence must identify the actual local action and its factorization maps.
The Calabi--Yau theorem supplies its categorical input and moduli form while leaving those comparisons to their stated hypotheses.

\chapter{Hall and BPS algebras}\label{cbc:hall:chapter}
An extension product counts the ways to assemble an object from two others.
For vector spaces over the finite field $\mathbb F_q$, let $u_n$ denote the class of an $n$-dimensional space.
An extension with quotient of dimension $a$ and subobject of dimension $b$ amounts to a $b$-dimensional subspace of $\mathbb F_q^{a+b}$.
Thus the Hall product is
\begin{equation}\label{cbc:hall:gaussian}
 u_a*u_b={a+b\brack b}_q u_{a+b},
 \qquad
 {n\brack b}_q=
 \prod_{i=0}^{b-1}\frac{q^n-q^i}{q^b-q^i}.
\end{equation}
The numerator counts ordered independent $b$-tuples, and the denominator counts ordered bases of a fixed subspace.
In particular, $u_1*u_1=(q+1)u_2$ and $u_0$ is the unit.

Multiplying three classes counts flags.
Choosing the smaller subspace first or the larger one first gives
\[
 {a+b\brack b}_q{a+b+c\brack c}_q
 = {b+c\brack c}_q{a+b+c\brack b+c}_q.
\]
For $a=b=c=1$ and $q=2$, both sides count the $21$ complete flags in $\mathbb F_2^3$.
This identity is the elementary form of Hall associativity.
Geometric Hall theories replace counting by operations on cohomology, retaining the same flag comparison.
Their BPS generators and Yangian presentations require further geometry, which determines the coefficient objects and the resulting algebra.

\section{The extension correspondence}
Let $\mathcal A$ be an exact category with a moduli stack
$\mathfrak M=\coprod_{\alpha\in\Gamma}\mathfrak M_\alpha$.
For a stable infinity-category, use cofiber sequences and derived flag stacks in the same construction.
Here $\Gamma$ is an additive charge monoid, and the charge of an extension is the sum of the charges of its factors.
Write $\mathfrak E_{\alpha,\beta}$ for the stack of exact sequences
\[
 0\longrightarrow E_\beta\longrightarrow E
   \longrightarrow E_\alpha\longrightarrow0.
\]
The maps recording the two factors and the middle term give
\[
 \mathfrak M_\alpha\times\mathfrak M_\beta
 \xleftarrow{\ p\ }\mathfrak E_{\alpha,\beta}
 \xrightarrow{\ q\ }\mathfrak M_{\alpha+\beta}.
\]

Fix a coefficient theory $H$, an external product $\boxtimes$, a pushforward $q_*$ and a normalized pullback $p^\sharp$ for these maps.
In a Borel--Moore theory, $p^\sharp$ is generally a refined or virtual Gysin map.
Its shifts, orientation and coefficient class form part of the definition.
Set
\begin{equation}\label{cbc:hall:product}
 a*b=q_*p^\sharp(a\boxtimes b).
\end{equation}

\begin{proposition}[Associativity from flags]\label{cbc:hall:flags}
Assume that these operations are functorial and compatible with external products.
Assume the base-change comparisons for both iterated extension diagrams, with the same shifts and orientations.
Then~\eqref{cbc:hall:product} is associative.
If extension by the zero object gives the identity correspondence with its identity coefficient map, its class is a unit.
\end{proposition}
\begin{proof}
Let $\mathfrak F_{\alpha,\beta,\gamma}$ classify flags
\[
 0\subset F_1\subset F_2\subset E,
 \qquad
 [E/F_2]=\alpha,\quad [F_2/F_1]=\beta,\quad [F_1]=\gamma.
\]
There are canonical equivalences of stacks
\[
 \begin{aligned}
 \mathfrak F_{\alpha,\beta,\gamma}
 &\simeq
 \mathfrak E_{\alpha,\beta}
  \mathop{\times}_{\mathfrak M_{\alpha+\beta}}
 \mathfrak E_{\alpha+\beta,\gamma},\\
 \mathfrak F_{\alpha,\beta,\gamma}
 &\simeq
 \mathfrak E_{\alpha,\beta+\gamma}
  \mathop{\times}_{\mathfrak M_{\beta+\gamma}}
 \mathfrak E_{\beta,\gamma}.
 \end{aligned}
\]
In the derived setting these are homotopy fiber products.

In the first, the inner extension is $0\to F_2/F_1\to E/F_1\to E/F_2\to0$.
Its inverse takes the preimage of the subobject in $E/F_1$.
In the second, the inner extension is $0\to F_1\to F_2\to F_2/F_1\to0$.
These constructions respect isomorphisms and families.

Let $r$ record the three successive factors and let $s$ record $E$.
Base change and functoriality identify each parenthesization with
\[
 s_*r^\sharp(a\boxtimes b\boxtimes c).
\]
The equality uses the stipulated agreement of the two Gysin normalizations.
The two zero-object correspondences prove the unit assertion.
\end{proof}

Normalization matters already for finite vector spaces.
Multiply the coefficient in~\eqref{cbc:hall:gaussian} by $-1$ only when $a=b=1$.
The unit is preserved, but at $q=2$ the two products of charges $(1,1,2)$ have coefficients $-105$ and $105$.
The modified normalization fails the flag comparison.

The proposition concerns the product on $H(\mathfrak M)$.
A product on complexes with coherent higher associativity also requires the compatible operations for flags of every length.
For a critical Hall theory, the vanishing-cycle coefficients and their orientation must participate in those comparisons.

One may form $\bigoplus_\alpha H(\mathfrak M_\alpha)$ when the product has finite charge support.
If every charge has finitely many decompositions, the same formula extends to the product completion $\prod_\alpha H(\mathfrak M_\alpha)$.
Each output then contains a finite convolution sum.
The supported surface algebra below uses a different topology, determined by its moduli stack.

\section{The BPS algebra and its Lie generators}
We now work over $\mathbb C$ with rational mixed Hodge coefficients.
Let $\mathcal A$ be an admissible finite-length Abelian subcategory of the homotopy category of a dg category $\mathcal C$.
Assume its objects form the classical truncation of an open $1$-Artin substack of the derived moduli of $\mathcal C$.
The geometric hypotheses needed here are the following.

\begin{enumerate}
\item The middle-term map $q$ is proper and representable.
The object stack has the resolution property, and the shifted relative $\operatorname{RHom}[1]$ complex on each pair of components has a global three-term vector-bundle presentation.
\item There is a good moduli map $\operatorname{JH}:\mathfrak M\to M$, where $M$ is a separated scheme with finite-type connected components.
The direct-sum map $\oplus:M\times M\to M$ is finite.
\item For every finite collection of simple objects, the full dg subcategory on those objects has a right $2$-Calabi--Yau structure.
The relative Hall construction is associative and unital with its specified virtual pullbacks.
\end{enumerate}
These are Assumptions~1--5 and~7 of Davison--Hennecart--Schlegel Mejia~\cite[\S\S4,6]{DHS}.
They include more than the existence of a good moduli space.
Choose a surjective charge map $\pi_0(M)\to\Gamma$ with finite fibers, where $\Gamma$ embeds in a lattice.
Assume the Euler form of $\mathcal C$ descends to a symmetric form $\chi$ on that lattice.

For complexes of mixed Hodge modules on $M$, define convolution by
\[
 F\boxdot G=\oplus_*(F\boxtimes G).
\]
The finiteness of $\oplus$ makes this tensor exact on mixed Hodge modules.
The skyscraper sheaf $\mathbb Q$ at the zero object is its unit.

Let $\mathbb L=H_c^*(\mathbb A^1,\mathbb Q)$ be the Tate complex, and use the normalization
\[
 \mathbb Q^{\mathrm{vir}}_{\mathfrak M_\alpha}
   =\mathbb Q_{\mathfrak M_\alpha}\otimes\mathbb L^{\chi(\alpha,\alpha)/2},
 \qquad
 \mathscr H=\operatorname{JH}_*\mathbb D\mathbb Q^{\mathrm{vir}}_{\mathfrak M}.
\]
Here $\mathbb D$ denotes Verdier duality.
The displayed pushforward denotes the defined complex on $M$, constructed by the stack approximations in the relative Hall theory.
It does not require a category of mixed Hodge modules on an arbitrary Artin stack.
The Hall product makes $\mathscr H$ an algebra in $D^+(\operatorname{MHM}(M))$.

The $2$-Calabi--Yau purity theorem gives $\mathcal H^i(\mathscr H)=0$ for $i<0$ and makes $\mathcal H^0(\mathscr H)$ semisimple~\cite[Lemma~6.3]{DHS}.
Define the relative BPS algebra and its absolute algebra by
\[
 \mathscr B_{\mathrm{Alg}}=\mathcal H^0(\mathscr H),
 \qquad
 B_{\mathrm{Alg}}=\mathbb H^*(M,\mathscr B_{\mathrm{Alg}}).
\]
The degree here is the mixed-Hodge-module degree, corresponding to perverse degree on the underlying sheaves.
It is taken before global sections.
By the decomposition theorem, $\mathscr B_{\mathrm{Alg}}\to\mathscr H$ has a retraction as a map of complexes.
Consequently $B_{\mathrm{Alg}}\to\mathbb H^*(M,\mathscr H)$ is an injective algebra map~\cite[Proposition~6.4]{DHS}.

The enveloping-algebra statement uses a sign twist.
Choose a bilinear form $\psi:\Gamma\times\Gamma\to\mathbb Z/2$ such that
\[
 \psi(\alpha,\beta)+\psi(\beta,\alpha)
       =\chi(\alpha,\beta)\pmod2,
 \qquad
 a*_{\psi}b=(-1)^{\psi(\alpha,\beta)}a*b.
\]
Such a choice exists because the diagonal values of $\chi$ are even.
Bilinearity gives the cocycle identity that preserves associativity.
Superscript $\psi$ will indicate this changed multiplication.
The cohomological Koszul signs remain in force as well.

We describe the generators before stating the comparison.
Write $p(\alpha)=2-\chi(\alpha,\alpha)$.
Let $\Sigma$ consist of the nonzero charges with $p(\alpha)\geq0$ such that
\[
 p(\alpha)>\sum_{i=1}^r p(\alpha_i)
 \quad\text{for every decomposition }\quad
 \alpha=\sum_{i=1}^r\alpha_i,
 \quad r\geq2,\quad\alpha_i\ne0.
\]
The simple-root set $\Phi^+$ is the union of $\Sigma$ with all positive multiples of its isotropic elements.
Here isotropic means $\chi(\alpha,\alpha)=0$.

For $\alpha\in\Sigma$, set $G_\alpha=\operatorname{IC}(M_\alpha)$.
We normalize the intersection complex to have weight zero, restricting on a smooth component of dimension $d$ to $\mathbb Q[d](d/2)$.
For $\alpha=n\alpha_0$, with $\alpha_0\in\Sigma$ isotropic and $n\geq2$, set
\[
 G_\alpha=(\delta_n)_*\operatorname{IC}(M_{\alpha_0}),
 \qquad
 \delta_n(x)=x^{\oplus n}.
\]
Thus a multiple isotropic root uses the small diagonal, rather than the intersection complex of the entire component.
The definitions are taken componentwise before a finite-fiber charge map combines components.

Let $\mathfrak n^+_{M,G}$ be the Lie algebra object generated by $G=\bigoplus_{\alpha\in\Phi^+}G_\alpha$ under convolution, with the generalized Serre relations.
For distinct simple-root generators $x$ of charge $\alpha$ and $y$ of charge $\beta$, these include
\[
 \begin{aligned}
 (\operatorname{ad}x)^{1-\chi(\alpha,\beta)}y&=0
       &&\text{if }\chi(\alpha,\alpha)=2,\\
 [x,y]&=0&&\text{if }\chi(\alpha,\beta)=0.
 \end{aligned}
\]
The orthogonal relation also applies within an isotropic generator space.
After global sections, brackets use the sign $(-1)^{|x||y|}$.
The resulting graded Lie algebra is the positive half of the generalized Kac--Moody Lie superalgebra with this Euler form and these generator spaces.

\begin{theorem}[Geometric BPS and PBW comparisons]\label{cbc:hall:bps}
Under the hypotheses above, the canonical generator map induces
\[
 U_{\boxdot}(\mathfrak n^+_{M,G})
       \xrightarrow{\sim}\mathscr B_{\mathrm{Alg}}^\psi,
 \qquad
 U(\mathfrak n^+_{\mathrm{BPS}})
       \xrightarrow{\sim}B_{\mathrm{Alg}}^\psi,
 \qquad
 \mathfrak n^+_{\mathrm{BPS}}=\mathbb H^*(M,\mathfrak n^+_{M,G}).
\]
These are algebra isomorphisms.
Suppose additionally that $\mathfrak M$ has a positive determinant line bundle: scalar automorphisms of every nonzero object act on its fiber with positive weight.
With $u=c_1(\mathcal O_{B\mathbb G_m}(1))$ of degree $2$, there is a PBW isomorphism
\begin{equation}\label{cbc:hall:pbw}
 \operatorname{Sym}\bigl(\mathfrak n^+_{\mathrm{BPS}}\otimes\mathbb Q[u]\bigr)
   \xrightarrow{\sim}\mathbb H^*(M,\mathscr H^\psi)
\end{equation}
of charge-graded mixed Hodge structures with cohomological grading.
The symmetric algebra uses the Koszul symmetry.
The map is symmetrized Hall multiplication, not an asserted algebra isomorphism from a commutative algebra.
\end{theorem}
\begin{proof}
The canonical inclusions of the intersection-complex summands give
$G\to\mathscr B_{\mathrm{Alg}}^\psi$.
Their Hall products define a map from the free convolution algebra.
Davison--Hennecart--Schlegel Mejia prove that it annihilates the Serre ideal and induces an isomorphism~\cite[Proposition~7.13 and Theorem~1.1]{DHS}.
This proves the first assertion.

Its proof uses the local Ext-quiver identification of the generator and BPS sheaves, including their products.
The induction on dimension vector reduces the kernel and cokernel to the strictly seminilpotent quiver calculation~\cite[Proposition~7.7, Corollary~7.14, \S7.7.2]{DHS}.
It therefore identifies the actual generator map, beyond its graded character.
Taking global sections gives the second isomorphism by the same convolution comparison~\cite[Theorem~1.2]{DHS}.

For the PBW map, compose
$\mathfrak n^+_{M,G}\to\mathscr B_{\mathrm{Alg}}^\psi\to\mathscr H^\psi$.
The determinant line bundle lets $\mathbb Q[u]=H^*(B\mathbb G_m,\mathbb Q)$ act, giving
$\mathfrak n^+_{M,G}\otimes\mathbb Q[u]\to\mathscr H^\psi$.
Apply iterated multiplication and average over permutations with their Koszul signs.
This constructs the relative symmetric-tensor map.

The relative PBW theorem proves it is an isomorphism~\cite[Theorem~1.7]{DHS}.
Concretely, a nonzero cone would have a nonzero costalk at a closed point.
The local Ext-quiver comparison identifies that costalk with the corresponding preprojective PBW comparison, whose cone vanishes.
Global sections give~\eqref{cbc:hall:pbw}, as in~\cite[Corollary~1.8]{DHS}.
\end{proof}

The relative BPS Lie algebra is $\mathfrak n^+_{M,G}$.
Its construction precedes any choice of a geometric coproduct on the full Hall algebra.
The enveloping-algebra isomorphism transports a cocommutative Hopf structure to $B_{\mathrm{Alg}}^\psi$.
Identifying that structure with geometric cooperation, or extending it to the full CoHA, requires a separate comparison.

\section{Point sheaves on a K3 surface}
Let $S$ be a projective K3 surface and let $\mathcal A=\operatorname{Coh}_0(S)$ be its category of zero-dimensional coherent sheaves.
The charge is length, and $M_n=\operatorname{Sym}^n(S)$.
The geometric hypotheses of Theorem~\ref{cbc:hall:bps} hold for this surface category~\cite[\S4]{DHS}.
The Euler form on point charges is zero, so $\Sigma=\{1\}$ and $\Phi^+=\mathbb Z_{>0}$.
Take $\psi=0$.

Every simple root is isotropic.
Therefore
\[
 G_n=(\delta_n)_*\operatorname{IC}(S),
 \qquad
 \mathfrak n^+_{\mathrm{BPS}}
   =\bigoplus_{n\geq1}H^*(S,\operatorname{IC}(S))\,z^n
\]
is abelian.
Here $z$ records length and $\operatorname{IC}(S)=\mathbb Q_S[2](1)$.
The brackets vanish because every pair of charges is orthogonal.

The generator in each length has dimensions $1,22,1$ in degrees $-2,0,2$.
All these degrees are even.
Forgetting the cohomological grading, the BPS dimension series is consequently
\begin{equation}\label{cbc:hall:k3}
 \sum_{n\geq0}\dim(B_{\mathrm{Alg},n})z^n
   =\prod_{r\geq1}(1-z^r)^{-24}
   =1+24z+324z^2+3200z^3+\cdots.
\end{equation}

The length-two coefficient consists of $\operatorname{Sym}^2(\mathbb Q^{24})$, of dimension $300$, and the new isotropic generator space of dimension $24$.
Removing the latter would change the BPS algebra.
In a category with odd generators, the corresponding symmetric factors are exterior factors, so this dimension formula cannot be reused without its parity hypothesis.
For an odd isotropic generator $\theta$, the relation $[\theta,\theta]=0$ gives $2\theta^2=0$ in its enveloping algebra.

The determinant of the vector bundle of global sections is a line bundle on the object stack.
A scalar acts with weight $n$ in length $n$, so this determinant is positive.
Writing $t$ for cohomological degree, Theorem~\ref{cbc:hall:bps} gives
\[
 \operatorname{ch}_t B_{\mathrm{Alg},1}=t^{-2}+22+t^2,
 \qquad
 \operatorname{ch}_t\mathbb H^*(M,\mathscr H)_1
       =\frac{t^{-2}+22+t^2}{1-t^2}.
\]
For instance, degree $4$ contributes $24$ classes to the full length-one CoHA and none to its BPS subalgebra.
Thus the BPS enveloping theorem and the full Hall PBW theorem have different targets even in this basic symplectic-surface example.

\section{Quiver and surface Yangian comparisons}
Let $Q$ be a finite quiver, possibly with loops, and let $\Pi_Q$ be its preprojective algebra.
It is the path algebra of the doubled quiver modulo $\sum_{a\in Q_1}[a,a^*]$.
Use the arrow-scaling torus
$T=(\mathbb G_m)^{Q_1}\times\mathbb G_m$, where the last factor scales the symplectic form.
Its coefficient ring is
\[
 R=H_T^*(\mathrm{pt},\mathbb Q)
     =\mathbb Q[t_a,\hbar\mid a\in Q_1].
\]

Let $\mathcal H_R(Q)$ be the equivariant Borel--Moore CoHA of the stacks of finite-dimensional $\Pi_Q$-modules, with the virtual grading and Hecke-convolution normalization of Schiffmann--Vasserot~\cite[\S2]{SVYangians}.
Let $\mathcal H_R^\vee(Q)$ use their Lusztig-nilpotent substacks instead.
The superscript here distinguishes two coefficient carriers.

Their Yangian $\mathbf Y_R(Q)$ is defined by stable-envelope operators on Nakajima quiver varieties.
The nilpotent version $\mathbf Y_R^\vee(Q)$ acts on the corresponding central fibers.
Schiffmann--Vasserot construct algebra isomorphisms
\begin{equation}\label{cbc:hall:sv}
 \mathcal H_R(Q)\xrightarrow{\sim}\mathbf Y_R^-(Q),
 \qquad
 \mathcal H_R^\vee(Q)\xrightarrow{\sim}\mathbf Y_R^{-,\vee}(Q)
\end{equation}
intertwining every such geometric action~\cite[Theorem~4.13]{SVYangians}.
The superscript minus is their convention for the half corresponding to positive dimension vectors in the Hall algebra.
The comparison holds over $R$.
No completion or passage to the fraction field is required in its statement.

The map first compares the Hecke classes with stable-envelope operators over $\operatorname{Frac}(R)$.
Generation and the faithful family of vacuum evaluations make this comparison injective.
The integral action lattices and the triangular decomposition then identify its image over $R$~\cite[Proposition~4.1, Theorem~4.8 and Proposition~4.10]{SVYangians}.
These operations distinguish~\eqref{cbc:hall:sv} from a character identity.
The central polynomial algebras of framing-tautological classes used in that proof are additional extensions of the halves.

For the surface comparison, let $\Gamma_0\subset\operatorname{SL}_2(\mathbb C)$ be a nontrivial finite subgroup and let
$X\to\mathbb C^2/\Gamma_0$ be the minimal resolution.
Let $Z$ be its exceptional divisor and let $T_0$ be a diagonal torus commuting with $\Gamma_0$.
The source is the stack of coherent sheaves with set-theoretic support in $Z$.
It includes zero-dimensional sheaves and sheaves on infinitesimal thickenings of $Z$.
Write $\mathcal H^{T_0}_{X,Z}$ for its nilpotent equivariant CoHA with rational coefficients.
Its underlying topological space is the inverse limit of equivariant Borel--Moore homology over an admissible quasi-compact open exhaustion, taken in each grading~\cite[\S III.7.2]{DPPSSV}.

Fix an orientation of the finite ADE diagram.
In the Grothendieck group, write $\alpha_i$ for its simple-root classes and $\delta$ for the point class.
The comparison twists the Hall multiplication by $(-1)^{\Theta(d,e)}$, where
\[
 \Theta(\alpha_i,\alpha_j)=\langle\alpha_i,\alpha_j\rangle_Q\pmod2,
 \qquad \Theta(\delta,d)=\Theta(d,\delta)=0.
\]
The brackets on the right denote the Euler form of the oriented finite quiver.
We include this twist in the notation $\mathcal H^{T_0}_{X,Z}$.

The target half has a specific topology.
Let $\mathbf Y_Q^-$ be the standard negative half of the affine ADE Yangian, specialized to $H^*_{T_0}(\mathrm{pt},\mathbb Q)$.
Choose a coweight $\vartheta$ with positive coefficients at the finite vertices and $\vartheta(\delta)=0$.
For a nonzero dimension vector $d$, define its slope by
$\mu_\vartheta(d)=\vartheta(d)/\sum_i d_i$.
Let $\vartheta_f$ be its finite coweight and let $T_\lambda$ denote the affine braid operator attached to a lattice translation $\lambda$.

Define the left ideal and translated quotients
\[
 J_{\vartheta,0}
   =\sum_{\mu_\vartheta(d)>0}\mathbf Y_Q^-(\mathbf Y_Q^-)_{-d},
 \qquad
 V_{2\ell}
   =\frac{T_{4\ell\vartheta_f}(\mathbf Y_Q^-)}
          {T_{4\ell\vartheta_f}(J_{\vartheta,0})},
 \qquad \ell\in\mathbb Z.
\]
These are the quotients for the slope interval $(-\infty,0]$ and its translates in~\cite[(III.3.6), \S III.7.4]{DPPSSV}.

For the transition from level $0$ to level $2$, first restrict the slope interval to $(-1/(4h),0]$, where $h=\sum_i\delta_i$.
Equivalently, quotient further by
$\sum_{\mu_\vartheta(d)\leq-1/(4h)}(\mathbf Y_Q^-)_{-d}\mathbf Y_Q^-$.
The braid and PBW comparison identifies this quotient with $V_2$.
Translating this map gives the restriction maps $V_{2\ell-2}\to V_{2\ell}$ of~\cite[(III.7.4)]{DPPSSV}, and hence
\begin{equation}\label{cbc:hall:surface-limit}
 \mathbf Y^+_\infty=\varprojlim_{\ell\to+\infty}V_{-2\ell}.
\end{equation}

The quotients are vector-space quotients by left ideals, so their products cannot be assumed.
For a fixed output level, multiply lifts at a sufficiently deeper level in the translated Yangian and project.
The stabilization theorem makes the result independent of those lifts and defines the continuous multiplication on~\eqref{cbc:hall:surface-limit}.

\begin{theorem}[The supported surface half]\label{cbc:hall:surface}
With this coefficient ring, sign twist and topology, there is an isomorphism of graded topological algebras
\[
 \Theta_{X,Z}:\mathcal H^{T_0}_{X,Z}
       \xrightarrow{\sim}\mathbf Y^+_\infty.
\]
It intertwines tensoring sheaves by line bundles with the corresponding affine braid operators.
\end{theorem}
\begin{proof}
The derived McKay equivalence identifies perverse coherent sheaves supported in $Z$ with nilpotent preprojective modules.
Tensoring by increasingly dominant negative line bundles gives a sequence of hearts whose limiting heart is the category of coherent sheaves supported in $Z$.
The continuity theorem for their flag stacks identifies the surface CoHA with the limit of the corresponding slope quotients~\cite[Theorem~III.7.5]{DPPSSV}.

The quiver Yangian comparison identifies each quotient with $V_{2\ell}$.
Compatibility of derived reflection functors with braid operators identifies the transition maps and the multiplications of sufficiently deep lifts.
Taking inverse limits therefore gives mutually inverse continuous maps preserving multiplication~\cite[Theorem~III.7.8]{DPPSSV}.
The same compatibility identifies the line-bundle action~\cite[Proposition~III.7.14]{DPPSSV}.
\end{proof}

For the trivial torus, the target is a completed enveloping algebra.
If $\mathfrak g_f$ is the finite ADE Lie algebra, with positive nilpotent part $\mathfrak n_f^+$ and Cartan $\mathfrak h_f$, its Lie algebra is
\[
 \mathfrak n^+_{\mathrm{ell}}
 =\mathfrak n_f^+[s^{\pm1},t]
   \oplus s^{-1}\mathfrak h_f[s^{-1},t]\oplus K^-.
\]
Here $K^-$ is the negative-$s$-degree part of the center in the universal central extension of $\mathfrak g_f[s^{\pm1},t]$.
The completion of $U(\mathfrak n^+_{\mathrm{ell}})$ is the slope completion corresponding to~\eqref{cbc:hall:surface-limit}~\cite[Theorem~III.7.25]{DPPSSV}.
In particular, the nonstandard half is not specified merely by choosing the usual positive roots of the affine Lie algebra.

\section{Cooperation and geometric reconstruction}
Even the vector-space Hall algebra illustrates the tensor data needed for a coproduct.
Keep the prime power $q>1$ fixed and work over $\mathbb Q$.
In charge-graded vector spaces, use the braiding
\[
 v_b\otimes w_c\longmapsto q^{-bc}w_c\otimes v_b.
\]
The induced product on two copies of the Hall algebra is
\[
 (u_a\otimes u_b)(u_c\otimes u_d)
     =q^{-bc}(u_a*u_c)\otimes(u_b*u_d).
\]

\begin{proposition}[The vector-space Hall coproduct]\label{cbc:hall:braided}
For this braided tensor product,
\[
 \Delta u_n=\sum_{a+b=n}q^{-ab}u_a\otimes u_b,
 \qquad \epsilon(u_n)=\delta_{n,0}
\]
defines a graded bialgebra.
Its antipode is
$S(u_n)=(-1)^nq^{-n(n-1)/2}u_n$.
The same formulas are continuous in the positive charge completion.
\end{proposition}
\begin{proof}
Put $[n]_q!=\prod_{j=1}^n(1+q+\cdots+q^{j-1})$ and $x=u_1$.
Equation~\eqref{cbc:hall:gaussian} gives $x^n=[n]_q!u_n$.
In the braided tensor algebra, $Y=1\otimes x$ and $X=x\otimes1$ satisfy $YX=q^{-1}XY$.
The $q^{-1}$-binomial identity therefore gives
\[
 (X+Y)^n
   =\sum_{a+b=n}{n\brack a}_{q^{-1}}X^aY^b,
 \qquad
 {n\brack a}_{q^{-1}}=q^{-ab}{n\brack a}_q.
\]
Dividing by $[n]_q!$ gives the displayed coproduct.
It is multiplicative because it was defined by $x\mapsto X+Y$ on the polynomial algebra generated by $x$.

The coefficient of $u_a\otimes u_b\otimes u_c$ in either iterated coproduct is $q^{-ab-ac-bc}$.
This proves coassociativity, and the endpoint terms prove the counit identities.
The connected graded bialgebra has the recursively determined antipode.
Braided reversal of $n$ copies of $x$ contributes $q^{-n(n-1)/2}$, giving its stated formula.
Every fixed-charge sum is finite, so all operations extend continuously.
\end{proof}

If one deletes the crossing factor while keeping this coproduct, multiplicativity fails.
The coefficient of $u_1\otimes u_1$ in $\Delta(u_1*u_1)$ is $1+q^{-1}$, whereas the ordinary tensor square of $\Delta u_1$ has coefficient $2$.
This calculation fixes both the carrier and the failed operation.
The shuffle obstruction~\eqref{cbc:l027:obstruction-relation} makes the corresponding tensor distinction on its own collision carrier.
Its rational exchanged recipient remains a positive construction with the stated removed divisors.

To obtain an ordered chiral Hall object, one must construct the coefficient objects over configurations, their separated factorization maps and their collision operations.
The extension-flag comparisons must commute with those maps and with the chosen charge-zero unit.
For a critical Hall theory, this includes descent of the vanishing-cycle coefficients and their orientation.
Supply the interval family, compatible augmentation and strong monoidal coefficient functor required by Theorem~\ref{cbc:ord:recognition}.
These data transport the Hall multiplication to an ordered chiral multiplication.
The Hall algebra on global cohomology alone does not give these objects or their descent.

A geometric coproduct must further specify its supported or localized recipient and the interchange governing that recipient.
A completed double requires compatible negative and neutral halves, Hopf operations and a continuous Hopf pairing.
A universal $R$-tensor additionally requires convergence of the pairing's canonical tensor.
The BPS theorem and the Yangian-half comparisons preserve their full algebraic content while leaving those geometric constructions separate.
For the supported surface algebra, even the geometric realization of its double is a further problem~\cite[Introduction, ``Representations'']{DPPSSV}.

\chapter{The Igusa denominator as a bar character}
Let $\g_\Delta$ be the Gritsenko--Nikulin Borcherds--Kac--Moody superalgebra and let $\nilp$ be its positive nilpotent part. Root spaces are finite-dimensional and the positive cone is locally finite. Fix a proper integral height $h$ that is positive on the effective root cone, and set
\[
 F^{>H}\nilp=\bigoplus_{h(\alpha)>H}\g_{\Delta,\alpha},
 \qquad
 \nilp^{(H)}=\nilp/F^{>H}\nilp.
\]
Height additivity makes $F^{>H}\nilp$ an ideal; properness makes $\nilp^{(H)}$ finite-dimensional and nilpotent. Every bar, Chevalley, enveloping, and character formula in this chapter is interpreted coefficientwise in the separated height completion.

\begin{theorem}[Bar--Chevalley identification]
For each height $H$ there is a natural quasi-isomorphism of conilpotent coalgebras
\[
 \Barc\bigl(U(\nilp^{(H)})\bigr)\simeq C_*(\nilp^{(H)}).
\]
The quotient maps are compatible with these quasi-isomorphisms, and therefore induce a quasi-isomorphism in the separated height-complete category
\[
 \Barc_h^\wedge\bigl(U(\nilp)\bigr):=\varprojlim_H\Barc\bigl(U(\nilp^{(H)})\bigr)
 \simeq
 C_{*,h}^\wedge(\nilp):=\varprojlim_H C_*(\nilp^{(H)}).
\]
The coefficientwise Euler supercharacter of the completed Chevalley complex is
\[
 \chi\,C_{*,h}^\wedge(\nilp)
 =\prod_{\alpha>0}(1-e^{-\alpha})^{\sdim\g_{\Delta,\alpha}}.
\]
\end{theorem}
\begin{proof}
For every $H$, the PBW filtration on $U(\nilp^{(H)})$ identifies the associated graded with $\Sym(\nilp^{(H)})$. The bar complex of the enveloping algebra and the Chevalley complex have the same associated-graded Koszul resolution. The standard antisymmetrization map is a filtered quasi-isomorphism. Naturality for the quotient maps gives the inverse-limit comparison; proper height makes every homogeneous coefficient stabilize. For an even root vector the exterior Chevalley factor contributes $1-e^{-\alpha}$; for an odd root vector it contributes its inverse. Multiplication over the root spaces gives the product.
\end{proof}

With the positive-root convention $e^{-\rho}=q^{1/2}r^{1/2}s^{1/2}$, the denominator theorem reads
\[
 e^{-\rho}\,\chi\,C_{*,h}^\wedge(\nilp)=D_5(2\Omega)=2^{-6}\Delta_5(2\Omega)
\]
\cite{GN97,GN98b}. This identity is the exact meeting point of ordered bar theory and automorphic products.

\chapter{Finite-height currents and an exact physical model}
Retain the finite-height quotients $\nilp^{(H)}$ fixed in the preceding chapter.

The current Lie conformal superalgebra is
\[
 \operatorname{Cur}(\nilp^{(H)})=\kfield[\partial]\tensor\nilp^{(H)},
 \qquad [a_\lambda b]=[a,b].
\]
Its universal enveloping vertex superalgebra is $V_H$. The current generators carry the cocommutative vertex-bialgebra coproduct
\[
 \Delta\bigl(a(-1)\mathbbm1\bigr)
 =a(-1)\mathbbm1\tensor\mathbbm1
 +\mathbbm1\tensor a(-1)\mathbbm1,
 \qquad a\in\nilp^{(H)},
\]
extended uniquely as a vertex-algebra morphism.

\begin{theorem}[Height-complete current realization]
Let $\mathfrak n=\bigoplus_{\alpha>0}\mathfrak n_\alpha$ be a root-graded Lie superalgebra with finite-dimensional root spaces and a proper additive height function.  Put
\[
 \mathfrak n^{(H)}=\mathfrak n/\!\bigoplus_{h(\alpha)>H}\mathfrak n_\alpha.
\]
Then every $\mathfrak n^{(H)}$ is finite-dimensional and nilpotent, the quotient maps form an inverse system, and
\[
 \widehat U_h(\mathfrak n)=\varprojlim_HU(\mathfrak n^{(H)}),\quad
 \widehat C_{*,h}(\mathfrak n)=\varprojlim_HC_*(\mathfrak n^{(H)}),\quad
 \widehat V_h(\mathfrak n)=\varprojlim_HU^{\ch}(\Cur_X\mathfrak n^{(H)})
\]
are separated complete algebras, coalgebras, and current algebras.  Their associated graded objects are the corresponding completed symmetric objects, and the bar--Chevalley comparison is compatible with the inverse system.
\end{theorem}
\begin{proof}
Properness gives finitely many roots below a fixed height.  Additivity implies that the roots of height $>H$ form a Lie ideal and that the quotient is nilpotent.  PBW and Chevalley antisymmetrization are natural under quotient maps.  Taking inverse limits in each finite total height therefore gives the displayed complete objects and preserves the coefficientwise quasi-isomorphism.
\end{proof}

Apply the finite BF model of Chapter 19 to every $\nilp^{(H)}$. The inverse system gives a height-complete holomorphic--topological classical field theory with completed boundary algebra and interval complex
\[
 U_h^\wedge(\nilp):=\varprojlim_H U(\nilp^{(H)}),
 \qquad
 C_{*,h}^\wedge(\nilp):=\varprojlim_H C_*(\nilp^{(H)}).
\]
This gives an exact algebraic realization, a local classical holomorphic--topological realization, and an exact quantum realization of the finite-dimensional zero-mode sectors of the positive enveloping algebra and its denominator character. The compact $K3\times E$ identification is the geometric comparison studied in the companion volumes.

\chapter{Character, bracket, parity, and orientation}
A character is an Euler invariant.  Faithful reconstruction combines it with the bracket, parity decomposition, and orientation line carried by the intended object.

\begin{theorem}[Three reconstruction fibres]
The following fibres contain distinct structured objects with the same scalar image.
\begin{enumerate}
\item The PBW-character fibre of a graded vector space with basis $x,y,z$ in degrees $a,b,a+b$ contains both the abelian Lie algebra and the Heisenberg Lie algebra $[x,y]=z$.
\item The superdimension-zero fibre contains both the zero super vector space and $\kfield^{1|1}$.
\item The squaring fibre of a line-valued section contains $s$ and $s\otimes\chi$ for every sign local system $\chi$.
\end{enumerate}
Consequently a faithful reconstruction datum consists respectively of a bracket, a parity-refined dimension, and an oriented square-root line.
\end{theorem}
\begin{proof}
PBW identifies the graded vector-space character of an enveloping algebra with the symmetric-superalgebra character of its underlying graded Lie space, so the two displayed brackets have the same PBW character.  The two super vector spaces in the second item both have superdimension zero and have parity dimensions $(0,0)$ and $(1,1)$.  Finally $\chi^{\otimes2}\cong\mathbbm1$, hence
\[
 (s\otimes\chi)^{\otimes2}\cong s^{\otimes2}.
\]
The extra structures listed in the conclusion distinguish each pair.
\end{proof}

\begin{recognition}[Minimal PBW recognition]
Let $\phi:\mathfrak g\twoheadrightarrow\mathfrak h$ be a surjective morphism of positively graded Lie superalgebras with finite-dimensional homogeneous pieces.  If the parity-refined Hilbert series of $U(\mathfrak g)$ and $U(\mathfrak h)$ agree, then $\phi$ is an isomorphism.  Equality of signed PBW characters suffices when the parity of every primitive degree is prescribed in advance.
\end{recognition}
\begin{proof}
If the kernel is nonzero, choose its first positive degree.  In that degree all decomposable PBW contributions agree by minimality, while the primitive even or odd dimension in $U(\mathfrak h)$ is strictly smaller.  This contradicts equality of parity-refined Hilbert series.  With prescribed parity, the formal logarithm of the PBW product recovers the signed primitive dimension inductively in height, and prescribed parity recovers the ordinary dimension.
\end{proof}

\chapter{Exact dependency theorem}
\begin{theorem}[Ordered-chiral closure]
Fix:
\begin{enumerate}
\item a six-functor Ran category and its factorization subcategory.
\item a transversely ordered algebra, an associative chiral algebra, or a chiral Lie algebra, as appropriate.
\item every mixed distributive, completion, duality, trace, or analytic datum named by the intended construction.
\item for a formal-disc field realization, the state space, vacuum and field map, together with the operation comparisons of Proposition~\ref{cbc:two:field-realization}.
\end{enumerate}
The three bars, their complete cobars and their module/comodule Koszul transforms are defined functorially under their stated hypotheses.
The same holds for the Hochschild center, ribbon theory and finite-height current realization under the hypotheses of those constructions.
They satisfy the identities proved in this volume.
The supplied formal-disc realization satisfies the corresponding conditional field-recognition statement.
Every passage between the pointwise and chiral-convolution monoidal products uses the displayed mixed distributive datum.
\end{theorem}
\begin{proof}
Each output is one of the constructions in the preceding chapters, in its specified source and target categories.
The comparison statements use the supplied distributive, PBW, completion, duality, trace, or realization datum.
Composition of the constructed functors gives the asserted functorial system.
The formal-disc assertion follows from Proposition~\ref{cbc:two:field-realization}.
The reverse geometric functor remains the construction sought in Problem~\ref{cbc:two:reverse}.
\end{proof}

\chapter{Synthesis}
The invariant language has three algebraic axes.  Transverse order is an $\Eone$ structure for $\tensor^!$.  Associative chirality is an $\Eone$ structure for $\starCh$.  Chiral Lie theory is a Lie structure for $\starCh$.  Their bars are $\BarPerp$, $\BarAssCh$, and $\BarLieCh$.  The aligned projector supplies the exact sector on which a mixed distributive law is defined.  Every comparison in the theorem system factors through one of these typed maps.

The constructed operations form the two functorial rows
\[
 \text{HT boundary theory}
 \xrightarrow{\ \mathrm{Obs}_{\partial}\ }
 \Alg_{\Eone}(\Fact_D(X))
 \xrightarrow{\ \Barc^{\ord}\ }
 \Coalg_{\mathrm{conil,cpl}}(\Fact_D(X)),
\]
\[
 \Alg_{\Eone}(\Fact_D(X))
 \xrightarrow{\ Z_{\ord}\ }
 \Alg_{\Etwo}(\Fact_D(X)).
\]
A primitive Lie algebra enters when a separate Chevalley or Hall comparison identifies the bar coalgebra with the Chevalley chains of that Lie algebra. Hall geometry supplies this comparison through extension correspondences. Quantum groups require an additional coproduct and exchange structure. Automorphic forms arise from graded Euler characters after a grading and completion have been fixed. Each arrow therefore preserves the geometric operation named at its source.

\part{Technical appendices to the preceding book}
\chapter{Sign calculation for the length-three bar}
For homogeneous $a,b,c\in\bar A$, write $[a|b|c]=sa\tensor sb\tensor sc$. With the standard suspension convention,
\[
 b_m[a|b|c]
 =(-1)^{|a|}[ab|c]
 +(-1)^{|a|+|b|+1}[a|bc].
\]
Applying $b_m$ again gives
\[
 b_m^2[a|b|c]
 =(-1)^{|b|}\bigl[(ab)c-a(bc)\bigr],
\]
which vanishes by associativity.

\chapter{Rational \texorpdfstring{$R$}{R}-matrix verification}
Let $P_{ij}$ permute tensor factors $i$ and $j$. Expanding
\[
 R_{12}(u)R_{13}(u+v)R_{23}(v)
\]
and its reversed product shows equality coefficientwise. The cubic coefficient uses
\[
 P_{12}P_{13}P_{23}=P_{23}P_{13}P_{12},
\]
the quadratic coefficients use the symmetric-group relations, and the linear coefficients agree termwise.

\chapter{Carriers and comparison maps}
The five-carrier geometry, transverse bar construction, mixed-forest signs, and noncommutative examples use four kinds of categorical data:
\begin{enumerate}
\item a factorization carrier is paired explicitly with its operation;
\item the Ran category is fixed before its tensor products are used;
\item the block-ordered operad is introduced after a representing chiral operad is chosen;
\item ribbon and stable-curve sewing are derived from nondegenerate Calabi--Yau and modular data.
\end{enumerate}
A comparison between these constructions therefore specifies its source and target categories and the operations preserved by its map.
